% Generated by scripts/gen_blueprint.py from `paper/core.tex`; do not edit.

\chapter{Introduction}\label{sec:intro}

\begin{theoremA}[Every dimension]\label{thm:A}
\lean{Z3NetA.A_sep_rowScale, Z3NetA.A_hom, Z3NetA.A_net, Z3NetA.A_covering, Z3NetA.A_sep, Z3NetA.A_sep_real, Z3NetA.A_sep_eucl, Z3NetA.covering_of_net, Z3NetA.Ustar_adm, Z3NetA.Ustar_reach, Z3NetA.Ustar_data, Z3NetA.Ustar_babai, Z3NetA.Ustar_growth, Z3NetA.aStar_le, Z3NetA.aStar_three, Z3NetA.Ainf_fin_le, Z3NetA.prodAp_fin_le, Z3NetA.Ainf_inf_le, Z3NetA.prodAp_inf_le, Z3NetA.hS_le_Hs, Z3NetA.aStar_bound, Z3NetA.aStar_bound_log, Z3NetA.A_sharp, Z3NetA.A_sharp_real, Z3NetA.thmA, Z3NetA.Apoint_indep, Z3NetA.placeA, Z3NetA.AdmBase, Z3NetA.UstarR, Z3NetA.UrS, Z3NetA.betaS, Z3NetA.Apoint}\leanok
\uses{def:Ns, def:Ustar}
Let $s\ge3$, let $q\ge2$ be an integer all of whose prime factors are at least $s-1$, and let $k\ge1$. Let $\mathcal P^{U^*_s}_{k,q}$ be the digital net over $\Z/q\Z$ of Definition~\ref{def:Ns} for the places $\infty,0,1,\dots,s-2$ and the design $U^*_s$ of Definition~\ref{def:Ustar}, with $N=q^{sk}$ points. Then:
\begin{enumerate}
\item $\mathcal P^{U^*_s}_{k,q}$ is a digital $(0,sk,s)$-net over $\Z/q\Z$, every cell $\prod_i[n_iq^{-k},(n_i+1)q^{-k})$ holds exactly one point, and $h\le\sqrt s\,N^{-1/s}$;
\item if $q\ge5a^*_s+4$, its toroidal (hence also Euclidean) minimum distance satisfies $\delta\ge(1-a^*_s/q)\,N^{-1/s}$, for every $k\ge1$; for $k\ge2$ the constant $a^*_s$ is first-order sharp for this net: $\delta N^{1/s}=1-a^*_s/q+O(q^{-2})$ as $q\to\infty$.
\end{enumerate}
Here $a^*_3=2$, and $a^*_s\le2(s-1)^{3/2}(3/2)^{s-2}(s-2)!=e^{s\log s+O(s)}$ for $s\ge4$.
\end{theoremA}

\begin{proof}[Proof of Theorem~\ref{thm:A}]
\leanok
\uses{lem:box, lem:Ss, prop:Gs, def:Ns, cor:tsk, prop:F, prop:D, lem:sat, lem:adelic, lem:babai}
\emph{Reachability.} By Lemma~\ref{lem:sat}(iii) every column $E_{d,j}$ is admissible at carry depth $j$, with polynomial $g_{d,j}\in R_S[X]$ of degree $j$ and leading coefficient $\gamma_j$, an $S$-unit; so by Proposition~\ref{prop:F} the off-diagonal part of $U^*_s$ is reachable with $\lambda'\equiv1$ over $R'=R_S$, and the diagonal is free (Lemma~\ref{lem:box}(a)).

(1) Every prime factor of $q$ is at least $s-1>\nu$, so $\gcd(q,\nu!)=1$, and Corollary~\ref{cor:tsk} gives a digital $(0,sk,s)$-net with one point in every cell, for every $k\ge1$. A cell has diameter $\sqrt s\,q^{-k}=\sqrt s\,N^{-1/s}$, which bounds $h$.

(2) \emph{Lower bound.} (H1) and (H2) of Lemma~\ref{lem:Ss} hold by Proposition~\ref{prop:Gs} and Definition~\ref{def:Ns}, with $L=\min(k,s-1)$. $U^*_s$ is diagonally dominant by construction, with the maximum over all carry depths, so Proposition~\ref{prop:D} gives $\delta q^k\ge1-a^*_s/q$ for $q\ge5a^*_s+4$ and every $k\ge1$.

\emph{Sharpness} ($k\ge2$). Let $i$ attain $a^*_s=\alpha_i$. For $q>2a^*_s$ the entries of $F_0$ are centred, so $\tilde M=F_0$. By (H2) the own-axis step $e_i$ from a cell whose last digit of $n_i$ is not $q-1$ changes coordinate $i$ by exactly $1-a^*_s/q$ cells if its first offset digit is at least $a^*_s$ (no wrap, no jitter), and every other coordinate $u$ by less than $(|\tilde M_{ui}|+1)/q$ if $r_u$ avoids the $|\tilde M_{ui}|$ values from which adding $\tilde M_{ui}$ wraps. For each $u\ne i$ the $t=0$ box with depths $k+1$ on $i$ and $u$ and $0$ elsewhere (total $2k+2\le sk$) makes $(n_i,r_i,n_u,r_u)$ jointly uniform (that box fixes exactly these $2k+2$ digits, so every value of them occurs equally often), so among the cells satisfying the conditions on $n_i$ and $r_i$ the fraction with a wrapping $r_u$ is $|\tilde M_{ui}|/q$; for $q>\sum_{u\ne i}|\tilde M_{ui}|$ a union bound leaves a cell. Its pair of points lies inside the unit cube, so
\[
\delta N^{1/s}<\Bigl(\bigl(1-\tfrac{a^*_s}q\bigr)^2+\sum_{u\ne i}\frac{(|F_0[u][i]|+1)^2}{q^2}\Bigr)^{1/2}=1-\frac{a^*_s}q+O(q^{-2}).
\]

\emph{The bound on $a^*_s$.} Let $s\ge4$. By Lemma~\ref{lem:babai} it suffices to show $\eta_{d,j}\le \Xi_s:=2\sqrt{s-1}\,(s-2)!$ for every column $d$ and carry depth $j$; then every entry of $U^*_s$ is at most $\max_{d,j}\|E_{d,j}\|_2\le(3/2)^{s-2}\Xi_s$, and $a^*_s\le(s-1)(3/2)^{s-2}\Xi_s=2(s-1)^{3/2}(3/2)^{s-2}(s-2)!$. By Lemma~\ref{lem:adelic} it suffices to give, at each place of $\mathbb Q$, one monic $g$ of degree $j$.

\emph{Finite column $d$.} Archimedean place: all nodes $y_c$ lie in $[-1,1]$ and $|w_c|\le1$, so the monic Chebyshev polynomial $2^{1-j}T_j$ ($j\ge1$; $g=1$ for $j=0$) gives $A_{j,\R}\le2\sqrt{s-1}$. At $p\in S$: let $\mathcal R_p$ be the finite rows $c\ne d$ with $p\mid c-d$, $e_c=v_p(c-d)\ge1$ for $c\in \mathcal R_p$, let $\mathcal R'\subset \mathcal R_p$ consist of $\min(j,|\mathcal R_p|)$ rows with the largest $e_c$, and $g_p=\prod_{r\in \mathcal R'}(X-y_r)\cdot X^{j-|\mathcal R'|}$. For $c\in \mathcal R'$ the entry vanishes. For $c,r\in \mathcal R_p$, $|y_c-y_r|_p=|(r-c)/((c-d)(r-d))|_p\le p^{\max(e_c,e_r)}$; so for $c\in \mathcal R_p\setminus \mathcal R'$ (then $|\mathcal R'|=j$ and $e_c\le e_r$ for $r\in \mathcal R'$) the entry has $p$-adic size at most $p^{e_c+\sum_{r\in \mathcal R'}e_r}\le p^{\sum_{r\in \mathcal R_p}e_r}$. For a finite row $c\notin \mathcal R_p$, and for the row $\infty$ ($y=0$, $w=1$), $|w_c|_p=1$, $|y_c|_p\le1$ and $|y_c-y_r|_p\le p^{e_r}$, so the size is at most $p^{\sum_{r\in \mathcal R'}e_r}$. Hence $A_{j,p}\le p^{\sum_{r\in \mathcal R_p}e_r}=p^{v_p(d!\,(\nu-d)!)}$, and $\prod_{p\in S}A_{j,p}\le d!\,(\nu-d)!\le \nu!$, since every prime factor of $d!(\nu-d)!$ lies in $S$. So $\eta_{d,j}\le2\sqrt{s-1}\,\nu!$.

\emph{Column $\infty$.} The nodes are $y_c=c\in\{0,\dots,\nu\}$ with $w_c=1$. At the archimedean place the monic Chebyshev polynomial of $[0,\nu]$ gives $A_{j,\R}\le2\sqrt{s-1}\,(\nu/4)^j$. At $p\in S$ the falling factorial $X(X-1)\cdots(X-j+1)$ takes the values $j!\binom cj$, so $A_{j,p}\le|j!|_p$ and $\prod_pA_{j,p}\le1/j!$. So $\eta_{\infty,j}\le2\sqrt{s-1}\,(\nu/4)^j/j!\le2\sqrt{s-1}\,e^{\nu/4}\le2\sqrt{s-1}\,\nu!$ for $\nu\ge2$.

Finally, with $\nu=s-2$ the elementary bounds $\nu^\nu e^{-\nu}\le \nu!\le \nu^\nu$ give $\bigl|\log\bigl(2(s-1)^{3/2}(3/2)^{s-2}(s-2)!\bigr)-s\log s\bigr|\le3s$ for $s\ge4$, so the bound is $e^{s\log s+O(s)}$; and $a^*_3=2$ is the display before Table~\texttt{tab:astar}. (At the archimedean place the Chebyshev polynomial even gives $A_{j,\R}\le\sqrt{s-1}$ for a finite column, since $|2^{1-j}T_j|\le1$ on $[-1,1]$ for $j\ge1$; the factor $2$ in $\Xi_s$ is slack.)
\end{proof}


\part{Existence}\label{part:existence}


\chapter{Graph forms and the one-digit relaxation}\label{sec:prelim}

\begin{lemma}[Box criterion; {Niederreiter \cite[Thm.~6.1]{Nie87}}]\label{lem:nie}
\lean{Z3NetF.nie_box, Z3NetF.nie_net, Z3NetF.nie_net_field, Z3NetF.nie_zmod}\leanok
For each $d$, every elementary box with $d_i$ digits in coordinate $i$ contains exactly one point if and only if the box minor for $d$ is a unit of $R$. In particular the point set is a $(0,m,s)$-net in base $q$ if and only if every box minor is a unit of $R$, i.e., over a field, nonzero.
\end{lemma}

\begin{proof}
\leanok
The point with index $v$ lies in a given box of shape $d$ if and only if the box matrix maps $v$ to the prescribed leading digits. So the boxes of shape $d$ correspond bijectively to the values of the linear map $v\mapsto(\text{box matrix})\,v$ on $R^m$. There are $q^m$ points and $q^m$ boxes, so every box holds exactly one point if and only if this map is bijective, i.e. if and only if the box matrix is invertible, i.e. its determinant is a unit. This is \cite[Thm.~6.1]{Nie87}.
\end{proof}


\section{Graph forms}

\begin{lemma}[Box condition and gauge moves]\label{lem:box}\label{lem:gauge}
\lean{Z3NetF.box_minor, Z3NetF.box_minor_self, Z3NetF.gauge_diag, Z3NetF.gauge_LU, Z3NetF.gauge_scale3, Z3NetF.gauge_net, Z3NetF.gauge_net3, Z3NetF.t3s4_reach, Z3NetPI.ppinv_window}\leanok
\uses{lem:nie}
The box minor of a graph form for $d$ equals, up to sign, the determinant of the square matrix formed by the top $d_c-k_c$ offset rows of the over coordinates $c$, restricted to the last $k_e-d_e$ levels of the under coordinates $e$ (it is $1$ if $d=\mathbf k$). In particular the diagonal blocks $\Phi_{cc}$ never enter. Consequently the following \emph{gauge moves} multiply every box minor by a unit of $R$, and so, by Lemma~\ref{lem:nie}, map $(0,m,s)$-nets over $\Z/q\Z$ or $\F_q$ to $(0,m,s)$-nets: (a) replacing the diagonal blocks arbitrarily; (b) for $s=3$ only: $\Phi_{cd}\mapsto\mu_{cd}\Phi_{cd}$ ($c\ne d$) with units $\mu_{cd}$ of $R$ (\emph{block scalings}); (c) $\Phi_{cd}\mapsto L_c\Phi_{cd}U_d$ ($c\ne d$) with $L_c$, $U_d$ lower-triangular whose diagonal entries are units of $R$.
\end{lemma}

\begin{proof}
\leanok
The first $\min(d_c,k_c)$ rows of $C_c$ are unit rows selecting the first $\min(d_c,k_c)$ levels of $n_c$. Laplace expansion along all these unit rows removes the selected columns; what remains are the top $d_c-k_c$ offset rows of the over coordinates, on the unselected columns, which are the levels $\ge d_e$ of the under coordinates. A coordinate is not both over and under, so no diagonal block occurs. (c) Left multiplication by $L_c$ preserves the span of the top rows of $\Phi_{c\cdot}$, and right multiplication by $U_d$ the span of the last columns of $\Phi_{\cdot d}$; the determinant changes by the product of the corresponding diagonal entries. (b) For $d\ne \mathbf k$ both over and under coordinates exist; for $s=3$ there is therefore a single over or a single under coordinate, so the matrix is $[\Phi_{cd}\,|\,\Phi_{cd'}]$ or $[\Phi_{cd};\Phi_{c'd}]$ restricted in this way, and scaling one block multiplies its determinant by a power of $\mu_{cd}$. For $s\ge4$ a box can have two over and two under coordinates, for instance $(k+1,k+1,k-1,k-1)$, and single-block scalings are no longer gauge moves.
\end{proof}

\begin{definition}[Geometric digit and geometric layer]\label{def:geom}
Let a graph form have the depths $(k,\dots,k)$. The first offset digit of a coordinate $c$ (its digit with index $k$) is its \emph{geometric digit} $r_c$. The \emph{geometric layer} is the integer matrix $\tilde M$ of the centred representatives $\tilde M_{cd}\in(-q/2,q/2]$ of the level-$(k-1)$ entries $(\Phi_{cd})_{0,k-1}$ of the first offset rows, diagonal included.
\end{definition}


\section{The one-digit relaxation}

\begin{lemma}[One-digit relaxation]\label{lem:Ss}
\lean{Z3NetF.Ss_coord, Z3NetF.Ss_sep, Z3NetU.U_row0, Z3NetF.Ss_sep_rowScale}\leanok
Let $q\ge3$, $k\ge1$, and let a digital net over $\Z/q\Z$ be given in graph form with depths $(k,\dots,k)$ such that (H1) row $0$ of every block $\Phi'_{cd}$, diagonal included, equals a design $(F_0,\dots,F_{L-1})$, $1\le L\le k$, on the levels $k-1,\dots,k-L$ and vanishes on the levels $0,\dots,k-1-L$; (H2) the rows $\ge1$ of the diagonal blocks vanish. In cell units write coordinate $i$ of a point as $n_i+(r_i+f_i)/q$ with $r_i$ its first offset digit and $f_i\in[0,1)$. For a step $e\in\{-1,0,1\}^s\setminus\{0\}$ to a toroidally adjacent cell and carry depths $j_d$ on the moving axes ($j_d$ is the number of further levels that change, by a carry if $e_d=1$ and by a borrow if $e_d=-1$), put
\[
\sigma_i=\sum_de_d\,M^{(\min(j_d,L-1))}[i][d],\qquad \bar\sigma_i=\sigma_i\bmod q\in[0,q),\qquad \mathrm{jit}_i=\begin{cases}0&e=\pm e_i,\\1&\text{otherwise,}\end{cases}
\]
and $b_i=\min_t\max(0,|e_iq+t|-\mathrm{jit}_i)$, the minimum over $t\in\{\bar\sigma_i,\bar\sigma_i-q\}$ if $\bar\sigma_i\ne0$ and over $t=0$ if $\bar\sigma_i=0$. Then the two points in cells $n$ and $n+e$ satisfy $q|\Delta_i|\ge b_i$ for every $i$, where $\Delta_i$ is coordinate $i$ of the point in cell $n+e$ minus that of the point in cell $n$, in cell units ($n+e$ taken as the toroidal neighbour), and the toroidal minimum distance satisfies
\[
\delta\,q^k\ \ge\ \mathrm{LB}:=\min\Bigl(1,\ \min_{e,j}\bigl(\textstyle\sum_ib_i^2\bigr)^{1/2}/q\Bigr).
\]
\end{lemma}

\begin{proof}
\leanok
A carry of depth $j$ on axis $d$ changes the levels $k-1,\dots,k-1-j$ of $n_d$ by $e_d$ modulo $q$; the torus wrap has $j=k-1$ (at $k=1$ every step has $j=0$). By (H1), $r_i$ changes by $\sigma_i$ modulo $q$, and capping the carry depth at $L-1$ is exact because row $0$ vanishes on the levels $0,\dots,k-1-L$. Since $r'_i-r_i\in(-q,q)$ and $r_i'-r_i\equiv\sigma_i$, we have $r'_i-r_i\in\{\bar\sigma_i,\bar\sigma_i-q\}$, and $r'_i-r_i=0$ if $\bar\sigma_i=0$. Always $|f'_i-f_i|<1$, and $f'_i=f_i$ if only axis $i$ moves: by (H2) the deeper digits of $i$ do not depend on $n_i$. So $q\Delta_i=e_iq+(r'_i-r_i)+(f'_i-f_i)$ gives $q|\Delta_i|\ge b_i$.

It remains to compare with the toroidal distance; this is the only place where the size of $q^k$ enters. In cell units the torus has period $G=q^k\ge3$. Let $n\ne n'$ be two cells, let $E\in\Z^s$ be the representative of $n'-n$ modulo $G$ with $|E_c|\le G/2$, and let $\theta_c\in(-1,1)$ be coordinate $c$ of $(r'-r+f'-f)/q$. The toroidal distance in coordinate $c$ is $\operatorname{tor}_c=\min_{z\in\Z}|E_c+zG+\theta_c|$. If some $|E_c|\ge2$, then $|E_c+zG|\ge|E_c|$ for every $z$, so $\operatorname{tor}_c\ge|E_c|-|\theta_c|>1$. Otherwise $E=e\in\{-1,0,1\}^s\setminus\{0\}$ and $\Delta=e+\theta$ with $|\Delta_c|<2$; if $|\Delta_c|\le G/2$, then $\operatorname{tor}_c=|\Delta_c|\ge b_c/q$, and otherwise $\operatorname{tor}_c\ge G-|\Delta_c|>G-2\ge1$. So either some coordinate has toroidal distance at least one cell, or the distance is at least $(\sum_ib_i^2)^{1/2}/q$ cells; as a cell has side $q^{-k}$, this is the claim.
\end{proof}

\begin{remark}[First-order expansion]\label{rem:K}
\lean{Z3NetF.remK_bad, Z3NetF.remK_good}\leanok
Call a moving coordinate \emph{backward} if $e_i\sigma_i\le0$ and \emph{forward} otherwise. For a fixed design and $q\to\infty$: if some configuration with at least two moving axes has no backward coordinate, or some diagonal partial sum $M^{(j)}[a][a]$ is positive, then $\mathrm{LB}=O(1/q)$; otherwise $\mathrm{LB}=1-K/q+O(q^{-2})$ with $K$ the maximum of $|M^{(j)}[a][a]|$ and of $|\sigma_b|+1$ over the configurations with at least two moving axes and a single backward coordinate $b$. (Two backward coordinates cost nothing at first order, and the jitter costs the $+1$.) In dimension three, $M=-I+S$ with $S$ skew with entries $\pm1$ gives $K=1$.
\end{remark}


\section{The lattice picture}

\begin{remark}[A two-dimensional example]\label{rem:gk2d}
The permutation $(0,m,2)$-nets of Gr\"unschlo{\ss} and Keller \cite[Sec.~3]{GK09} for even $m$ show this mechanism in the smallest case, base~$2$ and dimension two (Figure~\texttt{fig:arcgk}). For even $m=2r\ge4$, $Q=2^r$, $N=Q^2$ and $0\le i,j<Q$ they are the points $(X,Y)/N$ with $X=\bigl(Q\varrho(i)+\lfloor i/4\rfloor+\tfrac Q4(j+1)\bigr)\bmod N$, $Y=i+Qj$, where $\varrho$ reverses $r$ bits; distances are taken modulo the lattice $\langle(N,0),(N/4,N)\rangle$ of the row-shifted tiling for which they were designed \cite{GK09,Rud26} (the tiling metric).
(i) \emph{Cells and geometric digit.} Since $\lfloor i/4\rfloor<Q/4$, the binary digits of $X$ are the fields $[c_x;\,(j+1)\bmod4;\,\lfloor i/4\rfloor]$ with $c_x=\varrho(i)+\lfloor(j+1)/4\rfloor\bmod Q$, and those of $Y$ are $[j;\,i]$. So the cells of depths $(r,r)$ hold one point each, and the first two offset digits of $X$, one base-$4$ geometric digit, are $(j+1)\bmod4$, a function of the last two cell digits of $Y$. The occupied cells of depths $(r+2,r)$ are exactly the coset $\{(a,b):a\equiv b+1\pmod4\}$ of an index-$4$ lattice.
(ii) \emph{Lattice and jitter.} With $e=(Q,0)$, $v=(Q/4,Q)$ and $u=\varrho(i)$,
\[
(X,Y)\equiv u\,e+\Bigl(j+\frac{\varrho(u)}{Q}\Bigr)v+\Bigl(\frac Q4+\Bigl\lfloor\frac i4\Bigr\rfloor-\frac i4,\;0\Bigr)\pmod N:
\]
the points lie on the lines of the lattice $\Lambda=\Z e+\Z v$ in direction $v$, exactly $v$ apart along each line, and line $u$ is slid along itself by $\varrho(u)/Q$ of one step (the \emph{slide} of line $u$, a van der Corput value), up to less than one grid unit. The slides are the jitter. The tiling lattice is $Q\Lambda$, so the lines close up and there are no slips.
(iii) \emph{Separation.} The functional $X-Y/4$ separates distinct lines by at least $Q-\frac34$, so $\delta\sqrt N\ge(4/\sqrt{17})(1-1/Q)$ for every choice of the slides; the van der Corput slides give exactly $\delta^2N=241/256$ in the tiling metric for every even $m\ge8$ \cite{Rud26}. On these lines the slides that keep $t=0$ are exactly the nested scramblings of the van der Corput slides, and all of them are equally far from lattice slides (both computed for $m\le8$), so only their arrangement is free. Flipping the third binary digit of the slide on the lines $u\equiv2,3\pmod4$ is again a nested scrambling, so it gives a $(0,m,2)$-net, by (iv) and \cite{Owe95}. In the tiling metric it is worse at $m=8$ ($\delta^2N=229/256$ against $241/256$), and for every even $m$ with $10\le m\le20$ it has $\delta^2N=977/1024$ against $241/256=964/1024$, i.e.\ $\delta\sqrt N=0.9767\ldots$ against $0.9702\ldots$ (computed). The jitter is not reduced, only moved off the direction across the lines, as in Section~\ref{sec:drift}.
(iv) \emph{The net property.} The slides carry $t=0$: the net is a nested (Owen) scrambling of the Hammersley net, the frame of the two places $\infty$, $0$, and nested scramblings of a $(t,m,s)$-net are $(t,m,s)$-nets \cite{Owe95}; this gives a short proof of the net property proved in \cite{Rud26}.
The net is not digital (it does not contain the origin \cite{GK09,Rud26}), its scramblings are not linear, and its metric (the row-shifted torus) is adapted to its lattice of lines. The nets of this paper are digital: they use linear scrambles of frames with more places, keep the standard torus, and pay with a slip plane wherever a geometric digit wraps around modulo $q$; a nonlinear permutation can traverse a sub-cell (a box of side $1/q$ of a cell) without such a jump, and linear offsets cannot (Section~\ref{sec:whylat}).
\end{remark}


\section{Notation}\label{sec:conv}

\begin{definition}[Digital shift {\cite[Def.~2.2]{DGS2}}]\label{def:dshift}
\lean{Z3NetF.digPointShift_eq_dshiftPt, Z3NetF.digPointShift_zero, QUNetsM3.digitalPtS_eq_dshiftPt}\leanok
For a point set given by digit vectors $x_c\in(\Z/q\Z)^m$ and $\tau=(\tau_1,\dots,\tau_s)\in(\Z/q\Z)^{s\times m}$, the \emph{digitally shifted} point set $P\oplus\tau$ has the digit vectors $x_c+\tau_c$, added digit by digit modulo $q$ without carries. For prime $q$ this is \cite[Def.~2.2]{DGS2}. A \emph{digital net} is unshifted ($\tau=0$) unless stated otherwise.
\end{definition}

\begin{remark}[Formalized statements of this section without a result label]\label{bp:sec:prelim}
\lean{Z3NetF.graphForm_of_box}\leanok
An invertible box matrix makes a digital point set a graph form (Section~\ref{sec:prelim}, subsection "Graph forms").
\end{remark}


\chapter{Frames: \texorpdfstring{$t=0$}{t=0} for free}\label{sec:frames}

\begin{lemma}[Hermite interpolation]\label{lem:hermite}
\lean{Z3NetF.hermite_det, Z3NetF.hermite_ne_zero, Z3NetF.hermite_rows}\leanok
Let $x_1,\dots,x_n\in R$ and $\alpha_1,\dots,\alpha_n,\alpha_\infty\ge0$ with $\sum_P\alpha_P+\alpha_\infty=m$, and let $H$ be the matrix of the map that sends $V\in R[t]_{<m}$ to its $\alpha_P$-jets at the $x_P$ and its $\alpha_\infty$-jet at $\infty$. Then $\det H=\pm\prod_{P<Q}(x_Q-x_P)^{\alpha_P\alpha_Q}$. In particular, over a field, $H$ is invertible if the $x_P$ with $\alpha_P\ge1$ are pairwise distinct, and $H$ has two equal rows if two of them coincide.
\end{lemma}

\begin{proof}
\leanok
The value is the confluent Vandermonde determinant \cite[Sec.~6]{Nie87}. The consequences used below also follow by counting zeros: if over a field the points are distinct and all jets of $V$ vanish, the product $\prod_P(t-x_P)^{\alpha_P}$ of pairwise coprime factors, of degree $m-\alpha_\infty$, divides $V$, while $\deg V<m-\alpha_\infty$ because the jet at $\infty$ vanishes; so $V=0$. If $x_P=x_Q$ and $\alpha_P,\alpha_Q\ge1$, the rows $V(x_P)$ and $V(x_Q)$ coincide. Over a field this zero count is Skriganov's proof that the hyperderivative (Rosenbloom--Tsfasman Reed--Solomon) codes \cite{RT97} are MDS in that metric \cite[Prop.~5.1, Thms.~5.1, 5.2]{Skr01}.
\end{proof}


\section{The M\"obius frame}

\begin{proposition}[The frame is the Faure--Hammersley net]\label{prop:faure}
\lean{Z3Net.faure_Lam, Z3Net.faure_tau, Z3Net.faure_tau2, Z3Net.faure_DP, Z3Net.faure_reindex}\leanok
Let $P_{ji}=\binom{i}{j}$ be the upper Pascal matrix, $J$ the reversal, $D_{\pm}=\operatorname{diag}((-1)^i)$, and $(\Lambda_m)_{ij}=(-1)^{i-j}\binom{m-1-j}{i-j}$ (lower-unitriangular). Then $\Psi=\Lambda_m P$ and $\Psi^2=D_{\pm}JP$. Consequently, for each $m$, any net with generator matrices $(LT^2,L,LT)$, $L$ lower-triangular, is (after re-indexing $v\mapsto D_{\pm}Pv$) the net with generator matrices
\[
\bigl((-1)^{m-1}LJ,\; (LD_{\pm})\,P,\; (s_kL\Lambda_m D_{\pm})\,I\bigr),
\]
the Hammersley net of the two-dimensional generalized Faure sequence $(A_1,A_2P)$ of Tezuka \cite{Tez95} with $A_1=s_kL\Lambda_m D_{\pm}$, $A_2=LD_{\pm}$, and its Hammersley coordinate scrambled by $\pm L$. Since $A_1,A_2$ depend on $m$, the family is not the set of prefixes of one sequence.
\end{proposition}

\begin{proof}
\leanok
Expand $V=\sum_j w_j(t-1)^j$ with $w=Pv$. Then $\Psi V=\sum_j w_j t^j(1-t)^{m-1-j}$, giving $\Psi=\Lambda_m P$, and $\Psi^2 V=\sum_j w_j(-t)^{m-1-j}$, giving $\Psi^2=D_{\pm}JP$. Finally use $P^{-1}=D_{\pm}PD_{\pm}$ and $D_{\pm}JD_{\pm}=(-1)^{m-1}J$.
\end{proof}


\section{The frame of \texorpdfstring{$s$}{s} places}\label{sec:dim4}

\begin{lemma}[Frame]\label{lem:Ts}
\lean{Z3NetF.frame_minor, Z3NetF.frame_net_iff, Z3NetU.U_frameMinor}\leanok
Every box minor of the frame is, up to sign, $\prod_{\{c,d\}\subset F}(a_c-a_d)^{d_cd_d}$. In particular, for $m\ge2$ and $q$ coprime to the denominators of the places, the frame is a $(0,m,s)$-net over $\Z/q\Z$ if and only if the finite places are pairwise distinct modulo every prime factor of $q$.
\end{lemma}

\begin{proof}
\leanok
\uses{lem:nie, lem:hermite}
The box matrix for the depths $(d_c)$ is the Hermite matrix of Lemma~\ref{lem:hermite} for the places with multiplicities $d_c$, which gives the value, a unit modulo a prime $p$ if the finite places are distinct modulo $p$. Conversely, if $a_c\equiv a_d\pmod p$, a box with $d_c=d_d=1$ (it exists as $m\ge2$) contains the two equal rows $V(a_c)$, $V(a_d)$ modulo~$p$. The net statement then follows from Lemma~\ref{lem:nie}.
\end{proof}

\begin{lemma}[First offset rows of the frame]\label{lem:Rs}
\lean{Z3NetF.frameGraph_spec, Z3NetF.Rs_row0}\leanok
For every $k\ge1$ and $c\ne d$, row $0$ of $\Phi^F_{cd}$ has the generating function
\[
\frac{Q_c^k}{\Pi_d(d+x)\,(c-d-x)}\quad(c,d\in F),\qquad
\frac{Q_c^k\prod_{b\in F}(1-bx)^{-k}}{1-cx}\quad(c\in F,\ d=\infty),
\]
and $1/\Pi_d(d+x)$ for $c=\infty$, $d\in F$.
\end{lemma}

\begin{proof}
\leanok
The $(t-b)^k$, $b\in F$, are pairwise comaximal in $R[t]$, so partial fractions give $V/\Pi=H+\sum_{b\in F}N_b/(t-b)^k$ with $\deg H<k$ (as $\deg V<sk=\deg\Pi+k$) and $\deg N_b<k$; that is, $V=\Pi H+\sum_b\Pi_bN_b$.

\emph{Cells.} For $b\in F$, $V\equiv\Pi_bN_b\pmod{(t-b)^k}$ with $\Pi_b(b)=Q_b^k$ a unit, so the cell of $b$ and $N_b$ determine each other. As $\deg\sum_b\Pi_bN_b<m-k$, the cell of $\infty$ (the top $k$ coefficients of $V$) is that of $\Pi H$, which determines $H$ as $\Pi$ is monic.

\emph{Digit $k$ at $c\in F$.} $V=\Pi_cN_c+(t-c)^k\Pi_cG_c$ with $G_c=H+\sum_{b\ne c}N_b/(t-b)^k$, regular at $c$. So digit $k$ is a term depending on $N_c$ only (the own cell) plus $Q_c^kG_c(c)=Q_c^k\bigl(H(c)+\sum_{b\ne c}\mathrm{pp}_b(c)\bigr)$, where $\mathrm{pp}_b=N_b/(t-b)^k$ is the principal part at $b$.

\emph{Digit $k$ at $\infty$}, the coefficient of $t^{m-1-k}$, is a term from $\Pi H$ (the own cell) plus, from each $\Pi_bN_b$, the top coefficient of $N_b$, i.e. $\operatorname{Res}_b(V\,dt/\Pi)$.

\emph{Column $(d,j)$, $d\in F$.} All other cells vanish, so $H=0$, $N_b=0$ for $b\ne d$, and $\mathrm{pp}_d$ is the principal part at $d$ of $(t-d)^{j-k}/\Pi_d$. For a function $f$ whose only pole near $d$ has principal part $\mathrm{pp}$, $\mathrm{pp}(c)=\operatorname{Res}_{t=d}f(t)/(c-t)$; hence $\mathrm{pp}_d(c)=[u^{k-1-j}]\,1/(\Pi_d(d+u)(c-d-u))$ and $\operatorname{Res}_d(V/\Pi)=[u^{k-1-j}]\,1/\Pi_d(d+u)$. Summing over $j$ with weights $x^{k-1-j}$ gives the first and the third entry.

\emph{Column $(\infty,j)$.} The finite cells vanish, so $V=\Pi H\equiv t^{m-1-j}$ modulo degree $<m-k$, and $H$ is the polynomial part of $t^{k-1-j}\prod_b(1-b/t)^{-k}$, and $\sum_jx^{k-1-j}H_j(c)=\prod_b(1-bx)^{-k}/(1-cx)\bmod x^k$; multiply by $Q_c^k$.
\end{proof}

\begin{proposition}[Normal form]\label{prop:Ns}
\lean{Z3NetF.Ns_row0, Z3NetU.U_propN}\leanok
\uses{lem:box}
Apply the move of Lemma~\ref{lem:box}(c) with $L_c=\operatorname{diag}(Q_c^{-k},1,\dots,1)$, $U_d=\operatorname{Toep}(\Pi_d(d+x))$ for $d\in F$ and $U_\infty=\operatorname{Toep}(\prod_{b\in F}(1-bx)^k)$; their diagonal entries are units of $R$. Then row $0$ of every off-diagonal block becomes
\[
N_{cd}(x)=\frac{w_{cd}}{1-y_{cd}\,x},\qquad\text{i.e. the entry at carry depth $n$ is }w_{cd}\,y_{cd}^{\,n},
\]
independently of $k$, where $w_{cd}=y_{cd}=1/(a_c-a_d)$ for $c,d\in F$; $w=1$, $y=0$ for $c=\infty$; and $w=1$, $y=a_c$ for $d=\infty$.
\end{proposition}

\begin{proof}
\leanok
\uses{lem:Rs}
Substitute into Lemma~\ref{lem:Rs}: for instance $Q_c^{-k}\cdot Q_c^k/(\Pi_d(d+x)(c-d-x))\cdot\Pi_d(d+x)=1/(c-d-x)$.
\end{proof}


\section{Designs that work for every \texorpdfstring{$k$}{k}}

\begin{proposition}[Scrambles for every $k$]\label{prop:Gs}
\lean{Z3NetF.Gs_beta, Z3NetF.Gs_gauge, Z3NetF.Gs_row0, Z3NetF.PGen_indep, Z3NetU.U_propG, Z3NetU.U_gam_unique, Z3NetF.skel_Gs, Z3NetF.D4D2_data, Z3NetF.Gsplit_beta, Z3NetF.Gsplit_gauge, Z3NetF.Gsplit_row0}\leanok
\uses{prop:Ns}
Suppose the residual problem is solvable. For $k\ge1$ put $\lambda_c=\lambda'_cQ_c^{-k}$ and $U_d=U^0_dU'^{(k)}_d$, where $U^0_d$ is the Toeplitz column move of Proposition~\ref{prop:Ns}, and $U'^{(k)}_d$ agrees with $U'_d$ at the carry depths $<\min(s-1,k)$ and has, for $s-1\le n<k$, the column $e_n-\sum_{i\le s-2}\beta_{n,i}\cdot(\text{column }i\text{ of }U'_d)$, where $\beta_n\in R^{s-1}$ solves $\lambda'\odot v_n=\sum_i\beta_{n,i}F_i[\cdot][d]$ (rows $c\ne d$). Then row $0$ of $L_c\Phi^F_{cd}U_d$ equals the design at the carry depths $<s-1$ (truncated if $k<s-1$) and vanishes at the larger carry depths; all unit factors and all entries lie in $R$.
\end{proposition}

\begin{proof}
\leanok
\uses{prop:Ns}
For $n<s-1$ this is the residual problem (Proposition~\ref{prop:Ns}). The $y_{cd}$, $c\ne d$, are pairwise distinct with unit differences ($1/(a_c-a_d)-1/(a_{c'}-a_d)=(a_{c'}-a_c)/((a_c-a_d)(a_{c'}-a_d))$, and $0$ resp. the $a_c$ in the other cases), and the $w_{cd}$ are units, so the matrix $(v_0,\dots,v_{s-2})$ is invertible over $R$, and so is $(F_i[\cdot][d])_{i\le s-2}=\Lambda'(v_0,\dots,v_{s-2})U'_d$. Hence $\beta_n$ exists over $R$, and the new column $n$ has row-$0$ entries $\lambda'_cv_n[c]-\sum_i\beta_{n,i}F_i[c][d]=0$ and diagonal entry $1$.
\end{proof}

\begin{definition}\label{def:Ns}
\lean{Z3NetF.PGen_indep, Z3NetF.tsk_unit, Z3NetF.tsk_net, Z3NetF.away_hom, Z3NetU.U_defN, Z3NetF.Gsplit_tsk_unit, Z3NetF.Gsplit_tsk_net}\leanok
\uses{prop:Gs}
Let $(F_n)$ be a design whose residual problem is solvable over $R'=R[1/D']$, where $D'$ collects the primes in the units and denominators of $\lambda'$ and $U'$, and let $\gcd(q,DD')=1$. The digital net $\mathcal P^F_{k,q}$ over $\Z/q\Z$ is the graph form with depths $(k,\dots,k)$ and blocks $\Phi'_{cd}=L_c\Phi^F_{cd}U_d$ ($c\ne d$; Proposition~\ref{prop:Gs}), $\Phi'_{cc}$ with row $0$ equal to $(F_{s-2}[c][c],\dots,F_0[c][c])$ on the levels $(k-s+1,\dots,k-1)$ (truncated if $k<s-1$) and all other entries $0$, all entries reduced modulo $q$.
\end{definition}

\begin{corollary}\label{cor:tsk}
\lean{Z3NetF.tsk_unit, Z3NetF.tsk_net, Z3NetF.away_hom, Z3NetF.tsk_net_rowScale, Z3NetF.Gsplit_tsk_unit, Z3NetF.Gsplit_tsk_net}\leanok
Every box minor of $\mathcal P^F_{k,q}$ over $R'$ is a unit of $R'$; so $\mathcal P^F_{k,q}$ is a digital $(0,sk,s)$-net over $\Z/q\Z$ for every $k\ge1$, with one point in every cell.
\end{corollary}

\begin{proof}
\leanok
\uses{lem:nie, lem:box, lem:Ts, prop:Gs}
The box minors of $\Phi^F$ are those of the frame divided by its $(k,\dots,k)$ minor, units of $R$ (Lemma~\ref{lem:Ts}); Lemma~\ref{lem:box} and Proposition~\ref{prop:Gs} multiply them by units of $R'$; then Lemma~\ref{lem:nie}.
\end{proof}

\begin{remark}[Two-dimensional projections]\label{rem:projlines}
\lean{Z3NetF.projlines_frame, Z3NetF.projlines, Z3NetF.projlines_zmod, Z3NetF.projlines_box, Z3NetF.rowScale_box, Z3NetF.rowScale_net, Z3NetF.tsk_net_rowScale, Z3NetF.rowScale_rowRel, Z3NetF.Ss_sep_rowScale, Z3NetF.D_sep_rowScale, Z3Net.projlines_P, Z3Net.P_net_rowScale, Z3Net.P_sep_rowScale, Z3NetA.A_sep_rowScale}\leanok
\uses{thm:A, lem:Ss, def:Ns, cor:tsk, prop:D, def:net, thm:sep}
Let $m=sk$ be odd with $(s,k)\ne(3,1)$, and $r^*=(m-1)/2$. For any two coordinates of $\mathcal P^F_{k,q}$ (Definition~\ref{def:Ns}), and of $\mathcal P_{k,q}$ (Definition~\ref{def:net}), the digits with index $r^*$ agree up to a sign and a linear combination of more significant digits; for $\mathcal P^F_{k,q}$ the sign is $+1$ if one of the two places is $\infty$, and $(-1)^{r^*}$ otherwise. For the places $\infty$ and $0$ this is the reversal symmetry of the two-dimensional Hammersley set, which reads digit $r^*$ as $[t^{m-1-r^*}]V=[t^{r^*}]V$; a M\"obius change of variable transports it to every pair of places, and the lower-triangular scrambles of Corollary~\ref{cor:tsk} keep it, since $r^*\ge k+1$ and their diagonal entries in the rows below the first offset row are~$1$. This digit relation is proved. Its visible effect is computed, not proved: in the computed cases every two-dimensional projection lies on parallel lines, at a dual frequency of length about $\sqrt{2/q}\,\sqrt N$ (Figure~\texttt{fig:projlines}), which the $t$-parameter does not see. The phenomenon is akin to the known diagonal and rank-$1$-lattice clustering of generalized Faure sequences whose generator matrices coincide up to a scalar \cite{LF09,L15}; there the coincidence is in the top digits, here in the middle one. The remedy is of the classical kind: a $t$-preserving left multiplication by a nonsingular lower-triangular, here diagonal, matrix, as in Tezuka's generalized Faure sequences \cite{Tez93,Tez95} and Matou\v{s}ek's linear scrambling \cite{Mat98}, with per-coordinate multipliers chosen on two-dimensional projections as in \cite{FL09,LF09}. What is specific here is that only row $r^*$ is rescaled, by units $\rho_c$. This keeps $t=0$, as it multiplies every box minor by a unit, and, since Lemma~\ref{lem:Ss} and Theorem~\ref{thm:sep} use that row only through (H2), it keeps the proved separation bounds (of Lemma~\ref{lem:Ss}, Proposition~\ref{prop:D}, Theorems~\ref{thm:A}(2) and~\ref{thm:sep}); both facts are proved. For a two-dimensional projection with points $x$ let $S(\mathbf h)=N^{-1}\sum_xe^{2\pi i\,\mathbf h\cdot x}$, $\mathbf h\in\Z^2$, and call the least $|\mathbf h|/\sqrt N$ over the frequencies with $|S(\mathbf h)|\ge\frac12$ its \emph{spectral value} (about $1$ for a good two-dimensional lattice; lines at spacing $N^{-1/2}$ divided by this value appear when it is small). With $\rho$ found per $q$ by an exhaustive search that keeps an estimate of $|S(\mathbf h)|$ small for $|\mathbf h|<0.8\sqrt N$ in every pair of coordinates, the least spectral value over the pairs of coordinates at $s=5$, $q=13$--$23$ rises from $0.22$--$0.30$ to about $0.80$ (computed).
\end{remark}


\section{Reachable designs}

\begin{proposition}[Reachable designs]\label{prop:F}
\lean{Z3NetF.F_reach_iff, Z3NetF.F_poly_iff, Z3NetF.F_finite, Z3NetU.U_design, Z3NetF.skel_reach_iff, Z3NetF.Gsplit_reach_iff, Z3NetF.Gsplit_poly_iff}\leanok
\uses{lem:box}
The off-diagonal part of a design is reachable from the frame by the gauge moves of Lemma~\ref{lem:box}(c) with row scalars $\lambda'$ (that is, the residual problem is solvable with these $\lambda'$) if and only if, for every column $d$ and every $j\le s-2$,
\[
M^{(j)}[c][d]=\lambda'_c\,w_{cd}\,g_{d,j}(y_{cd})\qquad(c\ne d)
\]
for a polynomial $g_{d,j}\in R[X]$ of degree $j$ whose leading coefficient is a unit. For finite $c,d$ this says $M^{(j)}[c][d]=\lambda'_c\tilde g_{d,j}(z_{cd})$ with $z_{cd}=1/(a_c-a_d)$ and $\tilde g_{d,j}(z)=z\,g_{d,j}(z)$ of degree $j+1$ without constant term.
\end{proposition}

\begin{proof}
\leanok
\uses{lem:box}
A lower-triangular column move with unit diagonal is a Toeplitz move times an upper-triangular matrix in the carry depths, so the residual problem captures all of Lemma~\ref{lem:box}(c) on row $0$. Column $n$ of $\lambda'^{-1}\odot F_{\cdot d}$ must be $\sum_{n'\le n}v_{n'}(U'_d)_{n'n}$, i.e. $(w_{cd}f_n(y_{cd}))_c$ with $\deg f_n\le n$ and the coefficient of $X^n$ a unit; as the $y_{cd}$ are distinct with unit differences, interpolation is unique over $R$, and this is exactly the condition. Partial sums of such polynomials have the same form, and conversely.
\end{proof}


\section{Cauchy identities}

\begin{lemma}[Cauchy identities]\label{lem:CH}
\lean{Z3NetF.CH_C, Z3NetF.CH_C_neg, Z3NetF.CH_C_classes, Z3NetF.CH_P, Z3NetF.CH_H, Z3NetF.CH_X}\leanok
Let $R$ be a commutative ring, $\alpha,\beta\in R^s$, $D_{cd}=\alpha_c\beta_d-\alpha_d\beta_c$, and $W_{cd}=\lambda_c\mu_d/D_{cd}$ for $c\ne d$, where all $\lambda_c$, $\mu_d$ and $D_{cd}$ are units of $R$; let $\nu_c=\lambda_c\mu_c$ and $\pi_{cd}=W_{cd}W_{dc}$.
\begin{enumerate}
\item[(C)] For distinct $c,d,e$: $W_{cd}W_{de}W_{ec}=-W_{dc}W_{ed}W_{ce}$. Hence, if $R$ is an ordered field, $\pi_{cd}\pi_{de}\pi_{ec}<0$, and ``$c\sim d$ iff $\pi_{cd}<0$'' is an equivalence relation with at most two classes.
\item[(P)] For distinct $a,b,c,d$ let $\chi=D_{ac}D_{bd}/(D_{ab}D_{cd})$ and $\psi=D_{ad}D_{bc}/(D_{ab}D_{cd})$. Then $\chi-\psi=1$, $\chi^2=\pi_{ab}\pi_{cd}/(\pi_{ac}\pi_{bd})$ and $\psi^2=\pi_{ab}\pi_{cd}/(\pi_{ad}\pi_{bc})$; consequently $(\chi^2-\psi^2-1)^2=4\psi^2$.
\item[(H)] If $R\subset\R$: for distinct $i,j,k,l$, the largest of $|\pi_{ij}\pi_{kl}|^{-1/2}$, $|\pi_{ik}\pi_{jl}|^{-1/2}$, $|\pi_{il}\pi_{jk}|^{-1/2}$ is the sum of the other two.
\item[(X)] For distinct $i,j,l,d$:
\[
\frac{W_{ji}W_{ld}}{W_{jd}W_{li}}+\frac{W_{ij}W_{ld}}{W_{id}W_{lj}}=1 .
\]
\end{enumerate}
\end{lemma}

\begin{proof}
\leanok
Everything follows from $D_{dc}=-D_{cd}$ and the Pl\"ucker relation $D_{ab}D_{cd}-D_{ac}D_{bd}+D_{ad}D_{bc}=0$. (C) The $\lambda$, $\mu$ cancel, and the $D$ cancel up to $(-1)^3$; then $\pi_{cd}\pi_{de}\pi_{ec}=-(W_{dc}W_{ed}W_{ce})^2$. If $c\sim d$ and $d\sim e$, two of the three factors are negative, so the third is too; if $c\not\sim d$ and $d\not\sim e$, then $c\sim e$. (P) Dividing the Pl\"ucker relation by $D_{ab}D_{cd}$ gives $\chi-\psi=1$. Since $\pi_{xy}=-\nu_x\nu_y/D_{xy}^2$, we have $\pi_{ab}\pi_{cd}=\nu_a\nu_b\nu_c\nu_d/(D_{ab}D_{cd})^2$ and similarly for the other pairings, which gives $\chi^2$ and $\psi^2$; and $\chi=1+\psi$ gives $2\psi=\chi^2-\psi^2-1$. (H) By the same formula $|D_{ij}D_{kl}|$, $|D_{ik}D_{jl}|$, $|D_{il}D_{jk}|$ are proportional to the three numbers, and by the Pl\"ucker relation the products sum to zero with signs. (X) The $\lambda$ and $\mu$ cancel, the two terms are $D_{jd}D_{il}/(D_{ij}D_{ld})$ and $-D_{id}D_{jl}/(D_{ij}D_{ld})$, and $D_{il}D_{jd}-D_{id}D_{jl}=D_{ij}D_{ld}$ is the Pl\"ucker relation for $(i,j,l,d)$.
\end{proof}

\begin{remark}[Formalized statements of this section without a result label]\label{bp:sec:frames}
\lean{Z3NetF.wN_cauchy}\leanok
The frame layer $\lambda_c\mu_dw_{cd}$ is a Cauchy layer in homogeneous coordinates (Section~\ref{sec:frames}, subsection "Cauchy identities").
\end{remark}


\chapter{Cubic cells, \texorpdfstring{$m\equiv0\pmod s$}{m = 0 mod s}: designs from short lattice vectors}\label{sec:dims}\label{sec:small}


\section{Diagonal dominance}

\begin{proposition}[Diagonal dominance]\label{prop:D}
\lean{Z3NetF.D_sep_rowScale, Z3NetF.D_LB, Z3NetF.D_sep}\leanok
Let $F$ be a reachable design whose diagonal vanishes at every carry depth $n\ge1$, with $F_0[c][c]=-\alpha_c$ and
\[
\alpha_c\ \ge\ \sum_{d\ne c}\ \max_{0\le j\le s-2}\bigl|M^{(j)}[c][d]\bigr|,\qquad a=\max_c\alpha_c.
\]
Then $\mathrm{LB}\ge1-a/q$ for every $q\ge5a+4$ (and every $k\ge1$).
\end{proposition}

\begin{proof}
\leanok
\uses{lem:Ss}
Since the diagonal vanishes at every carry depth $n\ge1$, the own term of every moving coordinate $i$ is $-\alpha_i$ for every carry depth $j_i$, and dominance gives $e_i\sigma_i\le-\alpha_i+\alpha_i=0$ and $|\sigma_i|\le2a<q$: every moving coordinate is backward, so $b_i\ge q-|\sigma_i|-\mathrm{jit}_i$. An own-axis step along $c$ has $b_c=q-\alpha_c\ge q-a$ (no jitter). A step with at least two moving axes has $\sum b_i^2\ge2(q-2a-1)^2$, and $\sqrt2(q-2a-1)\ge q-a$ holds for $q\ge(2\sqrt2-1)a/(\sqrt2-1)+\sqrt2/(\sqrt2-1)$, in particular for $q\ge5a+4$. A truncated design is dominant with the same $a$, and at $k=1$ the configurations are a subset of those for $k\ge2$ (after Lemma~\ref{lem:Ss}).
\end{proof}


\section{Admissible design columns form a lattice}\label{sec:lattice}

\begin{lemma}[The lattice of admissible design columns]\label{lem:sat}
\lean{Z3NetA.sat_poly, Z3NetA.sat_gamma, Z3NetA.sat_adm, Z3NetA.sat_basis, Z3NetA.sat_coset, Z3NetA.gamS_unique, Z3NetA.RS, Z3NetA.InRS, Z3NetA.SUnit, Z3NetA.wQ, Z3NetA.yQ, Z3NetA.colPoly, Z3NetA.InLam, Z3NetA.kappa, Z3NetA.gamS, Z3NetA.AdmCol}\leanok
\begin{enumerate}
\item[(i)] For $E\in\Lambda_j$ the polynomial of $E$ has coefficients in $R_S$.
\item[(ii)] $\operatorname{lc}_j(\Lambda_j)=\gamma_j\Z$ for a positive $S$-unit $\gamma_j$.
\item[(iii)] An integer column is admissible at carry depth $j$ if and only if it lies in $\Lambda_j$ and $\operatorname{lc}_j(E)$ is an $S$-unit; in particular every $E\in\Lambda_j$ with $\operatorname{lc}_j(E)=\gamma_j$ is admissible.
\item[(iv)] If $E_i\in\Lambda_i$ and $\operatorname{lc}_i(E_i)=\gamma_i$ for $i=0,\dots,j$, then $E_0,\dots,E_j$ is a basis of $\Lambda_j$, and $\{E\in\Lambda_j:\operatorname{lc}_j(E)=\gamma_j\}=E_j+\Lambda_{j-1}$.
\end{enumerate}
\end{lemma}

\begin{proof}
\leanok
(i) Lagrange interpolation at any $j+1$ rows $J$ recovers $g=\sum_{b\in J}(E_b/w_b)\prod_{b'\in J\setminus b}(X-y_{b'})/(y_b-y_{b'})$, and all denominators are $S$-units. (ii) $\operatorname{lc}_j(\Lambda_j)$ is a finitely generated subgroup of $\mathbb Q$, nonzero because $\Lambda_j$ spans $V_j$, so it is $\gamma_j\Z$ with $\gamma_j>0$; by (i), $\gamma_j\in R_S$. The vector $(w_cy_c^j)_c$ of $g=X^j$ has $S$-unit or zero entries, so some product $P$ of primes in $S$ makes $P(w_cy_c^j)_c$ integral, with $\operatorname{lc}_j=P$; hence $P\in\gamma_j\Z$, and $\gamma_j$ has no prime factor outside $S$ in its numerator either. (iii) is (i) together with the definition. (iv) If $v\in\Lambda_j$ has $\operatorname{lc}_j(v)=t\gamma_j$, then $v-tE_j\in\Lambda_j\cap\ker\operatorname{lc}_j=\Lambda_{j-1}$; induction on $j$.
\end{proof}

\begin{lemma}[Adelic Chebyshev identity]\label{lem:adelic}
\lean{Z3NetA.adelic_p, Z3NetA.adelic_gamma, Z3NetA.adelic_inf, Z3NetA.adelic_proj, Z3NetA.adelic_norm, Z3NetA.hS_le_monic, Z3NetA.colE, Z3NetA.polyVecR, Z3NetA.polyVecQ, Z3NetA.Vlt, Z3NetA.lineA, Z3NetA.piA, Z3NetA.AinfSet, Z3NetA.Ainf, Z3NetA.ApSet, Z3NetA.Ap, Z3NetA.primesS, Z3NetA.hS}\leanok
Let $\pi_j$ be the orthogonal projection of $\R^{s-1}$ onto the line $V_j\cap V_{j-1}^\perp$. Then $\pi_j(\Lambda_j)$ is generated by a vector of length
\[
\eta_j=A_{j,\R}\prod_{p\in S}A_{j,p},\qquad\text{and}\qquad \gamma_j=\prod_{p\in S}A_{j,p}.
\]
Every $E\in\Lambda_j\setminus V_{j-1}$ has $\|E\|_2\ge \eta_j$.
\end{lemma}

\begin{proof}
\leanok
\emph{The $p$-adic places, $p\in S$.} If $E\in V_j$ has $\operatorname{lc}_j(E)=1$ and $\max_c|E_c|_p=p^e$, then $p^eE$ is $p$-integral, and $D'p^eE\in\Lambda_j$ for an integer $D'$ prime to $p$; so $D'p^e\in\gamma_j\Z$ and $v_p(\gamma_j)\le e$. Conversely, for $v\in\Lambda_j$ with $\operatorname{lc}_j(v)=\gamma_j$ the vector $v/\gamma_j$ has $\operatorname{lc}_j=1$ and $\max_c|v_c/\gamma_j|_p\le p^{v_p(\gamma_j)}$. So $A_{j,p}=p^{v_p(\gamma_j)}$, and as $\gamma_j>0$ is an $S$-unit, $\gamma_j=\prod_{p\in S}A_{j,p}$. \emph{The archimedean place.} On $V_j$, $\operatorname{lc}_j(E)=\langle\theta,E\rangle$ for a unique $\theta\in V_j$, and $\theta\perp V_{j-1}$. Hence $\|\pi_j(E)\|=|\operatorname{lc}_j(E)|/\|\theta\|$, and the least norm with $\operatorname{lc}_j=1$ is $A_{j,\R}=1/\|\theta\|$, attained at $\theta/\|\theta\|^2$. So $\pi_j(\Lambda_j)=\gamma_j\Z\,\theta/\|\theta\|^2$, generated by a vector of length $\gamma_jA_{j,\R}=\eta_j$, and $\|E\|\ge\|\pi_j(E)\|\ge \eta_j$ if $\operatorname{lc}_j(E)\ne0$.
\end{proof}

\begin{lemma}[Babai rounding along the flag]\label{lem:babai}
\lean{Z3NetA.Ustar_babai, Z3NetA.Ustar_growth, Z3NetA.aStar_le, Z3NetA.aStar_three, Z3NetA.babai, Z3NetA.babai_growth, Z3NetA.babai_flag}\leanok
Let $E_i\in\Lambda_i$ with $\operatorname{lc}_i(E_i)=\gamma_i$ for $i<j$, and let $v\in\Lambda_j$ with $\operatorname{lc}_j(v)=\gamma_j$. Write the orthogonal projection of $v$ onto $V_{j-1}$ as $\sum_{i<j}c_iE_i$ and put $E_j=v-\sum_{i<j}\lfloor c_i+\frac12\rfloor E_i$. Then $E_j\in\Lambda_j$, $\operatorname{lc}_j(E_j)=\gamma_j$, and
\[
\|E_j\|_2\le \eta_j+\tfrac12\textstyle\sum_{i<j}\|E_i\|_2 .
\]
Consequently, if $E_0,\dots,E_j$ are all obtained in this way and $\eta_i\le H$ for every $i\le j$, then $\|E_i\|_2\le(3/2)^iH$ for every $i\le j$.
\end{lemma}

\begin{proof}
\leanok
\uses{lem:adelic}
$E_j-v\in\Lambda_{j-1}=\ker\operatorname{lc}_j\cap\Lambda_j$, and $E_j=\pi_j(v)+\sum_{i<j}\bigl(c_i-\lfloor c_i+\frac12\rfloor\bigr)E_i$ with $\|\pi_j(v)\|=\eta_j$ (Lemma~\ref{lem:adelic}) and every coefficient in $[-\frac12,\frac12]$. For the consequence let $T_i=\sum_{i'\le i}\|E_{i'}\|$. Then $T_0=\eta_0\le H$ and $T_i\le\frac32T_{i-1}+H$, so $T_i\le2H((3/2)^{i+1}-1)$ by induction, and $\|E_i\|\le H+\frac12T_{i-1}\le(3/2)^iH$.
\end{proof}


\section{The design \texorpdfstring{$U^*_s$}{U*s} and the proof of Theorem~A}\label{sec:Ustar}

\begin{definition}[The design $U^*_s$]\label{def:Ustar}
\lean{Z3NetA.Ustar_cols, Z3NetA.Ustar_cols_unique, Z3NetA.Ustar_round, Z3NetA.Ustar_Msum, Z3NetA.Ustar_dominant, Z3NetA.Ustar_span, Z3NetA.Ustar_proj, Z3NetA.Ustar3_display, Z3NetA.ProjCoef, Z3NetA.IsUstarCols, Z3NetA.Ecol, Z3NetA.alphaS, Z3NetA.aStar, Z3NetA.UstarF}\leanok
Let $s\ge3$, $\nu=s-2$, and take the places $\infty,0,1,\dots,\nu$ and $\lambda'\equiv1$. For each column $d$ let $E_{d,0}$ be the generator of $\Lambda_0$ with $\operatorname{lc}_0>0$, and for $1\le j\le \nu$ let $E_{d,j}$ be the unique $E\in\Lambda_j$ with $\operatorname{lc}_j(E)=+\gamma_j$ whose orthogonal projection onto $V_{j-1}$ has all its coefficients in the basis $E_{d,0},\dots,E_{d,j-1}$ in $[-\frac12,\frac12)$. The design $U^*_s=(F_0,\dots,F_\nu)$ has the partial sums $M^{(j)}[c][d]=E_{d,j}[c]$ ($c\ne d$, $j\le \nu$), the diagonal $F_0[c][c]=-\alpha_c$ at carry depth $0$ with $\alpha_c=\sum_{d\ne c}\max_j|E_{d,j}[c]|$, and zero diagonal at the carry depths $n\ge1$. Let $a^*_s=\max_c\alpha_c$.
\end{definition}

\begin{remark}[The monomial design $U_s$]\label{rem:Us}
\lean{Z3NetU.U_net, Z3NetU.U_boxMinor, Z3NetU.U_covering, Z3NetU.U_sep, Z3NetU.U_sep_real, Z3NetU.U_aBound, Z3NetU.U_a3, Z3NetU.U_a4, Z3NetU.U_a5, Z3NetU.U_aLower, Z3NetU.U_growth, Z3NetU.U_sharp, Z3NetU.U_sharp_real, Z3NetU.U_logA_psi, Z3NetU.U_growth_pnt, Z3NetU.U_bound_growth_pnt}\leanok
\uses{prop:D}
The every-dimension theorem also holds with the monomial design $U_s$: put $L_s=\operatorname{lcm}(1,\dots,s-2)^{s-1}$, $\lambda'_c=L_s$ for finite $c$, $\lambda'_\infty=1$, and $F_n[c][d]=L_s/(c-d)^{n+1}$ for finite $c,d$, $F_n[\infty][d]=[n=0]$, $F_n[c][\infty]=L_s\binom cn$, with a diagonal that is dominant in the sense of Proposition~\ref{prop:D}. Its increments are monomials in the nodes, and the proof of Theorem~A applies verbatim with its constant $a_s=\max_c\alpha_c$ in place of $a^*_s$ (for $k\ge2$ the net property, the separation bound for $q\ge5a_s+4$ and the sharpness are formally verified for this design). But $\log a_s=(1+o(1))s^2$: for finite $c$, $\alpha_c\ge|F_0[c][\infty]|=L_s$, and $\log L_s=(s-1)\psi(s-2)$ with Chebyshev's function $\psi$, while $a_s\le L_s(s^2+2^{s-2})$; so $a_s=e^{\Theta(s^2)}$ already by Chebyshev's bounds $\psi(x)\asymp x$, and the constant $1$ uses the prime number theorem \cite{Had96,dlVP96}. In the language above the reason is $p$-adic: a monomial has all its roots at $0$, so at a collision of nodes it is as large as possible, $|y^j|_p=p^{j\,v_p(c-d)}$, at every carry depth and in every column, and the single row scalar $\lambda'$ must clear all of it. Theorem~A puts the roots near the colliding nodes, separately for each prime, and so replaces $e^{\Theta(s^2)}$ by the upper bound $e^{s\log s+O(s)}$ (Table~\texttt{tab:astar}: $a_5=17407$ against $a^*_5=16$).

\textbf{Lean: assuming classical theorems.} The growth statement $\log a_s=(1+o(1))s^2$ uses the prime number theorem as an explicit hypothesis (\texttt{PNTpsi}); everything else is unconditional.
\end{remark}


\chapter{Lattice skeletons: \texorpdfstring{$m=sk+r$}{m=sk+r}}\label{sec:skel}


\section{Fine cells and the separation reduction}\label{sec:fine}

\begin{lemma}[Separation reduction]\label{lem:skelred}\label{lem:fine}
\lean{Z3NetF.skelred_coord, Z3NetF.skelred_sep, Z3NetF.skelred_sep_real, Z3NetF.skelred_carry}\leanok
Let $v\ne w$ be two points, and let $a_c$ be the centred residue ($-p^{k+1}/2<a_c\le p^{k+1}/2$) of $P_c(w)-P_c(v)$ modulo $p^{k+1}$; we call $a$ the \emph{fine difference} of $v$ and $w$.
\begin{enumerate}
\item[(i)] The toroidal distance of $X_c(w)$ and $X_c(v)$ modulo $N$ is at least $(|a_c|-1)\varpi$. Hence, if $|a|_\infty\ge M$ for every pair of distinct points, the toroidal minimum distance satisfies $\delta N^{1/s}\ge(M-1)\,p^{-\ell/s}$.
\item[(ii)] If $|a_c|<p$ and $1\le j\le k$, then $\lfloor P_c(w)/p^j\rfloor-\lfloor P_c(v)/p^j\rfloor\equiv\xi_{c,j}\pmod{p^{k+1-j}}$ with $\xi_{c,j}\in\{0,\sgn a_c\}$, and $\xi_{c,j}\ne0$ implies $\xi_{c,j-1}\ne0$ for $j\ge2$.
\end{enumerate}
\end{lemma}

\begin{proof}
\leanok
(i) $N=p^{k+1}\varpi$, so $X_c(w)-X_c(v)\equiv x:=a_c\varpi+(\theta'-\theta)\pmod N$ with $|\theta'-\theta|<\varpi$. If $|x|\le N/2$, the distance is $|x|\ge|a_c|\varpi-(\varpi-1)$. Otherwise it is $N-|x|\ge N-(|a_c|+1)\varpi+1>(p^{k+1}-|a_c|-1)\varpi\ge(|a_c|-1)\varpi$, as $|a_c|\le p^{k+1}/2$. So $\delta N\ge(M-1)\varpi$, and $\varpi N^{1/s}/N=p^{m/s-k-1}=p^{-\ell/s}$. (ii) Write $P_c(w)=P_c(v)+a_c+\varepsilon p^{k+1}$ with $\varepsilon\in\{-1,0,1\}$. The difference of the floors is $\lfloor(P_c(v)+a_c)/p^j\rfloor-\lfloor P_c(v)/p^j\rfloor+\varepsilon p^{k+1-j}$. As $|a_c|<p\le p^j$, the first part lies in $\{0,\sgn a_c\}$, and it is nonzero only if a multiple of $p^j$, hence of $p^{j-1}$, lies between $P_c(v)$ and $P_c(v)+a_c$. Neither part uses that $p$ is odd or prime; $p\ge2$ suffices.
\end{proof}


\section{Necessary conditions, forced faults and completion}

\begin{lemma}[Necessary conditions on the skeleton digits]\label{lem:skelnec}
\lean{Z3NetF.skelnec_N1, Z3NetF.skelnec_N2, Z3NetF.skelnec_N3_full, Z3NetF.skelnec_N3_skel}\leanok
Let a graph form with the depths above be a $(0,m,s)$-net over $\F_p$. Let $A\in\F_p^{\ell\times r}$ hold the coefficients of the skeleton digits on the last levels $x_{d,k}$, and, for $k\ge1$ and $i\ne j$, let $g_{ij}$ be the coefficient of $S_{z_i}$ on $z_{j,k-1}$.
\begin{enumerate}
\item[(N1)] Every square submatrix of $A$ is invertible.
\item[(N2)] \emph{(Forced faults.)} If $\ell\ge2$ and $k\ge1$, then $g_{ij}\ne0$ for all $i\ne j$.
\item[(N3)] For every full coordinate $x_d$ and every set $I$ of $j$ skeleton coordinates, the coefficients of the $S_{z_i}$, $i\in I$, on the levels $x_{d,k},\dots,x_{d,k-j+1}$ form an invertible matrix; for a skeleton coordinate $z_e\notin I$ the same holds on the levels $z_{e,k-1},\dots,z_{e,k-j}$ (whenever these levels exist).
\end{enumerate}
\end{lemma}

\begin{proof}
\leanok
\uses{lem:box}
All three follow from Lemma~\ref{lem:box} for suitable boxes. For (N1) take depth $k+1$ on the skeleton coordinates of a set $I$, depth $k$ on the full coordinates of a set $J$ with $|J|=|I|$, and $k_c$ elsewhere: the over coordinates are those of $I$, by one digit (their skeleton digits), the under coordinates those of $J$, by one level (the level $k$), and the minor is $\det A[I,J]$. (N2) is the box with depth $k+1$ on $z_i$ and $k-1$ on $z_j$; (N3) the box with depth $k+1$ on $I$ and $k_e-j$ on the coordinate $e$.
\end{proof}

\begin{proposition}[Completion]\label{prop:skelcomp}
\lean{Z3NetF.skelcomp_normal, Z3NetF.skel_Gs, Z3NetF.skel_reach_iff, Z3NetF.skelcomp_ii, Z3NetF.skelcomp_iii, Z3NetF.frameGraphK_spec, Z3NetF.skel_graphGen, Z3NetF.skelreach_iff, Z3NetS1.S1gen_net}\leanok
\uses{lem:box, prop:Ns, prop:Gs, prop:F}
Let the frame and $\Phi^F$ be as above.
\begin{enumerate}
\item[(i)] \emph{(Normal form.)} After the moves of Proposition~\ref{prop:Ns}, applied to the restricted frame as described in the proof below, row $0$ of every block $(z_i,d)$, $d\ne z_i$, of $\Phi^F$ has the generating function $w_{id}/(1-y_{id}x)$ of Proposition~\ref{prop:Ns}, where $x^n$ marks carry depth $n$, the level $k_d-1-n$; this holds for every $k$. For a column block $d$ let $\ell_d$ be the number of skeleton coordinates other than $d$ ($\ell_d=\ell$ if $d$ is full, $\ell_d=\ell-1$ if $d$ is a skeleton coordinate). Then Propositions~\ref{prop:F} and~\ref{prop:Gs} hold for the skeleton rows, with the $\ell_d$ skeleton rows $z_i\ne d$ of block $d$ in place of the rows $c\ne d$, and $\ell_d-1$, $\ell_d$ in place of $s-2$, $s-1$ (the carry depths $n\le\ell_d-1$).
\item[(ii)] \emph{($\ell=1$.)} For every $k\ge0$ and all units $\kappa_1,\dots,\kappa_r$, the moves of Lemma~\ref{lem:box} turn $\Phi^F$ into a graph form whose skeleton digit is $S_z=\sum_d\kappa_dx_{d,k}$ (all other entries of row $0$ of $z$, own block included, are $0$). It is a $(0,m,s)$-net over $\F_p$.
\item[(iii)] \emph{($\ell=2$.)} Let $k\ge1$ and prescribe skeleton digits
\[
S_{z_i}=\sum_dA_{id}\,x_{d,k}+\sum_d\beta_{id}\,x_{d,k-1}+g_{ij}\,z_{j,k-1}+o_i\,z_{i,k-1}\qquad(\{i,j\}=\{1,2\})
\]
satisfying (N1)--(N3), with arbitrary own coefficients $o_i$. For a suitable choice of the places they are reachable for every $k\ge1$, and they complete to a $(0,m,s)$-net over $\F_p$.
\end{enumerate}
\end{proposition}

\begin{proof}
\leanok
\uses{lem:nie, lem:box, prop:Ns, prop:Gs, prop:F, def:B}
(i) \emph{Restriction.} Let $m'=s(k+1)$ and $V\in\F_p[t]_{<m'}$. The first digit of the reading at a finite skeleton place $a_z$ vanishes iff $(t-a_z)\mid V$, and at $\infty$ iff $\deg V<m'-1$. So the points of the $m'$-frame whose skeleton coordinates have first digit $0$ are given by $V=\prod_{z\text{ finite}}(t-a_z)\,V'$ with $\deg V'<m$, and their readings are those of $V'$ times a unit power series, shifted by one digit at a skeleton place. Hence this point set, with the skeleton coordinates rescaled by $p$, is the $m$-frame up to per-coordinate lower-triangular scrambles, i.e. up to the moves of Lemma~\ref{lem:box}; each new elementary box is an old one with one more leading digit on each skeleton coordinate. In the graph form of the $m'$-frame with depths $(k+1,\dots,k+1)$ the restriction is the vanishing of level $0$ of the skeleton cells, and the old first offset row of $z$ becomes its skeleton digit. The move of Proposition~\ref{prop:Ns} is lower-triangular in the levels and only scales level $0$, so it commutes with deleting the level-$0$ columns of the skeleton blocks; this gives the normal form, with carry depth $n$ at level $k_d-1-n$. The proofs of Propositions~\ref{prop:F} and~\ref{prop:Gs} use only this normal form and the distinctness, with unit differences, of the $y_{cd}$ on the rows involved, so they apply to the $\ell_d$ skeleton rows of a column.

(ii) For each full $x_d$, the entry of $S_z$ on $x_{d,k}$ is the box minor with depth $k+1$ on $z$ and $k$ on $x_d$, hence a unit. Subtract multiples of level $k$ from the levels $l<k$ (move (c) with $U_{x_d}$ lower-triangular) so that $S_z$ vanishes there, scale level $k$ to $\kappa_d$, and set all diagonal blocks to $0$ (move (a)). By Lemma~\ref{lem:box} every box minor stays a unit, and Lemma~\ref{lem:nie} applies. This is the recipe of Definition~\ref{def:B}($\mathcal B^2$) in $s$ dimensions; it needs no normal form.

(iii) By (i), the only constraints are these. On a full column $d$ ($\ell_d=2$): at carry depth $0$, $A_{id}=\lambda_i\mu_dw_{id}$ (a generalized Cauchy column); at carry depth $1$ Proposition~\ref{prop:F} only asks that the interpolating polynomial have a unit leading coefficient, and for two rows this is the determinant $\det[A_{\cdot d},\beta_{\cdot d}]$, which is (N3). On a skeleton column ($\ell_d=1$): $g_{ij}\ne0$, which is (N2). Larger carry depths are set to $0$ by Proposition~\ref{prop:Gs}, for every $k$, and the own terms by move (a). It remains to realize the ratios $A_{2d}/A_{1d}$. Put $z_1$ at $\infty$ and $z_2$ at a finite place $a_2$; then $w_{1d}=1$ and $w_{2d}=1/(a_2-b_d)$, with $b_d$ the place of $x_d$, and the condition is $b_d=a_2-\Lambda A_{1d}/A_{2d}$ with $\Lambda=\lambda_2/\lambda_1$. Take $\Lambda=1$, $\lambda_1=1$ and $\mu_d=A_{1d}$. The $b_d$ are finite, differ from $a_2$ because $A_{1d}\ne0$, and are pairwise distinct because the $2\times2$ minors of $A$ are nonzero (N1). The net property follows as in (ii).
\end{proof}

\begin{corollary}[Reachable skeletons with two skeleton coordinates]\label{cor:skelreach}
\lean{Z3NetF.skelreach_iff, Z3NetF.skelreach_nec}\leanok
\uses{prop:skelcomp}
Let $\ell=2$, and fix skeleton digits of the form in Proposition~\ref{prop:skelcomp}(iii), with coefficients that do not depend on $k$. For a suitable choice of the places they are reachable from the frame for every $k\ge1$ if and only if they satisfy (N1)--(N3). The necessity holds for each single $k\ge1$: every graph form with these skeleton digits that is a $(0,m,s)$-net over $\F_p$ satisfies (N1)--(N3).
\end{corollary}

\begin{proof}
\leanok
\uses{lem:nie, lem:box, lem:Ts, lem:skelnec, prop:skelcomp}
For skeleton digits of this form, (N1)--(N3) say: every entry and every $2\times2$ minor of $A$ is nonzero (N1, which contains (N3) for $j=1$ on a full coordinate); $g_{12}\ne0$ and $g_{21}\ne0$ (N2, which is also (N3) for $j=1$ on a skeleton coordinate); and $\det[A_{\cdot d},\beta_{\cdot d}]\ne0$ for every full coordinate $x_d$ ((N3) for $j=2$; for $\ell=2$ no skeleton coordinate lies outside a set $I$ of size $2$). These conditions involve only the coefficients, not $k$. \emph{If:} Proposition~\ref{prop:skelcomp}(iii). \emph{Only if:} let $k\ge1$, and let a graph form with the given skeleton digits be reachable from the frame for some choice of the places. All box minors of $\Phi^F$ are units (Lemma~\ref{lem:Ts}, over $\F_p$), and the moves multiply every box minor by a unit (Lemma~\ref{lem:box}), so the graph form is a $(0,m,s)$-net over $\F_p$ (Lemma~\ref{lem:nie}). Lemma~\ref{lem:skelnec} applies to every such net and gives (N1)--(N3).
\end{proof}


\section{Lattices}

\begin{lemma}[Short vectors]\label{lem:skellat}\label{lem:L12}
\lean{Z3NetF.skellat_L1, Z3NetF.skellat_L1_attained, Z3NetF.skellat_T, Z3NetF.skellat_G, Z3NetS1.L2s_min, Z3NetS1.L2s_attained, Z3NetS1.S1_cells, Z3NetS1.S1lat_index, Z3NetS1.S1lat_min, Z3NetS1.S1_opt}\leanok
Let $\Lambda_A$ be as above.
\begin{enumerate}
\item[(L)] ($\ell=1$.) If $\rho\ge2$, $\rho^s<p$ and $A=(\rho,\rho^2,\dots,\rho^{s-1})$, then $\lambda_\infty(\Lambda_A)=\rho$. For $\rho=\lfloor p^{1/s}\rfloor$ this is the largest minimum of any lattice of index $p$.
\item[(L1)] ($r=1$.) If $\rho\ge2$, $\rho^\ell(\rho+1)\le p$ and $A=(\rho,\rho^2,\dots,\rho^\ell)^{\!\top}$, then $\lambda_\infty(\Lambda_A)=\rho^\ell$.
\item[(T)] ($\ell=r=2$.) If $\rho\ge2$, $\rho^2-4<p$ and $A=\left(\begin{smallmatrix}\rho&1\\1&\rho\end{smallmatrix}\right)$, then $\lambda_\infty(\Lambda_A)\ge\rho-1$.
\item[(G)] ($\ell=2$, $r\ge3$.) Let $B\ge1$ and let $\nu_1>\sqrt2\,B$ and $\nu_2,\dots,\nu_{r-1}$ be integers with $\nu_l^2+1>B$. Put $\zeta_l=\nu_l+i$ and $\Omega_l=\zeta_1\cdots\zeta_l\in\Z[i]$, and define $A\in\Z^{2\times r}$ by
\[
A_1x+i\,A_2x=\sum_{l=1}^{r-2}x_l\,\Omega_l+(x_{r-1}+ix_r)\,\Omega_{r-1}.
\]
If $B(1+\|A_j\|_1)<p$ for $j=1,2$, then $\lambda_\infty(\Lambda_A)\ge B+1$.
\end{enumerate}
\end{lemma}

\begin{proof}
\leanok
(L) If $0<|u|_\infty\le\rho-1$, then $|z-\sum_d\rho^dx_d|\le(\rho-1)(1+\rho+\dots+\rho^{s-1})=\rho^s-1<p$, so $z=\rho(x_1+\rho x_2+\cdots)$ exactly. Then $\rho\mid z$ and $|z|<\rho$ give $z=0$, then $|x_1|<\rho$ and $\rho\mid x_1$ give $x_1=0$, and so on: $u=0$. The vector $x=e_1$, $z=\rho$ shows equality, and $\rho=\lfloor p^{1/s}\rfloor$ attains the pigeonhole bound. (L1) Let $u=(x,z_1,\dots,z_\ell)\ne0$ and $B=|u|_\infty$. If $B(\rho+1)\ge p$, then $B\ge p/(\rho+1)\ge\rho^\ell$. Otherwise $|z_1-\rho x|\le B(1+\rho)<p$ gives $z_1=\rho x$, and for $i<\ell$, $z_{i+1}-\rho z_i\equiv\rho^{i+1}x-\rho^{i+1}x=0$ with $|z_{i+1}-\rho z_i|<p$ gives $z_{i+1}=\rho z_i$; so $z_\ell=\rho^\ell x$. As $u\ne0$, $x\ne0$, so $B\ge|z_\ell|\ge\rho^\ell$; the vector $(1,\rho,\dots,\rho^\ell)$ shows equality. (T) If $0<|u|_\infty\le\rho-2$, then $|z_1-\rho x_1-x_2|\le(\rho-2)(\rho+2)<p$, so $z_1=\rho x_1+x_2$, and likewise $z_2=x_1+\rho x_2$. If $x=0$ then $z=0$. Otherwise, say $|x_1|\ge|x_2|$, and $|z_1|\ge(\rho-1)|x_1|\ge\rho-1$, a contradiction. (G) Let $0<|u|_\infty\le B$. Since $|z_j-A_jx|\le B(1+\|A_j\|_1)<p$, the relations are exact. Then $Z=z_1+iz_2$ is divisible by $\zeta_1$ in $\Z[i]$, and $|Z|\le\sqrt2B<|\zeta_1|$, so $Z=0$. Dividing by $\zeta_1$ shows that $\zeta_2$ divides the rational integer $x_1$. For a rational integer $n$, $n/\zeta_2=n(\nu_2-i)/(\nu_2^2+1)$, so $\zeta_2\mid n$ iff $\nu_2^2+1\mid n$; as $|x_1|\le B<\nu_2^2+1$, $x_1=0$. Continuing with $\zeta_3,\dots,\zeta_{r-1}$ gives $x_1=\dots=x_{r-2}=0$ and then $x_{r-1}+ix_r=0$, so $u=0$.
\end{proof}

\begin{definition}\label{def:G2rule}
\lean{Z3NetF.Bp_mono}\leanok
\uses{lem:skellat}
Let $s\ge5$, $r=s-2$ and $B\ge1$. Let $A(B)$ be the matrix of Lemma~\ref{lem:skellat}(G) with $\nu_1=\lfloor\sqrt2B\rfloor+1$ and $\nu_l=\lfloor\sqrt B\rfloor+1$ for $2\le l\le r-1$, and, for a prime $p$, let $B(p)$ be the largest integer $\beta\ge1$ with $\beta\,(1+\max_j\|A_j(\beta)\|_1)<p$. This $B(p)$ is nondecreasing in $p$.
\end{definition}

\begin{lemma}[(N1) for the Gaussian matrix]\label{lem:G2N1}
\lean{Z3NetF.G2N1_a, Z3NetF.G2N1_b, Z3NetF.G2N1_c}\leanok
Let $s\ge5$, $B\ge1$, $A=A(B)$ and $B(1+\max_j\|A_j\|_1)<p$. Then:
\begin{enumerate}
\item[(a)] $\nu_j^2+1<p$ for every $j$;
\item[(b)] if $\sum_{j=1}^{r-1}\arctan(1/\nu_j)<\pi/2$, every entry and every $2\times2$ minor of $A$ is nonzero modulo $p$, i.e. $A$ satisfies (N1);
\item[(c)] the hypothesis of (b) holds if $s\le8$ and $B\ge g^*+1$, where $g^*=1+\lceil r/2\rceil$, and, for every $s$, if $B>(2(r-1)/\pi)^2$.
\end{enumerate}
\end{lemma}

\begin{proof}
\leanok
\uses{def:G2rule}
(a) Columns $r-1$ and $r$ of $A$ are $(\operatorname{Re}\Omega_{r-1},\operatorname{Im}\Omega_{r-1})$ and $(-\operatorname{Im}\Omega_{r-1},\operatorname{Re}\Omega_{r-1})$, so $\|A_j\|_1\ge|\Omega_{r-1}|=\prod_j(\nu_j^2+1)^{1/2}\ge\sqrt5\,(\nu_1^2+1)^{1/2}$, as $r\ge3$ and $\nu_l\ge2$. Hence $p>B\sqrt5\,(\nu_1^2+1)^{1/2}$. Since $\nu_1\le\sqrt2B+1$, we have $\nu_1^2+1\le2B^2+2\sqrt2B+2\le5B^2$ for $B\ge2$, and $\nu_1^2+1=5=5B^2$ for $B=1$. So $p>\nu_1^2+1\ge \nu_j^2+1$.
(b) The arguments of all $\Omega_l$ lie in $(0,\pi/2)$, so the entries are nonzero integers, of absolute value $<p$ by the hypothesis. The $2\times2$ minor of the columns $l<l'\le r-1$ is $\operatorname{Im}(\overline{\Omega_l}\,\Omega_{l'})=|\Omega_l|^2\operatorname{Im}(\zeta_{l+1}\cdots\zeta_{l'})$, and that of the columns $l\le r-1$ and $r$ is $|\Omega_l|^2\operatorname{Re}(\zeta_{l+1}\cdots\zeta_{r-1})$. Here $|\Omega_l|^2$ is a product of the factors $\nu_j^2+1<p$, so it is prime to $p$. The second factor is a nonzero integer of absolute value at most $|\zeta_2\cdots\zeta_{r-1}|<|\Omega_{r-1}|<p$.
(c) The $\nu_j$ are nondecreasing in $B$. At $B=g^*+1$ they are $(6,3)$, $(6,3,3)$, $(8,3,3,3)$, $(8,3,3,3,3)$ for $s=5,6,7,8$, with arctangent sums below $0.487$, $0.809$, $1.090$, $1.412$, all $<\pi/2$. For every $B$, $\nu_j>\sqrt B$ (Definition~\ref{def:G2rule}), so the sum is $<(r-1)/\sqrt B$.
\end{proof}

\begin{remark}[Codes]\label{rem:codes}
\lean{Z3NetF.codes_linf, Z3NetF.codes_l2, Z3NetF.codes_mds}\leanok
\uses{prop:F, lem:CH, lem:skelnec, prop:skelcomp, thm:S1, def:S1, lem:frameid, cor:frameid, prop:S1E, rem:S1Elim}
The lattices of this section are Construction~A lattices \cite[Ch.~5]{CS99} (stated there for binary codes; the $p$-ary version is the same): with the systematic $[s,r]$ code $C_A=\{(x,Ax):x\in\F_p^r\}\subset\F_p^s$ of the skeleton digits, $\Lambda_A=C_A+p\Z^s$. For every prime $p$ the nonzero vectors of $p\Z^s$ have maximum and Euclidean norm at least $p$, and the centred lift of a codeword is a shortest lift in every coordinate. So $\lambda_\infty(\Lambda_A)$ is the least largest centred entry of a nonzero codeword, and $\lambda_2(\Lambda_A)^2=\min(p^2,w_E(C_A))$, with $w_E$ the least Euclidean (squared Lee) weight. We call skeleton digits \emph{pure} if they depend only on the last levels $x_{d,k}$ of the full coordinates (as in Definition~\ref{def:S1}). For pure skeleton digits, $P_{x_d}\equiv x_{d,k}$ and $P_z\equiv S_z\pmod p$ constrain only the lowest fine digit, so the occupied fine cells lie in $\Lambda_A$ (for $\ell=1$ they fill it, Theorem~\ref{thm:S1}(2)); the faults shift $C_A$ to cosets that vary between coarse blocks. In this language:

(a) Condition (N1) of Lemma~\ref{lem:skelnec} says that $C_A$ is MDS \cite[Ch.~11, Thm.~8]{MS77}, which says nothing about the Lee or Euclidean minimum: $A=(1,\dots,1)$ is MDS with $\lambda_\infty(\Lambda_A)=1$.

(b) Inside the frame orbit the layer at carry depth $0$ is a generalized Cauchy matrix in the places (Proposition~\ref{prop:skelcomp}(i)), so $C_A$ is a generalized Reed--Solomon code \cite{RS85} (doubly extended if $s=p+1$). This restricts the code only if $r\ge3$ and $\ell\ge3$, since every MDS code of dimension or codimension at most two is generalized Reed--Solomon; for $r\le2$ or $\ell\le2$ the restriction comes instead from the faults: the code and the faults must be realized with the same row scalars $\lambda'$ (Proposition~\ref{prop:F}), which ties their values together. Lemma~\ref{lem:frameid} is the Pl\"ucker relation (Lemma~\ref{lem:CH}(X)) applied jointly to the code and the faults; Corollary~\ref{cor:frameid} concerns the Korobov column with bounded faults, not other codes.

(c) For $\ell=1$, $\Lambda_A=\{a\in\Z^s:h\cdot a\equiv0\pmod p\}$ with the (parity-)check vector $\mathbf h=(\kappa_1,\dots,\kappa_r,-1)$ is the dual of the rank-$1$ lattice $\Z^s+\Z\mathbf h/p$ \cite{SJ94}, and $\lambda_2(\Lambda_A)$ is its spectral test \cite{CM67}, \cite[Sec.~3.3.4]{Knu97}; for $\kappa_d=\rho^d$ (Definition~\ref{def:S1}) it is Knuth's spectral-test lattice of the multiplier $\rho$ in dimension $s$, after $a_z\mapsto-a_z$. For $\rho=\lfloor p^{1/s}\rfloor$ the normalized minimum $\lambda_2(\Lambda_A)p^{-1/s}$ lies between $\rho p^{-1/s}$ and $\sqrt{\rho^2+1}\,p^{-1/s}$ and tends to $1$, against the upper bound $\sqrt{\gamma_s}$ valid for every $\mathbf h$ (Remark~\ref{rem:S1Elim}). So $\kappa_d=\rho^d$ is optimal in the maximum norm but not for the spectral test, which Proposition~\ref{prop:S1E} (Part~II) exploits.
\end{remark}


\section{One skeleton coordinate}

\begin{theorem}[One skeleton coordinate]\label{thm:S1}
\lean{Z3NetS1.S1_net, Z3NetS1.S1_sep, Z3NetS1.S1_sep_real, Z3NetS1.S1_const, Z3NetS1.S1_main, Z3NetS1.S1gen_net, Z3NetS1.S1_const_tendsto, Z3NetS1.S1_cells, Z3NetS1.S1lat_index, Z3NetS1.S1lat_min, Z3NetS1.S1_opt, Z3NetS1.S1_mu, Z3NetS1.S1_cover, Z3NetS1.S1_mesh, Z3NetS1.S1_mesh_tendsto, Z3NetKN.B2_net, Z3NetKN.B2_sep}\leanok
\uses{def:S1}
Let $s\ge3$, let $p$ be a prime, let $\rho\ge2$ be an integer with $\rho^s<p$ (one exists iff $p>2^s$), and let $k\ge0$. The net $\mathcal S^1_{k,p,\rho}$ of Definition~\ref{def:S1}, whose skeleton digit is $S_z=\sum_d\rho^dx_{d,k}$, with $N=p^{sk+s-1}$ points, is a digital $(0,sk+s-1,s)$-net over $\F_p$, and:
\begin{enumerate}
\item its toroidal (hence also Euclidean) minimum distance satisfies $\delta N^{1/s}\ge(\rho-1)\,p^{-1/s}$; for $\rho=\lfloor p^{1/s}\rfloor$ the bound exceeds $(\rho-1)/(\rho+1)$, so it is $>1/3$ for every prime $p>2^s$, $>1/2$ for $p>3^s$, and tends to $1$;
\item its occupied fine cells are exactly the cells of a lattice of index $p$ with minimum $\lambda_\infty=\rho$, which for $\rho=\lfloor p^{1/s}\rfloor$ is the largest minimum of any lattice of index $p$;
\item for $\rho=\lfloor p^{1/s}\rfloor$ its covering radius satisfies $hN^{1/s}\le\sqrt s\,(\mu+1)\,p^{-1/s}$ with $\mu=\lceil p/\rho^{s-1}\rceil<\rho(1+1/\rho)^s+1$, so the mesh ratio $h/\delta\le\sqrt s\,(\mu+1)/(\rho-1)$ is bounded uniformly in $p$ and $k$, and is at most $\sqrt s\,(1+o(1))$ as $p\to\infty$.
\end{enumerate}
\end{theorem}

\begin{proof}[Proof of Theorem~\ref{thm:S1}]
\leanok
\uses{lem:skelred, prop:skelcomp, lem:skellat}
\emph{Net.} Proposition~\ref{prop:skelcomp}(ii): $p>2^s\ge s-1$, and the prescribed coefficients $\rho^d$ are units since $\rho^{s-1}<p$.

(1) Since $x_{d,k}\equiv P_{x_d}$ and $P_z\equiv S_z\pmod p$, every point satisfies $P_z\equiv\sum_d\rho^dP_{x_d}\pmod p$. This congruence passes to the differences modulo $p^{k+1}$, so every difference $a\ne0$ lies in $\Lambda_A$ of Lemma~\ref{lem:skellat}(L), and $|a|_\infty\ge\rho$. Lemma~\ref{lem:skelred}(i) gives $\delta N^{1/s}\ge(\rho-1)p^{-1/s}$. There are no faults. With $\rho=\lfloor p^{1/s}\rfloor$ we have $p<(\rho+1)^s$, so the bound exceeds $(\rho-1)/(\rho+1)$, which is $\ge1/3$ for $\rho\ge2$ ($p>2^s$) and $\ge1/2$ for $\rho\ge3$ ($p>3^s$).

(2) There are $N=p^{s(k+1)-1}$ points with distinct $P$, all in $\Lambda_A$, and $\Lambda_A$ has $p^{s(k+1)}/p=N$ classes modulo $p^{k+1}\Z^s$. So the occupied fine cells are exactly the cells of $\Lambda_A$. Optimality is Lemma~\ref{lem:skellat}(L).

(3) Let $Q\in\{0,\dots,p^{k+1}-1\}^s$ be the fine cell of a point $y\in[0,1)^s$. We look for a cell $a=Q-v\in\Lambda_A$ in the same range with $|v|_\infty\le\mu:=\lceil p/\rho^{s-1}\rceil$. Prescribe the sign $\epsilon_c$ of $v_c$: $-1$ if $Q_c<p^{k+1}/2$, else $+1$. Since $\mu\le p/4+1\le p^{k+1}/2$, any $v$ with these signs and $|v|_\infty\le\mu$ keeps $a$ in range. With $T\equiv Q_z-\sum_d\rho^dQ_{x_d}\pmod p$, the vector $a$ lies in $\Lambda_A$ iff $v_z-\sum_d\rho^dv_{x_d}\equiv T$. Choose the representative $T'$ of $T$ in $[0,p)$ or in $(-p,0]$, and expand $T'=\sum_{i=0}^{s-1}t_i\rho^i$ in base $\rho$ with signed digits: $t_0,\dots,t_{s-2}$ are taken one at a time from the window $[0,\rho-1]$ or $[-(\rho-1),0]$, according to the sign of $T'$, and $t_{s-1}$ is the rest. Put $v_z=t_0$ and $v_{x_i}=-t_i$; the windows give the required signs of $v_z,v_{x_1},\dots,v_{x_{s-2}}$. The lower digits contribute less than $\rho^{s-1}$ in absolute value, so $t_{s-1}$ has the sign of $T'$ (which we choose) and $|t_{s-1}|<p/\rho^{s-1}+1$. So $|v|_\infty\le\mu$, since also $\rho-1<p/\rho^{s-1}$. The point in the cell $a$ differs from $y$ by less than $(\mu+1)p^{-(k+1)}$ in every coordinate, so
\[
hN^{1/s}\le\sqrt s\,(\mu+1)\,p^{-1/s},\qquad \frac h\delta\le\frac{\sqrt s\,(\mu+1)}{\rho-1}.
\]
For $\rho=\lfloor p^{1/s}\rfloor$, $\mu<\rho(1+1/\rho)^s+1$, so both bounds are bounded in $p$ for fixed $s$, and $h/\delta\le\sqrt s\,(1+o(1))$ as $p\to\infty$.
\end{proof}

\begin{definition}\label{def:S1}
\lean{Z3NetS1.S1_net, Z3NetS1.S1_sep, Z3NetS1.S1_sep_real, Z3NetS1.S1gen_net}\leanok
\uses{prop:skelcomp}
Let $s\ge3$, $\rho\ge2$ with $\rho^s<p$, $k\ge0$ and $m=sk+s-1$. Take the frame that assigns the places $\infty,0,1,\dots,s-2$ to the coordinates $x_1,\dots,x_{s-1},z$, respectively, the depths $(k+1,\dots,k+1,k)$ with the skeleton coordinate $z$ last, and the recipe of Proposition~\ref{prop:skelcomp}(ii) with $\kappa_d=\rho^d$. The result is the digital net $\mathcal S^1_{k,p,\rho}$ over $\F_p$. Its skeleton digit is
\[
S_z=\sum_{d=1}^{s-1}\rho^d\,x_{d,k}.
\]
\end{definition}


\section{Two skeleton coordinates}

\begin{theorem}[Two skeleton coordinates]\label{thm:S2}
\lean{Z3NetF.S2_net4, Z3NetF.S2_sep4, Z3NetF.S2_sep4_real, Z3NetF.S2_const4, Z3NetF.S2_net, Z3NetF.S2_sep, Z3NetF.S2_sep_real, Z3NetF.S2_thr, Z3NetF.S2_main58, Z3NetF.S2_eventually, Z3NetF.Bp_asymp}\leanok
\uses{def:G2rule, def:S2}
Let $s\ge4$, let $p$ be a prime and $k\ge1$, and let $\mathcal S^2_{k,p}$ be the net of Definition~\ref{def:S2}, with $N=p^{sk+s-2}$ points.
\begin{enumerate}
\item ($s=4$.) For every prime $p\ge5$, $\mathcal S^2_{k,p}$ is a digital $(0,4k+2,4)$-net over $\F_p$ with $\delta N^{1/4}\ge(\rho-4)\,p^{-1/2}$, where $\rho=\max\{\rho:\rho^2<p+4\}$; the bound is positive for every prime $p\ge23$ and tends to $1$.
\item ($5\le s\le8$.) For every prime $p\ge137$, $443$, $2179$, $6101$ ($s=5,6,7,8$), $\mathcal S^2_{k,p}$ is a digital $(0,sk+s-2,s)$-net over $\F_p$ with $\delta N^{1/s}\ge(B(p)-g^*)\,p^{-2/s}>0$, where $g^*=1+\lceil(s-2)/2\rceil$ and $B(p)\asymp p^{2/s}$ is defined in Definition~\ref{def:G2rule}.
\item ($s\ge9$.) The same holds for every prime $p$ with $B(p)>g^*$ and $\sum_{j=1}^{s-3}\arctan(1/\nu_j)<\pi/2$ for the $\nu_j$ of Definition~\ref{def:G2rule} at $B=B(p)$; this holds if $B(p)>\max\{g^*,(2(s-3)/\pi)^2\}$, in particular for all sufficiently large $p$.
\end{enumerate}
\end{theorem}

\begin{proof}[Proof of Theorem~\ref{thm:S2}]
\leanok
\uses{lem:skelred, prop:skelcomp, lem:skellat, lem:G2N1}
\emph{(N1).} For $s=4$ the entries $\rho$, $1$ are units, and $\det A=\rho^2-1=(\rho-1)(\rho+1)$ with $0<\rho-1<\rho+1<p$, since $\rho^2<p+4$ and $p\ge5$. For $s\ge5$ use Lemma~\ref{lem:G2N1}: in (2), $B(p)\ge g^*+1$ exactly for $p>132$, $440$, $2170$, $6095$ ($s=5,\dots,8$), which gives the primes $137$, $443$, $2179$, $6101$. These four numbers are $\Gamma(g^*+1)$ for $\Gamma(B)=B(1+\max_j\|A_j(B)\|_1)$, and $\Gamma$ is increasing on $B\ge g^*+1$ (below), so some $B\ge g^*+1$ has $\Gamma(B)<p$ iff $\Gamma(g^*+1)<p$. In (3) the hypothesis is that of Lemma~\ref{lem:G2N1}(b), and the closed-form sufficient condition is that of Lemma~\ref{lem:G2N1}(c). \emph{Net.} Proposition~\ref{prop:skelcomp}(iii).

\emph{Monotonicity of $\Gamma$.} Let $5\le s\le8$ and $g^*+1\le B<B'$. The $\nu_j$ do not decrease from $B$ to $B'$, and $\sum_j\arctan(1/\nu_j)<\pi/2$ at $B$ by Lemma~\ref{lem:G2N1}(c). The absolute values of the entries of $A(B)$ are the real and imaginary parts of $\Omega_1,\dots,\Omega_{r-1}$ (proof of Lemma~\ref{lem:G2N1}(a)). Raise the $\nu_j$ one at a time from their values at $B$ to those at $B'$. If $\nu_j$ is a factor of $\Omega_l=P\,(\nu_j+i)$, the other factors have arctangents no larger than at $B$, so $P=a+ib$ with $\arg P\in[0,\pi/2)$, i.e. $a>0$, $b\ge0$; and $P(\nu_j+i)=(a\nu_j-b)+i(b\nu_j+a)$, so raising $\nu_j$ by $\Delta\ge0$ raises the real part by $a\Delta$ and the imaginary part by $b\Delta$. All parts stay positive, so every $|A_{jd}|$ is nondecreasing, $\|A_j(B')\|_1\ge\|A_j(B)\|_1$, and $\Gamma(B')>\Gamma(B)$. (The thresholds are also verified in Lean, by a different route: $\Gamma$ is evaluated exactly for $g^*+1\le B<B_1$, $B_1=7,7,9,9$, and $(B\max_j\|A_j(B)\|_1)^2\ge2B^s$, from $\|A_j\|_1\ge|\Omega_{r-1}|$, covers $B\ge B_1$.)

\emph{Separation.} Let $a\ne0$ be a difference and $M=|a|_\infty$. If $M\ge p-1$, then $M\ge\lambda_\infty(\Lambda_A)-g^*$ below holds, since $\lambda_\infty(\Lambda_A)\le\lfloor p^{2/s}\rfloor\le p-1$. Otherwise let $\xi_c=\xi_{c,1}$ be the floor changes of Lemma~\ref{lem:skelred}(ii); they give the changes of the level-$(k-1)$ digits modulo $p$. Put
\[
u=\Bigl(a_x,\ a_{z_1}-\sum_{d\in I_1}\xi_{x_d}-\xi_{z_2},\ a_{z_2}-\sum_{d\in I_2}\xi_{x_d}+\xi_{z_1}\Bigr).
\]
By the skeleton digits, the change of $S_{z_i}$ is $A_ia_x$ plus the contributions of the level-$(k-1)$ changes $\xi$; subtracting these, as in $u$, leaves a vector of $\Lambda_A$. So $u\in\Lambda_A$, and $|u|_\infty\le M+g^*$ with $g^*=1+\lceil r/2\rceil$. If $u\ne0$, then $M\ge\lambda_\infty(\Lambda_A)-g^*$. If $u=0$, then $a_x=0$, so every $\xi_{x_d}=0$, and $a_{z_1}=\xi_{z_2}$, $a_{z_2}=-\xi_{z_1}$. If $\xi_{z_2}=1$ then $a_{z_2}>0$ and $a_{z_1}=1$, so $\xi_{z_1}\in\{0,1\}$ and $a_{z_2}=-\xi_{z_1}\le0$, a contradiction. The case $\xi_{z_2}=-1$ is symmetric, and $\xi_{z_2}=0$ gives $a_{z_1}=0$, $\xi_{z_1}=0$, $a_{z_2}=0$, contradicting $a\ne0$. Here the opposite signs of $g_{12}$ and $g_{21}$ give the contradiction, so no dependence of $S_{z_i}$ on its own cell is needed. So $M\ge\lambda_\infty(\Lambda_A)-g^*$, and Lemma~\ref{lem:skelred}(i) gives $\delta N^{1/s}\ge(\lambda_\infty(\Lambda_A)-g^*-1)p^{-2/s}$. With $\lambda_\infty\ge\rho-1$ and $g^*=2$ ($s=4$), resp. $\lambda_\infty\ge B+1$ (Lemma~\ref{lem:skellat}), this is the claim. For $s=4$, $\rho\ge5$ iff $p\ge23$, and $(\rho-4)/\sqrt p\to1$.
\end{proof}

\begin{definition}\label{def:S2}
\lean{Z3NetF.IsS2skel, Z3NetF.S2A}\leanok
\uses{prop:skelcomp, lem:skellat, def:G2rule}
Let $s\ge4$, $r=s-2$, $k\ge1$, $m=sk+s-2$. For $s=4$ let $A$ be the matrix of Lemma~\ref{lem:skellat}(T) with $\rho$ maximal such that $\rho^2-4<p$, and put $I_1=\{2\}$, $I_2=\{1\}$. For $s\ge5$ let $A=A(B(p))$ (Definition~\ref{def:G2rule}), and let $I_1\sqcup I_2=\{1,\dots,r\}$ with $|I_1|=\lfloor r/2\rfloor$. The net $\mathcal S^2_{k,p}$ is the completion (Proposition~\ref{prop:skelcomp}(iii)) of the skeleton digits
\[
S_{z_1}=A_1x_{\cdot,k}+\sum_{d\in I_1}x_{d,k-1}+z_{2,k-1},\qquad S_{z_2}=A_2x_{\cdot,k}+\sum_{d\in I_2}x_{d,k-1}-z_{1,k-1},
\]
with no dependence of $S_{z_i}$ on the cell of $z_i$ itself (own coefficients $0$), and $x_{\cdot,k}=(x_{1,k},\dots,x_{r,k})$.
\end{definition}

\begin{remark}[Comparison with $m\equiv0$]\label{rem:skelr0}
\uses{thm:A, lem:skelred}
Seen one digit deeper, the case $m=sk$ fits this picture as the case $\ell=s$: every coordinate is a skeleton coordinate whose skeleton digit is its geometric digit, the occupied fine cells are a translate of all of $p\Z^s$, and the geometric layer $\tilde M$ plays the role of the faults; so $1-K/q$ is the analogue of ``$\lambda_\infty$ minus faults'' (up to the one fine cell more that Lemma~\ref{lem:skelred} loses). Since $p\Z^s$ attains the pigeonhole bound, only the scale of the faults differs: $K/q$ at $r=0$ against $g^*\,p^{-\ell/s}$ for $r\ne0$. This is why the thresholds ($p>2^s$ for $r=s-1$; $p\ge23$, $137$, $443$, $2179$, $6101$ for $r=s-2$, $s=4,\dots,8$) are far below the threshold $5a^*_s+4\le e^{s\log s+O(s)}$ of Theorem~\ref{thm:A}, while at a fixed large $p$ the constants for $r=0$ are larger: the skeletons gain in the range of $p$, not in the constant (heuristic comparison).
\end{remark}


\section{Three or more skeleton coordinates}

\begin{lemma}[A Pl\"ucker identity for skeleton digits]\label{lem:frameid}
\lean{Z3NetF.frameid}\leanok
\uses{lem:box}
Let $\ell\ge3$, let a graph form be obtained from a frame over $\F_p$ by the moves of Lemma~\ref{lem:box}, and let $d$ be a full coordinate. Let $\kappa_i$ be the coefficient of $S_{z_i}$ on $x_{d,k}$, $g_{ij}$ that on $z_{j,k-1}$, and $\psi_i=1/\kappa_i$. Then, for distinct $i,j,l$,
\[
g_{ji}\,g_{lj}\,\psi_j+g_{ij}\,g_{li}\,\psi_i=g_{li}\,g_{lj}\,\psi_l .
\]
\end{lemma}

\begin{proof}
\leanok
\uses{lem:CH, prop:skelcomp}
The moves only scale the entries of the first offset rows on the last cell level, so by Proposition~\ref{prop:skelcomp}(i) these entries form a Cauchy layer $W_{ce}=\lambda_c\mu_e/D_{ce}$ in the places (Section~\ref{sec:frames}), with $W_{z_iz_j}=g_{ij}$ and $W_{z_ix_d}=\kappa_i$. Lemma~\ref{lem:CH}(X) for the coordinates $z_i,z_j,z_l,x_d$ reads $g_{ji}\kappa_l/(\kappa_jg_{li})+g_{ij}\kappa_l/(\kappa_ig_{lj})=1$; multiplying by $g_{li}g_{lj}/\kappa_l$ gives the identity.
\end{proof}

\begin{corollary}[Korobov skeleton digits need large faults]\label{cor:frameid}
\lean{Z3NetF.frameid_cor}\leanok
\uses{lem:box, lem:frameid}
Let $\ell\ge3$, $k\ge1$, let $d$ be a full coordinate, and let $\rho,G\ge1$ be integers with $G^2<\rho$ and $p>G^2(2\rho^2+1)$. Then no graph form obtained from a frame over $\F_p$ by the moves of Lemma~\ref{lem:box} has, in the notation of Lemma~\ref{lem:frameid}, $\kappa_1,\kappa_2,\kappa_3\equiv\rho,\rho^2,\rho^3$ and centred $g_{ij}$ with $|g_{ij}|\le G$ for $i\ne j\le3$.
\end{corollary}

\begin{proof}
\leanok
Suppose such a graph form exists. The moves multiply every box minor by a unit, so it is a $(0,m,s)$-net, and $g_{ij}\ne0$ by (N2). Multiply the identity for $(i,j,l)=(1,2,3)$ by $\rho^3$: $g_{21}g_{32}\rho+g_{12}g_{31}\rho^2\equiv g_{31}g_{32}\pmod p$. The difference of the two sides is at most $G^2(\rho^2+\rho+1)<p$ in absolute value, so the congruence is an equality of integers. Hence $\rho\mid g_{31}g_{32}$. But $0<|g_{31}g_{32}|\le G^2<\rho$, a contradiction.
\end{proof}

\begin{proposition}[A Korobov skeleton with faults $\pm1$; computer-assisted]\label{prop:korobov4}
\lean{Z3NetF.t3s4_reach, Z3NetF.t3s4_net, Z3NetF.t3s4_sep, Z3NetF.t3s4_sep_real, Z3NetF.t3s4_M, Z3NetF.t3s4_main, Z3NetF.t3s4_table, Z3NetF.t3s4_lat, Z3NetF.t3s4_tab}\leanok
Let $p=13$, $s=4$ and $m=4k+1$ with $k\ge1$, with one full coordinate $x$ of depth $k+1$ (we write $x_l$ for its level $l$) and three skeleton coordinates $z_1,z_2,z_3$ of depth $k$, and take the frame with the place $\infty$ for $x$ and the places $0$, $1$, $3$ for $z_1$, $z_2$, $z_3$. The skeleton digits
\begin{align*}
S_{z_1}&=4x_k-2x_{k-1}+x_{k-2}-z_{1,k-1}-z_{2,k-1}+z_{2,k-2}-z_{3,k-1}+z_{3,k-2},\\
S_{z_2}&=7x_k-x_{k-1}-z_{1,k-1}+z_{1,k-2}-z_{3,k-1},\\
S_{z_3}&=8x_k-x_{k-1}-z_{1,k-1}+z_{2,k-1}
\end{align*}
(terms with nonexistent levels omitted) are reachable from the frame for every $k\ge1$, so they complete to digital $(0,4k+1,4)$-nets over $\F_{13}$. Every graph form over $\F_{13}$ with the depths $(k+1,k,k,k)$ and these skeleton digits, net or not, satisfies $\delta N^{1/4}\ge3\cdot13^{-3/4}>0.438$ (toroidal).
\end{proposition}

\begin{proof}
\leanok
\uses{lem:nie, lem:box, lem:Ts, prop:Gs, prop:F, lem:skelred, prop:skelcomp}
Write $\kappa=(4,7,8)$, $\beta=(-2,-1,-1)$ and $\eta=(1,0,0)$ for the coefficients of the $S_{z_i}$ on $x_k$, $x_{k-1}$, $x_{k-2}$, $g_{ij}$ and $H_{ij}$ for those of $S_{z_i}$ on $z_{j,k-1}$ and $z_{j,k-2}$ ($i\ne j$), and $\alpha=(0,1,3)$ for the places of $z_1,z_2,z_3$; all faults $g_{ij}$ are $\pm1$.

\emph{Reachability.} By Proposition~\ref{prop:skelcomp}(i), Propositions~\ref{prop:F} and~\ref{prop:Gs} hold for the skeleton rows: column $x$ carries the three rows $S_{z_i}$ and the levels $k$, $k-1$, $k-2$ (the carry depths $0,1,2$), and column $z_j$ carries the two rows $S_{z_i}$, $i\ne j$, and the levels $k-1$, $k-2$; one vector $\lambda'$ of row scalars is shared by all columns. The coefficients of $S_{z_i}$ on its own cell are diagonal entries, free by move (a). By Proposition~\ref{prop:F} (with $w=1$, $y=\alpha_i$ on the column at $\infty$, and $w=y=1/(\alpha_i-\alpha_j)$ on the column $z_j$), the digits are reachable, with all deeper levels $0$, if there are units $\lambda'_i$, $\mu_x$, $\mu_{z_j}$ such that, modulo $13$:
(R1) $\kappa_i=\lambda'_i\mu_x$; $\kappa_i+\beta_i=\lambda'_i(c_0+c_1\alpha_i)$ with $c_1\ne0$; and $\kappa_i+\beta_i+\eta_i=\lambda'_iG(\alpha_i)$ with $\deg G=2$;
(R2) $g_{ij}=\lambda'_i\mu_{z_j}/(\alpha_i-\alpha_j)$, and $g_{ij}+H_{ij}=\lambda'_i(\alpha_i-\alpha_j)^{-1}\bigl(c'_0+c'_1/(\alpha_i-\alpha_j)\bigr)$ with $c'_1\ne0$, for each $j$ and the two $i\ne j$.
Take $\lambda'=\kappa$ and $\mu_x=1$. For (R2) at level $k-1$ each column must give one value of $\mu_{z_j}=g_{ij}(\alpha_i-\alpha_j)/\kappa_i$ for both $i\ne j$: column $z_1$ needs $7\mu=-1$ and $8\mu=-3$, column $z_2$ needs $4\mu=1$ and $8\mu=2$, column $z_3$ needs $4\mu=3$ and $7\mu=2$, and indeed $(\mu_{z_1},\mu_{z_2},\mu_{z_3})=(-2,-3,4)$ solves all six. For (R1), $(\kappa_i+\beta_i)/\kappa_i=(2/4,6/7,7/8)=(7,12,9)=c_0+c_1\alpha_i$ with $(c_0,c_1)=(7,5)$. The remaining conditions say that leading coefficients of interpolation polynomials are nonzero: three values $\lambda'_iG(\alpha_i)$ at distinct $\alpha_i$ always come from a $G$ of degree at most $2$, and since $\det[\lambda',\lambda'\odot(c_0+c_1\alpha),\lambda'\odot G(\alpha)]$ is $\lambda'_1\lambda'_2\lambda'_3\,c_1$ times the leading coefficient of $G$ times a Vandermonde determinant, $\deg G=2$ is equivalent to $\det[\kappa,\beta,\eta]\ne0$; here $\det[\kappa,\beta,\eta]=7\cdot(-1)-8\cdot(-1)=1$. Likewise $c'_1\ne0$ in column $z_j$ is equivalent to $g_{ij}H_{lj}-H_{ij}g_{lj}\ne0$ for $\{i,l\}=\{1,2,3\}\setminus\{j\}$, and these are $1$, $-1$, $1$ for $j=1,2,3$. The places $\infty,0,1,3$ are distinct modulo $13$. So Proposition~\ref{prop:Gs} gives, for every $k\ge1$, moves of Lemma~\ref{lem:box} that turn the graph form $\Phi^F$ of the frame into one with these skeleton digits (truncated for $k=1$). The box minors of $\Phi^F$ are units (Lemma~\ref{lem:Ts}), the moves multiply them by units, and Lemma~\ref{lem:nie} gives the net property.

\emph{Separation.} Let $a$ be the centred difference of the fine positions of two distinct points (Lemma~\ref{lem:skelred}). If some $|a_c|\ge13/2$, then $|a|_\infty\ge7$. Otherwise, as $P_x\equiv x_k$ and $P_{z_i}\equiv S_{z_i}\pmod{13}$, Lemma~\ref{lem:skelred}(ii) gives
\[
a_{z_i}\equiv\kappa_ia_x+\beta_i\xi_{x,1}+\eta_i\xi_{x,2}+\sum_{j\ne i}\bigl(g_{ij}\xi_{z_j,1}+H_{ij}\xi_{z_j,2}\bigr)-[i=1]\,\xi_{z_1,1}\pmod{13},
\]
with sign-consistent prefix patterns $(\xi_{c,1},\xi_{c,2})\in\{(0,0),(\sgn a_c,0),(\sgn a_c,\sgn a_c)\}$ (and $\xi=0$ for missing levels). We enumerate $|a_x|\le6$ and all patterns ($4625$ combinations, with both signs for the patterns of the $z_i$ before the sign check; rerun by the scripts of Appendix~\texttt{app:verify}), take the centred $a_{z_i}$ and keep the sign-consistent solutions: the least $|a|_\infty$ of a nonzero solution is $M_{13}=4$. Lemma~\ref{lem:skelred}(i) gives $\delta N^{1/4}\ge(M_{13}-1)\,13^{-3/4}=0.4381\ldots$.
\end{proof}


\chapter{Dimension three, every resolution}\label{sec:dim3}


\section{\texorpdfstring{$m\equiv0$}{m = 0}: the scramble}\label{sec:construction}

\begin{definition}\label{def:net}
\lean{Z3Net.projlines_P, Z3Net.tau, Z3Net.T, Z3Net.K, Z3Net.E, Z3Net.L, Z3Net.PGen, Z3Net.Ppoint}\leanok
$\mathcal P_{k,q}$ is the digital net over $\Z/q\Z$ with generator matrices $(LT^2,\,L,\,LT)$. For $V\in(\Z/q\Z)[t]_{<3k}$ the point has coordinates $\phi_q(d(W))$ for $W=T^2V,V,TV$, where $d(W)=L\,[W]_m$:
\begin{align*}
d_j&=\coef{j}\,W(1-t)^{-k}\quad(j\notin\{k-2,k-1,k\}),&
d_{k-2}&=\coef{k-2}\,W(1-t)^{-k-1},\\
d_{k-1}&=-\coef{k-1}\,W(1-t)^{-k},&
d_k&=\coef{k}\,(1+t)\,W(1-t)^{-k}.
\end{align*}
\end{definition}


\section{The net property}

\begin{theorem}\label{thm:t0}
\lean{Z3Net.P_net_rowScale, Z3Net.P_boxMinor, Z3Net.P_net}\leanok
For every $k\ge1$, every box minor of $(LT^2,L,LT)$ equals $\pm1$. Hence $\mathcal P_{k,q}$ is a $(0,3k,3)$-net in base $q$ for every $q\ge2$.
\end{theorem}

\begin{proof}
\leanok
\uses{lem:nie}
From $(\Psi V)(t)=\sum_iv_i(1-t)^{m-1-i}$ one reads off, over any commutative ring:
(F1) if $t^j\mid V$ then $\deg\Psi V<m-j$; (F2) if $\deg V<m-j$ then $(1-t)^j\mid\Psi V$.
Let $d_1+d_2+d_3=m$ and suppose that, over a field, the first $d_1,d_2,d_3$ digits of $T^2V$, $V$, $TV$ vanish: $t^{d_1}\mid\Psi^2V$, $t^{d_2}\mid V$, $t^{d_3}\mid\Psi V$. Since $V=\pm\Psi(\Psi^2V)$, (F1) gives $\deg V<m-d_1=d_2+d_3$; and (F1) applied to $\Psi V$ gives $\deg\Psi^2V<m-d_3$, whence $(1-t)^{d_3}\mid V$ by (F2). As $1=t+(1-t)$, the factors $t^{d_2}$ and $(1-t)^{d_3}$ are coprime, so their product, of degree $d_2+d_3$, divides $V$; hence $V=0$. So every box minor of the frame is nonzero modulo every prime, i.e. $\pm1$; the scramble multiplies it by $\pm1$ (\texttt{eq:scr}), and Lemma~\ref{lem:nie} concludes.
\end{proof}


\section{Cells and the geometric digit}\label{sec:sep}

\begin{proposition}\label{prop:cell}
\lean{Z3Net.cell_digits, Z3Net.cell_layer}\leanok
The geometric digit of each coordinate depends only on the last two cell digits of the three coordinates, and the deeper digits of a coordinate do not depend on its own cell. Precisely, for every $V$, writing $d^x,d^y,d^z$ for the digit vectors of the three coordinates:
(a) $d^x_k=-d^x_{k-1}-d^y_{k-1}+d^z_{k-1}-d^z_{k-2}$;
(b) for $j>k$, $d^x_j$ is a function of the cells $(n_y,n_z)$ alone;
(c) the same holds after the cyclic shift $(x,y,z)\mapsto(y,z,x)$.
Consequently $r\equiv Mu+Jv\pmod q$ with
\[
M=\begin{pmatrix}-1&-1&1\\1&-1&-1\\-1&1&-1\end{pmatrix},\qquad J=\begin{pmatrix}0&0&-1\\-1&0&0\\0&-1&0\end{pmatrix},
\]
and $f_c$ does not depend on $n_c$.
\end{proposition}

\begin{proof}
\leanok
For $\deg F\le N$ put $\Psi_NF=(1-t)^NF(1/(1-t))$, so $\Psi=\Psi_{m-1}$, and $\Psi_{a+b}(FG)=\Psi_aF\cdot\Psi_bG$ whenever $\deg F\le a$, $\deg G\le b$ (on binary forms, $\Psi_N$ is the $N$-th symmetric power of the matrix $g$ of Section~\ref{sec:frames}). Moreover $\Psi_kt^k=1$, $\Psi_k(t-1)^k=t^k$, $\Psi_k(1-t)^k=(-t)^k$, $\Psi_k1=(1-t)^k$, and for $\deg F<k$, with $\Psi'=\Psi_{k-1}$: $(\Psi'^2F)(0)=F_{k-1}=(\Psi'F)(1)$. Partial fractions of $V/(t^k(t-1)^k)$ (over $\Z$: divide by the monic $t^k(t-1)^k$, then use a B\'ezout identity $1=\alpha t^k+\beta(t-1)^k$ in $\Z[t]$, obtained by expanding $1=(t-(t-1))^{2k-1}$) give
\[
V=A\,(t-1)^k+B\,t^k+Q\,t^k(t-1)^k,\qquad\deg A,\deg B,\deg Q<k,
\]
and multiplicativity gives
\begin{align*}
\Psi V&=\Psi'A\cdot t^k(1-t)^k+\Psi'B\cdot(1-t)^k+\Psi'Q\cdot t^k,\\
\Psi^2V&=(-t)^k\,\Psi'^2A+(-t)^k(1-t)^k\,\Psi'^2B+(1-t)^k\,\Psi'^2Q .
\end{align*}
Write $\psi_j(W)=(\mathcal K[W]_m)_j=\coef{j}W(1-t)^{-k}$, so that $d(W)=\mathcal E\,(\psi_j(W))_j$.

\emph{Cells and (b).} Modulo $t^k$, $V\equiv A(t-1)^k$ and $\Psi V\equiv\Psi'B\,(1-t)^k$, so the $y$-cell determines $A$ and the $z$-cell $B$ ($L$ is lower-triangular with diagonal $\pm1$, $(t\mp1)^k$ are units modulo $t^k$, and $\Psi'$ is invertible). For $j<k$, $\psi_j(\Psi^2V)=\coef{j}\Psi'^2Q$, so the $x$-cell determines $Q$, and the three cells determine $V$. For $j\ge k$, $\psi_j(\Psi^2V)=\coef{j}(-t)^k\bigl(\Psi'^2A\,(1-t)^{-k}+\Psi'^2B\bigr)$ depends on $A,B$ only; as $\mathcal E$ changes only rows $k-2,k-1,k$, this is (b).

\emph{(a).} From the displays, $\psi_k(\Psi^2V)=(-1)^k(A_{k-1}+B_{k-1})$, $\psi_{k-1}(V)=(-1)^kA_{k-1}$ and, as $\Psi V(1-t)^{-k}\equiv\Psi'B\pmod{t^k}$, $\sum_{j<k}\psi_j(\Psi V)=(\Psi'B)(1)=B_{k-1}$. By the definition of $\mathcal E$ (and $T^2=\Psi^2$, $T=s_k\Psi$), $d^x_k=-d^x_{k-1}+(-1)^k(A_{k-1}+B_{k-1})$, $d^y_{k-1}=-(-1)^kA_{k-1}$ and $d^z_{k-2}-d^z_{k-1}=s_kB_{k-1}$; with $(-1)^ks_k=-1$ this is (a). Part (c) follows from $T^3=1$.
\end{proof}

\begin{remark}[Lattice reading]\label{rem:lattice3}
\lean{Z3Net.lat3_cross, Z3Net.lat3_rot, Z3Net.lat3_decomp, Z3Net.lat3_det, Z3Net.lat3_norm, Z3Net.lat3_min, Z3Net.lat3_min_attained, Z3Net.lat3_cP}\leanok
\uses{thm:sep}
Figure~\texttt{fig:arcsmall} shows the lattice picture for $\mathcal P_{k,q}$. On a patch where no geometric digit wraps, the points form, up to the jitter, a piece of a translate of the lattice $\Lambda_q=(I+M/q)\Z^3$ (in cell units). Since $M=-I+[\mathbf 1]_\times$, where $[\mathbf 1]_\times E=\mathbf 1\times E$ is the cross product with $\mathbf 1=(1,1,1)$, $\Lambda_q$ is $\Z^3$ rotated by $\arctan(\sqrt3/(q-1))$ about $\mathbf 1$, scaled by $1-1/q$ along $\mathbf 1$ and by $((1-1/q)^2+3/q^2)^{1/2}$ across it; $\det(I+M/q)<1$. Its shortest vectors are $\pm(I+M/q)e_i$, e.g. $(I+M/q)e_x=(1-\tfrac1q,\tfrac1q,-\tfrac1q)$, of length $\lambda_1(\Lambda_q)=(1-2/q+3/q^2)^{1/2}=1-1/q+O(q^{-2})$: for $E\in\Z^3$, $|(I+M/q)E|^2=(1-1/q)^2|E|^2+(3/q^2)|E_\perp|^2$ with $E_\perp$ the component orthogonal to $\mathbf 1$, which is at least $2(1-1/q)^2>1-2/q+3/q^2$ when $|E|^2\ge2$ and $q\ge3$. The own-axis term $1-1/q$ of $c(q)$ is its component along the moving axis. The proof of Theorem~\ref{thm:sep} uses only this component: the transverse parts $\pm1/q$ are of the size of the jitter, which may cancel them, so the certified $c(q)$ is below $\lambda_1(\Lambda_q)$. $\Lambda_q$ has covolume $\det(I+M/q)<1$, so it has slightly more than one site per cell; reducing $Mu$ modulo $q$ skips the extra sites at the slip planes, so that each cell keeps exactly one point. Theorem~\ref{thm:sep} bounds the steps coordinatewise and does not use $\lambda_1(\Lambda_q)$.
\end{remark}


\section{The separation bound}

\begin{theorem}[Dimension three, $m\equiv0$]\label{thm:main}
\lean{Z3Net.P_boxMinor, Z3Net.P_net, Z3Net.P_covering, Z3Net.P_sep, Z3Net.P_sep_ge8, Z3Net.P_sep_real}\leanok
\uses{def:net, thm:t0}
Let $q\ge 4$ and $k\ge 2$, and let $\mathcal P_{k,q}$ be the digital net of Definition~\ref{def:net}, with $N=q^{3k}$ points. Then:
\begin{enumerate}
\item $\mathcal P_{k,q}$ is a $(0,3k,3)$-net in base $q$ (Theorem~\ref{thm:t0}; this holds for every $q\ge2$);
\item every cell $\prod_i[n_i q^{-k},(n_i+1)q^{-k})$ contains exactly one point, so $h_{\mathcal P}\le\sqrt3\,N^{-1/3}$;
\item the toroidal (hence also the Euclidean) minimum distance satisfies $\delta_{\mathcal P}\ge c(q)\,N^{-1/3}$, where
\[
c(q)^2=\min\Bigl\{\bigl(1-\tfrac1q\bigr)^2,\ \bigl(1-\tfrac2q\bigr)^2+\bigl(1-\tfrac4q\bigr)^2,\ 3\bigl(1-\tfrac3q\bigr)^2\Bigr\}.
\]
In particular $c(q)=1-1/q$ for all $q\ge8$ (Theorem~\ref{thm:sep}).
\end{enumerate}
\end{theorem}

\begin{theorem}\label{thm:sep}
\lean{Z3Net.P_sep_rowScale, Z3Net.P_sep, Z3Net.P_sep_ge8, Z3Net.P_sep_real}\leanok
\uses{thm:main}
For $q\ge4$ and $k\ge2$, the toroidal minimum distance of $\mathcal P_{k,q}$ is at least $c(q)\,q^{-k}$, with $c(q)$ as in Theorem~\ref{thm:main}.
\end{theorem}

\begin{proof}
\leanok
\uses{prop:cell}
Work in cell units, with two points in cells $n$ and $n+e$ (toroidally). If some $|e_i|\ge2$, the $i$-th coordinates differ by at least $1\ge c(q)$; $e=0$ is impossible. So let $e\in\{-1,0,1\}^3\setminus\{0\}$, and let $\nu$ be the number of moving axes.

\emph{The step.} For a moving axis $d$, $u_d$ changes by $e_d$ modulo $q$, and $v_d$ changes by $e_d$ if $u_d$ wraps ($\omega_d=1$; this includes the step across the torus boundary, $n_d$ from $q^k-1$ to $0$ or back, where all cell digits wrap; since $k\ge2$, $v_d$ then still changes by exactly $e_d$) and is unchanged otherwise ($\omega_d=0$). Hence $r'-r\equiv\sigma\pmod q$ with
\begin{equation}\label{eq:sigma}
\sigma_i=-e_i-e_{i+1}+(1-\omega_{i+2})\,e_{i+2}\qquad(\text{indices mod }3,\ x\to y\to z),
\end{equation}
$|\sigma_i|\le3$, and $\Delta_i=e_i+((r'_i-r_i)+(f'_i-f_i))/q$, where $r'_i-r_i\in\{\sigma_i,\ \sigma_i-q\sgn\sigma_i\}$ ($\{0\}$ if $\sigma_i=0$), $|f'_i-f_i|<1$, and $f'_i=f_i$ if only axis $i$ moves (Proposition~\ref{prop:cell}(b)).

\emph{Per-coordinate bound.} Call a moving coordinate \emph{backward} if $e_i\sigma_i\le0$ and \emph{forward} otherwise. For a backward coordinate, $r'_i-r_i=\sigma_i$ gives $|\Delta_i|\ge1-(|\sigma_i|+1)/q$ and $r'_i-r_i=\sigma_i+qe_i$ gives $|\Delta_i|\ge2-(|\sigma_i|+1)/q$; as $q\ge4>|\sigma_i|$, both give $|\Delta_i|\ge c_{|\sigma_i|+1}$ with $c_j:=1-j/q\ge0$. Forward coordinates, and coordinates that do not move, are only bounded by $0$.

\emph{Signs.} From \eqref{eq:sigma}:
(i) If $e=\pm e_i$ then $\sigma_i=-e_i$ and there is no jitter, so $\Delta_i\in\{e_i(1-1/q),\ e_i(2-1/q)\}$ and $|\Delta|\ge c_1$.
(ii) $e_i\sigma_i=-1-e_ie_{i+1}+(1-\omega_{i+2})e_ie_{i+2}$ is positive iff $e_{i+1}=-e_i$, $e_{i+2}=e_i$ and $\omega_{i+2}=0$. So a forward coordinate exists only if $\nu=3$, and it is unique; then $\sigma_{i+1}=(1-\omega_i)e_i$ (backward, $|\sigma|\le1$) and $\sigma_{i+2}=-(3-\omega_{i+1})e_i$ (backward, $|\sigma|\le3$).
(iii) $|\sigma_i|=3$ forces $e_{i+1}=e_i$, $e_{i+2}=-e_i$, $\omega_{i+2}=0$, and then $\sigma_{i+2}=(1-\omega_{i+1})e_i$: backward with $|\sigma|\le1$.
(iv) If $\nu=2$ with moving axes $\{i,i+1\}$, both are backward by (ii); either $e_{i+1}=e_i$ and $|\sigma_{i+1}|\le1$, or $e_{i+1}=-e_i$ and $\sigma_i=0$; the other value is at most $2$ in absolute value.

\emph{Cases.} $\nu=1$: $|\Delta|\ge c_1$ by (i). $\nu=2$: $|\Delta|^2\ge c_2^2+c_3^2$ by (iv). $\nu=3$ with some backward $|\sigma_i|=3$: $|\Delta|^2\ge c_4^2+c_2^2$ by (iii). $\nu=3$ with a forward coordinate and all backward $|\sigma|\le2$: $|\Delta|^2\ge c_2^2+c_3^2$ by (ii). Otherwise three backward coordinates with $|\sigma|\le2$: $|\Delta|^2\ge3c_3^2$. As $c_3\ge c_4$, the minimum is $c(q)^2$. For $q\ge8$, $c_2^2+c_4^2\ge c_1^2$ (as $1-10x+19x^2\ge0$ for $x=1/q\le1/8$) and $3c_3^2\ge c_1^2$, so $c(q)=1-1/q$.
\end{proof}


\section{Resolution \texorpdfstring{$m=3$}{m=3}: column nets}\label{sec:m3}

\begin{definition}[Column nets]\label{def:col}
\lean{Z3NetC.colW_column, Z3NetC.col_column}\leanok
Let $q\ge5$ with $\gcd(q,6)=1$, let $w=(-2,1,1)$ if $q\equiv2\pmod3$ and $w=(-1,-1,2)$ if $q\equiv1\pmod3$, and let $a=\frac13(q\mathbf 1+w)$, an integer vector with $a_x+a_y+a_z=q$ ($w$ is forced modulo $3$, $w\equiv-q\mathbf 1$, by the integrality of $a$, and $w\cdot\mathbf 1=0$). The \emph{column net} $\mathcal C_q$ is the digital net over $\Z/q\Z$ whose point with index $n\in(\Z/q\Z)^3$ has, in every coordinate $c$, the three digits
\[
n_c,\qquad r(n)=a\cdot n,\qquad j(n)=n_y-n_x,
\]
i.e. the generator matrices with the rows $\e_c$, $a$ and $b:=(-1,1,0)$.
\end{definition}

\begin{theorem}[Minimum distance of column point sets]\label{thm:colmin}
\lean{Z3NetC.colmin, Z3NetC.colmin_real}\leanok
Let $q\ge2$, let $X$ be a column point set, and for a direction $i\in\{x,y,z\}$ and a cell $n$ let $\Delta_i([n])=\varphi([n+\e_i])-\varphi([n])$. Then the toroidal minimum distance $d$ of $X$ at integer scale $q^3$ satisfies
\[
d^2=\frac23q^4+3g^2,\qquad g=\min_{i,[n]}\operatorname{dist}\Bigl(\Delta_i([n]),\ \frac23q^2+q^2\Z\Bigr),
\]
and it is attained by two points in neighbouring columns $[n]$, $[n+\e_i]$. In particular $\delta N^{1/3}\ge\sqrt{2/3}$.
\end{theorem}

\begin{proof}
\leanok
Let $n\ne n'$ be cells and $D$ a representative of $X(n')-X(n)$ modulo $q^3$. Then $D=q^2E+\Delta\mathbf 1$ with an integer vector $E\equiv n'-n\pmod q$ and $\Delta=\varphi([n'])-\varphi([n])$; every such $E$ occurs, and the toroidal distance of the pair is the least $|D|$. With $\Sigma E=E_x+E_y+E_z$,
\[
|q^2E+\Delta\mathbf 1|^2=q^4|PE|^2+\tfrac13\bigl(q^2\Sigma E+3\Delta\bigr)^2,\qquad
|PE|^2=\tfrac13\bigl((E_x-E_y)^2+(E_y-E_z)^2+(E_z-E_x)^2\bigr).
\]
The three differences $E_x-E_y$, $E_y-E_z$, $E_z-E_x$ are integers with sum $0$. If they all vanish, $E=t\mathbf 1$ with $t\in\Z$; then $n'\equiv n+t\mathbf 1$ lies in the column of $n$, $\Delta=0$, $t\not\equiv0\pmod q$, and $|D|^2=3t^2q^4\ge3q^4$. If one of them has absolute value at least $2$, the other two sum to its negative, so the three squares sum to at least $6$ and $|D|^2\ge2q^4$. Otherwise they are $\pm(1,-1,0)$ up to order, i.e. $E=\pm\e_i+t\mathbf 1$ for some $i$ and $t\in\Z$. For $E=\e_i+t\mathbf 1$ we get $D=q^2\e_i+U\mathbf 1$ with $U=\Delta+tq^2$, and
\[
|q^2\e_i+U\mathbf 1|^2=(q^2+U)^2+2U^2=\tfrac23q^4+3\bigl(U+\tfrac13q^2\bigr)^2 .
\]
As $t$ runs over $\Z$, $U$ runs over $\Delta+q^2\Z$, so the least such value is $\frac23q^4+3\operatorname{dist}(\Delta,-\frac13q^2+q^2\Z)^2$, and $-\frac13q^2\equiv\frac23q^2\pmod{q^2}$; it is at most $\frac23q^4+3(\frac12q^2)^2<2q^4$. A pair with $E=-\e_i+t\mathbf 1$ is, read backwards, a pair with $E=\e_i-t\mathbf 1$. Every pair of neighbouring columns occurs (take any cell $n$ and $n'=n+\e_i$), so the minimum over all pairs is $\frac23q^4+3g^2$.
\end{proof}

\begin{theorem}[Resolution $m=3$]\label{thm:col}
\lean{Z3NetC.col_net, Z3NetC.col_covering, Z3NetC.col_mindist, Z3NetC.col_closest, Z3NetC.col_expansion, Z3NetC.col_sep_real, Z3NetC.col_mono, Z3NetC.col_values, Z3NetC.col_boxMinor_int}\leanok
\uses{def:col, thm:colmin}
Let $q\ge5$ with $\gcd(q,6)=1$, and let $\mathcal C_q$ be the column net of Definition~\ref{def:col}, with $N=q^3$ points.
\begin{enumerate}
\item $\mathcal C_q$ is a digital $(0,3,3)$-net over $\Z/q\Z$; up to sign, its box minors are $1$, $a_x$, $a_y$, $a_z$ and $a_x+a_y$ (the last two modulo $q$; as integers the minors of the boxes with depth $3$ on $x$, $y$ and $z$ are, up to sign, $a_z$, $a_z$ and $a_x+a_y$). Every cell of side $q^{-1}$ holds one point, so $h\le\sqrt3\,N^{-1/3}$.
\item Its toroidal minimum distance $d$ at integer scale $q^3$, and also its Euclidean one, is
\[
d^2=\frac23q^4+3g^2,\qquad 3g=\begin{cases}q^2-q-3&\text{if }q\equiv2\pmod3,\\ q^2-2q&\text{if }q\equiv1\pmod3,\end{cases}
\]
with $g$ as in Theorem~\ref{thm:colmin},
that is, $d^2=q^4-\frac23q^3-\frac53q^2+2q+3$ resp.\ $q^4-\frac43q^3+\frac43q^2$, and
\[
\delta N^{1/3}=1-\frac1{3q}-\frac8{9q^2}+O(q^{-3})\quad\text{resp.}\quad\delta N^{1/3}=1-\frac2{3q}+\frac4{9q^2}+O(q^{-3}).
\]
Two points are closest on the torus if and only if they differ by $\pm D$ modulo $q^3$, where $D=\frac13\bigl(-(2q^2-q-3),\,q^2+q+3,\,-(2q^2-q-3)\bigr)$ for $q\equiv2$ and $D=\frac13\bigl(-(2q^2-2q),\,-(2q^2-2q),\,q^2+2q\bigr)$ for $q\equiv1\pmod3$.
\item Within each residue class of $q$ modulo $3$, $\delta N^{1/3}$ increases with $q$. It is $0.9059\ldots$ at $q=5$ and $0.9147\ldots$ at $q=7$, so $\delta N^{1/3}>0.9059$ for every such $q$, and $\delta N^{1/3}>0.9823$ for every such $q\ge37$.
\end{enumerate}
\end{theorem}

\begin{proof}
\leanok
\uses{lem:nie, lem:box, thm:colmin}
(1) The first rows are $\e_x,\e_y,\e_z$, so $\mathcal C_q$ is a graph form with the depths $(1,1,1)$, and we use Lemma~\ref{lem:box}. The box of shape $(1,1,1)$ has determinant $1$ and puts one point in every cell. A box with depth $2$ on $c$, $1$ on $c'$ and $0$ on $c''$ has, up to sign, the determinant $a_{c''}$; as $3a_{c''}=q+w_{c''}\equiv w_{c''}\in\{\pm1,\pm2\}$ and $\gcd(q,6)=1$, it is a unit. A box with depth $3$ on $c$ has, up to sign, the $2\times2$ minor of the rows $a$ and $b$ on the two columns other than $c$. Both rows sum to $0$ modulo $q$, so in each of them column $c$ is minus the sum of the other two modulo $q$, and modulo $q$ all three such minors equal $\pm\mu$ with $\mu=a_xb_y-a_yb_x=a_x+a_y$ (as integers they are $a_x+a_y$ and $\pm a_z$); and $3\mu=2q+w_x+w_y\equiv-w_z\in\{-1,-2\}\pmod q$, a unit. Lemma~\ref{lem:nie} concludes.

(2) We apply Theorem~\ref{thm:colmin}. Let $\bar b_i\in\{0,\dots,q-1\}$ be the residue of $b\cdot\e_i$, so $\bar b=(q-1,1,0)$. A step from the cell $n$ to $n+\e_i$ changes $r$ by $a_i$ and $j$ by $\bar b_i$ modulo $q$, so
\[
\Delta_i([n])=q\bigl(a_i-q\,[r(n)+a_i\ge q]\bigr)+\bigl(\bar b_i-q\,[j(n)+\bar b_i\ge q]\bigr).
\]
The map $[n]\mapsto(r(n),j(n))$ is a bijection from $\mathbb T$ onto $(\Z/q\Z)^2$, because $\mu$ is a unit; so the two wraps occur independently, in all four combinations (two if $\bar b_i=0$). Modulo $q^2$ the values of $\Delta_i$ are therefore $s_i$ and, if $\bar b_i\ne0$, $s_i-q$, where $s_i=qa_i+\bar b_i$, and $3\operatorname{dist}(\Delta_i,-\frac13q^2+q^2\Z)$ is the least value of $|3U+q^2|$ over $U\in s_i+\{0,-q\}+q^2\Z$. Here $3s_i+q^2$ lies between $2q^2-q+3$ and $2q^2+2q$, so the least value is attained at $U=s_i-q^2$ or $U=s_i-q^2-q$, where $3U+q^2=-(2q^2-3s_i)$ resp.\ $-(2q^2-3s_i+3q)$; every other $U$ gives at least $2q^2-4q+3>q^2+q$. The values $2q^2-3s_i$ for $i=x,y,z$ are
\[
q^2-q+3,\quad q^2-q-3,\quad q^2-q\quad(q\equiv2),\qquad\qquad q^2-2q+3,\quad q^2+q-3,\quad q^2-2q\quad(q\equiv1),
\]
all positive, and adding $3q$ only increases them. So $3g$ is the least of them, $q^2-q-3$ resp.\ $q^2-2q$, attained only at $i=y$ (resp.\ $i=z$) with $U=s_i-q^2$, and Theorem~\ref{thm:colmin} gives $d^2=\frac23q^4+3g^2$. The closest differences are then $\pm(q^2\e_i+U\mathbf 1)=\pm D$ modulo $q^3$. The Euclidean minimum distance is at least $d$, and it is attained without a wrap: $\varphi([\mathbf 1])=0$ and $\varphi([\e_i])=qa_i+\bar b_i=s_i$, so the points of the cells $\mathbf 1$ and $\e_i$ differ by $q^2(\e_i-\mathbf 1)+s_i\mathbf 1=q^2\e_i+(s_i-q^2)\mathbf 1=D$ exactly. The expansions follow from $(\delta N^{1/3})^2=\frac23+\frac13(3g/q^2)^2$.

(3) $3g/q^2=1-\frac1q-\frac3{q^2}$ resp.\ $1-\frac2q$ is positive and increases with $q$. The values at $q=5$ and $7$ are $\sqrt{513}/25$ and $\sqrt{2009}/49$, and at $q=37$ (the least $q\ge37$ coprime to $6$, in the class $q\equiv1$; the least in the other class is $41$, with $0.9913\ldots$) it is $0.9823\ldots$.
\end{proof}

\begin{lemma}[Shared digits]\label{lem:shared}
\lean{Z3NetC.shared_count, Z3NetC.shared_fun}\leanok
Let $P$ be a $(0,m,s)$-net in base $q$, digital or not, $j\ge0$, and $a\ne b$ two coordinates. If $2(j+1)\le m$, then the pair of the digits with index $j$ of $a$ and of $b$ takes every value in $\{0,\dots,q-1\}^2$ at exactly $q^{m-2}$ points of $P$. In particular, if $f$ is a map of digits such that for every point of $P$ the digit with index $j$ of coordinate $b$ equals $f$ of the digit with index $j$ of coordinate $a$, then $2(j+1)>m$.
\end{lemma}

\begin{proof}
\leanok
Let $2(j+1)\le m$ and choose the depths $d_a=j+1$, $d_b=m-j-1\ge j+1$ and $d_c=0$ for the other coordinates. Prescribing the digits $0,\dots,j$ of $a$ and of $b$ fixes the coordinate $a$ of an elementary box of this shape and the first $j+1$ digits of its coordinate $b$, and leaves $q^{m-2(j+1)}$ choices for the other digits of $b$; by $t=0$ each such box holds exactly one point. So each prescription is met by $q^{m-2j-2}$ points, and each value of the pair of digits with index $j$ by $q^{2j}\cdot q^{m-2j-2}=q^{m-2}$ points. A map $f$ as stated would leave the values $(\alpha,\beta)$, $\beta\ne f(\alpha)$, unmet.
\end{proof}

\begin{proposition}[Arbitrary third digits]\label{prop:colfree}
\lean{Z3NetC.colfree}\leanok
Let $q\ge2$, let $w\in\Z^3$ with $w_x+w_y+w_z=0$ and $w_c\equiv-q\pmod3$, $a=\frac13(q\mathbf 1+w)$, and $\beta\in\Z$. Let $P$ be a set of $q^3$ points of the torus $(\Z/q^3\Z)^3$, one for each cell $n\in\{0,\dots,q-1\}^3$, at
\[
X(n)=q^2n+q\,r(n)\mathbf 1+j(n),\qquad r(n)\in[0,q),\ \ r(n)\equiv a\cdot n+\beta\pmod q,\qquad j(n)\in[0,q)^3
\]
(integers; $j$ arbitrary). Then any two distinct points of $P$ have toroidal distance at least $q^2-\bigl((\|w\|_1+3)q-3\bigr)/\sqrt3$.
\end{proposition}

\begin{proof}
\leanok
Let $n\ne n'$ and let $D=q^2E+q\,\Delta r\,\mathbf 1+\Delta j$ be a representative of $X(n')-X(n)$ modulo $q^3$, with an integer vector $E\equiv n'-n\pmod q$, $\Delta r\in(-q,q)$ and $\Delta j\in(-q,q)^3$. If some $|E_c|\ge2$, then $|D_c|\ge2q^2-q(q-1)-(q-1)>q^2$. So let $E\in\{-1,0,1\}^3\setminus\{0\}$ and $\Sigma E=E_x+E_y+E_z$. The constant $\beta$ cancels: $\Delta r\equiv a\cdot E\pmod q$, and $3a\cdot E=q\Sigma E+w\cdot E$, so $3\Delta r=w\cdot E+zq$ with an integer $z\equiv\Sigma E\pmod3$. Hence $v=E-\frac{z+2\Sigma E}3\mathbf 1$ is an integer vector with $Hv=E+\frac z3\mathbf 1$, and
\[
3D=3q^2E+q\,(w\cdot E+zq)\mathbf 1+3\Delta j=3q^2Hv+R,\qquad R=q\,(w\cdot E)\mathbf 1+3\Delta j,
\]
with $|R_c|\le q\|w\|_1+3(q-1)$, since $|w\cdot E|\le\|w\|_1$. If $v\ne0$, then $|D|\ge q^2|Hv|-|R|/3\ge q^2-\bigl((\|w\|_1+3)q-3\bigr)/\sqrt3$, as $H$ is orthogonal. If $v=0$, then $E=t\mathbf 1$ with $3t=z+2\Sigma E=z+6t$, so $z=-3t$ and, as $w\cdot\mathbf 1=0$, $\Delta r=-tq$; then $|\Delta r|<q$ forces $t=0$, i.e. $E=0$, a contradiction.
\end{proof}

\begin{proposition}[Every digital shift at $m=3$]\label{prop:colshift}
\lean{Z3NetC.colshift_eq, Z3NetC.colshift_agree2, Z3NetC.colshift_lower, Z3NetC.colshift_upper, Z3NetC.colshift_eq_eucl}\leanok
\uses{def:dshift, thm:col}
Let $q\ge5$ with $\gcd(q,6)=1$, let $\tau$ be a digital shift of $\mathcal C_q$ (Definition~\ref{def:dshift}) with the digits $\tau_{c,1},\tau_{c,2},\tau_{c,3}$ in coordinate $c$ (the first, second and third digit), and put $W=2$ if $q\equiv2$ and $W=4$ if $q\equiv1\pmod3$ ($W/(3q)$ is the largest distance of $a_i/q$ or $2a_i/q$ from $\frac13$ or $\frac23$ that occurs; see the proof), and $a_{\max}=\max_ca_c$. The toroidal minimum distance of $\mathcal C_q\oplus\tau$ satisfies:
\begin{enumerate}
\item if $\tau_{x,2}=\tau_{y,2}=\tau_{z,2}$ and $\tau_{x,3}=\tau_{y,3}=\tau_{z,3}$, it equals that of $\mathcal C_q$ (Theorem~\ref{thm:col}(2));
\item if $\tau_{x,2}=\tau_{y,2}=\tau_{z,2}$, then $\delta N^{1/3}\ge1-(7q-3)/(\sqrt3q^2)$;
\item for every $\tau$,
\[
\delta N^{1/3}\ \ge\ \frac1{\sqrt3}\Bigl(1-\frac{W+3}q+\frac3{q^2}\Bigr),
\]
which exceeds $0.4693$ for every $q\ge37$; and if $\tau_{x,2},\tau_{y,2},\tau_{z,2}$ are not all equal, then $\delta N^{1/3}\le\sqrt3\,(q\,a_{\max}+q-1)/q^2=\frac1{\sqrt3}\bigl(1+\frac{3+w_{\max}}q-\frac3{q^2}\bigr)$, $w_{\max}=\max_cw_c$.
\end{enumerate}
The lower bounds of (2) and (3), and the value of (1) as a lower bound, hold for the Euclidean minimum distance too, which is at least the toroidal one; so does the upper bound of (3), which is witnessed by a pair that does not wrap.
\end{proposition}

\begin{proof}
\leanok
\uses{thm:colmin, thm:col, prop:colfree, lem:shift}
\emph{Shifted digits.} By Lemma~\ref{lem:shift}, after re-indexing by the cells, the point of cell $n$ of $\mathcal C_q\oplus\tau$ has, in coordinate $c$, the digits $n_c$, $r_c(n)\equiv a\cdot n+\beta_c$ and $j_c(n)\equiv b\cdot n+\beta'_c$ with $\beta_c=\tau_{c,2}-a\cdot\tau_1$ and $\beta'_c=\tau_{c,3}-b\cdot\tau_1$, $\tau_1=(\tau_{x,1},\tau_{y,1},\tau_{z,1})$. So the $\beta_c$ are congruent if and only if the $\tau_{c,2}$ are equal, and likewise for $\beta'_c$ and $\tau_{c,3}$.

(1) Then $X(n)=q^2n+\varphi'([n])\mathbf 1$ with $\varphi'=q\,r+j$, $r=(a\cdot n+\beta)\bmod q$, $j=(b\cdot n+\beta')\bmod q$: a column point set. The map $[n]\mapsto(r,j)$ is still a bijection onto $(\Z/q\Z)^2$, so in the proof of Theorem~\ref{thm:col}(2) the same values $\Delta_i$ occur, and Theorem~\ref{thm:colmin} gives the same $d$.

(2) Proposition~\ref{prop:colfree} with $\|w\|_1=4$.

(3) \emph{Lower bound.} Let $n\ne n'$ and $D=q^2E+q\,\Delta r+\Delta j$ a representative, now with $\Delta r=(\Delta r_c)_c\in(-q,q)^3$. If some $|E_c|\ge2$, then $|D_c|>q^2$. So let $E\in\{-1,0,1\}^3\setminus\{0\}$, and let $\vartheta\in[0,q)$ be the residue of $a\cdot E$. The constants $\beta_c$ cancel: $\Delta r_c\equiv a\cdot E$ and $|\Delta r_c|<q$, so $\Delta r_c=\vartheta-q\omega_c$ with $\omega_c\in\{0,1\}$ ($\omega=0$ if $\vartheta=0$). Hence $D=q^2u+\Delta j$ with $u=E-\omega+\theta\mathbf 1$, $\theta=\vartheta/q$, and $|D|\ge q^2|u|-\sqrt3(q-1)$. Every $u_c$ lies in $\theta+\Z$, so $|u_c|\ge\|\theta\|$, the distance from $\theta$ to the nearest integer, and $\theta\equiv(a\cdot E)/q\pmod1$. If $\Sigma E=0$, then $E=\pm(\e_c-\e_d)$, the two moving coordinates of $u$ are $\pm(\theta+\{0,1\})$ and $\pm(\theta-\{1,2\})$, and $|u|^2\ge\theta^2+(1-\theta)^2\ge\frac12$. If $E=\pm\mathbf 1$, then $\vartheta=0$ and $|u|=\sqrt3$. Otherwise, up to sign, $E=\e_i$ or $E=\mathbf 1-2\e_i$ or $E=\mathbf 1-\e_i$, with $a\cdot E\equiv a_i$, $-2a_i$, $-a_i$ respectively. As $3a_i=q+w_i$ with $0<a_i<\frac q2$, $\|a_i/q\|=\frac{q+w_i}{3q}$ and $\|2a_i/q\|=\min\bigl(\frac{q-2w_i}{3q},\frac{2q+2w_i}{3q}\bigr)$; over the entries $w_i\in\{-2,1\}$ resp.\ $\{-1,2\}$ the least of these is $\frac{q-W}{3q}$. So $|u|\ge\sqrt3\,\|\theta\|\ge(q-W)/(\sqrt3q)$ in these cases, which is also less than $\frac1{\sqrt2}$ and $\sqrt3$. Therefore $|D|\ge q(q-W)/\sqrt3-\sqrt3(q-1)=\bigl(q^2-(W+3)q+3\bigr)/\sqrt3$ for every pair, which is the bound. It increases with $q$ within each residue class; for $q\ge37$ its least value is $0.4693\ldots$ at $q=37$ ($W=4$), and $0.5079\ldots$ at $q=41$ in the class $W=2$.

\emph{Upper bound.} Put $b_c\equiv-\beta_c\pmod q$; they are not all equal. Along $\e_c$ the shared digit changes by $a_c$, and $r_d$ wraps if and only if $\vartheta=a\cdot n\bmod q$ lies in the cyclic interval $J^{(c)}_d=\{b_d-a_c,\dots,b_d-1\}$ of $a_c$ consecutive residues of $\Z/q\Z$. We find $c$ and a residue $\vartheta$ that lies in $J^{(c)}_c$ and in no $J^{(c)}_d$, $d\ne c$. If the $b_d$ are distinct, let $g_c\ge1$ be the distance from $b_c$ forward around $\Z/q\Z$ to the next of them. The gaps sum to $q=a_x+a_y+a_z$, so $g_c\ge a_c$ for some $c$; then $\vartheta=b_c-1$ lies in $J^{(c)}_c$, and it lies in $J^{(c)}_d$ only if $b_d$ is at forward distance at most $a_c-1$ from $b_c$, which no $b_d$ is. If exactly two of them are equal, $b_d=b_{d'}\ne b_c$, then $J^{(c)}_c$ and $J^{(c)}_d=J^{(c)}_{d'}$ are two different cyclic intervals of the same length $a_c<q$, so neither contains the other. Take a cell $n$ with $n_c=0$ and $a\cdot n\equiv\vartheta$ (possible, as the entries $a_d$, $d\ne c$, are units) and $n'=n+\e_c$, so that the pair does not wrap in the cell digits. Then $\Delta r_c=a_c-q$ and $\Delta r_d=a_c$ for $d\ne c$, so the representative with $E=\e_c$ is $D=q^2\e_c+q(a_c\mathbf 1-q\e_c)+\Delta j=q\,a_c\mathbf 1+\Delta j$, and $|D|\le\sqrt3\,(q\,a_{\max}+q-1)$.
\end{proof}


\section{\texorpdfstring{$m\equiv1,2$}{m = 1, 2}: lattice skeletons}\label{sec:allm}

\begin{definition}\label{def:B}
\lean{Z3NetKN.B2gen_net, Z3NetKN.B1gen_net, Z3NetKN.B2_net, Z3NetKN.B2_sep, Z3NetKN.B1_net, Z3NetKN.B1_sep}\leanok
Let $\rho$ be an integer with $\rho\not\equiv0\pmod p$.
\begin{enumerate}
\item[($\mathcal B^2$)] $\mathbf k=(k+1,k+1,k)$; the recipe of Proposition~\ref{prop:skelcomp}(ii) with $\kappa=(\rho,\rho^2)$. In column block $x$: $Z0|x_k\ne0$ (box $(k,k+1,k+1)$); subtract multiples of level $x_k$ from every level $x_l$, $l<k$, so that $Z0|x_l=0$, then scale level $x_k$ so that $Z0|x_k=\rho$. In column block $y$ the same, with pivot $Z0|y_k$ (box $(k+1,k,k+1)$) and target $\rho^2$. Set all diagonal blocks to $0$. The result is the digital net $\mathcal B^2_{k,p,\rho}$ over $\F_p$; its skeleton digit is
\[
S_z\equiv\rho\,x_k+\rho^2\,y_k\pmod p.
\]
\item[($\mathcal B^1$)] $\mathbf k=(k+1,k,k)$. The goal is a graph form whose skeleton digits are $S_y\equiv\rho x_k+z_{k-1}$ and $S_z\equiv\rho^2x_k+x_{k-1}+y_{k-1}-z_{k-1}\pmod p$, with every box minor a unit; the steps are as follows. Column block $x$: the matrix $\left(\begin{smallmatrix}Y0|x_k&Y0|x_{k-1}\\ Z0|x_k&Z0|x_{k-1}\end{smallmatrix}\right)$ is invertible (box $(k-1,k+1,k+1)$); call its determinant $\Delta$; use the levels $x_k$, $x_{k-1}$ to make $Y0$ and $Z0$ vanish on the levels $l<k-1$; then add a multiple of $x_k$ to $x_{k-1}$ so that $Y0|x_{k-1}=0$. Now $Y0|x_k\ne0$ (box $(k,k+1,k)$), $Z0|x_k\ne0$ (box $(k,k,k+1)$) and $Z0|x_{k-1}=\Delta/Y0|x_k\ne0$. Scale the block $\Phi_{zx}$ so that $Z0|x_{k-1}=1$, the level $x_k$ so that $Z0|x_k=\rho^2$, and the block $\Phi_{yx}$ so that $Y0|x_k=\rho$. Column block $y$: with the pivot $Z0|y_{k-1}\ne0$ (box $(k+1,k-1,k+1)$) make $Z0|y_l=0$ for $l<k-1$ and scale $\Phi_{zy}$ so that $Z0|y_{k-1}=1$. Column block $z$: the same for $Y0|z$ (box $(k+1,k+1,k-1)$). Diagonal blocks: $Z0|z_{k-1}=-1$ and all other entries $0$. The result is $\mathcal B^1_{k,p,\rho}$; its skeleton digits are
\[
S_y\equiv\rho\,x_k+z_{k-1},\qquad S_z\equiv\rho^2\,x_k+x_{k-1}+y_{k-1}-z_{k-1}\pmod p.
\]
\end{enumerate}
\end{definition}

\begin{proposition}\label{prop:Bnet}
\lean{Z3NetKN.B2gen_net, Z3NetKN.B1gen_net, Z3NetKN.B2_net, Z3NetKN.B2_sep, Z3NetKN.B1_net, Z3NetKN.B1_sep}\leanok
For every prime $p$, every $k\ge1$ and every $\rho\not\equiv0\pmod p$, $\mathcal B^2_{k,p,\rho}$ is a $(0,3k+2,3)$-net and $\mathcal B^1_{k,p,\rho}$ is a $(0,3k+1,3)$-net over $\F_p$.
\end{proposition}

\begin{proof}
\leanok
\uses{lem:nie, lem:box}
Every box minor of $\Phi^F$ is $\pm1$, and each move multiplies every box minor by a unit of $\F_p$ (Lemma~\ref{lem:box}; the targets $\rho$, $\rho^2$, $1$ and the pivots are units). Lemma~\ref{lem:nie} completes the proof.
\end{proof}

\begin{remark}[Dimension three, $m\equiv2\pmod3$]\label{rem:B2}
\lean{Z3NetKN.B2_net, Z3NetKN.B2_sep, Z3NetKN.B2_skel_iff, Z3NetKN.B2_rem_sep, Z3NetKN.B2_rem_sep_real, Z3NetKN.B2_relax, Z3NetKN.B2_const}\leanok
\uses{lem:skelred, prop:skelcomp, thm:S1, def:S1, def:B, prop:Bnet}
For $s=3$ the recipes of Definitions~\ref{def:S1} and~\ref{def:B}($\mathcal B^2$) coincide (Proposition~\ref{prop:skelcomp}(ii) with $\kappa=(\rho,\rho^2)$, skeleton digit $S_z=\rho x_k+\rho^2y_k$), and Theorem~\ref{thm:S1} uses only the skeleton digit. So for every prime $p$, every $k\ge1$ and every $\rho\ge1$ with $\rho^3<p$, $\mathcal B^2_{k,p,\rho}$ is a $(0,3k+2,3)$-net over $\F_p$ (Proposition~\ref{prop:Bnet}) with toroidal $\delta N^{1/3}\ge(\rho-1)\,p^{-1/3}$ (Theorem~\ref{thm:S1}(1) for $\rho\ge2$; the bound is $0$ for $\rho=1$); its occupied fine cells are exactly the cells of the index-$p$ lattice $\Lambda_{\mathcal B^2}=\{a:a_z\equiv\rho a_x+\rho^2a_y\pmod p\}$, of minimum $\rho$ (Theorem~\ref{thm:S1}(2), Figure~\texttt{fig:Blattice}(a),~(b)); and for $\rho=\lfloor p^{1/3}\rfloor$ its mesh ratio is bounded uniformly in $p$ (Theorem~\ref{thm:S1}(3)). With $\rho$ maximal, $p<(\rho+1)^3$, so the bound exceeds $(\rho-1)/(\rho+1)$ and tends to $1$. It is a sawtooth in $p$ (it drops each time $\lfloor p^{1/3}\rfloor$ increases), first exceeds $0.9$ at $p=1361$, and stays above $0.9$ for every prime $p\ge6863$ ($\rho\ge19$). The bound is sharp for the method: the bound of Lemma~\ref{lem:skelred}(i), which allows a jitter of one fine cell in every coordinate, equals $(\rho-1)p^{-1/3}$ exactly, attained at the step $(-\rho,1,0)$.
\end{remark}

\begin{theorem}[Dimension three, $m\equiv1\pmod3$]\label{thm:allm}
\lean{Z3NetKN.B2_net, Z3NetKN.B2_sep, Z3NetKN.B1_net, Z3NetKN.B1_sep, Z3NetKN.B1_rho_spec, Z3NetKN.B1_sep_real, Z3NetKN.B1_tendsto, Z3NetKN.B1_rho27, Z3NetKN.B1_gt09, Z3NetKN.B1_g_lt, Z3NetKN.gB1_strictMonoOn, Z3NetKN.gB1_tendsto, Z3NetKN.gB1_27}\leanok
\uses{def:B}
Let $p$ be a prime, $k\ge1$, and let $\rho\ge2$ with $\rho^2(\rho+1)\le p$. The digital net $\mathcal B^1_{k,p,\rho}$ of Definition~\ref{def:B} is a $(0,3k+1,3)$-net over $\F_p$ whose toroidal minimum distance satisfies $\delta N^{1/3}\ge(\rho^2-4)\,p^{-2/3}$. With $\rho$ the largest integer such that $\rho^2(\rho+1)\le p$, the bound tends to $1$ as $p\to\infty$; it exceeds $0.9$ once $\rho\ge27$, i.e.\ $p\ge27^2\cdot28=20412$, so for every prime $p\ge20431$.
\end{theorem}

\begin{proof}[Proof of Theorem~\ref{thm:allm}]
\leanok
\uses{lem:skelred, lem:skellat, prop:Bnet}
The net property is Proposition~\ref{prop:Bnet}; the hypotheses give $2\le\rho<p$, so $\rho\not\equiv0$. Let $v\ne w$ be two points, $a\ne0$ the centred difference of their fine positions, and $M=|a|_\infty$. If $M\ge p-1$ we are done, since $p-1\ge\rho^2-3$. Otherwise let $\xi_c=\xi_{c,1}$ be the floor changes of Lemma~\ref{lem:skelred}(ii) and $u=(a_x,\ a_y-\xi_z,\ a_z-\xi_x-\xi_y+\xi_z)$. By the displayed relations, $u\in\Lambda_{\mathcal B^1}$, and $|u|_\infty\le M+3$. If $u\ne0$, Lemma~\ref{lem:skellat}(L1) gives $M+3\ge\rho^2$. If $u=0$, then $a_x=0$, hence $\xi_x=0$, $a_y=\xi_z$ and $a_z=\xi_y-\xi_z$. If $\xi_z=1$ then $a_z>0$, so $\xi_y\ge2$, impossible; if $\xi_z=-1$ then $a_z<0$ and $\xi_y\le-2$, impossible. So $\xi_z=0$, $a_y=0$, $\xi_y=0$ and $a_z=0$, contradicting $a\ne0$. Hence $M\ge\rho^2-3$, and Lemma~\ref{lem:skelred}(i) with $\ell=2$ gives $\delta N^{1/3}\ge(\rho^2-4)p^{-2/3}$.

\emph{Constants.} Let $c_1(p)=(\rho^2-4)p^{-2/3}$ be the bound with $\rho$ maximal. Then $p<(\rho+1)^2(\rho+2)$, hence
\[
c_1(p)^3>G(\rho)=\frac{(\rho-2)^3(\rho+2)}{(\rho+1)^4}.
\]
$G$ increases for $\rho\ge3$, since $(\log G)'=9/((\rho-2)(\rho+1))-1/((\rho+1)(\rho+2))$ is positive, and tends to $1$. Since $\rho\ge27$ iff $p\ge20412$, and $G(27)^{1/3}=0.9033\ldots$, the bound exceeds $0.9$ from this prime on.
\end{proof}


\section{Every resolution}\label{sec:everym}

\begin{corollary}[Every prime $p\ge37$, every $m\ge3$]\label{cor:allm}
\lean{Z3NetKN.B2_net, Z3NetKN.B2_sep, Z3NetKN.cor_allm, Z3NetKN.bndDim3_gt, Z3NetKN.bndDim3_tendsto, Z3NetKN.cor_allm_mesh, Z3NetKN.cor_allm_mesh2, Z3NetKN.cor_allm_consts}\leanok
\uses{thm:S1, thm:main, thm:col, rem:B2, thm:allm}
For every prime $p\ge37$ and every $m\ge3$ the nets of Theorems~\ref{thm:col} ($m=3$), \ref{thm:main} ($m\equiv0$, $m\ge6$), \ref{thm:allm} ($m\equiv1$) and~\ref{thm:S1} with $s=3$ ($m\equiv2$; the nets $\mathcal B^2$ of Remark~\ref{rem:B2}), with $\rho$ maximal, are explicit digital $(0,m,3)$-nets over $\F_p$ with
\[
\delta N^{1/3}>0.2715,\quad\text{and}\quad\delta N^{1/3}>0.4250\ \text{if } p\ge83,
\]
uniformly in $m$, and the lower bound tends to $1$ as $p\to\infty$. Their covering radius satisfies $hN^{1/3}\le\sqrt3$ for $m\equiv0$, $hN^{1/3}\le p^{1/3}\sqrt{2+p^{-2}}$ for $m\equiv1$, and $hN^{1/3}\le\min\{p^{2/3}\sqrt{1+2p^{-2}},\,4.264\}$ for $m\equiv2\pmod 3$. So for each such $p$ the mesh ratio $h/\delta$ is bounded uniformly in $m$: the nets are quasi-uniform along every $N=p^m$, $m\ge3$. For $m\equiv2$ the mesh ratio is moreover at most $7\sqrt3<12.2$ for every prime $p\ge11$ and every $m\ge5$ (Theorem~\ref{thm:S1}(3)).
\end{corollary}

\begin{proof}[Proof of Corollary~\ref{cor:allm}]
\leanok
\uses{thm:S1, thm:col, rem:B2}
Take $\mathcal C_p$ for $m=3$ (Theorem~\ref{thm:col}; for $p\ge37$, $\delta N^{1/3}>0.9823$, and its covering radius is at most $\sqrt3N^{-1/3}$ because every cell holds a point), $\mathcal P_{k,p}$ for $m=3k\ge6$, $\mathcal B^1_{k,p,\rho}$ for $m=3k+1\ge4$ and $\mathcal B^2_{k,p,\rho}$ for $m=3k+2\ge5$, with $\rho$ maximal: $\rho_1$ and $\rho_2$, the largest integers with $\rho_1^2(\rho_1+1)\le p$ and $\rho_2^3<p$. For $p\ge37$, $\rho_1\ge3$ and $\rho_2\ge3$. If $37\le p\le79$ then $\rho_1=3$ and the bound of $\mathcal B^1$ is $5p^{-2/3}\ge5\cdot79^{-2/3}=0.27157\ldots$, the bound of $\mathcal B^2$ is $\ge2\cdot61^{-1/3}>0.508$, and $1-1/p>0.97$. If $p\ge83$ then $\rho_1\ge4$ and $\rho_2\ge4$, and the bounds exceed $G(4)^{1/3}=(48/625)^{1/3}>0.4250$ and $(\rho-1)/(\rho+1)=3/5$ at $\rho=4$ (Remark~\ref{rem:B2}). All bounds tend to $1$. For the covering radius: by $t=0$ the elementary box with depths $\mathbf k$ (resp. $(k,k,k)$) around any point of the torus contains a net point, and its diagonal is $p^{-k}\sqrt{2+p^{-2}}=p^{1/3}\sqrt{2+p^{-2}}\,N^{-1/3}$ for $m\equiv1$ and $p^{-k}\sqrt{1+2p^{-2}}=p^{2/3}\sqrt{1+2p^{-2}}\,N^{-1/3}$ for $m\equiv2$. For $m\equiv2$, Theorem~\ref{thm:S1}(3) with $s=3$ gives $hN^{1/3}\le\sqrt3(\mu+1)p^{-1/3}$ and $h/\delta\le\sqrt3(\mu+1)/(\rho-1)$, $\mu=\lceil p/\rho^2\rceil$; over all primes $p\ge11$ the maxima are $4.2633\ldots<4.264$ and $7\sqrt3$, both at $p=23$ (for $\rho\ge4$ use $\mu<\rho(1+1/\rho)^3+1$).
\end{proof}

\begin{remark}[Per-prime designs]\label{rem:perprime}
\lean{Z3NetKN.KN1_net, Z3NetKN.KN1_sep, Z3NetKN.KN2_net, Z3NetKN.KN2_sep}\leanok
\uses{lem:box, lem:Ss, lem:ppinv, prop:pp2, prop:pp1, prop:ppsep, rem:ppshift}
For small primes, computer-found designs give larger constants. They are graph forms with the depths $\mathbf k$, built by the same moves from the frame, in which the skeleton digits and the next offset digit of each coordinate (which controls the jitter) are prescribed on the last three or four levels; their skeletons are index-$p^2$ Korobov-type ($m\equiv1$) and index-$p$ congruence lattices ($m\equiv2$). A finite computation analogous to Lemma~\ref{lem:Ss} (it enumerates the steps between nearby fine cells, all carry patterns and both representatives of each digit difference) certifies the constants of Table~\texttt{tab:perprime} for every $m=3k+1\ge10$ resp. $m=3k+2\ge8$, since the prescribed rows do not depend on $k$; they exceed $1$ for $p=47$, $53$, $59$, $61$, $71$, $73$, $79$ ($m\equiv1$) and $p=41$, $43$, $59$, $61$, $79$ ($m\equiv2$). This separation bound is Proposition~\ref{prop:ppsep}, and it holds for every digital shift as well (Remark~\ref{rem:ppshift}). The existence of such a $(0,m,3)$-net for every $k$ (Propositions~\ref{prop:pp2} and~\ref{prop:pp1}, Appendix~\ref{app:perprime}) rests on quantities that the moves of Lemma~\ref{lem:box} do not change and that do not depend on $k$ (Lemma~\ref{lem:ppinv}: five for $m\equiv1$, two for $m\equiv2$), and on finitely many exact checks per design (proofs in Appendix~\ref{app:perprime}; formalization status in Section~\texttt{sec:formal}).
\end{remark}


\chapter{Bases two and three: the deep-shift lift}\label{sec:base23}

\begin{remark}[Why $s=3$, and base two]\label{rem:base2}
\lean{Z3Net.base2_net}\leanok
\uses{lem:Ts}
A $(0,m,s)$-net in base $2$ with $m\ge2$ requires $s\le3$ \cite[Cor.~5.11]{Nie87}; for digital nets over rings see \cite[Thm.~2]{LNS96}. So $s=3$ is the largest dimension in which one set of integer generator matrices can give $t=0$ in every base (for dimension $s$ see the text after Lemma~\ref{lem:Ts}). In the frame language of Section~\ref{sec:frames}: $t=0$ in every base needs places whose pairwise differences are $\pm1$; $\{0,1\}$ is the largest set of integers with this property (in the language of exceptional units \cite{Len77}), and the place $\infty$ gives the third coordinate.
\end{remark}

\begin{definition}[Deep-shift lift]\label{def:lift}
\lean{Z3NetB23.Zsub, Z3NetB23.pairSub, Z3NetB23.Admissible, Z3NetB23.VanishTop, Z3NetB23.liftGen, Z3NetB23.liftShift, Z3NetB23.LiftChain}\leanok
Let $P$ be as above, $0\le D$, and let $s_A,s_B\in\F_b^{3m_0}$ vanish on the first $D$ digits of each coordinate, i.e.\ change only the \emph{deep} digits, those of index $>D$. They are \emph{separating} if, for every $1\le e\le m_0$,
\[
\lambda^{yz}_e\bigl(s_A^{yz}\bigr)\ne0,\qquad \lambda^{xy}_e\bigl(s_B^{xy}\bigr)\ne0,\qquad \lambda^{xz}_e\bigl((s_B-s_A)^{xz}\bigr)\ne0,
\]
where $s^{cc'}$ is the restriction of $s$ to the coordinates $c,c'$. The \emph{lift} $L_D(P)$ is the set of $b^{m_0+3}$ points with index $(n,\alpha,\beta_1,\beta_2)\in\F_b^{m_0+3}$ whose coordinates have the $m_0+3$ digits
\[
\begin{aligned}
x&:\ \alpha\ \bigm|\ (\xi(n)+s_\beta)_x\ \bigm|\ \beta_2,\ \beta_1,\\
y&:\ \alpha+\beta_1\ \bigm|\ (\xi(n)+s_\beta)_y\ \bigm|\ \beta_1+\beta_2,\ \beta_1,\\
z&:\ \alpha+\beta_1+\beta_2\ \bigm|\ (\xi(n)+s_\beta)_z\ \bigm|\ \beta_1,\ \beta_2,
\end{aligned}
\qquad s_\beta=\beta_1s_A+\beta_2s_B,
\]
one \emph{top} digit, $m_0$ \emph{block} digits and two \emph{jitter} digits. The shift $\tau$ of $P$ stays in the block digits, so $L_D(P)$ is again a digitally shifted digital net over $\F_b$.
\end{definition}

\begin{lemma}[Lift]\label{lem:lift}
\lean{Z3NetB23.lift_net, Z3NetB23.lift_exists}\leanok
Let $P$ be a digitally shifted digital $(0,m_0,3)$-net over $\F_b$.
\begin{enumerate}
\item[(A)] If $s_A,s_B$ are separating, then $L_D(P)$ is a digitally shifted digital $(0,m_0+3,3)$-net over $\F_b$.
\item[(B)] If $2D\le m_0$, separating $s_A,s_B$ exist.
\end{enumerate}
\end{lemma}

\begin{proof}
\leanok
\uses{lem:nie}
(A) Let $(d_x,d_y,d_z)$ be a box shape with $\sum_cd_c=m_0+3$; by Lemma~\ref{lem:nie} it suffices that the box map, from the index to the prescribed digits, is injective. \emph{All $d_c\ge1$.} The top digits $(\alpha,\alpha+\beta_1,\alpha+\beta_1+\beta_2)$ determine $(\alpha,\beta)$. The remaining depths $d_c-1$ sum to $m_0$ and are at most $m_0$ each, so no jitter digit is prescribed, and what remains is an $m_0$-box of the shifted net $P+s_\beta$, which is injective. \emph{Exactly one $d_c=0$}, say $d_x=0$ (the cases $d_y=0$, $d_z=0$ are the same with the pairs $xz$ and $xy$). The top digits of $y$ and $z$ fix $\alpha+\beta_1$ and $\beta_2$, and the remaining index is $(n,t)$ with $t=\beta_1$; the block digits of $(y,z)$ are those of $P_{yz}+\text{const}+t\,s_A^{yz}$. If $d_y-1=e$ and $d_z-1=m_0+1-e$ with $1\le e\le m_0$, injectivity on $(n,t)$ is injectivity of the projection to the first $e$ resp.\ $m_0+1-e$ digits on $P_{yz}+\F_bs_A^{yz}$; it holds on $P_{yz}$ by the $(e,m_0-e)$ box, and for $t\ne0$ a collision would put $s_A^{yz}$ in $H_e=\ker\lambda^{yz}_e$, which the separating condition excludes. If $d_y=1$ or $d_z=1$, the other coordinate reads all of its block and one jitter digit, $\beta_1$ for $z$ and $\beta_1+\beta_2$ for $y$, which determines $t$, and then the block determines $n$. For $d_y=0$ the pair is $(x,z)$ with the line $s_B-s_A$ and the jitter digits $\beta_2$, $\beta_1$; for $d_z=0$ it is $(x,y)$ with $s_B$ and $\beta_2$, $\beta_1+\beta_2$. \emph{Two $d_c=0$.} One coordinate carries all $m_0+3$ digits: its top digit, its block (a bijection in $n$ for fixed $\beta=(\beta_1,\beta_2)$) and its two jitter digits, which are a bijection of $\F_b^2$ onto $(\beta_1,\beta_2)$ in every coordinate.

(B) The functional $\lambda_e=\lambda^{cc'}_e$ vanishes on $Z_e$, so it is supported on the first $e$ digits of $c$ and the first $m_0+1-e$ digits of $c'$. Its coefficient on digit $e$ of $c$ is nonzero: otherwise $\lambda_e$ would live on a window of $e-1$ plus $m_0+1-e$ digits, onto which $P_{cc'}$ projects bijectively (the $(e-1,m_0+1-e)$ box), and vanish on $P_{cc'}$, hence be $0$. Likewise its coefficient on digit $m_0+1-e$ of $c'$ is nonzero. If $e>D$, digit $e$ of $c$ is \emph{deep} (index $>D$); if $e\le D$, digit $m_0+1-e\ge m_0+1-D>D$ of $c'$ is deep. So each $\lambda_e$ has a nonzero coefficient on a deep digit, its own: digit $e$ of $c$ if $e>D$, digit $m_0+1-e$ of $c'$ if $e\le D$. In any vanishing linear combination, the term with the largest $e>D$ (or, if there is none, the smallest $e$) is the only one that sees its own digit, so its coefficient is $0$; by induction all coefficients are $0$. Hence the restrictions of $\lambda_1,\dots,\lambda_{m_0}$ to the deep digits are linearly independent, and the deep parts of $s_A^{yz}$ can be chosen with $\lambda^{yz}_e(s_A^{yz})=1$ for all $e$, those of $s_B^{xy}$ with $\lambda^{xy}_e(s_B^{xy})=1$; the conditions on $(s_B-s_A)^{xz}$ involve the still free deep digits of $(s_A)_x$ and $(s_B)_z$, and are solved in the same way.
\end{proof}

\begin{lemma}[Freezing and separation]\label{lem:liftsep}
\lean{Z3NetB23.liftsep}\leanok
Let $P_L\to P_{L+3}\to\cdots$ be lifts whose depth parameters satisfy $D_{L+3j'}\ge\lfloor L/2\rfloor+j'$ for every $j'\ge0$ (for instance the canonical continuation). Then for every $j\ge0$ the toroidal (hence also Euclidean) minimum distance of $P_{L+3j}$ satisfies
\[
\delta N^{1/3}\ \ge\ \min\bigl(g_L,\ b^{L/3}-W\,b^{-2L/3}\bigr).
\]
\end{lemma}

\begin{proof}
\leanok
\emph{Freezing.} By induction on $j$, $P_{L+3j}$ is the set of points whose coordinate $c$ has the digits $\pi_c(u)\,|\,(p+\eta(u))_c\,|\,\mathrm{low}_c(u)$, where $u\in\F_b^{3j}$ collects the new index digits, $u\mapsto\pi(u)=(\pi_c(u))_c\in\F_b^{3j}$ is a bijection ($j$ top digits per coordinate), $p$ runs over $P_L$, $\eta(u)$ is a digitwise shift of the $L$ block digits, and $\mathrm{low}(u)$ has $2j$ digits per coordinate. The lift at resolution $\ell=L+3j'$ shifts only digits of index $>D_\ell$ of its own block, in which the block of $P_L$ occupies the indices $j'+1,\dots,j'+L$; so it touches block digits of index $>D_\ell-j'\ge\lfloor L/2\rfloor$ only, and never the $j'$ top digits before them. Hence $\eta(u)$ vanishes on the top $\lfloor L/2\rfloor$ digits of every coordinate: these are \emph{frozen}, equal to those of $p$.

\emph{Separation.} Measure in units of $b^{-(L+3j)}b^{2j}$, so that the torus has period $b^{L+j}$ and a point has coordinates $\pi_c(u)b^L+t_c(p)W+\varepsilon_c$ with $\varepsilon_c\in[0,W)$. Two points with different $p\ne p'$: reducing modulo $b^L$, which divides the period, does not increase toroidal distances, and modulo $b^L=AW$ the difference in coordinate $c$ is $(t_c(p)-t_c(p'))W+(\varepsilon_c-\varepsilon'_c)$, of toroidal size more than $W(|t_c(p)-t_c(p')|_A-1)$; so their distance is at least $W\sqrt{v^*_L}$. Two points with the same $p$ and different $u$: some $\pi_c$ differ, while $t_c$ agree, so coordinate $c$ differs by $(\pi_c-\pi'_c)b^L+(\varepsilon_c-\varepsilon'_c)$ with $\pi_c\not\equiv\pi'_c\pmod{b^j}$, of toroidal size more than $b^L-W$. A distance of $x$ units is $x\,b^{-L-j}$ in $[0,1)^3$, and $N^{1/3}=b^{L/3+j}$, so $\delta N^{1/3}$ is the distance in units times $b^{-2L/3}$.
\end{proof}

\begin{lemma}[Every digital shift]\label{lem:shiftcert}
\lean{Z3NetB23.shiftcert_a, Z3NetB23.vcirc_eq_iInf, Z3NetB23.shiftcert_b}\leanok
\uses{def:dshift, lem:liftsep}
For $\mathbf d\in\F_b^n$ let $\mu_n(\mathbf d)$ be the least toroidal distance $|X'-X|_{b^n}$ (the distance to the nearest multiple of $b^n$) over the integers $X\in[0,b^n)$, where $X'$ has the digits of $X$ plus $\mathbf d$, digit by digit modulo $b$. Let $H=\lfloor L/2\rfloor$.
\begin{enumerate}
\item[(a)] Under the hypothesis of Lemma~\ref{lem:liftsep}, for every $j\ge0$ and every digital shift $\tau$ of $P_{L+3j}$ (Definition~\ref{def:dshift}),
\[
\delta(P_{L+3j}\oplus\tau)\,N^{1/3}\ \ge\ \min\bigl(W\sqrt{v^\circ_L}\,b^{-2L/3},\ b^{L/3}-W\,b^{-2L/3}\bigr),
\]
where $v^\circ_L=\min_\Delta\sum_c\max\bigl(0,\mu_H(\Delta_c)-1\bigr)^2$, the minimum over the differences $\Delta=(\Delta_x,\Delta_y,\Delta_z)$ of the top $H$ digits of two distinct points of $P_L$ ($\Delta=0$ allowed, with weight $0$).
\item[(b)] For every digitally shifted digital $(0,m,3)$-net $P$ over $\F_b$,
\[
\min_\tau\delta(P\oplus\tau)^2=b^{-2m}\min_\Delta\textstyle\sum_c\mu_m(\Delta_c)^2,
\]
the minimum over the digital shifts $\tau$ and over the nonzero digit differences $\Delta$ of two points of $P$ (toroidal distances).
\end{enumerate}
\end{lemma}

\begin{proof}
\leanok
\uses{lem:liftsep}
For a pair of points with digit difference $\Delta$, coordinate $c$ of the shifted pair has the digits $X_c$ and $X_c+\Delta_c$, where $X_c$ is coordinate $c$ of the first point plus $\tau_c$. As $\tau$ varies, $(X_x,X_y,X_z)$ runs over all triples of digit vectors, independently in the three coordinates, so the minimum over $\tau$ of the squared distance of the pair is $\sum_c\min_{X_c}|X'_c-X_c|^2$. This gives (b). Since $t\mapsto\max(0,t-1)^2$ is nondecreasing, the same argument gives $v^\circ_L=\min_{\tau'}v^*_L(P_L\oplus\tau')$ over the digital shifts $\tau'$ of $P_L$. For (a), write the points of $P_{L+3j}$ as in the proof of Lemma~\ref{lem:liftsep}. A digital shift $\tau$ keeps $u\mapsto\pi(u)+\tau_\pi$ a bijection, where $\tau_\pi$ are its digits on the top $j$ digits, and adds the same vector $\tau_{\mathrm{blk}}$ to the block digits of every point. So $P_{L+3j}\oplus\tau$ has the structure used in that proof, with $P_L$ replaced by the digitally shifted copy $P_L\oplus\tau_{\mathrm{blk}}$: the top $H$ block digits of every point are those of a point of $P_L\oplus\tau_{\mathrm{blk}}$, and the proof applies verbatim. The first term is then $W\sqrt{v^*_L(P_L\oplus\tau_{\mathrm{blk}})}\,b^{-2L/3}\ge W\sqrt{v^\circ_L}\,b^{-2L/3}$, and the second term does not involve the shift.
\end{proof}

\begin{theorem}[Bases two and three; computer-assisted]\label{thm:base23}
\lean{Z3NetB23.b23_seed, Z3NetB23.b23_lift, Z3NetB23.b23_vcirc, Z3NetB23.base23_net, Z3NetB23.base23_exists, Z3NetB23.base23_everyshift, Z3NetB23.base23_everyshift_base, Z3NetB23.base23_witcur, Z3NetB23.base23_shifted, Z3NetDGS.dgs_family_b23}\leanok
\uses{def:dshift}
Let $b\in\{2,3\}$. For every $m\ge1$ there is an explicit digital $(0,m,3)$-net $P_m$ over $\F_b$ and a digital shift $\tau_m\in(\F_b^m)^3$ (Definition~\ref{def:dshift}) such that the toroidal (hence also Euclidean) minimum distances satisfy
\begin{enumerate}
\item $\delta(P_m\oplus\tau_m)\,N^{1/3}\ge0.4129$ for $b=2$ and $\ge0.3538$ for $b=3$;
\item $\delta(P_m)\,N^{1/3}\ge0.1718$ for $b=2$ and $\ge0.0770$ for $b=3$;
\item $\delta(P_m\oplus\tau)\,N^{1/3}\ge0.1718$ for $b=2$ and $\ge0.0770$ for $b=3$, for every digital shift $\tau\in(\F_b^m)^3$.
\end{enumerate}
Per residue $m\equiv0,1,2\pmod3$ the constants are those of Table~\texttt{tab:base23}. For $m\ge9$ ($b=2$) resp.\ $m\ge6$ ($b=3$) the nets come from a finite seed by deep-shift lifts: along a fixed chain up to the certificate resolution of Table~\texttt{tab:base23}, then by any continuation with $D_\ell=\lfloor\ell/2\rfloor$. For smaller $m$ they are explicit nets with exactly computed minimum distances.
\end{theorem}

\begin{proof}
\leanok
\uses{lem:lift, lem:liftsep, lem:shiftcert}
For the lifted resolutions: Lemma~\ref{lem:lift}(A) gives the net property at every resolution, Lemma~\ref{lem:lift}(B) the existence of the continuation, and Lemma~\ref{lem:liftsep} the bound $\delta N^{1/3}\ge g_L$ for every $m\ge L$ in the residue class of a certificate resolution $L$, once the finite certificate $v^*_L$ is known and the proviso $g_L\le b^{L/3}-Wb^{-2L/3}$ holds. The finite part is an exact computation on explicitly given nets: the $t=0$ property of the seeds and of the nets at the certificate resolutions, the integers $v^*_L$ (with the proviso, which holds by a wide margin), the exact minimum distances at the finitely many resolutions between the seed and $L$, and, for $m$ below the seed resolutions, explicit digital $(0,m,3)$-nets with their exact minimum distances (all above the stated constants). The constants are the minima of these exact values, truncated. Without the shift the same matrices are used: Lemmas~\ref{lem:lift} and~\ref{lem:liftsep} never use the shift, and only the certificates change. For (3), Lemma~\ref{lem:shiftcert}(a) gives the bound for every digital shift and every $m\ge L$ from one integer $v^\circ_L$, computed at the same certificate resolution $L$ (the proviso holds again); at the finitely many resolutions below $L$, and for the explicit nets below the seed, the exact minimum of $\delta(P_m\oplus\tau)N^{1/3}$ over all $\tau$ is computed by Lemma~\ref{lem:shiftcert}(b).
\end{proof}

\begin{remark}[Why every-shift certificates vanish at some resolutions]\label{rem:shiftparity}
\lean{Z3NetB23.shiftparity_parity, Z3NetB23.shiftparity_first_even, Z3NetB23.shiftparity_dichotomy, Z3NetB23.shiftparity_floor_iff, Z3NetB23.shiftparity_tendsto}\leanok
\uses{def:lift, lem:liftsep, thm:base23}
Along the chains the every-shift certificate $v^\circ_L$ vanishes at many resolutions (for $b=3$, $m\equiv0$, at every resolution below $18$; computed), and it is positive at the six certificate resolutions of Table~\texttt{tab:base23}, all of them even. Two facts explain the pattern. (i) \emph{Parity.} If $L$ is even and $v^\circ_L=0$, then $v^\circ_{L+3}=0$ for every lift $P_{L+3}$ of $P_L$. Indeed, the indices $(n,0,0,0)$ of the lift (Definition~\ref{def:lift}) give points with top digit $0$, block digits those of $P_L$ and jitter digits $0$, so the top $\lfloor(L+3)/2\rfloor=H+1$ digits of their differences are $0$ followed by the top $H$ digits of the differences of $P_L$. For such differences the toroidal $\mu_{H+1}$ of the extended vector is at most the plain (non-toroidal) value of the original one, and a vector $\pm(0,\dots,0,1,\dots,1)$ has plain value $1$; so a vanishing certificate at $L$ gives one at $L+3$. For odd $L$ the window of $P_{L+3}$ gains one more block digit, so a certificate can first become positive only at an even resolution. (ii) \emph{Dichotomy.} If $v^\circ_L=0$, then along every lift chain from $P_L$ every $m\ge L$ in the class has $\min_\tau\delta(P_m\oplus\tau)N^{1/3}<2\sqrt3\,b^{\lceil L/2\rceil-2L/3}$. Indeed, the indices whose new digits are all $0$ embed $P_L$ into $P_m$ with top and jitter digits $0$, and a difference whose top window is $0$ or $\pm(0,\dots,0,1,\dots,1)$ in each coordinate can be shifted to have each coordinate below $2W$ in the units $b^{-(L+j)}$ of the proof of Lemma~\ref{lem:liftsep}, $m=L+3j$. So, for a chain frozen from every resolution, a floor $\kappa N^{-1/3}$, $\kappa>0$, that holds for every shift from some resolution on exists if and only if $v^\circ_L>0$ for some $L$; if the certificate vanished at infinitely many $L$, the every-shift separation would tend to $0$. For the nets of Theorem~\ref{thm:base23} the certificate is positive at every certificate resolution (Table~\texttt{tab:base23}), so no shift makes their separation decay.
\end{remark}


\chapter{The question of Dick, Goda and Suzuki}\label{sec:dgs}

\begin{remark}[Low discrepancy]\label{rem:lowdisc}
\uses{lem:nie, lem:shift}
All nets in this paper are $(0,m,s)$-nets in base $q$: every elementary box of volume $q^{-m}$ contains exactly one point (Lemma~\ref{lem:nie}). This is a property of the point set alone, whatever ring the net is digital over, so Niederreiter's discrepancy bound for $(t,m,s)$-nets in an integer base \cite{Nie87} applies: $D^*_N\le C_{s,q}(\log N)^{s-1}/N+O_{s,q}\bigl((\log N)^{s-2}/N\bigr)$ for $N=q^m$, with $C_{s,q}$ independent of $m$. So all our nets have low discrepancy in the sense of \cite{DGS2,DGLPS1}, and the question of Dick, Goda and Suzuki reduces, for these nets, to separation. The same holds for digitally shifted nets, which are again $(0,m,s)$-nets (Lemma~\ref{lem:shift}).
\end{remark}


\section{Digital shifts}

\begin{lemma}[Shifted graph forms]\label{lem:shift}
\lean{Z3NetF.shift_net, Z3NetF.digPointShift_net, Z3NetF.shift_graph, Z3NetF.shift_graph_cells}\leanok
\uses{def:dshift}
Let $P$ be a graph form over $\Z/q\Z$ with depths $\mathbf k$ and blocks $\Phi$, and $\tau=(\tau^{\mathrm{cell}}_c,\tau^{\mathrm{off}}_c)_c$ a digital shift (Definition~\ref{def:dshift}), split into its digits on the cell digits and on the offset digits of each coordinate. Then $P\oplus\tau$ is the set of points whose coordinate $c$ has the digits $n_c$ followed by $\sum_d\Phi_{cd}n_d+\gamma_c$, $n\in\prod_c(\Z/q\Z)^{k_c}$, where $\gamma_c=\tau^{\mathrm{off}}_c-\sum_d\Phi_{cd}\tau^{\mathrm{cell}}_d$. In particular $n$ is the cell of the point, every cell holds one point, $P\oplus\tau$ is a $(0,m,s)$-net if $P$ is, and for two points in the cells $n,n'$ every offset digit difference is congruent to the corresponding entry of $\Phi(n'-n)$ modulo $q$, as for $P$.
\end{lemma}

\begin{proof}
\leanok
Substitute $n=\tilde n-\tau^{\mathrm{cell}}$ and index the shifted point by $\tilde n$. The shift maps every elementary box of shape $d$ bijectively onto one of the same shape, since it shifts the first $d_c$ digits of coordinate $c$.
\end{proof}

\begin{proposition}[Every digital shift]\label{prop:shift}
\lean{Z3NetF.Ss_coord_shift, Z3NetF.Ss_sep_shift, Z3NetF.D_sep_shift, Z3NetF.S2_sep4_shift, Z3NetF.S2_sep4_shift_real, Z3NetF.S2_sep_shift, Z3NetF.S2_sep_shift_real, Z3NetA.Apoint_eq_AGen, Z3NetA.A_net_shift, Z3NetA.A_sep_shift, Z3NetA.A_sep_shift_real, Z3NetA.A_sep_shift_eucl, Z3NetA.A_sharp_shift, Z3NetA.A_sharp_shift_real, Z3Net.P_sep_shift, Z3Net.P_sep_shift_ge8, Z3Net.P_sep_shift_real, Z3NetS1.S1_sep_shift, Z3NetS1.S1_sep_shift_real, Z3NetS1.S1_cells_shift, Z3NetS1.S1_cover_shift, Z3NetKN.B1_sep_shift, Z3NetKN.B1_sep_shift_real, Z3NetKN.B2_rem_sep_shift_real, Z3NetHL.GF_relax_shift, Z3NetHL.GF_dom_shift, Z3NetHL.GF_dom_shift_real, Z3NetHL.HL_relax_shift, Z3NetHL.HL_dom_shift, Z3NetHL.HL_dom_shift_real, Z3NetHL.HL_domL_shift_real, Z3NetHL.HL_main_shift_real, Z3NetHL.HL_main_shift, Z3NetHL.GFpoint_shift_zero, Z3NetHL.HLpoint_shift_zero}\leanok
\uses{thm:A, lem:Ss, prop:D, thm:S1, thm:S2, thm:main, thm:allm, prop:Hk, thm:HL}
Lemma~\ref{lem:Ss} holds verbatim, with the same bound, for every digital shift of a graph form that satisfies (H1) and (H2). Consequently the lower bounds on the separation in Theorems~\ref{thm:A}, \ref{thm:main}, \ref{thm:S1}, \ref{thm:S2}, \ref{thm:allm} and \ref{thm:HL} and in Propositions~\ref{prop:D} and~\ref{prop:Hk} hold for every digital shift of the respective nets (in Theorem~\ref{thm:S1}(2) the occupied fine cells form a coset of the lattice), and so does the first-order sharpness in Theorem~\ref{thm:A}(2).
\end{proposition}

\begin{proof}
\leanok
\uses{lem:Ss, thm:S2, lem:shift}
By Lemma~\ref{lem:shift} the shift only adds the constant $\gamma$ to the offsets. The proofs use the cells (the step and its carries), the residues of digit differences, the facts that a digit difference lies in $(-q,q)$ and vanishes if its residue does, and $|f'-f|<1$; the constant cancels in every residue, and in Lemma~\ref{lem:Ss} the deeper digits of a coordinate whose own axis alone moves are still equal, by (H2). Theorem~\ref{thm:S2} does not even need the re-indexing by the shifted cells: its proof uses only differences of the skeleton digits. The sharpness argument prescribes the cell digits and the first offset digits of a pair; in the shifted net these are digits of the unshifted net prescribed in translated sets of the same size, so the joint uniformity of the digits of the unshifted net, given by its $t=0$, suffices.
\end{proof}

\begin{corollary}[Corollary~\ref{cor:allm} under digital shifts]\label{cor:allmshift}
\lean{Z3NetKN.allmshift, Z3NetC.allmshift_m3}\leanok
\uses{thm:col, cor:allm}
Let $p\ge37$ be prime. For every $m\ge3$ and every digital shift, the shifted nets of Corollary~\ref{cor:allm} satisfy its separation bounds, with the same constants $0.2715$ and $0.4250$, and its covering bounds; so they are quasi-uniform along every $N=p^m$, $m\ge3$, as the Larcher--Pillichshammer nets in \cite[Thm.~3.11]{DGS2}. For $m\ge4$ the lower bound still tends to $1$ as $p\to\infty$. At $m=3$ the nets are the column nets $\mathcal C_p$, whose three coordinates share the second and third digits, and the bound depends on whether the shift keeps them shared:
\begin{enumerate}
\item[(a)] for every shift, $\delta N^{1/3}>0.4693$;
\item[(b)] if the shift digits of the three coordinates agree at the second and at the third digit, the exact value of Theorem~\ref{thm:col}(2) holds;
\item[(c)] if they agree at the second digit, $\delta N^{1/3}\ge1-(7p-3)/(\sqrt3p^2)$;
\item[(d)] otherwise $1/\sqrt3-7/(\sqrt3p)\le\delta N^{1/3}\le1/\sqrt3+5/(\sqrt3p)$.
\end{enumerate}
\end{corollary}

\begin{proof}
\leanok
\uses{thm:S1, thm:main, prop:colshift, thm:allm, cor:allm, lem:shift, prop:shift}
A digital shift keeps one point in every cell and keeps $t=0$ (Lemma~\ref{lem:shift}), which is all the covering bounds use, except for the mesh-ratio constant $4.264$ of Corollary~\ref{cor:allm}, which comes from Theorem~\ref{thm:S1}(3); at $m=3$ this gives $h\le\sqrt3N^{-1/3}$. In the shifted net the occupied fine cells form a coset $g+\Lambda_A$ (Proposition~\ref{prop:shift}), and the proof of Theorem~\ref{thm:S1}(3) runs with its target residue moved by $g_z-\sum_d\rho^dg_{x_d}$; this is not a translation of the unshifted construction, since the target cell must stay in range. For $m\ge4$ the separation bounds of Theorems~\ref{thm:main}, \ref{thm:allm} and~\ref{thm:S1} used in Corollary~\ref{cor:allm} hold for every shift by Proposition~\ref{prop:shift}. For $m=3$ Proposition~\ref{prop:colshift} gives the lower bound $0.4693>0.4250$, the exact value and the bound for agreeing shift digits, and otherwise the two bounds $\frac1{\sqrt3}(1-\frac{W+3}p+\frac3{p^2})$ and $\frac1{\sqrt3}(1+\frac{3+w_{\max}}p-\frac3{p^2})$ with $W\le4$ and $w_{\max}\le2$, which lie in the stated interval.
\end{proof}

\begin{remark}[Shared digits under shifts]\label{rem:shiftshared}
\lean{Z3NetC.shiftshared_col}\leanok
\uses{thm:col, prop:colshift, rem:glue}
So the exact value of Theorem~\ref{thm:col} does not extend to all shifted nets: if the shift digits of the three coordinates of $\mathcal C_q$ at the second digit are not all equal, then for every $q\ge11$, prime or not, the shifted net has $\delta N^{1/3}\le\frac1{\sqrt3}(1+\frac{3+w_{\max}}q-\frac3{q^2})$, below the value of Theorem~\ref{thm:col} (Proposition~\ref{prop:colshift}(3)). For each coordinate, the set of values of the shared digit at which a step wraps (its \emph{wrap window}) is translated by that coordinate's shift digit; unequal shift digits give distinct windows, so some cell wraps in exactly one coordinate, and its step along that axis gives the pair $qa_c\mathbf 1+\Delta j$ at integer scale. In general a digital shift translates the wrap windows independently in each coordinate. Bounds certified with arbitrary wraps (Remark~\ref{rem:glue}) therefore survive every shift; bounds that use locked wraps, wraps that are correlated across coordinates, need not. The nets $\mathcal T_p$ and $\mathcal W_p$ of Part~II are of the second kind (Remark~\ref{rem:wshift}).
\end{remark}


\section{Well-separated families}

\begin{corollary}[Well-separated families in the sense of Dick, Goda and Suzuki]\label{cor:dgs}
\lean{Z3NetDGS.lemma37, Z3NetDGS.lemma37_strict, Z3NetDGS.isTMSNet_of_isNet0, Z3NetDGS.dgsNet_eq, Z3NetDGS.dgs_main, Z3NetDGS.dgs_P, Z3NetDGS.dgs_A, Z3NetDGS.dgs_S1, Z3NetDGS.dgs_HL, Z3NetDGS.dgs_allm, Z3NetDGS.dgs_S1_floor, Z3NetDGS.dgs_family_b23}\leanok
\uses{thm:A, def:dshift, thm:S1, thm:main, cor:allm, thm:HL}
Let $s\ge2$ and $q\ge2$ be integers and $I\subseteq\N$ infinite, and for $m\in I$ let $\mathcal P_m\subset[0,1)^s$ be a $(0,m,s)$-net in base $q$ whose toroidal minimum distance satisfies $\delta\ge c\,q^{-m/s}$, with $c>0$ independent of $m$. Then every $\mathcal P_m$ is $\kappa_m$-separated and toroidally $\kappa_m$-separated in base $q$ \cite[Def.~3.3]{DGS2}, with
\[
\kappa_m=\Bigl\lceil\frac ms+\log_q\frac{\sqrt s}{c}\Bigr\rceil\ \le\ \frac ms+\log_q\frac{\sqrt s}{c}+1,
\]
so $(\mathcal P_m)_{m\in I}$ is a well-separated and toroidally well-separated family on $I$ \cite[Def.~3.8]{DGS2}. In particular:
\begin{enumerate}
\item the nets $\mathcal P_{k,q}$ of Theorem~\ref{thm:main}, for every integer $q\ge4$, are $(k+1)$-separated, i.e.\ $\kappa_m=m/3+1$ on $I=\{3k:k\ge2\}$;
\item the nets of Theorem~\ref{thm:A}(2) are $(k+1)$-separated on $I=\{sk:k\ge1\}$;
\item the nets $\mathcal H^s_{k,p}$ of Theorem~\ref{thm:HL} are $(k+1)$-separated for every prime $p\ge\max(2s,2B_s+4)$ (two more than in Theorem~\ref{thm:HL}; see Step~3 of the proof) and every $k\ge1$;
\item the nets $\mathcal S^1_{k,p,\rho}$ of Theorem~\ref{thm:S1}, with $m=sk+s-1$, are $(k+1)$-separated, i.e.\ $\kappa_m=\lceil m/s\rceil$, if $(\rho-1)^2\ge s$ (for $\rho=\lfloor p^{1/s}\rfloor$ this holds when $p>(\lceil\sqrt s\rceil+1)^s$), and $(k+2)$-separated otherwise;
\item for every prime $p\ge37$ the nets of Corollary~\ref{cor:allm} are $(\lfloor m/3\rfloor+1)$-separated for every $m\ge3$.
\end{enumerate}
All of these hold toroidally as well, and for every digital shift (Definition~\ref{def:dshift}) of the respective nets. For prime $q$ the nets are digital nets over $\F_q$ in the sense of \cite[Def.~2.1]{DGS2}; this covers (3)--(5) entirely and the prime instances of (1) and~(2). For composite $q$ the nets of (1) and (2) are digital over $\Z/q\Z$, which is not covered by \cite[Def.~2.1]{DGS2}; Definitions~3.1 and~3.2 of \cite{DGS2} are stated for every integer base $b\ge2$, and we read Definitions~3.3 and~3.8 in the same generality.
\end{corollary}

\begin{proof}
\leanok
\uses{thm:S1, thm:sep, thm:colmin, prop:colshift, thm:allm, prop:shift}
\emph{Step 1 (\cite[Lemma~3.7]{DGS2}, with strict inequality).} Let $P\subset[0,1)^s$ be finite and $\kappa\ge0$, and suppose every two distinct points of $P$ have toroidal $\ell^\infty$-distance at least $q^{-\kappa}$. A toroidal shifted elementary interval $\tilde J_{\mathbf c,a,e}$ with $\kappa$ digits in every coordinate, $\mathbf c=(\kappa,\dots,\kappa)$ \cite[Def.~3.2]{DGS2}, is contained in a product of sets $I_j\bmod1$, where $I_j=[\alpha_j,\alpha_j+q^{-\kappa})$ is half-open. If $x,y\in\tilde J_{\mathbf c,a,e}$, then for each $j$ there are $z,z'\in\Z$ with $x_j+z,\ y_j+z'\in I_j$, so $\min(|x_j-y_j|,1-|x_j-y_j|)\le|(x_j+z)-(y_j+z')|<q^{-\kappa}$, and hence $x=y$. So $P$ is toroidally $\kappa$-separated. A non-toroidal interval $J_{\mathbf c,a,e}$ gives the same bound for $|x_j-y_j|$, which is at least the toroidal distance, so $P$ is also $\kappa$-separated. This sharpens the contrapositive of \cite[Lemma~3.7]{DGS2} to the strict bound $q_\infty(P)<1/(2q^\kappa)$ for non-separated $P$, where $q_\infty(P)$ is the separation radius of \cite{DGS2} in the maximum norm; the strictness comes from the intervals being half-open.

\emph{Step 2.} For distinct $x,y\in\mathcal P_m$ the toroidal coordinate distances $t_j$ satisfy $\sum_jt_j^2\ge\delta^2$, so $\max_jt_j\ge\delta/\sqrt s\ge(c/\sqrt s)\,q^{-m/s}\ge q^{-\kappa_m}$ by the choice of $\kappa_m$. Since toroidal distances are at most $\sqrt s/2$, we have $c\,q^{-m/s}\le\sqrt s/2$, so $\kappa_m\ge1$. Step~1 gives the separation, and $\kappa_m-m/s\le1+\log_q(\sqrt s/c)$ is bounded, which is the condition of \cite[Def.~3.8]{DGS2}.

\emph{Step 3 (the families).} By Step~2 a net with $\delta\ge\theta\,q^{-\kappa}$ is $\kappa$-separated as soon as $\theta\ge\sqrt s$.
(1) $m=3k$ and $\delta\ge c(q)q^{-k}$ (Theorem~\ref{thm:sep}); with $\kappa=k+1$ we need $q\,c(q)\ge\sqrt3$, i.e.\ $\min\{(q-1)^2,(q-2)^2+(q-4)^2,3(q-3)^2\}\ge3$, which holds for $q\ge4$ (with equality at $q=4$).
(2) $\delta\ge(1-a^*_s/q)q^{-k}$; with $\kappa=k+1$ we need $q-a^*_s\ge\sqrt s$. From $q\ge5a^*_s+4$ we get $q-a^*_s\ge\frac45(q+1)$, and from $\operatorname{lpf}(q)\ge s-1$ we get $q+1\ge s$, so $q-a^*_s\ge\frac45s\ge\sqrt s$ for $s\ge2$.
(3) $\delta\ge c_p\,p^{-k}$ with $c_p=\min(\sqrt s(\frac12-\epsilon_p),\sqrt2(1-\epsilon_p))$ and $\epsilon_p=(B_s+1)/p$; with $\kappa=k+1$ we need $p\,c_p\ge\sqrt s$. If $p\ge2B_s+4$, then $\sqrt s(\frac12-\epsilon_p)p=\frac{\sqrt s}2(p-2B_s-2)\ge\sqrt s$ and $\sqrt2(1-\epsilon_p)p\ge p/\sqrt2\ge\sqrt2\,s\ge\sqrt s$, since $\epsilon_p\le\frac12$ and $p\ge2s$.
(4) $m=sk+s-1$ and $\delta\ge(\rho-1)p^{-(k+1)}$ (Theorem~\ref{thm:S1}(1)); $\kappa=k+1$ needs $\rho-1\ge\sqrt s$, and otherwise $\kappa=k+2$ needs $(\rho-1)p\ge\sqrt s$, which holds as $p>2^s$. If $p>(\lceil\sqrt s\rceil+1)^s$, then $\rho=\lfloor p^{1/s}\rfloor\ge\lceil\sqrt s\rceil+1$.
(5) $m=3$: $\delta\ge\sqrt{2/3}\,p^{-1}$ (Theorem~\ref{thm:colmin}), and $\kappa=2$ needs $p\sqrt{2/3}\ge\sqrt3$, i.e.\ $p\ge3/\sqrt2$. $m=3k$, $k\ge2$: part~(1). $m=3k+1$: $\delta\ge(\rho^2-4)p^{-(k+1)}$ with $\rho$ maximal subject to $\rho^2(\rho+1)\le p$ (Theorem~\ref{thm:allm}), and $p\ge37$ gives $\rho\ge3$, so $\rho^2-4\ge5\ge\sqrt3$ and $\kappa=k+1$. $m=3k+2$: part~(4) with $s=3$ and $\rho=\lfloor p^{1/3}\rfloor\ge3$, so $(\rho-1)^2=4\ge3$ and $\kappa=k+1$. In every case $\kappa_m=\lfloor m/3\rfloor+1$.

\emph{Digital shifts.} By Proposition~\ref{prop:shift} the separation bounds used in (1)--(4) and, for $m\ge4$, in (5) hold for every digitally shifted net. For $m=3$, Proposition~\ref{prop:colshift}(3) gives $\delta\ge\frac1{\sqrt3}(1-\frac{W+3}p+\frac3{p^2})\,p^{-1}$ with $W\le4$, and $\kappa=2$ needs $p^2-(W+6)p+3\ge0$, which holds for $p\ge10$.
\end{proof}

\begin{remark}[Finitely many resolutions]\label{rem:dgsfinite}
\lean{Z3NetDGS.dgsfinite_sep, Z3NetDGS.frame3_net, Z3NetDGS.dgs_family, Z3NetDGS.dgs_pp, Z3NetDGS.dgs_family_pp, Z3NetDGS.dgs_family_b23, Z3NetDGS.dgs_family_all}\leanok
\uses{lem:nie, lem:Ts, thm:main, thm:col, cor:allm, rem:perprime, thm:base23, lem:shift, prop:shift, cor:allmshift, cor:dgs, rem:ppshift}
Definition~3.8 of \cite{DGS2} constrains only $m\in I$, and finitely many further resolutions can be added at no cost. If $\mathcal Q\subset q^{-m}\Z^s\cap[0,1)^s$ is a $(0,m,s)$-net in base $q$, the elementary boxes of shape $(m,0,\dots,0)$ each hold one point, so the first coordinates of the points are distinct elements of $q^{-m}\Z\cap[0,1)$ and any two of them have toroidal distance at least $q^{-m}$; by Step~1 of the proof of Corollary~\ref{cor:dgs}, $\mathcal Q$ is $m$-separated and toroidally $m$-separated. Adding finitely many such $m$ to $I$ changes the $O(1)$ of \cite[Def.~3.8]{DGS2} by a constant. In dimension three digital $(0,m,3)$-nets over $\Z/q\Z$ exist for every integer $q\ge2$ (over $\F_p$ for every prime $p$) and every $m\ge1$: by Lemma~\ref{lem:Ts}, whose proof does not use $m\equiv0\pmod s$, every box minor of the frame of the places $\infty,0,1$ is $\pm1$, a unit modulo every $q$, so the frame is one (Lemma~\ref{lem:nie}; see also \cite[Thm.~6.2, Cor.~6.3]{Nie87} and \cite[Sec.~8.2]{DP10}), and trivially for $m=0$. Together, Theorems~\ref{thm:main} and~\ref{thm:col}, Corollary~\ref{cor:allm}, Remark~\ref{rem:perprime} and Theorem~\ref{thm:base23} cover all but finitely many $m$ in each prime base with a constant independent of $m$; with the frame nets for the remaining $m$ they give well-separated families on all of~$\N$, in every prime base. The same holds under every digital shift: the constants survive every shift (Proposition~\ref{prop:shift}, Corollary~\ref{cor:allmshift}, Remark~\ref{rem:ppshift} and Theorem~\ref{thm:base23}(3)), and a digitally shifted $(0,m,s)$-net is again a $(0,m,s)$-net on the grid $q^{-m}\Z^s$ (Lemma~\ref{lem:shift}), so the argument above applies to it.
\end{remark}


\section{The maximum norm}

\begin{proposition}[Maximum norm]\label{prop:linf}
\lean{Z3NetDGS.linf_dist, Z3NetDGS.linf_mesh}\leanok
Let $q\ge4$, $k\ge2$. For every digital shift $\tau$, any two distinct points of $\mathcal P_{k,q}\oplus\tau$ have toroidal $\ell^\infty$-distance at least $(1-3/q)\,q^{-k}$. Hence the $\ell^\infty$ mesh ratio of $\mathcal P_{k,q}\oplus\tau$ in the sense of \cite{DGS2}, toroidal or not, is at most $2q/(q-3)$, and $\mathcal P_{k,q}\oplus\tau$ is toroidally $(k+1)$-separated \cite[Def.~3.3]{DGS2}.
\end{proposition}

\begin{proof}
\leanok
\uses{thm:sep, prop:shift, cor:dgs}
By Proposition~\ref{prop:shift} the proof of Theorem~\ref{thm:sep} applies to $\mathcal P_{k,q}\oplus\tau$. In its sign analysis every step has a coordinate with $|\Delta_i|\ge1-3/q$ ($=\beta_3$ in that proof): an own-axis step has $\beta_1$; for $\nu=2$, (iv) gives a backward coordinate with $|\sigma|\le1$ or $\sigma=0$; for $\nu=3$, a forward coordinate or a backward one with $|\sigma|=3$ forces a backward one with $|\sigma|\le1$ by (ii), (iii), and otherwise all three are backward with $|\sigma|\le2$. One point per cell gives $h_\infty\le q^{-k}$, and $(1-3/q)q^{-k}\ge q^{-k-1}$ for $q\ge4$; the separation statement is Step~1 of the proof of Corollary~\ref{cor:dgs}.
\end{proof}


\part{Separation constants}\label{part:constants}


\chapter{The first-order deficit: certificates and ceilings}\label{sec:deficit}


\section{The certified deficit of a design}

\begin{definition}[Certified first-order deficit]\label{def:K}
\lean{Z3NetF.Kdef_ne_top_iff, Z3NetF.Kdef_eq_Kexp, Z3NetF.K_expansion, Z3NetF.jit, Z3NetF.cost, Z3NetF.Kdef, Z3NetF.Sdef, Z3NetF.cfgs, Z3NetF.sigma, Z3NetF.bwd}\leanok
The \emph{cost} of a configuration is $|\sigma_b|+\mathrm{jit}_b$ if exactly one moving coordinate $b$ is backward, $0$ if at least two are, and $+\infty$ if none is. The \emph{certified deficit} $K(F)$ is the maximum of the cost over all configurations, and $\sigma_{\max}(F)$ is the maximum of $|\sigma_i|$ over all configurations and all $i$.
\end{definition}

\begin{proposition}[Certificate for all large $q$]\label{prop:Dprime}
\lean{Z3NetF.Dprime_LB, Z3NetF.Dprime, Z3NetF.Dprime_sep, Z3NetF.QD}\leanok
\uses{lem:Ss}
Let $F$ be a design with $K=K(F)<\infty$ and $\sigma_{\max}=\sigma_{\max}(F)$, and put
\[
q_{\mathrm{tail}}(K,\sigma_{\max})=\max\Bigl(\sigma_{\max}+1,\ \Bigl\lceil\frac{\sqrt2(\sigma_{\max}+1)-K}{\sqrt2-1}\Bigr\rceil\Bigr).
\]
Then $\mathrm{LB}\ge1-K/q$ for every $q\ge q_{\mathrm{tail}}(K,\sigma_{\max})$, where $\mathrm{LB}$ is the relaxation of Lemma~\ref{lem:Ss} for $F$ or for any truncation $(F_0,\dots,F_{L'-1})$, $L'\le L$.
\end{proposition}

\begin{proof}
\leanok
\uses{lem:Ss}
Fix a configuration and let $q\ge q_{\mathrm{tail}}$, so $q>\sigma_{\max}\ge|\sigma_i|$ for all $i$. For a moving backward coordinate $i$ the bound $b_i$ of Lemma~\ref{lem:Ss} is $q-|\sigma_i|-\mathrm{jit}_i\ge0$: for $e_i=1$ and $-q<\sigma_i<0$, $\bar\sigma_i=q+\sigma_i$ and the two options give $|q+t|\in\{2q-|\sigma_i|,\,q-|\sigma_i|\}$; for $\sigma_i=0$ the only option gives $q$; $e_i=-1$ is symmetric. If at least two coordinates are backward, $\sum_ib_i^2\ge2(q-\sigma_{\max}-1)^2\ge(q-K)^2$, because $\sqrt2(q-\sigma_{\max}-1)\ge q-K$ is the second condition on $q_{\mathrm{tail}}$ (and $K\le\sigma_{\max}+1\le q$). If exactly one coordinate $b$ is backward, $\sum_ib_i^2\ge(q-|\sigma_b|-\mathrm{jit}_b)^2\ge(q-K)^2$. No backward coordinate is excluded by $K<\infty$. So every configuration gives $(\sum_ib_i^2)^{1/2}/q\ge1-K/q$. The configurations of a truncation are among those of $F$, with the same $\sigma$ (Lemma~\ref{lem:Ss}, capped carry depths), so its $\mathrm{LB}$ is at least that of $F$.
\end{proof}

\begin{lemma}[The relaxation stops at one]\label{lem:Kone}
\lean{Z3NetF.Kone}\leanok
Let $F$ be a design with some nonzero off-diagonal entry $F_0[a][b]$. Then $K(F)\ge1$.
\end{lemma}

\begin{proof}
\leanok
Write $M=F_0$ and consider the carry-free steps $e=\e_a+\varepsilon\e_b$, $\varepsilon=\pm1$, along which only $a$ and $b$ move, with $\sigma_a=M_{aa}+\varepsilon M_{ab}$ and $\varepsilon\sigma_b=M_{bb}+\varepsilon M_{ba}$. If for some $\varepsilon$ at most one of $a$, $b$ is backward, that configuration costs at least $\mathrm{jit}=1$, or $+\infty$. Otherwise $M_{aa}+\varepsilon M_{ab}\le0$ for both signs, so $M_{aa}\le-|M_{ab}|\le-1$, and the one-axis step along $a$ without carry costs $-M_{aa}\ge1$.
\end{proof}


\section{The first-order ceiling}\label{sec:ceil}

\begin{proposition}\label{prop:ceil}
\lean{Z3Net.ceil_neg, Z3Net.ceil_pos, Z3Net.sign_boxMinor, Z3Net.sign_net, Z3Net.sign_graph, Z3Net.sign_cap}\leanok
Let $p$ be an odd prime, $k\ge2$, and let $\mathcal P$ be a digital $(0,3k,3)$-net over $\F_p$ with normalized minimum distance $c=\delta N^{1/3}$ (toroidal or not). Let $a$, $b$, $b'$ be the three coordinates $x$, $y$, $z$ in some order.
\begin{enumerate}
\item[(i)] If $\tilde M_{aa}\le-1$, then $c^2<(1-1/p+p^{1-k})^2+\bigl((|\tilde M_{ba}|+1)^2+(|\tilde M_{b'a}|+1)^2\bigr)/p^2$.
\item[(ii)] If $\tilde M_{aa}\ge1$, then $c^2<\bigl((\tilde M_{aa}+1)^2+(|\tilde M_{ba}|+1)^2+(|\tilde M_{b'a}|+1)^2\bigr)/p^2$.
\end{enumerate}
\end{proposition}

\begin{proof}
\leanok
Let $\mu\in\{1,\dots,p-1\}$ be the least nonnegative residue of $(\Phi_{aa})_{0,k-1}$, so $\mu=p+\tilde M_{aa}$ in (i) and $\mu=\tilde M_{aa}$ in (ii). Take a cell $n$ with $n_a=0$ and its neighbour $n+e_a$; only the last digit of $n_a$ changes, so each offset digit of each coordinate changes by a constant modulo $p$, the first one of coordinate $u$ by $(\Phi_{ua})_{0,k-1}$. Require
(A) the first $k-1$ offset digits of $a$ all equal $p-1$, and
(B) for $u\in\{b,b'\}$, the first offset digit of $u$ is not in the set $W_u$ of the $|\tilde M_{ua}|$ digit values on which adding $\tilde M_{ua}$ leaves $\{0,\dots,p-1\}$.
Under (A) the first offset digit of $a$ becomes $\mu-1$, the next $k-2$ (all $p-1$) can only decrease, and the remaining digits contribute less than $p^{1-k}$. So in cell units $\Delta_a=\mu/p+\theta$ with $-1/p<\theta<p^{1-k}$. Under (B), the first offset digit of $u$ changes by exactly $\tilde M_{ua}$, so $|\Delta_u|<(|\tilde M_{ua}|+1)/p$. In (i), $0<\Delta_a<1-1/p+p^{1-k}$; in (ii), $0<\Delta_a<\tilde M_{aa}/p+p^{1-k}\le(\tilde M_{aa}+1)/p$ since $k\ge2$.
It remains to find such a cell. The $p^{2k}$ indices with $n_a=0$ are parametrized by $(n_b,n_{b'})$. By the $t=0$ boxes $(2k-1,k+1,0)$ and $(2k-1,0,k+1)$ (order $a,b,b'$), exactly $p^{k+1}$ of them satisfy (A), and among these the first offset digits of $b$ and of $b'$ each take every value exactly $p^k$ times. Hence at most $(|W_b|+|W_{b'}|)p^k\le(p-1)p^k<p^{k+1}$ of them violate (B), since $|\tilde M_{ua}|\le(p-1)/2$. The pair $n$, $n+e_a$ lies inside the cube, so it bounds the toroidal and the Euclidean distance.
\end{proof}

\begin{corollary}[First-order ceiling in dimension three]\label{cor:ceil3}\label{prop:A}
\lean{Z3Net.propA_neg, Z3Net.propA_pos, Z3Net.propA_pos_k2}\leanok
Let $p$ be an odd prime and $k\ge2$. Suppose the first offset digit of some coordinate $a$ of a digital $(0,3k,3)$-net over $\F_p$ depends on the last cell digit of $a$ with a coefficient whose centred representative is nonzero, and the coefficients of the other two coordinates on that digit are bounded by $C\ge1$ in absolute value. If that representative is $\le-1$ and $k\ge3$, then $\delta N^{1/3}<1-1/p+7C^2/p^2$. If it lies in $[1,C]$, then $\delta N^{1/3}<4C/p$.
\end{corollary}

\begin{proof}
\leanok
\uses{prop:ceil}
The coefficient of the corollary is $\tilde M_{aa}$; let $|\tilde M_{ba}|,|\tilde M_{b'a}|\le C$, $C\ge1$. In (i) of Proposition~\ref{prop:ceil}, with $k\ge3$, put $\alpha=1-1/p+p^{1-k}\ge2/3$ and $\beta=\bigl((|\tilde M_{ba}|+1)^2+(|\tilde M_{b'a}|+1)^2\bigr)/p^2\le8C^2/p^2$; then $c<\sqrt{\alpha^2+\beta}\le\alpha+\beta/(2\alpha)\le\alpha+6C^2/p^2$, and $p^{1-k}\le1/p^2\le C^2/p^2$, so $c<1-1/p+7C^2/p^2$. In (ii), with $1\le\tilde M_{aa}\le C$, $c^2<3(C+1)^2/p^2\le12C^2/p^2$, so $c<4C/p$. (For $k=2$ the term $p^{1-k}=1/p$ cancels the $-1/p$; hence $k\ge3$.)
\end{proof}

\begin{corollary}[Ceiling in every dimension]\label{cor:ceils}
\lean{Z3NetF.ceils_neg, Z3NetF.ceils_pos, Z3NetF.ceils_pos_C}\leanok
Let $s\ge3$, $k\ge2$, and let $\mathcal P$ be a digital $(0,sk,s)$-net over $\Z/q\Z$, $q\ge3$, in graph form, with geometric layer $\tilde M$ and $c=\delta N^{1/s}$ (toroidal or not). Let $a$ be a coordinate with $\sum_{u\ne a}|\tilde M_{ua}|<q$.
\begin{enumerate}
\item[(i)] If $\tilde M_{aa}\le-1$, then
\[
c^2<\Bigl(1-\frac{|\tilde M_{aa}|}q+q^{1-k}\Bigr)^2+\sum_{u\ne a}\frac{(|\tilde M_{ua}|+1)^2}{q^2}.
\]
\item[(ii)] If $\tilde M_{aa}\ge1$, then
\[
c^2<\frac{(\tilde M_{aa}+1)^2+\sum_{u\ne a}(|\tilde M_{ua}|+1)^2}{q^2};
\]
in particular $c<\sqrt s\,(C+1)/q$ if $|\tilde M|\le C$.
\end{enumerate}
\end{corollary}

\begin{proof}
\leanok
\uses{prop:ceil}
The proof of Proposition~\ref{prop:ceil} applies with all counts multiplied by $q^{(s-3)k}$ and $\Delta_a=\mu/q+\theta$, $\mu=q+\tilde M_{aa}$ in (i), $\mu=\tilde M_{aa}$ in (ii) (where $0<\Delta_a<(\tilde M_{aa}+1)/q$ as $k\ge2$). The boxes used have depth $0$ off $a$ and the transverse coordinate considered, and total depth $\le3k\le sk$; no field structure is used; and the union bound over $u\ne a$ leaves a cell because $\sum_{u\ne a}|\tilde M_{ua}|<q$.
\end{proof}

\begin{remark}[The sign of the diagonal; diagonal zero]\label{rem:sign}\label{rem:diag0}
\lean{Z3Net.sign_boxMinor, Z3Net.sign_net, Z3Net.sign_graph, Z3Net.sign_cap}\leanok
\uses{lem:box, thm:sep, prop:ceil}
Case (i) of the sign analysis in the proof of Theorem~\ref{thm:sep} uses the diagonal entry $-1$ of $M$: on a step along its own axis a coordinate moves by $1-1/q$ cells, slightly less than a full cell, never slightly more. Replacing it by $+1$ (in $\mathcal P_{k,q}$: row $k$ of $E$ becomes $(Ed)_k=-d_{k-1}+d_k$, which changes only the entry $(\Phi_{cc})_{0,k-1}$) keeps $t=0$ (Lemma~\ref{lem:gauge}(a)), but for odd prime $q$ Proposition~\ref{prop:ceil}(ii) caps the minimum distance at $O(1/q)$. If $\tilde M_{aa}=0$ for all $a$, Proposition~\ref{prop:ceil} does not apply; heuristically such layers stay below $1-1/p$ at first order too, but the argument would need a joint condition on three digits that is not a single box, and we have no proof (Remark~\ref{rem:open}(1)).
\end{remark}


\section{Dimension four}\label{sec:dim4ex}

\begin{proposition}[Dimension four, deficit two; computer-assisted]\label{prop:4D2}
\lean{Z3NetF.D4D2_data, Z3NetF.D4D2_net, Z3NetF.D4D2_boxMinor, Z3NetF.D4D2_covering, Z3NetF.D4D2_KS, Z3NetF.D4D2_LB_small, Z3NetF.D4D2_LB, Z3NetF.D4D2_sep, Z3NetF.D4D2_sep_eucl, Z3NetF.D4D2_ceil, Z3NetF.D4D2_ceil_real}\leanok
\uses{def:Ns}
Let $F$ be the design of Appendix~\texttt{app:designs} on the places $\infty,0,1,2$ with $\lambda'\equiv1$, and for odd $q$ and $k\ge1$ let $\mathcal P^F_{k,q}$ be the net of Definition~\ref{def:Ns}, $N=q^{4k}$. Then:
\begin{enumerate}
\item $\mathcal P^F_{k,q}$ is a digital $(0,4k,4)$-net over $\Z/q\Z$ for every odd $q\ge3$ and every $k\ge1$; every box minor of its generator matrices over $\Z[\frac12]$ is $\pm2^i$, every cell holds exactly one point, and $h\le2N^{-1/4}$;
\item for every odd $q\ge17$ and every $k\ge1$ the toroidal (hence also Euclidean) minimum distance satisfies $\delta N^{1/4}\ge1-2/q$;
\item for every odd $q\ge17$ and every $k\ge2$, toroidal or Euclidean, $\delta N^{1/4}<1-2/q+8/q^2$.
\end{enumerate}
So for $k\ge2$ the first-order deficit of these nets is exactly $2$.
\end{proposition}

\begin{proof}
\leanok
\uses{lem:Ss, prop:Gs, def:Ns, cor:tsk, prop:F, prop:Dprime, cor:ceils}
(1) Exact interpolation shows that every column of every partial sum $M^{(j)}$ is admissible (Proposition~\ref{prop:F}): its polynomial $g_{d,j}$ has degree $j$, coefficients in $\Z[\frac12]$ and leading coefficient $\pm2^i$. So the residual problem is solvable over $R'=\Z[\frac12]$, and Corollary~\ref{cor:tsk} gives the net property with box minors in $\Z[\frac12]^\times=\{\pm2^i\}$, units modulo every odd $q$; a cell has diameter $2q^{-k}=2N^{-1/4}$.

(2) (H1) and (H2) of Lemma~\ref{lem:Ss} hold by Proposition~\ref{prop:Gs} and Definition~\ref{def:Ns}, with $L=\min(k,3)$. Full enumeration of the $2400$ configurations gives $K(F)=2$ and $\sigma_{\max}(F)=7$, so $q_{\mathrm{tail}}=23$ and Proposition~\ref{prop:Dprime} gives $\mathrm{LB}\ge1-2/q$ for $q\ge23$ and every $k\ge1$. For $q=17,19,21$ the exact relaxation (with $L=3$, which bounds those with $L<3$) is $\mathrm{LB}=(q-2)/q$. The bound is attained: the one-axis step along $\infty$ without carry has $\sigma=(-2,1,1,1)$, so $b=q-2$ for the moving coordinate and $b=0$ for the others.

(3) As $|F_0|\le2<q/2$, the geometric layer is $\tilde M=F_0$, and column $\infty$ of $F_0$ is $(-2,1,1,1)$. This is the sharpness argument of Theorem~A (Section~\ref{sec:Ustar}). Equivalently, it is Corollary~\ref{cor:ceils}(i) with $a=\infty$ without its term $q^{1-k}$: that term accounts for the deeper offset digits of the moving coordinate, and here they do not change, by (H2). So the own-axis step $e_\infty$ from a cell whose last digit of $n_\infty$ is not $q-1$ and whose first offset digit $r_\infty$ is at least $2$ moves coordinate $\infty$ by exactly $1-2/q$ cells, and as $\sum_{u\ne\infty}|F_0[u][\infty]|=3<q$, the union bound of that argument leaves such a cell at which every other coordinate $u$ moves by less than $(|F_0[u][\infty]|+1)/q$. Hence $c^2<(1-2/q)^2+12/q^2$ for every $k\ge2$, toroidal and Euclidean (the pair lies inside the unit cube). With $\alpha=1-2/q\ge15/17$ we get $c<\alpha+6/(\alpha q^2)\le1-2/q+\frac{102}{15}/q^2<1-2/q+8/q^2$.
\end{proof}

\begin{remark}[Why the entries have size two]\label{rem:pm1}
\lean{Z3NetF.pm1_ring, Z3NetF.pm1}\leanok
\uses{lem:CH}
\emph{For any four places with unit differences and every base $q\ge2$ other than $3$, no Cauchy layer $\lambda_c\mu_dw_{cd}$ ($c\ne d$) has all twelve off-diagonal entries in $\{\pm1\}$ modulo $q$.} Indeed, the two terms of Lemma~\ref{lem:CH}(X) would be $\pm1$, and $\pm1\pm1=1$ forces $1=0$ or $3=0$ in $\Z/q\Z$. So in dimension four no net in the frame orbit (the nets obtained from a frame by scrambles and gauge moves, Section~\ref{sec:frames}) has a $\pm1$ layer like the layer $M$ of dimension three; this does not decide whether $K=1$ is possible.
\end{remark}


\section{The case \texorpdfstring{$k=1$}{k=1}}\label{sec:k1}

\begin{proposition}[Separation at $k=1$]\label{prop:k1}
\lean{Z3NetF.Ss_sep_rowScale, Z3NetF.k1_sep, Z3NetF.k1_sep_eucl, Z3NetF.k1_tail, Z3NetF.k1_transfer, Z3NetF.k1_Ns_sep, Z3NetF.k1_ceil, Z3NetF.k1_ceil_eucl, Z3NetF.k1_ceil_real}\leanok
\uses{thm:A, lem:Ss, def:Ns, prop:D, def:K, prop:Dprime}
Let $q\ge3$ and let $\mathcal P$ be a digital net over $\Z/q\Z$ in graph form with depths $(1,\dots,1)$ ($m=s$, $N=q^s$) whose first offset rows form an integer matrix $F_0$ (diagonal included, reduced modulo $q$) and whose deeper rows vanish in their own column (the offset digits after the first do not depend on the point's own cell digit: (H2) of Lemma~\ref{lem:Ss}). Let $\mathrm{LB}_1(q)$ be the relaxation $\mathrm{LB}$ of Lemma~\ref{lem:Ss} for the design $(F_0)$ ($L=1$, $\sigma=F_0e$), and $K_1=K((F_0))$, $\sigma_{\max,1}=\sigma_{\max}((F_0))$ (Definition~\ref{def:K}). Then:
\begin{enumerate}
\item $\delta N^{1/s}\ge\mathrm{LB}_1(q)$ (toroidal, hence Euclidean), and $\mathrm{LB}_1(q)\ge1-K_1/q$ for $q\ge q_{\mathrm{tail}}(K_1,\sigma_{\max,1})$;
\item for a design $F$ realized by Definition~\ref{def:Ns}, $\mathrm{LB}_1(q)$ is at least the relaxation of $F$ at the same $q$ for every $k\ge2$; so the lower bounds of Propositions~\ref{prop:D} and~\ref{prop:Dprime} and of Theorem~\ref{thm:A}(2) hold at $k=1$;
\item let $a$ be a coordinate with $-F_0[a][a]=K_1$ and let $q>2\max|F_0|$. If $\gcd(\det F_0,q)=1$, or if $\mathcal P$ is a $(0,s,s)$-net with $s\ge4$ and $\sum_{u\ne a}|F_0[u][a]|<q$, then
\[
\delta N^{1/s}<\Bigl(\Bigl(1-\frac{K_1}q\Bigr)^2+\sum_{u\ne a}\frac{(|F_0[u][a]|+1)^2}{q^2}\Bigr)^{1/2}
\]
for the toroidal distance, in the second case also for the Euclidean one. With (1), $\delta N^{1/s}=1-K_1/q+O(q^{-2})$.
\end{enumerate}
\end{proposition}

\begin{proof}
\leanok
\uses{lem:nie, lem:Ss, def:Ns, prop:Dprime}
(1) The hypotheses are (H1) and (H2) of Lemma~\ref{lem:Ss} with $k=L=1$; then Proposition~\ref{prop:Dprime}. (2) At $k=1$ Definition~\ref{def:Ns} realizes the truncated design $(F_0)$ with vanishing deeper diagonal rows, and its configurations are those of $F$ with all $j_d=0$, with the same $\sigma$ and jitter (after Lemma~\ref{lem:Ss}). (3) As $q>2\max|F_0|$, the centred layer is $F_0$. Take the step $e_a$ from a cell with $r_a\ge K_1$ and $r_u+F_0[u][a]\in[0,q)$ for every $u\ne a$. Then $\Delta r_a=-K_1$ exactly and $\Delta f_a=0$ (the deeper rows of $a$ vanish on $n_a$), so coordinate $a$ moves by exactly $1-K_1/q$ cells, and every $u\ne a$ by $(F_0[u][a]+\Delta f_u)/q$, of absolute value less than $(|F_0[u][a]|+1)/q$. Such a cell exists: if $\gcd(\det F_0,q)=1$, $n\mapsto r=F_0n$ is a bijection of $(\Z/q\Z)^s$, so $r$ can be prescribed. In the second case, for each $u\ne a$ the box with depths $2$ on $a$ and on $u$, $0$ on two further coordinates and $1$ elsewhere (total $s$) makes $(n_a,r_a,r_u)$ uniformly distributed (Lemma~\ref{lem:nie}); among the cells with $n_a\ne q-1$ and $r_a\ge K_1$ the fraction on which $r_u$ wraps is $|F_0[u][a]|/q$, and the union bound leaves a cell as $\sum_{u\ne a}|F_0[u][a]|<q$. With $n_a\ne q-1$ the pair does not cross the boundary of the cube, so the bound holds for the Euclidean distance too.
\end{proof}

\begin{proposition}[Designs for $k=1$; computer-assisted]\label{prop:k1designs}
\lean{Z3NetF.k1designs_frame_K, Z3NetF.k1designs_frame6_cauchy, Z3NetF.k1designs_free_K, Z3NetF.k1designs_free5_net, Z3NetF.k1designs_free6a_net, Z3NetF.k1designs_free6b_net, Z3NetF.k1designs_frame5_cauchy, Z3NetA.k1designs_frame4_Ustar}\leanok
\begin{enumerate}
\item \emph{Frame layers.} With its diagonal re-chosen for $k=1$, the geometric layer $F_0$ of $U^*_4$ has $K_1=2$, and that of the $s=5$ design of Table~\texttt{tab:Kdesigns} has $K_1=3$; a choice of the row and column scalars of the Cauchy layer gives $K_1=6$ for $s=6$.
\item \emph{Frameless layers.} Integer layers $F_0$ with entries of size at most $3$, chosen directly, without a frame (Appendix~\texttt{app:designs}), completed by deeper rows chosen per prime factor of $q$, give digital $(0,s,s)$-nets over $\Z/q\Z$ with $K_1=2$ for $s=5$ at every $q$ with $\gcd(q,30)=1$, and with $K_1=3$ for $s=6$ at every $q$ with no prime factor in $\mathcal E_6=\{2,3,5,7,11,13,17,23,31\}$, with $\delta N^{1/s}\ge0.8$ for every such $q\ge13$ ($s=5$) resp.\ $q\ge19$ ($s=6$).
\end{enumerate}
\end{proposition}

\begin{proof}
\leanok
\uses{lem:nie, lem:box, prop:k1}
(1) is a finite computation of $K_1$. (2) At $k=1$ the frame is not needed. By Lemma~\ref{lem:box}, a box minor of a graph form with depths $(1,\dots,1)$ is the determinant of stacked offset rows restricted to the coordinates of depth $0$; those that use only the first offset rows are the minors $F_0[T,Z]$ with disjoint $T$, $Z$. If a prime $p$ divides none of them, every box minor is a nonzero polynomial over $\F_p$ in the entries of the deeper rows (it is not the zero polynomial: take the deeper rows it uses to be unit vectors that complete the independent rows of $F_0$), of degree at most the number of deeper rows it uses, and by the Schwartz--Zippel lemma deeper rows with all box minors nonzero exist modulo $p$ once $p$ exceeds the sum $\mathrm{SZ}(s)$ of these degrees ($\mathrm{SZ}(5)=105$, $\mathrm{SZ}(6)=504$); below it we constructed them explicitly for each admissible prime and certified $t=0$. Any integer lift of such a completion modulo $p$ works modulo $p^a$: every box minor is an integer polynomial in the entries, so it is congruent modulo $p$ to the minor modulo $p$, and a determinant prime to $p$ is a unit modulo $p^a$. The Chinese remainder theorem combines the lifts for the prime powers of $q$, and Lemma~\ref{lem:nie} gives the net property. The separation bounds are Proposition~\ref{prop:k1} with exact scans below its tail. Computer-assisted are the layers, $K_1$, the per-prime completions up to $\mathrm{SZ}(s)$ and the finite scans.
\end{proof}


\chapter{Why the deficit grows with the dimension}\label{sec:deficitgrowth}

\begin{proposition}[Pair ceiling]\label{prop:P2}
\lean{Z3NetF.P2_ceiling, Z3NetF.P2_forward, Z3NetF.P2_forward_C}\leanok
Let $s\ge4$, $3k+4\le sk$ (so that the boxes of the proof, of total depth $3k+4$, fit into $sk$ digits), $a\ne b$, $\varepsilon=\pm1$, and suppose $\sigma_a>0$, $\varepsilon\sigma_b<0$ and $\sum_c|\sigma_c|<q$. Then
\[
c^2<\Bigl(1-\frac{|\sigma_b|}q+\frac1{q^2}\Bigr)^2+\frac{(\sigma_a+1)^2+\sum_{u\ne a,b}(|\sigma_u|+1)^2}{q^2}.
\]
If instead $\sigma_a>0$ and $\varepsilon\sigma_b>0$, then $c^2<\sum_c(|\sigma_c|+1)^2/q^2\le s(2C+1)^2/q^2$, $C=\max|\tilde M|$.
\end{proposition}

\begin{proof}
\leanok
Let $e=\e_a+\varepsilon\e_b$ be the step of the proposition, from the cell $n$ to the cell $n+e$. Write $\rho'_b$ and $\rho''_b$ for the second offset digit of $b$ at $n$ and at $n+e$; along $e$ it changes by a constant $\sigma'_b$ modulo $q$. Count the cells $n$ with the following properties:
\begin{enumerate}
\item[(a)] the last digits of $n_a$ and $n_b$ cause no carry along $e$;
\item[(b)] $r_a+\sigma_a\ge q$, so that $a$ wraps ($\sigma_a$ values of $r_a$);
\item[(c)] $r_b$ does not wrap ($q-|\sigma_b|$ values);
\item[(d)] $\rho'_b$ lies in a nonempty set $\mathcal V_b$ of digit values chosen so that $\varepsilon(\rho''_b-\rho'_b)\le0$;
\item[(e)] $r_u+\sigma_u\in[0,q)$ for every coordinate $u\notin\{a,b\}$.
\end{enumerate}
For one $u$ at a time, the digits involved form a box with depths $k+1$ on $a$, $k+2$ on $b$, $k+1$ on $u$ and $0$ elsewhere, of total $3k+4\le sk$, so they are jointly uniform; the conditions (a)--(d) have positive density and the fraction of cells violating (e) at $u$ is $|\sigma_u|/q$, so some cell satisfies all of them, because $\sum_u|\sigma_u|<q$. For such a cell, $|\Delta_a|<(\sigma_a+1)/q$, $|\Delta_u|<(|\sigma_u|+1)/q$ and $\varepsilon\Delta_b<1-|\sigma_b|/q+1/q^2$. The second part of the proposition (the case $\varepsilon\sigma_b>0$) uses, in the same way, the boxes with depths $k+1$ on $a$, $b$ and $u$ and $0$ elsewhere (total $3k+3\le sk$), with $r_a$ and $r_b$ both wrapping, so that $|\Delta_a|<(\sigma_a+1)/q$ and $|\Delta_b|<(|\sigma_b|+1)/q$. The pair lies inside the cube, so the bound holds for the toroidal and the Euclidean distance.
\end{proof}

\begin{definition}\label{def:Ktrue}
\lean{Z3NetF.Ktrue_ceiling}\leanok
Let $\tilde M$ be a layer with nonzero diagonal. If some diagonal entry is positive, or some carry-free step $\e_a+\varepsilon\e_b$ has both moving coordinates forward, let $K_{\mathrm{pair}}(\tilde M)=\infty$. Otherwise $K_{\mathrm{pair}}(\tilde M)$ is the largest of
\begin{enumerate}
\item[(a)] $|\tilde M_{aa}|$ over the coordinates $a$, and
\item[(b)] $|\sigma_b|$ over the carry-free steps $\e_a+\varepsilon\e_b$ (as above) in which $a$ is forward and $b$ is the only backward coordinate, with $\varepsilon\sigma_b<0$.
\end{enumerate}
\end{definition}

\begin{theorem}[The deficit grows with $s$]\label{thm:L}
\lean{Z3NetF.L_Kpair, Z3NetF.L_ceiling, Z3NetF.L_pm1}\leanok
\uses{def:Ktrue}
Let $\mathcal P$ be a digital $(0,sk,s)$-net over $\Z/q\Z$ with $q\ge2$, $s\ge4$ and $k\ge1$, whose layer $\tilde M$ has no zero diagonal entry. Then $K_{\mathrm{pair}}(\tilde M)\ge\sqrt{(s-3)/8}$ (Definition~\ref{def:Ktrue}). Hence, by (\texttt{eq:Kpair}) (which includes the case $K_{\mathrm{pair}}=\infty$), a family of such nets with $k\ge3$ and $3k+4\le sk$ whose layers stay bounded, $|\tilde M|\le C$, has $\delta N^{1/s}<1-\sqrt{(s-3)/8}/q+6sC^2/q^2$ for $q>4sC$.
\end{theorem}

\begin{proof}
\leanok
\uses{lem:box}
Let $K=K_{\mathrm{pair}}<\infty$. Every diagonal entry is $-\alpha_a$ with $1\le\alpha_a\le K$. Every pair has a small entry: if $|x|=|\tilde M_{ab}|>K$, put $\varepsilon=\operatorname{sgn}x$; then $\sigma_a=-\alpha_a+|x|>0$, so $b$ is backward, $\varepsilon\sigma_b=\varepsilon y-\alpha_b\le0$ for $y=\tilde M_{ba}$, and if it is negative the definition gives $\alpha_b-\varepsilon y\le K$; so $|y|\le\max(\alpha_b,K-\alpha_b)\le K$. Orient $c\to d$ when $|\tilde M_{cd}|\le K$, and let $\mathrm{out}(c)$ be the number of $d\ne c$ with $c\to d$. Every pair carries at least one arc, so $\sum_c\mathrm{out}(c)\ge\binom s2$, and by convexity $\sum_c\binom{\mathrm{out}(c)}2\ge s\binom{(s-1)/2}2=s(s-1)(s-3)/8$. The left side counts the triples $(c,\{d,d'\})$ with $c\to d$ and $c\to d'$, so averaging over the $\binom s2$ pairs $\{d,d'\}$, some two columns $d,d'$ have at least $(s-3)/4$ common rows $c\notin\{d,d'\}$ with $|\tilde M_{cd}|,|\tilde M_{cd'}|\le K$. By $t=0$ and Lemma~\ref{lem:box}, the box with depths $k+1$ on two such rows $c,c'$, $k-1$ on $d,d'$ and $k$ elsewhere makes every $2\times2$ minor $\tilde M_{cd}\tilde M_{c'd'}-\tilde M_{cd'}\tilde M_{c'd}$ a unit, hence nonzero, and the $1\times1$ boxes make the entries nonzero. So these rows give pairwise non-proportional vectors in $([-K,K]\setminus0)^2$, of which there are at most $2K^2$ directions: $(s-3)/4\le2K^2$.
\end{proof}

\begin{theorem}[Frame designs have $K_{\mathrm{pair}}\ge2$ for $s\ge5$]\label{thm:O5}
\lean{Z3NetF.O5_Kpair, Z3NetF.O5_sep}\leanok
\uses{def:Ns}
Let $s\ge5$ and let $F$ be an integer design (Section~\ref{sec:prelim}) whose residual problem is solvable over $\mathbb{Q}$ for a frame of rational places, with nonzero diagonal of $F_0$. Then $K_{\mathrm{pair}}(F_0)\ge2$; hence for $k\ge3$, $3k+4\le sk$ and every $q>4sC$ coprime to $DD'$ (Definition~\ref{def:Ns}), $C=\max|F_0|$, $\delta N^{1/s}<1-2/q+6sC^2/q^2$ (for $q>2C$ the layer is $\tilde M=F_0$, and (\texttt{eq:Kpair}) applies, including the case $K_{\mathrm{pair}}=\infty$).
\end{theorem}

\begin{proof}
\leanok
\uses{lem:CH}
Suppose $K_{\mathrm{pair}}\le1$. Then every diagonal entry is $-1$, and the definition of $K_{\mathrm{pair}}$ forces the \emph{pair rule}: each pair $(\tilde M_{ab},\tilde M_{ba})$ is in $\{\pm1\}^2$, or one entry $x$ has $|x|\ge2$ and the other is $\operatorname{sgn}x$ (the step $e_a+\varepsilon e_b$ has $\sigma_a=-1+\varepsilon x$, $\sigma_b=y-\varepsilon$ for $(x,y)=(\tilde M_{ab},\tilde M_{ba})$; if $|x|\ge2$ and $\varepsilon=\operatorname{sgn}x$, $a$ is forward, so $b$ is backward with $|\sigma_b|\le1$ whenever $\varepsilon\sigma_b<0$, which with $y\ne0$ forces $y=\operatorname{sgn}x$). So pairs in one class of Lemma~\ref{lem:CH}(C) have $\pi=-1$, and pairs across classes have an integer $\pi=P\ge1$. A class of size $\ge4$ contradicts (H) ($1=1+1$), so for $s=5$ the classes are $X=\{x_1,x_2,x_3\}$ and $Y=\{y_1,y_2\}$. For $\{x_i,x_j,y_1,y_2\}$, (H) gives $1=1/u+1/v$ with integers $u^2=P_{i1}P_{j2}$, $v^2=P_{i2}P_{j1}$; squaring shows $u,v\in\mathbb{Q}$, hence integers, hence $u=v=2$; so $P_{i1}P_{j2}=4$ for all $i\ne j$, which forces $P_{i1}=\alpha$ for all $i$. Then the three pairings of $\{x_1,x_2,x_3,y_1\}$ all have magnitude $\alpha^{-1/2}$, and $1=2$ is false. For $s=6$ a class of size $\ge4$ again contradicts (H), so both classes have size $3$, and any five of the coordinates form classes of sizes $3$ and $2$, to which the case $s=5$ applies; for $s\ge7$ some class has size $\ge4$. Since $K_{\mathrm{pair}}$ is an integer, $K_{\mathrm{pair}}\ge2$.
\end{proof}


\section{Lower bounds for the deficit of frame designs}\label{sec:Lstar}

\begin{lemma}[Orientation]\label{lem:O}
\lean{Z3NetF.O_orient}\leanok
Let $F$ be a design with $K=K(F)<\infty$. Then for every pair $a\ne b$ of coordinates, $|M^{(j)}[a][b]|\le K$ for all $j$, or $|M^{(j)}[b][a]|\le K$ for all $j$.
\end{lemma}

\begin{proof}
\leanok
One-axis steps give $0\le\alpha^{(j)}_c\le K$ for all $c$ and $j$. Suppose $u=M^{(j_1)}[a][b]$ and $v=M^{(j_2)}[b][a]$ have $|u|,|v|>K$. Take the step along $a$ and $b$ with carry depths $j_a=j_2$, $j_b=j_1$ and $e_ae_b=\varepsilon:=\sgn u$. Then $e_a\sigma_a=-\alpha^{(j_2)}_a+|u|>0$, so $a$ is forward, and $e_b\sigma_b=\varepsilon v-\alpha^{(j_1)}_b$. If $\sgn v=\sgn u$, then $b$ is forward too and the cost is $+\infty$; otherwise $b$ is the only backward coordinate, of cost $\alpha^{(j_1)}_b+|v|+1>K$. Both contradict the definition of $K$.
\end{proof}

\begin{proposition}[Constant row scalars]\label{prop:E}
\lean{Z3NetF.E_lcm, Z3NetF.E_two_pow, Z3NetF.E_Ustar}\leanok
\uses{prop:F}
Let $s\ge4$ and let $F$ be an integer design with $L\ge1$ levels reachable over $\mathbb Q$ (Proposition~\ref{prop:F}) from the places $\infty,0,1,\dots,s-2$ with row scalars $\lambda'$ that take one value on all finite places ($\lambda'_\infty$ arbitrary). If $K(F)<\infty$, then $K(F)\ge\operatorname{lcm}(1,\dots,s-3)$.
\end{proposition}

\begin{proof}
\leanok
\uses{prop:F, lem:O}
Let $\nu=s-2$ and $\lambda$ the common value. By Proposition~\ref{prop:F} at carry depth $0$, $M^{(0)}[c][d]=\lambda'_cg_{d,0}w_{cd}$ with $g_{d,0}\ne0$; put $\mu_d=\lambda g_{d,0}$. In column $0$ the finite rows $c=1,\dots,\nu$ give the integers $\mu_0/c$, so $\operatorname{lcm}(1,\dots,\nu)$ divides $\mu_0$ and $|M^{(0)}[1][0]|=|\mu_0|\ge\operatorname{lcm}(1,\dots,\nu)$. In column $1$ the rows $c\in\{0,2,\dots,\nu\}$ give the integers $\mu_1/(c-1)$, $c-1\in\{-1,1,\dots,\nu-1\}$, so $|M^{(0)}[0][1]|=|\mu_1|\ge\operatorname{lcm}(1,\dots,\nu-1)$. Lemma~\ref{lem:O} for the pair $\{0,1\}$ gives $K(F)\ge\operatorname{lcm}(1,\dots,\nu-1)$.
\end{proof}

\begin{lemma}[A cross-ratio count]\label{lem:C}
\lean{Z3NetF.C_count, Z3NetF.C_count_sharp}\leanok
\uses{prop:Ns}
Take the places $\infty,0,1,\dots,\nu$, two distinct places $d,d'$ and a rational $\varrho\ne0$, and let $w_{cd}$ be as in Proposition~\ref{prop:Ns}. If $\mathcal C$ is a nonempty set of places $c\notin\{d,d'\}$ with $w_{cd}/w_{cd'}\in\varrho\,Q_K$ for all $c\in\mathcal C$, then $|\mathcal C|\le1+2K\operatorname{div}^*(\nu)\operatorname{div}^*(\nu K)^2$.
\end{lemma}

\begin{proof}
\leanok
A rational lies in $Q_K$ if and only if its numerator and denominator in lowest terms are at most $K$ in absolute value. Put $\varphi(c)=w_{cd}/w_{cd'}$.
\emph{One of $d,d'$ is $\infty$.} As $\varphi\mapsto1/\varphi$ and $Q_K=Q_K^{-1}$, we may take $d'=\infty$; then $c$ is finite and $\varphi(c)=1/t$ with $t=c-d\in[-\nu,\nu]\setminus0$, and the condition says $tQ'/P'\in Q_K$ for $\varrho^{-1}=P'/Q'$ in lowest terms. With $e=\gcd(t,P')$ the lowest terms of $tQ'/P'$ are $(tQ'/e)/(P'/e)$, so $t=et'$ with $e\mid P'$ and $|t'|\le K$: at most $2K\operatorname{div}(P')$ values. Some $c_0\in\mathcal C$ gives $\varrho^{-1}\in t_0Q_K$, so $|P'|\le \nu K$.
\emph{$d,d'$ finite.} Row $\infty$ has $\varphi=1$ and contributes at most one element. For finite $c$ put $u=c-d$ and $g=d-d'\ne0$, so $\varphi(c)=(u+g)/u$; with $h=\gcd(u,g)$, $u=hu'$, $g=hg'$, this is $(u'+g')/u'$ in lowest terms. Let $\varrho=P/Q$ in lowest terms, $A=\gcd(Q,u')$ and $B=\gcd(P,u'+g')$. Then $\varphi(c)/\varrho=\bigl(\frac QA\cdot\frac{u'+g'}B\bigr)\big/\bigl(\frac PB\cdot\frac{u'}A\bigr)$ is in lowest terms (each pair of factors is coprime, by $\gcd(P,Q)=\gcd(u',u'+g')=1$ and $\gcd(Q/A,u'/A)=\gcd(P/B,(u'+g')/B)=1$), so $u'+g'=B\beta$ with $1\le|\beta|\le K$. For fixed $(h,A,B)$, $\beta$ determines $u'=B\beta-g'$, hence $c=d+hu'$: at most $2K\operatorname{div}(g)\operatorname{div}(Q)\operatorname{div}(P)$ places, with $|g|\le \nu$ and, as above, $|P|,|Q|\le \nu K$.
\end{proof}

\begin{theorem}[Frame designs]\label{thm:Lstar}
\lean{Z3NetF.PGLImage_self, Z3NetF.Lstar_bound}\leanok
\uses{prop:F}
Let $s\ge5$ and let $F$ be an integer design with $L\ge1$ levels reachable over $\mathbb Q$ (Proposition~\ref{prop:F}) from the places $\infty,0,1,\dots,s-2$, or from their image under an element of $PGL_2(\mathbb Q)$ that contains $\infty$ (as the places of every frame do, Section~\ref{sec:dim4}), with arbitrary row scalars. If $K=K(F)<\infty$, then
\[
\frac{s-3}4\ \le\ 1+2K\,\operatorname{div}^*(s-2)\,\operatorname{div}^*\bigl((s-2)K\bigr)^2 .
\]
\end{theorem}

\begin{proof}
\leanok
\uses{prop:F, thm:L, lem:O, lem:C}
By Lemma~\ref{lem:O} at carry depth $0$ every pair $\{c,d\}$ has an entry $|M^{(0)}[c][d]|\le K$; orient $c\to d$ for such an entry. The convexity count in the proof of Theorem~\ref{thm:L} gives columns $d\ne d'$ and a set $\mathcal C$ of at least $(s-3)/4$ rows $c\notin\{d,d'\}$ with $1\le|M^{(0)}[c][d]|,|M^{(0)}[c][d']|\le K$ (entries at carry depth $0$ are nonzero integers). By Proposition~\ref{prop:F}, $M^{(0)}[c][d]=\lambda'_cg_{d,0}w_{cd}$, so $w_{cd}/w_{cd'}=(g_{d',0}/g_{d,0})\,M^{(0)}[c][d]/M^{(0)}[c][d']$ lies in $\varrho\,Q_K$ with $\varrho=g_{d',0}/g_{d,0}$, and Lemma~\ref{lem:C} with $n=s-2$ bounds $|\mathcal C|$. For the image of the places under $g\in PGL_2(\mathbb Q)$, in normalized homogeneous coordinates every $D_{cd}$, hence every $w_{cd}$, is multiplied by a factor depending only on $c$ times one depending only on $d$; the first cancels in $w_{cd}/w_{cd'}$ and the second changes only $\varrho$.
\end{proof}


\chapter{Drift and jitter}\label{sec:drift}


\section{Excess, drift and jitter}\label{sec:excess}

\begin{lemma}[Drift and jitter]\label{lem:drift}
\lean{Z3NetC.colW_column, Z3NetC.drift_step, Z3NetC.drift_cycle, Z3NetC.drift_range, Z3NetC.drift_eta, Z3NetC.drift_net, Z3NetC.drift_excess}\leanok
\uses{def:col}
Let $q\ge2$ with $3\nmid q$, let $w\in\Z^3$ with $w_x+w_y+w_z=0$ and $w_c\equiv-q\pmod3$, $a=\frac13(q\mathbf 1+w)$, and $b\in\Z^3$ with $b_x+b_y+b_z\equiv0\pmod q$ (the third-digit row, as in Definition~\ref{def:col}); let $X$ be the digital net over $\Z/q\Z$ whose point of cell $n$ has, in every coordinate $c$, the digits $n_c$, $a\cdot n$, $b\cdot n$ (a \emph{digital column net}; $\mathcal C_q$ is the case of Definition~\ref{def:col}), and let $\bar b_i\in\{0,\dots,q-1\}$ be the residue of $b_i$ (a third-digit step, not a geometric digit). Put $\varepsilon_i=\frac13qw_i+\bar b_i$.
\begin{enumerate}
\item A step along $\e_i$ has $\Delta_i\equiv\frac13q^2+\varepsilon_i$ if the third digit $b\cdot n$ does not wrap (does not pass from $q-1$ to $0$) and $\Delta_i\equiv\frac13q^2+\varepsilon_i-q$ if it wraps, modulo $q^2$; if $\bar b_i\ne0$ both occur.
\item Along each cycle of $q$ steps in the direction $\e_i$ the third digit wraps exactly $\bar b_i$ times. If moreover $|\varepsilon_i|<\frac12q^2$ and $|\varepsilon_i-q|<\frac12q^2$ (for instance if $|w_i|\le2$ and $q\ge4$), the excesses of the cycle are $\varepsilon_i$ and $\varepsilon_i-q$, and they sum to $\frac13q^2w_i$.
\item If $q\ge10$ and $|w_c|\le2$, then $\eta=\max_i\varepsilon_i$, so $d^2=q^4-2q^2\eta+3\eta^2$.
\item If $X$ is a $(0,3,3)$-net, then $\varepsilon_x+\varepsilon_y+\varepsilon_z\ge q$, and $\max_i\varepsilon_i\ge\frac13q$; if moreover $q\equiv1\pmod3$, then $\max_i\varepsilon_i\ge\frac23q$.
\end{enumerate}
\end{lemma}

\begin{proof}
\leanok
\uses{lem:box}
(1) The digit $a\cdot n$ changes by $a_i$ and $b\cdot n$ by $\bar b_i$ modulo $q$, so $\Delta_i\equiv qa_i+\bar b_i-q\omega=\frac13q^2+\frac13qw_i+\bar b_i-q\omega$ modulo $q^2$, with $\omega\in\{0,1\}$ the wrap of the third digit (a wrap of $a\cdot n$ changes $\Delta_i$ by $q^2$). (2) Along a cycle the third digit takes the values $b\cdot n+t\bar b_i\bmod q$, $t=0,\dots,q-1$, and returns after a total increase $q\bar b_i$, so it wraps exactly $\bar b_i$ times and does not wrap $q-\bar b_i\ge1$ times; this also shows that both cases of (1) occur if $\bar b_i\ne0$. Under the range condition the representatives $\frac13q^2+\varepsilon_i$ and $\frac13q^2+\varepsilon_i-q$ of (1) lie in the window $(-\frac16q^2,\frac56q^2]$, so the excesses are $\varepsilon_i$ and $\varepsilon_i-q$, and those of the cycle sum to $q\varepsilon_i-q\bar b_i=\frac13q^2w_i$. (Without it the identity can fail: for $q=5$, $w=(-8,4,4)$ and $b=(0,1,-1)$ the excesses of a cycle in the direction $\e_x$ sum to $\frac{175}3$, not $-\frac{200}3$.) (3) Then $-\frac53q\le\varepsilon_i-q<\varepsilon_i\le\frac53q$, and $\frac53q\le\frac16q^2$ for $q\ge10$, so every excess lies in $[-\frac16q^2,\frac13q^2]$ and the largest one is $\max_i\varepsilon_i$ (attained, since the third digit does not wrap on some step of each cycle, by (2)). (4) As $w_x+w_y+w_z=0$, $\sum_i\varepsilon_i=\sum_i\bar b_i\equiv\sum_ib_i\equiv0\pmod q$, and $\bar b\ne0$: otherwise $b\equiv0$ and the box with depth $3$ on a coordinate, whose minor is a $2\times2$ minor of $(a,b)$ (Lemma~\ref{lem:box}), would be singular. So $\sum_i\bar b_i\ge q$ and $\max_i\varepsilon_i\ge\frac13q$. For $q\equiv1\pmod3$, $w_c\equiv2\pmod3$ and $\sum_cw_c=0$, so some $w_c\ge2$ and $\varepsilon_c\ge\frac23q$.
\end{proof}


\section{The holonomy, and the optimum among column nets}

\begin{corollary}[Column ceiling]\label{cor:colceil}
\lean{Z3NetC.colceil}\leanok
\uses{thm:colmin}
Every column point set in base $q\ge2$ has $g\le\frac13q^2$ in Theorem~\ref{thm:colmin}, i.e. toroidal minimum distance at most $q^2$ at integer scale $q^3$: $\delta N^{1/3}\le1$.
\end{corollary}

\begin{proof}
\leanok
If $g>\frac13q^2$, every step value $\Delta_i([n])$ has distance more than $\frac13q^2$ from $\frac23q^2+q^2\Z$, so its residue modulo $q^2$ lies in the open interval $(0,\frac13q^2)$. The three steps $[n]\to[n+\e_x]\to[n+\e_x+\e_y]\to[n+\mathbf 1]=[n]$ have step values that sum to $0$, so their residues sum to a multiple of $q^2$; but that sum lies in $(0,q^2)$.
\end{proof}

\begin{theorem}[Column optimum]\label{thm:O}
\lean{Z3NetC.colO_net_iff, Z3NetC.colO_even, Z3NetC.colO_opt}\leanok
\uses{thm:col}
Let $q\ge2$.
\begin{enumerate}
\item A column point set, with the set of columns $\mathbb T$ and the heights $\varphi$ of Section~\ref{sec:m3}, is a $(0,3,3)$-net in base $q$ if and only if $\varphi$ is a bijection from $\mathbb T$ onto $\{0,\dots,q^2-1\}$ and $\lfloor\varphi/q\rfloor$ is injective on every line of $\mathbb T$ in the directions $[\e_x],[\e_y],[\e_z]$.
\item If $q$ is even, no column point set is a $(0,3,3)$-net in base $q$.
\item If $q\ge5$ is odd and $3\nmid q$, every column point set that is a $(0,3,3)$-net in base $q$ has toroidal minimum distance $d$ with $d^2\le d^2(\mathcal C_q)$, the value of Theorem~\ref{thm:col}(2); equivalently $\eta\ge\frac13q+1$ for $q\equiv2$ and $\eta\ge\frac23q$ for $q\equiv1\pmod3$, where $\eta=\frac13q^2-g$.
\end{enumerate}
\end{theorem}

\begin{proof}
\leanok
\uses{thm:colmin, thm:col, cor:colceil}
(1) The point of cell $n$ has the digits $n_c$, $\lfloor\varphi/q\rfloor$, $\varphi\bmod q$ in every coordinate. The box of shape $(1,1,1)$ holds one point per cell. A box with depth $3$ on $c$ prescribes $n_c$ and $\varphi$; the $q^2$ cells with a given $n_c$ meet every column exactly once ($n+t\mathbf 1$ for the one $t$ with $n_c+t$ given), so these boxes hold one point each if and only if $\varphi$ is a bijection. A box with depth $2$ on $c$, $1$ on $c'$ and $0$ on $c''$ prescribes $n_c$, $n_{c'}$ and $\lfloor\varphi/q\rfloor$; the $q$ cells with given $n_c,n_{c'}$ form a line in the direction $\e_{c''}$, whose columns form a line of $\mathbb T$ in the direction $[\e_{c''}]$, and every such line of $\mathbb T$ arises. These are all box shapes.

(2) Identify $\mathbb T$ with $(\Z/q\Z)^2$ by $[n]\mapsto(n_x-n_z,n_y-n_z)$; the lines in the directions $[\e_x]$, $[\e_y]$, $[\e_z]$ become the horizontal lines, the vertical lines and the lines of slope $1$ of $(\Z/q\Z)^2$. A level set $\{\lfloor\varphi/q\rfloor=v\}$ has $q$ elements and meets each of these lines once, by (1). So it is the graph of a permutation $\pi$ of $\Z/q\Z$ (once on every horizontal and every vertical line) such that $y\mapsto\pi(y)-y$ is a permutation too (once on every line of slope $1$). Then $\sum_y(\pi(y)-y)\equiv0$, while the sum over all residues is $\frac q2\not\equiv0\pmod q$ for even $q$. In other words, $y\mapsto\pi(y)-y$ is a complete mapping of $\Z/q\Z$, and this is the classical argument that a cyclic group of even order has none \cite[Thm.~1 and its corollary]{Pai51}; more generally, no finite group with a nontrivial cyclic Sylow $2$-subgroup has one \cite[Thm.~5]{HP55}. In the language of Latin squares, (1) says that each level set is a transversal of the cyclic Latin square $(a,b)\mapsto b-a$, and that the array $\lfloor\varphi/q\rfloor$ on $(\Z/q\Z)^2$ is a Latin square orthogonal to it.

(3) Put $\eta=\frac13q^2-g$. By Corollary~\ref{cor:colceil}, $0\le\eta\le\frac13q^2$, and by Theorem~\ref{thm:colmin} $d^2=q^4-2q^2\eta+3\eta^2$, which decreases in $\eta$ on this interval and equals $d^2(\mathcal C_q)$ at $\eta_0=\frac13q+1$ ($q\equiv2$) resp.\ $\eta_0=\frac23q$ ($q\equiv1$). So it suffices to show $\eta\ge\eta_0$. Suppose $\eta<\eta_0$.

\emph{Then $\eta<\frac1{12}q^2$.} For $q\ge8$ this holds because $\eta_0\le\frac1{12}q^2$. For $q=5$ and $7$: $3g=|3U+q^2|$ for an integer $U$, which is not divisible by $3$ since $q^2\equiv1\pmod3$, so $3\eta=q^2-3g\equiv0$ or $2\pmod3$; hence $\eta<\frac83$ gives $\eta\le2<\frac{25}{12}$ at $q=5$, and $\eta<\frac{14}3$ gives $\eta\le4<\frac{49}{12}$ at $q=7$.

\emph{Windows and triangles.} Every step value has distance at least $\frac13q^2-\eta$ from $\frac23q^2+q^2\Z$, so it has a unique representative $u$ in the window $[-\eta,\frac13q^2+\eta]$ modulo $q^2$ (narrower than the window of Section~\ref{sec:excess}). As $[\e_x]+[\e_y]+[\e_z]=0$ in $\mathbb T$, for every column and every order of the three directions the three steps around the triangle $[n]\to[n+\e_i]\to[n+\e_i+\e_{i'}]\to[n]$ have values summing to $0$, so their representatives sum to a multiple of $q^2$ in $[-3\eta,q^2+3\eta]$, i.e. to $0$ or $q^2$. In a triangle with sum $0$ every $u$ lies in $[-\eta,2\eta]$, and in one with sum $q^2$ every $u$ lies in $[\frac13q^2-2\eta,\frac13q^2+\eta]$; as $\eta<\frac1{12}q^2$ these ranges are disjoint, so two triangles that share a step have the same sum. The step $([n],i)$ shares a triangle with $([n+\e_i],i')$ for both $i'\ne i$, and two such moves give $([n+\e_i+\e_{i'}],i)=([n-\e_{i''}],i)$; and the classes $[\e_{i''}]$, $i''\ne i$, generate $\mathbb T$; so all steps are connected through shared triangles, and all triangles have the same sum. Sum $q^2$ means that every step is close to $\frac13q^2$, the configuration of the reflected cube $H\Z^3$; sum $0$ means that every step is close to $0$, the configuration of $\Z^3$.

\emph{Sum $q^2$ (the reflected cube).} Put $\xi=u-\frac13q^2\in[-2\eta,\eta]$, the excess; then $\operatorname{dist}(u,\frac23q^2+q^2\Z)=\frac13q^2-\xi$, so $\eta$ is the largest excess, and every triangle has excess sum $0$. The class $[\e_i]$ has order $q$ in $\mathbb T$, so the steps in the direction $\e_i$ form $q$ cycles $\gamma$ of length $q$. Along a cycle the step values telescope to $0$, so $\sum_{\gamma}u\in q^2\Z$, and the \emph{cycle excess} $\Xi_\gamma=\sum_\gamma\xi\in q^2\Z-\frac13q^3$. Summing the excess over the triangles $[n]\to[n+\e_x]\to[n+\e_x+\e_y]\to[n]$ for all $[n]\in \mathbb T$ counts every step once, so the $3q$ cycle excesses sum to $0$. Also $-2q\eta\le\Xi_\gamma\le q\eta$.
If $q\equiv1\pmod3$, then $\Xi_\gamma=q^2z_\gamma-\frac13q^2$ with integers $z_\gamma$ (as $\frac13q^3\equiv\frac13q^2$), and the cycle excesses sum to $0$ only if $\sum_\gamma z_\gamma=q$. So some $z_\gamma\ge1$, $\Xi_\gamma\ge\frac23q^2$, and $q\eta\ge\frac23q^2$: $\eta\ge\frac23q$, a contradiction.
If $q\equiv2\pmod3$, then $\Xi_\gamma=q^2z_\gamma-\frac23q^2$ and $\sum_\gamma z_\gamma=2q$. Every excess lies in $\Z+\frac23$ (as $u\in\Z$ and $\frac13q^2\in\Z+\frac13$), and so does $\frac13q$; hence $\eta<\frac13q+1$ gives $\eta\le\frac13q$, and then every $\Xi_\gamma\in[-\frac23q^2,\frac13q^2]$, i.e. $\Xi_\gamma\in\{-\frac23q^2,\frac13q^2\}$. A cycle excess $\frac13q^2\le q\eta$ forces $\xi=\eta=\frac13q$ on each step of its cycle, and a cycle excess $-\frac23q^2\ge-2q\eta$ forces $\xi=-2\eta=-\frac23q$ on each step. So every step has $u=\frac13q(q+1)$ or $\frac13q(q-2)$, a multiple of $q$, and $\varphi$ is constant modulo $q$ on $\mathbb T$. This contradicts (1): $\varphi$ takes every value in $\{0,\dots,q^2-1\}$.

\emph{Sum $0$ (the cube).} Now $u\in[-\eta,2\eta]$ and $\operatorname{dist}(u,\frac23q^2+q^2\Z)=u+\frac13q^2$, so $\min u=-\eta$. The cycle sums $\sum_\gamma u$ lie in $q^2\Z$ and add up to $0$. If one of them is nonzero, one is at most $-q^2$; it is at least $-q\eta$, so $\eta\ge q$. Otherwise every cycle closes: the partial sums $S_t$ of $u$ along a cycle are integers with $S_t\equiv\varphi(c_t)-\varphi(c_0)\pmod{q^2}$, where $c_t$ are its columns. By (1) the $q$ values $\varphi(c_t)$ lie in $q$ different blocks $[kq,kq+q)$, so the $q$ values $\varphi(c_0)+S_t$ meet all $q$ blocks modulo $q^2$; an interval of integers of length $\mathrm{rg}$ meets at most $\mathrm{rg}/q+2$ of them, so the range of the walk is $\mathrm{rg}\ge q(q-2)$. Here $\eta>0$, since $\eta=0$ would make every step value $0$ and $\mathrm{rg}=0<q(q-2)$. A closed walk whose up-steps are at most $2\eta$ and whose down-steps are at least $-\eta$ needs at least $\mathrm{rg}/(2\eta)+\mathrm{rg}/\eta$ steps, so $q\ge3\,\mathrm{rg}/(2\eta)$ and $\eta\ge\frac32(q-2)$. Both $q$ and $\frac32(q-2)$ are at least $\frac23q\ge\eta_0$ for $q\ge4$, a contradiction.

Equality holds for $\mathcal C_q$ by Theorem~\ref{thm:col}(2).
\end{proof}


\chapter{One point per cube: cube tilings and the value 1}\label{sec:whylat}

\begin{lemma}[Additive digit functionals]\label{lem:hom}
\lean{Z3NetF.hom_iff}\leanok
Let $q\ge2$, $j\ge1$, $\beta_0,\dots,\beta_{j-1}\in\Z/q\Z$, and $\phi(P)=\sum_{i<j}\beta_i\,d_i(P)$ for $P\in\Z/q^j\Z$ (digits of the representative in $[0,q^j)$). Then $\phi$ is a homomorphism $\Z/q^j\Z\to\Z/q\Z$ if and only if $\beta_i=0$ for every $i\ge1$.
\end{lemma}

\begin{proof}
\leanok
If $\beta_i=0$ for $i\ge1$, then $\phi(P)=\beta_0\,(P\bmod q)$, a homomorphism. Conversely, a homomorphism on the cyclic group $\Z/q^j\Z$ satisfies $\phi(P)=P\,\phi(1)=\beta_0P=\beta_0\,d_0(P)$ in $\Z/q\Z$; at $P=q^i$, $i\ge1$, this gives $\beta_i=\phi(q^i)=\beta_0\,d_0(q^i)=0$.
\end{proof}

\begin{corollary}[Linear offsets]\label{cor:linoff}
\lean{Z3NetF.linoff_a, Z3NetF.linoff_b, Z3NetF.linoff_c}\leanok
Let $q\ge2$.
\begin{enumerate}
\item[(a)] \emph{One level.} Let a set of integer vectors $Y\in\Z^s$ consist of all $Y$ whose lowest base-$q$ digits $(Y_c\bmod q)_c$ satisfy a system of linear equations over $\Z/q\Z$, with solution set $y_0+C$ for a linear code $C\subset(\Z/q\Z)^s$, all higher digits being free. Then the set is the coset $y_0+C+q\Z^s$ of the Construction~A lattice $C+q\Z^s$.
\item[(b)] \emph{Two levels.} For an $s\times s$ matrix $B$ over $\Z/q\Z$ let $\mathcal D_B=\{q\ell+(B\ell\bmod q):\ell\in\{0,\dots,q-1\}^s\}\subset(\Z/q^2\Z)^s$ (the \emph{two-digit set} of $B$; $\ell$ plays the part of the last cell digits), with $B\ell\bmod q$ taken in $[0,q)^s$. Then $\mathcal D_B$ is a subgroup if and only if $Bc\equiv0\pmod q$ for every carry vector $c=(\lfloor(v_i+v'_i)/q\rfloor)_i$ of two vectors $v,v'\in\{B\ell\bmod q\}$. If $B$ is invertible, $\mathcal D_B$ is not a subgroup.
\item[(c)] \emph{Locally a lattice.} Let a digital net over $\Z/q\Z$ be in graph form with depths $(k,\dots,k)$ and geometric layer $\tilde M$, and let the cells $n$, $n'$ differ only in their last level, by $n'-n=E\in\Z^s$ (no carry). If the geometric digits $r$ of $n$ satisfy $r+\tilde ME\in[0,q)^s$, then in cell units the points of $n'$ and $n$ differ by $(I+\tilde M/q)E+(f'-f)/q$, with $f,f'\in[0,1)^s$ the contributions of the deeper digits.
\end{enumerate}
\end{corollary}

\begin{proof}
\leanok
(a) The set of $Y\in\Z^s$ with $Y\bmod q\in y_0+C$ is $y_0+C+q\Z^s$. (b) Let $Y=q\ell+v$ and $Y'=q\ell'+v'$ be in $\mathcal D_B$, and let $c$ be the carry vector of $v,v'$. Then $Y+Y'=q(\ell+\ell'+c)+(v+v'-qc)$ with $v+v'-qc\in[0,q)^s$, and $v+v'\equiv B(\ell+\ell')$; so $Y+Y'\in\mathcal D_B$ if and only if $Bc\equiv0$. A nonempty finite subset of a group that is closed under addition is a subgroup. If $B$ is invertible, its image contains $v=v'=\lceil q/2\rceil\e_i$, whose carry vector is $\e_i$; then $B\e_i\equiv0$ would be needed, which is false. (c) Only the last digits of the cells change, so $r$ changes by $\tilde ME$ modulo $q$, and $r'=r+\tilde ME$ exactly since both sides lie in $[0,q)^s$; the cell digits change by $E$.
\end{proof}

\begin{lemma}[Wrap patterns; rank of the layer]\label{lem:phase}
\lean{Z3NetF.phase_i, Z3NetF.phase_i_free, Z3NetF.phase_ii, Z3NetF.phase_ii_rank}\leanok
\begin{enumerate}
\item[(i)] Let $p$ be a prime and consider a graph form over $\F_p$ with depths $(1,\dots,1)$: coordinate $c$ of the point with index $n\in\F_p^s$ has the digits $n_c$, $r_c(n)=(Bn)_c$ and deeper digits, with $B\in\F_p^{s\times s}$. For $E\in\Z^s$ let $\beta=BE\bmod p\in\{0,\dots,p-1\}^s$ and $n'\equiv n+E$. Then $r(n')=r(n)+\beta-p\omega$ with $\omega_c=[r_c(n)+\beta_c\ge p]$, and $\omega_c=0$ whenever $\beta_c=0$. If $B$ is invertible, then every $\omega\in\{0,1\}^s$ with $\omega_c=0$ whenever $\beta_c=0$ is the wrap vector of some cell $n$.
\item[(ii)] Let a digital $(0,sk,s)$-net over $\Z/q\Z$, $k\ge1$, be in graph form with depths $(k,\dots,k)$ and geometric layer $\tilde M$. For disjoint $T,Z\subset\{1,\dots,s\}$ with $|T|=|Z|$, $\det\tilde M_{T,Z}$ is a unit of $\Z/q\Z$. In particular $\tilde M$ has rank at least $\lfloor s/2\rfloor$ modulo every prime factor of $q$.
\end{enumerate}
\end{lemma}

\begin{proof}
\leanok
\uses{lem:nie, lem:box}
(i) $r(n')\equiv B(n+E)\equiv r(n)+\beta$, and $r(n')\in[0,p)^s$, which gives $\omega$; if $\beta_c=0$, then $r_c(n)+\beta_c<p$. If $B$ is invertible, take $n=B^{-1}y$ with $y_c=p-1$ if $\omega_c=1$ (then $\beta_c\ge1$ and $y_c+\beta_c\ge p$) and $y_c=0$ if $\omega_c=0$. (ii) The box with depths $k+1$ on $T$, $k-1$ on $Z$ and $k$ elsewhere has total depth $sk$; by Lemma~\ref{lem:box} its minor is, up to sign, the determinant of the first offset rows of $T$ on the last level of $Z$, which is $\det\tilde M_{T,Z}$ modulo $q$, and by Lemma~\ref{lem:nie} it is a unit. Disjoint $T$, $Z$ of size $\lfloor s/2\rfloor$ exist.
\end{proof}

\begin{remark}[Glue under free wraps]\label{rem:glue}
\uses{thm:col, lem:shared, prop:cube, prop:twocoset, prop:wnet, prop:Hk}
A coordinate with $E_c=0$ and $\beta_c\ne0$ contributes $\min(b_c,1-b_c)\le\frac12$ to the fully folded length, with equality only for $b_c=\frac12$: a step keeps its full length under free wraps only if its glue in each coordinate that does not move is $\frac12$, the only value that is its own negative modulo $1$. So half-integer layers do not need the wraps of different coordinates to be correlated: for odd $s$ the differences between the two cosets of $D_s^+$ fill $(\Z+\frac12)^s$, a set that free wraps (integer changes of single coordinates) map to itself, so these steps keep a length of at least $\sqrt s/2$, and the relaxation of Proposition~\ref{prop:Hk} at $k=1$, which allows every wrap, tends to $\min(\sqrt s/2,\sqrt2)$ for dominant layers. The glue $\frac13\mathbf 1$ of the reflected cube does not survive: with free wraps the step $E=\e_1$ of the layer $\frac p3\mathbf 1\mathbf 1^{\!\top}$ would have fully folded length $1/\sqrt3=.5773\ldots$, and bounded perturbations of the equal rows $\mathbf 1a^{\!\top}$ of the column nets make the wraps independent (computed instance: Figure~\texttt{fig:arcwhylat}(b)). Wraps that are correlated across coordinates because the layer is singular (\emph{locked} wraps, Section~\ref{sec:wnet}) are needed otherwise: the column nets (Theorem~\ref{thm:col}) share one geometric digit among all three coordinates, Proposition~\ref{prop:twocoset} between two, and in Proposition~\ref{prop:wnet} a layer of rank one locks all three, one coordinate reading half of the digit of the others; literally equal rows exist only at $(s,m)=(3,3)$ (Lemma~\ref{lem:shared}). Locked wraps are what produces the coherent phases defined before Proposition~\ref{prop:cube}. The half-integer rule concerns coordinates whose wraps are free: a coordinate whose wraps are locked to those of the others can carry other glue and a step longer than one cell along its own axis (in the net $\mathcal W_p$ of Section~\ref{sec:wnet}, the glue $\frac12+3\varepsilon$ and the step $1+2\varepsilon$ along $\e_z$, with $\varepsilon=\rho/(2p)$). The rule concerns the length of a single step under free wraps; it is a heuristic for the choice of glue, not an optimality statement for the minimum, even for free coordinates (the remark on $\sqrt{233}/15$ in Section~\ref{sec:wnet}).
\end{remark}

\begin{proposition}[One point per cube]\label{prop:cube}
\lean{Z3NetF.cube_i, Z3NetF.cube_ii, Z3NetF.cube_iii, Z3NetF.cube_iv_rat, Z3NetF.cube_iv_min, Z3NetF.cube_iv_vals}\leanok
Let $S$ be a periodic one-per-cube set in $\R^s$ with period lattice $\Pi$.
\begin{enumerate}
\item[(i)] If $S=L$ is a lattice and $\Z^s+L$ is dense in $\R^s$, then $L$ contains a unit vector $\e_c$; so $\mu(L)\le1$.
\item[(ii)] If $S=L\subset\mathbb Q^s$ is a lattice whose phase group $Y$ equals $\pi_1(Y)\times\dots\times\pi_s(Y)$, the product of its projections to the coordinate circles, then $L$ contains a unit vector $\e_c$; so $\mu(L)\le1$.
\item[(iii)] If $s\le7$ and $\Z^s+\Pi$ is dense in $\R^s$, then two points of $S$ differ by a unit vector; so $\mu(S)\le1$.
\item[(iv)] If $S-x\subset\mathbb Q^s$ for some $x$, all squared distances in $S$ are rational. Hence $\mu(S)<(4/3)^{1/4}$ for $s=2$ and $\mu(S)<2^{1/6}$ for $s=3$, the optimal packing values at density $1$; and for $4\le s\le7$, $\mu(S)$ differs from the value of the densest known packing at density $1$ ($D_4$, $D_5$, $E_6$, $E_7$, with squared minima $\sqrt2$, $2^{3/5}$, $2\cdot3^{-1/6}$, $2^{6/7}$).
\end{enumerate}

\textbf{Lean: assuming classical theorems.} Uses Haj\'os's theorem, Keller's conjecture for $s\le7$ and the optimal packing bounds in dimensions two and three as explicit hypotheses.
\end{proposition}

\begin{proof}
\leanok
Let $f(x)=\#\bigl(S\cap(x+[0,1)^s)\bigr)$. Then $f(x+\pi)=f(x)$ for $\pi\in\Pi$ and $f(n)=1$ for $n\in\Z^s$, so $f=1$ on $\Z^s+\Pi$. For every $x$ there is $\varepsilon>0$ such that $f$ is constant on $x-[0,\varepsilon)^s$: $S$ is locally finite, so we may choose $\varepsilon$ below every positive difference $x_c-\sigma_c$ and $x_c+1-\sigma_c$ for the finitely many $\sigma\in S$ near the cube, and then shifting the cube by $-t$, $t\in[0,\varepsilon)^s$, neither adds nor removes a point. If $f\equiv1$, the cubes $\sigma+(-1,0]^s$, $\sigma\in S$, tile $\R^s$, since $x\in\sigma+(-1,0]^s$ if and only if $\sigma\in x+[0,1)^s$.
(i) Density gives a point of $\Z^s+L$ in every $x-[0,\varepsilon)^s$, so $f\equiv1$, and the cubes form a lattice tiling. By Haj\'os's theorem (Minkowski's conjecture) \cite{Haj42} two of them share a facet, so $L$ contains some $\e_c$.
(ii) For each $c$, $\pi_c(L)+\Z$ is a finitely generated subgroup of $\mathbb Q$ containing $\Z$, hence equal to $\frac1{e_c}\Z$ for an integer $e_c\ge1$. For $x_c\in(\frac{j-1}{e_c},\frac j{e_c}]$ and $\sigma\in L$ the condition $x_c\le\sigma_c<x_c+1$ is equivalent to $\frac j{e_c}\le\sigma_c\le\frac{j-1}{e_c}+1$, which does not depend on $x_c$; so $f$ is constant on every box $\prod_c(\frac{j_c-1}{e_c},\frac{j_c}{e_c}]$, and its value there is that at the corner $(j_c/e_c)_c$. By hypothesis the corner lies in $\Z^s+L$, so $f\equiv1$, and Haj\'os's theorem applies as in (i).
(iii) As in (i), $f\equiv1$ and the cubes $\sigma+(-1,0]^s$ tile $\R^s$. Keller's conjecture holds for $s\le7$ (Perron \cite{Per40} for $s\le6$, Brakensiek, Heule, Mackey and Narv\'aez \cite{BHMN20} for $s=7$): two cubes of the tiling share a complete facet, so their corners differ by some $\e_c$.
(iv) The differences are rational vectors. As $S$ is a finite union of $\Pi$-cosets, $\mu(S)$ is a minimum over the finitely many $\Pi$-orbits of pairs, so it is attained and $\mu(S)^2$ is rational. A periodic set of density $1$ with minimum $\mu$ is a packing of balls of radius $\mu/2$ with density $\mathrm{vol}(B_{\mu/2})$, so $\mu^2$ is at most $4\vartheta_s^{2/s}$ with $\vartheta_s$ the optimal centre density; this is $2/\sqrt3$ for $s=2$ \cite[Ch.~1]{CS99} and $2^{1/3}$ for $s=3$ \cite{Hal05}, both irrational. The four squared minima for $4\le s\le7$ are $4\vartheta_s^{2/s}$ with the centre densities $\frac18$, $\frac1{8\sqrt2}$, $\frac1{8\sqrt3}$, $\frac1{16}$ \cite{CS99}, and they are irrational.
\end{proof}

\begin{remark}[What the cube costs; computed]\label{rem:cubecost}
\uses{lem:CH, cor:linoff, tab:S1E, prop:wnet}
The largest minimum of a periodic one-per-cube set is not known. In dimension three an exhaustive enumeration of the lattices of covolume $1$ in $\frac1e\Z^3$, $e\le11$, finds one-per-cube lattices with minimum up to $\sqrt{29/25}=1.0770\ldots$ (computed), for instance the lattice with basis $\tfrac1{5}(3,-4,-2)$, $\tfrac1{10}(6,2,-9)$, $\tfrac1{2}(1,0,2)$. It has covolume $1$; its Gram matrix $100^{-1}\left(\begin{smallmatrix}116&56&-10\\56&121&-60\\-10&-60&125\end{smallmatrix}\right)$ is Minkowski-reduced, so its minimum is $\sqrt{116/100}$; and a finite check over one period shows that every unit cube holds exactly one of its points. So the largest minimum lies between $1.0770$ and $2^{1/6}=1.1224\ldots$ \cite{Hal05}. The maximizers found have offsets that are not affine modulo $1$ in the cell, so no digital net realizes them through its geometric digit (Corollary~\ref{cor:linoff}(c)); the best digital value at $m=3$ in this paper is $\sqrt{281}/16=1.0476\ldots$ (Proposition~\ref{prop:wnet}). For $s\le9$, restricting to lattices is not known to lose anything: lattices hold the packing records there and are optimal for $s=3$, $8$ \cite{CS99,Hal05,Via17}. The constraints visible here are the \emph{cube constraint} (one point in every cube of side $q^{-k}$, at $m\equiv0\pmod s$), the linearity of the offsets and, in the frame orbit, Lemma~\ref{lem:CH}. The skeleton regime ($m\not\equiv0\pmod s$), whose cells are not cubes, is not bound by the cube constraint (Table~\ref{tab:S1E}).
\end{remark}


\chapter{Beyond the cubic value: Euclidean skeletons and half-integer layers}\label{sec:large}


\section{Euclidean skeletons}\label{sec:S1E}

\begin{proposition}[Euclidean form of Theorem~\ref{thm:S1}(1)]\label{prop:S1E}
\lean{Z3NetF.S1E_net, Z3NetF.S1E_sep, Z3NetF.S1E_sep_real, Z3NetF.S1E_lam2, Z3NetF.S1E_tab_s8_p10007_rho}\leanok
\uses{prop:skelcomp, def:S1}
Let $s\ge3$, let $p\ge s-1$ be a prime, $k\ge0$, $m=sk+s-1$, and let $\kappa_1,\dots,\kappa_{s-1}$ be units of $\F_p$. The recipe of Proposition~\ref{prop:skelcomp}(ii) on the frame of Definition~\ref{def:S1} gives a digital $(0,m,s)$-net over $\F_p$ with the skeleton digit $S_z=\sum_d\kappa_dx_{d,k}$, and every graph form with the depths $(k+1,\dots,k+1,k)$ and this skeleton digit satisfies, for the toroidal distance,
\[
\delta N^{1/s}\ \ge\ p^{-1/s}\lambda_{\mathrm{jit}}(\kappa)\ \ge\ p^{-1/s}\bigl(\lambda_2(\Lambda_\kappa)-\sqrt s\bigr),\qquad \lambda_{\mathrm{jit}}(\kappa)=\min_{0\ne a\in\Lambda_\kappa}\bigl\|(|a|-\mathbf 1)_+\bigr\|_2,
\]
where $\Lambda_\kappa=\{a\in\Z^s:a_z\equiv\sum_d\kappa_da_{x_d}\pmod p\}$ and $(|a|-\mathbf 1)_+$ is taken coordinatewise. (So $\lambda_{\mathrm{jit}}$ is the length of the shortest nonzero lattice vector after one fine cell of jitter is subtracted from each coordinate.)
\end{proposition}

\begin{proof}
\leanok
\uses{lem:skelred, prop:skelcomp, thm:S1, def:S1}
The net property is Proposition~\ref{prop:skelcomp}(ii); the $s$ places of Definition~\ref{def:S1} are distinct because $p\ge s-1$. As in the proof of Theorem~\ref{thm:S1}(1), $P_z\equiv\sum_d\kappa_dP_{x_d}\pmod p$ for every point, and the centred difference $a\ne0$ of the fine positions of two distinct points lies in $\Lambda_\kappa$. By Lemma~\ref{lem:skelred}(i) in each coordinate $c$, the two points are at toroidal distance at least $(|a_c|-1)_+\varpi$ in coordinate $c$, so their distance is at least $\varpi\|(|a|-\mathbf 1)_+\|_2\ge \varpi\lambda_{\mathrm{jit}}(\kappa)$, and $\varpi N^{1/s}/N=p^{-1/s}$. Finally $|a|\le(|a|-\mathbf 1)_++\mathbf 1$ coordinatewise gives $\|(|a|-\mathbf 1)_+\|_2\ge\|a\|_2-\sqrt s$.
\end{proof}

\begin{remark}[Table of Proposition~\ref{prop:S1E}]\label{tab:S1E}
\lean{Z3NetF.S1E_tab_s8_p10007_rho, Z3NetF.S1E_tab_s3_p1009, Z3NetF.S1E_tab_s3_p10007, Z3NetF.S1E_tab_s3_p1000003, Z3NetF.S1E_tab_s3_p100000007, Z3NetF.S1E_tab_s3_p996001, Z3NetF.S1E_tab_s3_p10002435263, Z3NetF.S1E_tab_s4_p1009, Z3NetF.S1E_tab_s4_p10007, Z3NetF.S1E_tab_s4_p1000003, Z3NetF.S1E_tab_s4_p100000007, Z3NetF.S1E_tab_s4_p1006301, Z3NetF.S1E_tab_s4_p10040205517, Z3NetF.S1E_tab_s5_p1009, Z3NetF.S1E_tab_s5_p10007, Z3NetF.S1E_tab_s5_p1000003, Z3NetF.S1E_tab_s5_p100000007, Z3NetF.S1E_tab_s5_p1025909, Z3NetF.S1E_tab_s5_p10123347709, Z3NetF.S1E_tab_s6_p1009, Z3NetF.S1E_tab_s6_p10007, Z3NetF.S1E_tab_s6_p1000003, Z3NetF.S1E_tab_s6_p100000007, Z3NetF.S1E_tab_s6_p1033493, Z3NetF.S1E_tab_s6_p10136355487, Z3NetF.S1E_tab_s7_p1009, Z3NetF.S1E_tab_s7_p10007, Z3NetF.S1E_tab_s7_p1000003, Z3NetF.S1E_tab_s7_p100000007, Z3NetF.S1E_tab_s7_p1138273, Z3NetF.S1E_tab_s7_p10660100951, Z3NetF.S1E_tab_s8_p1009, Z3NetF.S1E_tab_s8_p10007, Z3NetF.S1E_tab_s8_p1000003, Z3NetF.S1E_tab_s8_p100000007, Z3NetF.S1E_tab_s8_p1071337, Z3NetF.S1E_tab_s8_p11556622703}\leanok
Certified lower bounds $\lambda_{\mathrm{jit}}(\kappa)p^{-1/s}$ for $\delta N^{1/s}$ (Proposition~\ref{prop:S1E}; toroidal, every $k\ge0$), truncated to three digits; every entry except those of the last two columns is such a certified bound. Left: $\kappa$ found by a random and local search; compare Table~\texttt{tab:skel} at $p=1009$, $10007$, $10^6{+}3$. ${}^*$Here the search is weaker than $\kappa_d=\rho^d$, which gives $.632$. Right: $\kappa$ obtained by rounding a randomly rotated copy of the root lattice to an integer matrix of prime determinant $p'$; the primes are $p'=996001$, $1006301$, $1025909$, $1033493$, $1138273$, $1071337$ and $10002435263$, $10040205517$, $10123347709$, $10136355487$, $10660100951$, $11556622703$ for $s=3,\dots,8$. The last two columns give $\sqrt{\gamma_s}$ (truncated) and the lattice that attains it. The $\kappa$ are listed in Table~\texttt{tab:S1Ekappa} (Appendix~\texttt{app:tables}); they are also fixed data of the Lean formalization (Appendix~\texttt{app:formal}).

\begin{tabular}{@{}ccccccccc@{}}
\toprule
& \multicolumn{4}{c}{searched $\kappa$, at $p=$} & \multicolumn{2}{c}{rounded root lattice, $p'\approx$} & & \\
$s$ & $1009$ & $10007$ & $10^6{+}3$ & $10^8{+}7$ & $10^6$ & $10^{10}$ & $\sqrt{\gamma_s}$ & root lattice\\
\midrule
3 & .940 & 1.027 & 1.090 & 1.097 & 1.106 & 1.121 & 1.122 & $A_3$\\
4 & .813 & .948 & 1.034 & 1.110 & 1.128 & 1.183 & 1.189 & $D_4$\\
5 & .752 & .823 & .989 & 1.069 & 1.098 & 1.208 & 1.231 & $D_5$\\
6 & .631 & .714 & .938 & 1.036 & 1.047 & 1.235 & 1.290 & $E_6$\\
7 & .372 & .599 & .867 & .984 & .974 & 1.238 & 1.345 & $E_7$\\
8 & .421 & $.547^*$ & .814 & .953 & .881 & 1.219 & 1.414 & $E_8$\\
\bottomrule
\end{tabular}
\end{remark}


\section{Resolution \texorpdfstring{$m=3$}{m=3}: a layer of two cosets}\label{sec:twocoset}

\begin{proposition}[Two-coset nets at $m=3$]\label{prop:twocoset}
\lean{QUNetsM3.TC_params, QUNetsM3.TC_cells, QUNetsM3.TC_net, QUNetsM3.TC_freeJitter, QUNetsM3.TC_sep, QUNetsM3.TC_sep_real, QUNetsM3.TC_bound, QUNetsM3.TC_pair, QUNetsM3.TC_pair_lt, QUNetsM3.TC_delta, QUNetsM3.TC_limit, QUNetsM3.TC_shift}\leanok
Let $p\ge5$ be a prime, $h=(p-1)/2$, let $r$ be the integer nearest to $p/8$, and let $\theta=8r-p\in\{\pm1,\pm3\}$. Let $\mathcal T_p$ be the digital net over $\F_p$ whose generator matrices $C_x,C_y,C_z$ have the rows
\[
C_x:\ \e_x,\ (0,h,h),\ (0,1,2);\qquad C_y:\ \e_y,\ (r,-1,1),\ (3,0,1);\qquad C_z:\ \e_z,\ (r,-1,1),\ (1,5,0).
\]
\begin{enumerate}
\item $\mathcal T_p$ is a $(0,3,3)$-net if and only if $p\ne23$, and every cell of side $p^{-1}$ holds one point.
\item For $p\ge41$ its toroidal (hence also Euclidean) minimum distance satisfies
\[
\delta N^{1/3}\ \ge\ \sqrt{\frac{33}{32}-\frac{41}{16p}}-\frac{\sqrt3}p.
\]
The bound exceeds $\sqrt{66}/8-3.014/p$ for $p\ge83$, and it exceeds $1$ for $p\ge197$.
\item For $p\ge7$ the points of the cells $0$ and $(p-1)\mathbf 1$ have toroidal distance $d_*$ at integer scale $p^3$, where $d_*^2=(p^2-2p+3)^2+((r-1)p+4)^2+((r-1)p+6)^2$. So $\delta N^{1/3}\le d_*/p^2$, and $d_*/p^2<\sqrt{66}/8-2/p$ for $p\ge53$. In particular $\delta N^{1/3}\to\sqrt{66}/8=1.0155\ldots$ as $p\to\infty$.
\end{enumerate}
\end{proposition}

\begin{proof}
\leanok
\uses{lem:nie, lem:box}
\emph{Net property.} The first rows are $\e_x,\e_y,\e_z$, so $\mathcal T_p$ is a graph form with the depths $(1,1,1)$, and we use Lemma~\ref{lem:box}: the box of shape $(1,1,1)$ has determinant $1$ and puts one point in every cell. A box with depth $2$ on $a$, $1$ on $b$ and $0$ on $c$ has, up to sign, the determinant $(C_a)_{2,c}$, the entry of the second row of $C_a$ in column $c$; these off-diagonal entries of the second rows are $h,h,r,1,r,-1$. The box with depth $3$ on $a$ has, up to sign, the $2\times2$ minor of the second and third rows of $C_a$ off column $a$; it is $2h-h=h$, $r-3$ and $5r+1$ for $a=x,y,z$. Here $1\le r<p$. Since $r\le p/8+\frac12$, $0<5r+1<p$ for $p\ge11$, and $5r+1=6$ for $p=5,7$, a unit modulo $5$ and $7$. Finally $r-3\equiv0$ means $r=3$, as $-2\le r-3<p$, i.e.\ $p/8\in(\frac52,\frac72)$; among primes this holds only for $p=23$, where the box with depth $3$ on $y$ is singular. Lemma~\ref{lem:nie} concludes.

\emph{Separation: reduction.} At integer scale $p^3$ the point of cell $n\in\{0,\dots,p-1\}^3$ is $X(n)=p^2n+p\,g(n)+f(n)$, with the geometric digits $g(n)$ and the third digits $f(n)$ in $\{0,\dots,p-1\}^3$. Let $n\ne n'$, and let $D=p^2E+p\,\Delta g+\Delta f$ be a representative of $X(n')-X(n)$ modulo $p^3$, where $E\equiv n'-n\pmod p$ is an integer vector, $\Delta g=g(n')-g(n)$ and $\Delta f=f(n')-f(n)\in(-p,p)^3$; the toroidal distance is the least $|D|$ over these representatives. If some $|E_c|\ge3$, then $|D_c|\ge3p^2-p(p-1)-(p-1)>2p^2$. So let $E\in\{-2,\dots,2\}^3\setminus\{0\}$, and put $S=E_y+E_z$ and $T=E_y-E_z$ (so $|S|\le4$, and $|S|\le3$ if $T$ is odd). As $C_y$ and $C_z$ have the same second row, $g_y(n)=g_z(n)$ for every cell, so $\Delta g=(\xi,\eta,\eta)$ with integers $\xi,\eta\in(-p,p)$: \emph{the $y$- and $z$-digits wrap together}. Since $\xi\equiv hS$, $2h\equiv-1$, $\eta\equiv rE_x-T$ and $8r=p+\theta$, there are integers $i$, $j$ with
\[
2\xi=jp-S,\qquad 8\eta=ip+\theta E_x-8T.
\]
Reducing modulo $2$ and modulo $8$, with $p$ odd and $\theta\equiv-p\pmod8$, gives $j\equiv S\pmod2$ and $i\equiv E_x\pmod8$. As $|\xi|<p$ and $|S|<p$, $|j|\le2$, and $j=\pm2$ only if $\pm S>0$. If $|T|\le1$, then $|\theta E_x-8T|\le14<p$, so likewise $|i|\le8$, $i=\pm8$ only if $\mp(\theta E_x-8T)>0$, and $i=0$ if $E_x=T=0$.

In cell units $D/p^2=V+\Delta f/p^2$, with $|\Delta f|<\sqrt3\,p$ and $V=E+(\xi,\eta,\eta)/p$. So $|D|/p^2>|V|-\sqrt3/p$, and it suffices to show $|V|^2\ge\frac{33}{32}-\frac{41}{16p}$. With the half-integer $\alpha=E_x+\frac j2$ and the quarter-integer $\beta=S+\frac i4$ (the limits of $V_x$ and $V_y+V_z$ as $p\to\infty$),
\[
V_x=\alpha-\frac S{2p},\qquad V_y+V_z=\beta+\frac{\theta E_x-8T}{4p},\qquad V_y-V_z=T,\qquad |V|^2=V_x^2+\tfrac12\bigl((V_y+V_z)^2+T^2\bigr).
\]
The terms of order $1/p$ are the drift (Section~\ref{sec:drift}): the deviation of the geometric digits from their limit values. Each numerical estimate below increases with $p$ and is evaluated at $p=41$.

\emph{The cases.} (a) $|T|\ge2$: $|V|^2\ge T^2/2\ge2$.

(b) $T=0$, $E_x=0$: then $S\ne0$ is even and $i=0$, so $V_y=V_z=E_y$ and $|V|^2\ge S^2/2\ge2$.

(c) $|T|=1$, $E_x=0$: $j$ is odd, so $|\alpha|=\frac12$ and $|V_x|\ge\frac12-\frac3{2p}$; $i\in\{0,\pm8\}$, so $\beta$ is an odd integer and $|V_y+V_z|\ge1-\frac2p$. Hence $|V|^2\ge(\frac12-\frac3{2p})^2+\frac12(1-\frac2p)^2+\frac12>1.16$.

(d) $|E_x|=2$. If $|T|=1$, then $j=\pm1$, $|\alpha|\ge\frac32$ and $|V|^2\ge(\frac32-\frac3{2p})^2>2$. If $T=0$, then $S$ and hence $j$ are even, so $\alpha$ is an integer, and $\alpha\ne0$ (which would need $j=\mp4$); so $|V_x|\ge1-\frac2p$; and $4\beta=4S+i\equiv E_x\equiv2\pmod4$, so $|\beta|\ge\frac12$ and $|V_y+V_z|\ge\frac12-\frac3{2p}$. Hence $|V|^2\ge(1-\frac2p)^2+\frac12(\frac12-\frac3{2p})^2=\frac98-\frac{19}{4p}+\frac{41}{8p^2}$, which is at least $\frac{33}{32}-\frac{41}{16p}$ for $p\ge24$.

(e) $|E_x|=1$, $T=0$: $\alpha\in\Z$, $i\in\{E_x,-7E_x\}$, and $4\beta=4S+i\equiv E_x\pmod8$, so $|\beta|\in\{\frac14,\frac74,\frac94,\dots\}$; moreover $|V_x-\alpha|\le\frac2p$ and $|V_y+V_z-\beta|\le\frac3{4p}$. If $\alpha=0$, then $j=-2E_x$, so $E_xS\le-2$ and $E_x\beta=E_xS+E_xi/4\le-\frac74$; then $|V|^2\ge\frac12(\frac74-\frac3{4p})^2>1.4$. If $|\alpha|\ge2$, or $|\alpha|=1$ and $|\beta|\ge\frac74$, then also $|V|^2>1.4$. There remains $|\alpha|=1$, $|\beta|=\frac14$. Since $S=\beta-\frac i4$ is an even integer, $\beta=\frac{E_x}4$ and $S\in\{0,2E_x\}$, and if $S=2E_x$ then $\alpha=E_x$ (as $|j|\le2$). With $(a-b)^2\ge a^2-2ab$,
\[
|V|^2\ \ge\ \alpha^2-\frac{\alpha S}p+\frac12\Bigl(\beta^2+\frac{\beta\theta E_x}{2p}\Bigr)=\frac{33}{32}-\frac{\alpha S-\theta/16}p\ \ge\ \frac{33}{32}-\frac{35}{16p},
\]
as $\alpha S\in\{0,2\}$ and $\theta\ge-3$.

(f) $|E_x|=1$, $|T|=1$: $j=\pm1$, so $\alpha\in\{\pm\frac12,\pm\frac32\}$; $i\in\{E_x,-7E_x\}$, and $4\beta=4S+i\equiv4+E_x\pmod8$, so $|\beta|\in\{\frac34,\frac54,\frac{11}4,\dots\}$; moreover $|V_x-\alpha|\le\frac3{2p}$ and $|V_y+V_z-\beta|\le\frac{11}{4p}$, as $|\theta E_x-8T|\le11$. If $|\alpha|=\frac32$, then $|V|^2\ge(\frac32-\frac3{2p})^2>2$. If $|\beta|\ge\frac54$, then $|V|^2\ge(\frac12-\frac3{2p})^2+\frac12(\frac54-\frac{11}{4p})^2+\frac12>1.41$. There remains $|\alpha|=\frac12$, $|\beta|=\frac34$. Then $\alpha=\frac{E_x}2$ and, from $4\beta\equiv4+E_x\pmod 8$, $\beta=-\frac34E_x$; and $S=\beta-\frac i4=\pm E_x$. As in (e),
\[
|V|^2\ \ge\ \alpha^2-\frac{\alpha S}p+\frac12\Bigl(\beta^2+\frac{\beta(\theta E_x-8T)}{2p}+T^2\Bigr)=\frac{33}{32}-\frac{\frac12E_xS-\frac32E_xT+\frac3{16}\theta}{p}\ \ge\ \frac{33}{32}-\frac{41}{16p},
\]
as $E_xS,E_xT\in\{\pm1\}$ and $\theta\le3$.

So $|V|^2\ge\frac{33}{32}-\frac{41}{16p}$ in every case, which proves (2). The bound increases with $p$; it is $1.0002\ldots$ at $p=197$ and below $1$ at $p=193$. For $p\ge83$ we have $\frac{33}{32}-\frac{41}{16p}\ge1$, and $\sqrt a-\sqrt{a-x}=x/(\sqrt a+\sqrt{a-x})\le x/2$ gives the bound $\sqrt{66}/8-(\frac{41}{32}+\sqrt3)/p$, with $\frac{41}{32}+\sqrt3<3.014$.

\emph{The pair.} For $n'=(p-1)\mathbf 1\equiv-\mathbf 1$ one has $g(n')=(1,p-r,p-r)$ and, as the entries of the third rows sum to $3$, $4$, $6$, $f(n')=(p-3,p-4,p-6)$; also $g(0)=f(0)=0$. The representative with $E=-\mathbf 1$ is $D_*=-\bigl(p^2-2p+3,\,(r-1)p+4,\,(r-1)p+6\bigr)$. Since $|D_{*,c}|<p^2\le p^3/2$ in every coordinate, this representative is the toroidal one, so the toroidal distance is $d_*$ and $\delta N^{1/3}\le|D_*|/p^2=d_*/p^2$. With $r=(p+\theta)/8$,
\[
p^4\Bigl(\frac{\sqrt{66}}8-\frac2p\Bigr)^2-d_*^2=\Bigl(\frac92-\frac\theta{16}-\frac{\sqrt{66}}2\Bigr)p^3-\Bigl(\frac{21}2-\frac\theta2+\frac{\theta^2}{32}\Bigr)p^2+\Bigl(32-\frac{5\theta}2\Bigr)p-61,
\]
and for $|\theta|\le3$ the three coefficients exceed $\frac14$, $-12.3$ and $24$, so the difference is positive for $p\ge53$. With (2) this gives the limit.
\end{proof}


\section{Resolution \texorpdfstring{$m=3$}{m=3}: a rank-one layer}\label{sec:wnet}

\begin{proposition}[A rank-one layer at $m=3$; (4) computer-assisted]\label{prop:wnet}
\lean{QUNetsM3.W_shift, QUNetsM3.W_layer, QUNetsM3.W_net, QUNetsM3.W_cells, QUNetsM3.W_freeJitter, QUNetsM3.W_sep, QUNetsM3.W_sep_real, QUNetsM3.W_bound, QUNetsM3.W_pair, QUNetsM3.W_pair_lt, QUNetsM3.W_delta, QUNetsM3.W_limit, QUNetsM3.W_exact, QUNetsM3.W_four, QUNetsM3.W_least}\leanok
Let $p\ge5$ be a prime, $h=(p-1)/2$, $r=\lfloor p/48\rfloor$ and $\rho=2r+1$. Let $\mathcal W_p$ be the digital net over $\F_p$ whose generator matrices $C_x,C_y,C_z$ have the rows
\[
C_x:\ \e_x,\ \rho(-1,-1,2),\ (0,1,2);\qquad C_y:\ \e_y,\ \rho(-1,-1,2),\ (1,0,1);
\]
\[
C_z:\ \e_z,\ (h-r,h-r,\rho),\ (0,1,0).
\]
\begin{enumerate}
\item $\mathcal W_p$ is a $(0,3,3)$-net, and every cell of side $p^{-1}$ holds one point.
\item For $p\ge53$ its toroidal (hence also Euclidean) minimum distance satisfies, for any third rows,
\[
\delta N^{1/3}\ \ge\ \sqrt{\frac{281}{256}-\frac{345}{128p}}-\frac{\sqrt3}p .
\]
The bound increases with $p$; it exceeds $1$ for $p\ge67$ and $\sqrt{66}/8$ for $p\ge97$.
\item For $p\ge53$ the points of the cells $(1,0,0)$ and $(2,0,p-1)$ have toroidal distance $d_*$ at integer scale $p^3$, where $d_*^2=(p^2-(6r+2)p-2)^2+(6r+3)^2p^2+(h+3r+2)^2p^2$, and $d_*/p^2<\sqrt{281}/16+5/(2p)$. So $\delta N^{1/3}\to\sqrt{281}/16=1.0476\ldots$ as $p\to\infty$.
\item $\delta N^{1/3}>1$ for every prime $p\ge17$, and $\delta N^{1/3}>\sqrt{66}/8$ for every prime $p\ge53$.
\end{enumerate}
\end{proposition}

\begin{proof}[Proof of Proposition~\ref{prop:wnet}]
\leanok
\uses{lem:nie, lem:box, prop:twocoset, lem:wwrap}
(1) The first rows are $\e_x,\e_y,\e_z$, so $\mathcal W_p$ is a graph form with the depths $(1,1,1)$, and we use Lemma~\ref{lem:box}: the box of shape $(1,1,1)$ has determinant $1$ and puts one point in every cell. A box with depth $2$ on $a$, $1$ on $b$ and $0$ on $c$ has, up to sign, the determinant $(C_a)_{2,c}$; these off-diagonal entries of the second rows are $-\rho,2\rho,-\rho,2\rho,h-r,h-r$. The box with depth $3$ on $a$ has, up to sign, the $2\times2$ minor of the second and third rows of $C_a$ off column $a$: $\det\left(\begin{smallmatrix}-\rho&2\rho\\1&2\end{smallmatrix}\right)=-4\rho$, $\det\left(\begin{smallmatrix}-\rho&2\rho\\1&1\end{smallmatrix}\right)=-3\rho$ and $\det\left(\begin{smallmatrix}h-r&h-r\\0&1\end{smallmatrix}\right)=h-r$ for $a=x,y,z$. As $1\le\rho\le\frac p{24}+1<p$, $1\le h-r<p$ and $p\ge5$, all these are units, and Lemma~\ref{lem:nie} concludes.

(2) \emph{Reduction.} At integer scale $p^3$ the point of cell $n$ is $X(n)=p^2n+p\,g(n)+f(n)$, with the second digits $g(n)=(2\zeta(n)\bmod p,\,2\zeta(n)\bmod p,\,\zeta(n))$ and the third digits $f(n)$. As in the proof of Proposition~\ref{prop:twocoset}, a representative of $X(n')-X(n)$ is $D=p^2E+p\,\Delta g+\Delta f$ with $E\equiv n'-n$, $|\Delta f_c|<p$, and only $E\in\{-2,\dots,2\}^3\setminus\{0\}$ can give $|D|\le2p^2$; then $|D|/p^2>|V|-\sqrt3/p$ with $V=E+\Delta g/p$, and by Lemma~\ref{lem:wwrap} $\Delta g=(2\Delta\zeta-p\omega,\,2\Delta\zeta-p\omega,\,\Delta\zeta)$. It suffices to show $|V|^2\ge\frac{281}{256}-\frac{345}{128p}$.

\emph{Parametrization.} Put $S=E_x+E_y$, $Q=S-2E_z$ and $\varepsilon=\rho/(2p)$. As $2(h-r)=p-\rho$, $2u\cdot E=pS-\rho Q$, so $\Delta\zeta=p(\nu-Q\varepsilon)$ with $\nu\in\frac S2+\Z$, and
\[
V=A-Q\varepsilon\mathbf w,\qquad A=E+(2\nu-\omega)(1,1,0)+\nu\e_z,\qquad \mathbf w=(2,2,1),
\]
\[
|V|^2=|A|^2-2Q(A\cdot\mathbf w)\,\varepsilon+9Q^2\varepsilon^2 .
\]
With $\theta=48r-p$, an odd integer in $[-47,-1]$ (here $\theta$ is not that of Proposition~\ref{prop:twocoset}), we have $\rho=(p+\theta)/24+1$, so $\varepsilon=\frac1{48}+\frac{\theta+24}{48p}$ and $|\varepsilon-\frac1{48}|\le\frac{23}{48p}$; for $p\ge53$, $\varepsilon\le\frac{19}{636}<\frac3{100}$. The drift of the layer is this deviation of $\varepsilon$ from its limit.

\emph{Which $(\nu,\omega)$ occur.} Let $p\ge53$. Then $|Q|\le8$ and $4\rho\le\frac p6+4<h$. By Lemma~\ref{lem:wwrap}: if $\nu=0$, then $\Delta\zeta=-Q\rho/2$, $|\Delta\zeta|\le h$, and $\omega\in\{0,-\sgn Q\}$ (only $\omega=0$ if $Q=0$); if $\nu=\pm\frac12$, then $\Delta\zeta=\frac12(\pm p-Q\rho)$ has the sign of $\nu$, $\omega=\sgn \nu$ occurs, and $\omega=0$ occurs if and only if $\nu Q>0$; if $\nu=\pm1$, the step occurs only if $\nu Q>0$, and then $|\Delta\zeta|>h$ and $\omega=\sgn \nu$; and $|\nu|\ge\frac32$ gives $|\Delta\zeta|\ge\frac32p-4\rho>p$, which does not occur. None of these rules depends on $p$ or on its class.

\emph{Long options.} $A$ lies in the lattice $\Lambda_{\mathcal W}=\{A:A_x,A_y\in\Z,\ 2A_z\in\Z,\ A_x+A_y\equiv2A_z\pmod2\}$, as $A_x+A_y\equiv S\equiv2\nu\equiv2A_z$. By the rules $|2\omega-2\nu|\le2$, and $Q=(A_x+A_y-2A_z)-2\nu+2\omega$, so $|Q|\le\sqrt6\,|A|+2$. Hence $|V|\ge|A|-3|Q|\varepsilon\ge(1-\frac{9\sqrt6}{100})|A|-\frac9{50}>1.05$ if $|A|^2\ge3$, and $1.05^2>\frac{281}{256}$.

\emph{Short options.} Since $2A_z\in\Z$, $|A|^2$ is an integer or an integer plus $\frac14$, and $|A|^2=\frac94$ would need $A=(\pm1,\pm1,\pm\frac12)$ or $\pm\frac32\e_z$, which violate the parity condition; so the vectors of $\Lambda_{\mathcal W}$ with $0<|A|^2<3$ are those with $|A|^2\le2$: $\pm\e_z$, $(\pm1,0,\pm\frac12)$, $(0,\pm1,\pm\frac12)$ and $(\pm1,\pm1,0)$. For each of them $E=A-(2\nu-\omega)(1,1,0)-\nu\e_z$, with $\nu\equiv A_z\pmod1$ and $|\nu|\le1$, so at most three values of $\nu$, and checking the rules for these values and their $\omega$ leaves the options of Table~\texttt{tab:wcontacts}; $A=0$ does not occur (it would need $E=\omega(1,1,0)$ with $\nu=0$, $\omega=-\sgn Q=-\sgn(2\omega)$, or $E=-\nu\mathbf 1$ with $|\nu|=1$ and $Q=0$). The squared length $q(\varepsilon)=|V|^2$ of each row is a convex quadratic in $\varepsilon$, so it lies above its tangent line at $\frac1{48}$, and as $|\varepsilon-\frac1{48}|\le\frac{23}{48p}$, $q(\varepsilon)\ge q(\frac1{48})-|q'(\frac1{48})|\cdot\frac{23}{48p}=:T(p)$. For the diagonal contact (a), $T(p)=\frac{281}{256}-\frac{345}{128p}$, as $\frac{45}8\cdot\frac{23}{48}=\frac{345}{128}$. For the step (b) along $\e_z$, $T(p)=\frac{211}{192}-\frac{253}{96p}$, which is larger, as $\frac{211}{192}>\frac{281}{256}$ and $\frac{253}{96}<\frac{345}{128}$. For the rows (c)--(g), $q(\frac1{48})>\frac{281}{256}$, so $p\,(T(p)-\frac{281}{256})=p\,(q(\frac1{48})-\frac{281}{256})-\frac{23}{48}|q'(\frac1{48})|$ is affine in $p$ with positive slope, and at $p=53$ it is positive (for row (c), $53\cdot\frac5{96}-\frac{23}{48}\cdot\frac{37}8>0$, and the other rows have more room); so these options have $|V|^2>\frac{281}{256}$. This proves (2). The constant $\frac{345}{128}$ is the slope of the diagonal contact (a) times the largest drift $\frac{23}{48p}$, which is attained for $p\equiv1\pmod{48}$; exactly, with $x=\varepsilon-\frac1{48}=\frac{\theta+24}{48p}$,
\[
|V_{\mathrm a}|^2=\frac{281}{256}-\frac{45}8\,x+81x^2 .
\] The bound increases with $p$; it is $1.0024\ldots$ at $p=67$, below $1$ at $p=61$, $1.0164\ldots$ at $p=97$ and below $\sqrt{66}/8=1.01550\ldots$ at $p=89$.

(3) Here $\zeta(1,0,0)=h-r$ and $\zeta(2,0,p-1)=p-2-4r$, so the second digits of the two points are $(p-1-2r,\,p-1-2r,\,h-r)$ and $(p-4-8r,\,p-4-8r,\,p-2-4r)$ (as $p\ge8r+4$), and the third digits are $(0,1,0)$ and $(p-2,1,0)$. Hence $D=(p^2-(6r+2)p-2,\,-(6r+3)p,\,-(h+3r+2)p)$ is a representative of the difference with $|D_c|<p^3/2$, i.e. the toroidal one. As $6r+3=3\rho$ and $h+3r+2=\frac12(p+3\rho)$, $D/p^2=(1-6\varepsilon,-6\varepsilon,-\frac12-3\varepsilon)+(\frac1p-\frac2{p^2},0,0)$: the contact (a) with the step $E=(1,0,-1)$, plus a jitter. So $d_*/p^2\le|V_{\mathrm a}|+\frac1p$, where by the display in (2), with $|x|\le\frac{23}{48p}$, $|V_{\mathrm a}|^2\le\frac{281}{256}+\frac{345}{128p}+\frac{4761}{256p^2}<\frac{281}{256}+\frac{3.05}p$ for $p\ge53$; with $\sqrt{a+x}\le\sqrt a+x/(2\sqrt a)$ this gives $d_*/p^2<\sqrt{281}/16+2.46/p$, slightly more than (3) states. Together with (2), $\delta N^{1/3}\to\sqrt{281}/16$.

(4) By (2), $\delta N^{1/3}>1$ for every prime $p\ge67$ and $>\sqrt{66}/8$ for every prime $p\ge97$. The remaining primes are settled by the exact minimum distances (Appendix~\texttt{app:verify}): at every prime $17\le p\le61$, $\delta N^{1/3}\ge1.0009\ldots$ (least at $p=47$), and at every prime $53\le p\le89$, $\delta N^{1/3}\ge1.0257\ldots$ (least at $p=89$).
\end{proof}

\begin{lemma}[Wraps of a doubled digit]\label{lem:wwrap}
\lean{QUNetsM3.W_wrap, QUNetsM3.W_wrap_one}\leanok
Let $p\ge5$ be a prime, $h=(p-1)/2$, $u\in\F_p^3\setminus\{0\}$, and $\zeta(n)=u\cdot n\bmod p\in\{0,\dots,p-1\}$ for $n\in\F_p^3$. For $E\in\Z^3$ and $n'\equiv n+E$ put $\Delta\zeta=\zeta(n')-\zeta(n)$ and $\omega=[2\zeta(n')\ge p]-[2\zeta(n)\ge p]$, so that the digit $2\zeta\bmod p$ changes by $2\Delta\zeta-p\omega$. Let $\sigma\in\{0,\dots,p-1\}$ be the residue of $u\cdot E$. As $n$ runs over $\F_p^3$, the pairs $(\Delta\zeta,\omega)$ that occur are exactly: $(0,0)$ if $\sigma=0$; and if $\sigma\ne0$, for each $\Delta\zeta\in\{\sigma,\sigma-p\}$ the pair $(\Delta\zeta,\sgn\Delta\zeta)$, and also $(\Delta\zeta,0)$ if $|\Delta\zeta|\le h$.
\end{lemma}

\begin{proof}
\leanok
As $u\ne0$, $\zeta$ takes every value, and $\zeta(n')\equiv\zeta(n)+\sigma$; so $\Delta\zeta=\sigma$ for $\zeta(n)\le p-1-\sigma$ and $\Delta\zeta=\sigma-p$ otherwise, and both occur if $\sigma\ne0$. Reversing the pair exchanges $(\Delta\zeta,\omega)$ and $(-\Delta\zeta,-\omega)$, so let $\Delta\zeta=d>0$ and $\zeta=\zeta(n)\in[0,p-1-d]$. As $2\zeta\ge p$ if and only if $\zeta\ge h+1$, $\omega\in\{0,1\}$, and $\omega=1$ if and only if $\zeta\le h<\zeta+d$; such $\zeta$ form the interval $[\max(0,h+1-d),\min(h,p-1-d)]$, which is not empty, as $1\le d\le p-1$ and $h+2\le p$. And $\omega=0$ if and only if $\zeta+d\le h$, possible if and only if $d\le h$, or $\zeta\ge h+1$, possible if and only if $h+1\le p-1-d$, i.e. $d\le h-1$.
\end{proof}

\begin{proposition}[Optimum of the rank-one layers with image $(2,2,1)$]\label{prop:wfamily}
\lean{QUNetsM3.Wfam_star, QUNetsM3.Wfam_one, QUNetsM3.Wfam_two, QUNetsM3.Wfam_three_tau, QUNetsM3.Wfam_three}\leanok
Let $p\ge5$ be a prime, $\mathbf w=(2,2,1)$, and let $P$ be a digital $(0,3,3)$-net over $\F_p$ in graph form with the depths $(1,1,1)$ whose layer is $B=\mathbf w u^{\!\top}$ for some $u\in\F_p^3$, with any third rows: $z$ reads the second digit $\zeta(n)=u\cdot n\bmod p$, and $x$ and $y$ read $2\zeta(n)\bmod p$. For $\psi\in(\R/\Z)^3$ let $F(\psi)$ be the least value of $|E+\xi\mathbf w-\omega(1,1,0)|^2$ over the steps $E\in\Z^3\setminus\{0\}$ and the pairs $(\xi,\omega)$ with $\xi\in(-1,1)$, $\xi\equiv\psi\cdot E\pmod1$, and $\omega=\sgn\xi$, or $\omega=0$ if $|\xi|<\frac12$. Let $\psi^*=(\frac12-\varepsilon^*,\frac12-\varepsilon^*,2\varepsilon^*)$.
\begin{enumerate}
\item Let $F_p(u)$ be the squared toroidal minimum distance of the points $p^2n+p\,g(n)$ at integer scale $p^3$, divided by $p^4$, where $g(n)=(2\zeta(n)\bmod p,\,2\zeta(n)\bmod p,\,\zeta(n))$ are the second digits: the value of $(\delta N^{1/3})^2$ without the third digits. Then $F_p(u)=F(u/p)$, and $|\delta N^{1/3}-\sqrt{F_p(u)}|<\sqrt3/p$ for any third rows.
\item For $0<\varepsilon\le\frac3{100}$, $F(\frac12-\varepsilon,\frac12-\varepsilon,2\varepsilon)=\min(q_a(\varepsilon),q_b(\varepsilon))$.
\item $F(\psi)\le\mu^{*2}$ for every $\psi\in(\R/\Z)^3$, with equality only at $\psi=\psi^*$. As $\psi^*$ has irrational coordinates, $F_p(u)<\mu^{*2}$, and $\delta N^{1/3}<\mu^*+\sqrt3/p$.
\end{enumerate}
\end{proposition}

\begin{proof}
\leanok
\uses{prop:wnet, lem:wwrap}
Call a step $E$ with a pair $(\xi,\omega)$ as in the definition of $F$ an \emph{option} of $\psi$, with the vector $V=E+\xi\mathbf w-\omega(1,1,0)$. Every option has $|\xi|<1$ and $|2\xi-\omega|<1$, so $|V|>2$ if some $|E_c|\ge3$. The step $\e_z$ gives $V=\e_z$ if $\psi_z\equiv0$, and otherwise, with $\xi=\{\psi_z\}-1$ and $\omega=-1$, it gives $V=(2\xi+1,2\xi+1,1+\xi)$, whose entries lie in $(-1,1)$. Hence $F\le3$, and only the steps $E\in\{-2,\dots,2\}^3\setminus\{0\}$ matter.

(1) At integer scale $p^3$ the point of cell $n$ is $p^2n+p\,g(n)+f(n)$ with the third digits $f(n)$. As in the proof of Proposition~\ref{prop:wnet}(2), the representatives of the difference of the points of the cells $n$ and $n'$ are $D=p^2E+p\,\Delta g+\Delta f$ with $E\equiv n'-n\pmod p$ and $|\Delta f_c|<p$; without the third digits they are $p^2V$ with $V=E+\Delta g/p$, and $\bigl||D|/p^2-|V|\bigr|\le|\Delta f|/p^2<\sqrt3/p$. Comparing a pair that attains $\delta$ with one that attains $\sqrt{F_p(u)}$ gives $|\delta N^{1/3}-\sqrt{F_p(u)}|<\sqrt3/p$. By Lemma~\ref{lem:wwrap} (which holds for every $u$; part (1) uses neither $u\ne0$ nor the net property), $\Delta g=(2\Delta\zeta-p\omega,\,2\Delta\zeta-p\omega,\,\Delta\zeta)$, and the pairs $(\Delta\zeta,\omega)$ that occur for the step $E$ are those of the lemma. With $\xi=\Delta\zeta/p$ they are exactly the pairs of the definition of $F$ for $\psi=u/p$: $\Delta\zeta\in\{\sigma,\sigma-p\}$ (only $0$ if $\sigma=0$) means $\xi\in(-1,1)$ and $\xi\equiv u\cdot E/p\pmod1$, and, as $p$ is odd, $|\Delta\zeta|\le h$ means $|\xi|<\frac12$. Then $V=E+\xi\mathbf w-\omega(1,1,0)$. As $p\ge5$, every step $E\in\{-2,\dots,2\}^3\setminus\{0\}$ joins two distinct cells, so $F_p(u)=F(u/p)$.

(2) Let $\psi=(\frac12-\varepsilon,\frac12-\varepsilon,2\varepsilon)$ with $0<\varepsilon\le\frac3{100}$. With $S=E_x+E_y$ and $Q=S-2E_z$, $\psi\cdot E=\frac S2-Q\varepsilon$, so $\xi=\nu-Q\varepsilon$ with $\nu\in\frac S2+\Z$, and $V=A-Q\varepsilon\mathbf w$ with $A=E+(2\nu-\omega)(1,1,0)+\nu\e_z$. This is the parametrization of the proof of Proposition~\ref{prop:wnet}(2), where $u/p$ is of this form with $\varepsilon=\rho/(2p)$. As $|Q\varepsilon|\le8\varepsilon<\frac12$, the pairs that occur follow the rules stated there: for $\nu=0$, $|\xi|<\frac12$ and $\omega\in\{0,-\sgn Q\}$; for $\nu=\pm\frac12$, $\xi$ has the sign of $\nu$, $\omega=\sgn \nu$ occurs, and $\omega=0$ occurs if and only if $\nu Q>0$; $\nu=\pm1$ occurs only if $\nu Q>0$, and then $\omega=\sgn \nu$; and $|\nu|\ge\frac32$ does not occur. The rest of that proof uses only $\varepsilon\le\frac3{100}$: the options with $|A|^2\ge3$ have $|V|>1.05$, and the others are the rows of Table~\texttt{tab:wcontacts}. On $(0,\frac3{100}]$ the rows (c), (d) and (f) decrease, to $1.1081$, $1.1681$ and $1.5524$, and the rows (e) and (g) are at least $\frac54$; all of these, and $1.05^2=1.1025$, exceed $\mu^{*2}$, while $\min(q_a(\varepsilon),q_b(\varepsilon))\le\mu^{*2}$, as $q_a$ decreases and $q_b$ increases through $\varepsilon^*$. The rows (a) and (b) occur, since $(\nu,\omega)=(\frac12,1)$ and $(0,0)$ are always allowed. So $F(\psi)=\min(q_a(\varepsilon),q_b(\varepsilon))$.

(3) Let $\psi\in(\R/\Z)^3$ with $F(\psi)\ge\mu^{*2}$, and represent $\psi_x,\psi_y,\psi_z$ by numbers in $[0,1)$. We use only options with $\omega=\sgn\xi$, which always occur, and options with $\omega=0$ and $|\xi|<\frac12$. The numbers $\frac{246}{225}=1.0933\ldots$, $\frac{17}{16}$ and $\frac9{25}$ are below $\mu^{*2}=1.0983\ldots$.

\emph{The step $\e_z$.} If $\psi_z=0$, then $V=\e_z$ and $|V|^2=1$. Otherwise $\xi=\psi_z-1$, $\omega=-1$ gives $V=(2\psi_z-1,2\psi_z-1,\psi_z)$ with $|V|^2=9\psi_z^2-8\psi_z+2$, a convex function with the values $\frac{246}{225}$ at $\psi_z=\frac2{15}$ and $\frac{17}{16}$ at $\psi_z=\frac34$; so $\psi_z\notin[\frac2{15},\frac34]$. If $\psi_z>\frac34$, then $\xi=\psi_z-1$, $\omega=0$ gives $|V|^2=8(1-\psi_z)^2+\psi_z^2\le\frac{17}{16}$. So $0<\psi_z<\frac2{15}$, and $\xi=\psi_z$, $\omega=0$ gives $V=(2\psi_z,2\psi_z,1+\psi_z)$ with $|V|^2=q_b(\psi_z/2)$; as $q_b$ increases, $\psi_z\ge2\varepsilon^*$.

\emph{The steps $\e_x$ and $\e_y$.} If $\psi_x=0$, then $|V|^2=1$. Otherwise $\xi=\psi_x$, $\omega=1$ gives $V=(2\psi_x,2\psi_x-1,\psi_x)$ with $|V|^2=9\psi_x^2-4\psi_x+1$, a convex function with the values $1$ at $\psi_x=0$ and $\frac{246}{225}$ at $\psi_x=\frac7{15}$; so $\psi_x>\frac7{15}$. If $\psi_x>\frac12$, then $\xi=\psi_x-1$, $\omega=0$ gives $V=(1-2t,-2t,-t)$ with $t=1-\psi_x$, and $|V|^2=9t^2-4t+1$ is the same function, so $t>\frac7{15}$. Hence $\alpha=\frac12-\psi_x\in(-\frac1{30},\frac1{30})$, and likewise $\beta=\frac12-\psi_y\in(-\frac1{30},\frac1{30})$. Put $\gamma=\psi_z$.

\emph{The step $\mathbf 1$ (the lock).} Here $\psi\cdot\mathbf 1\equiv\gamma-\alpha-\beta=:\tilde\sigma\in(-\frac1{15},\frac15)$. If $\tilde\sigma>0$, then $\xi=\tilde\sigma-1$, $\omega=-1$ gives $V=(2\tilde\sigma,2\tilde\sigma,\tilde\sigma)$ with $|V|^2=9\tilde\sigma^2<\frac9{25}$. So $\alpha+\beta\ge\gamma$.

\emph{The diagonal steps $(1,0,-1)$ and $(0,1,-1)$.} For $E=(1,0,-1)$, $\psi\cdot E\equiv\psi_x-\psi_z=\frac12-\lambda$ with $\lambda=\alpha+\gamma\in(2\varepsilon^*-\frac1{30},\frac16)\subset(0,\frac16)$, and $\xi=\frac12-\lambda$, $\omega=1$ gives $V=(1-2\lambda,-2\lambda,-\frac12-\lambda)$ with $|V|^2=\frac54-3\lambda+9\lambda^2=q_a(\lambda/3)$: row (a) of Table~\texttt{tab:wcontacts} with $\varepsilon=\lambda/3$. Likewise $E=(0,1,-1)$ gives $q_a(\lambda'/3)$ with $\lambda'=\beta+\gamma$. As $q_a$ decreases on $[0,\frac1{18}]$ and $q_a(\varepsilon^*)=\mu^{*2}\le F(\psi)$, we get $\lambda\le3\varepsilon^*$ and $\lambda'\le3\varepsilon^*$. But $\lambda+\lambda'=\alpha+\beta+2\gamma\ge3\gamma\ge6\varepsilon^*$, by the lock and the step $\e_z$. So equality holds throughout: $\lambda=\lambda'=3\varepsilon^*$, $\alpha+\beta=\gamma$ and $\gamma=2\varepsilon^*$, hence $\alpha=\beta=\varepsilon^*$, i.e.\ $\psi=\psi^*$.

Hence $F(\psi)\ge\mu^{*2}$ only at $\psi=\psi^*$, where $F(\psi^*)=\mu^{*2}$ by (2): $F\le\mu^{*2}$ everywhere, with equality only at $\psi^*$. As $\varepsilon^*$ is irrational, $u/p\ne\psi^*$, so $F_p(u)=F(u/p)<\mu^{*2}$, and by (1) $\delta N^{1/3}<\sqrt{F_p(u)}+\sqrt3/p<\mu^*+\sqrt3/p$.
\end{proof}

\begin{remark}[Reaching $\mu^*$]\label{rem:wstar}
\lean{QUNetsM3.W_layer, QUNetsM3.Wfam_star, QUNetsM3.Wstar_net, QUNetsM3.Wstar_rho, QUNetsM3.Wstar_eps, QUNetsM3.Wstar_bound, QUNetsM3.Wstar_limit, QUNetsM3.Wstar_thresholds, QUNetsM3.Wstar_ahead}\leanok
\uses{prop:wnet, prop:wfamily}
Let $\mathcal W^*_p$ be the net of Proposition~\ref{prop:wnet} with $r=\lfloor p\varepsilon^*\rfloor$ in place of $\lfloor p/48\rfloor$. Then $\rho=2r+1$ is the odd integer nearest to $2p\varepsilon^*$, and $\varepsilon_p=\rho/(2p)$ satisfies $|\varepsilon_p-\varepsilon^*|<\frac1{2p}$ (both use that $p\varepsilon^*$ is never an integer, as $\varepsilon^*$ is irrational). $\mathcal W^*_p$ is a $(0,3,3)$-net for every prime $p\ge5$, since the proof of Proposition~\ref{prop:wnet}(1) needs only $1\le\rho<p$ and $1\le h-r<p$. Its layer is $\mathbf w u^{\!\top}$ with $u/p=(\frac12-\varepsilon_p,\frac12-\varepsilon_p,2\varepsilon_p)$, and $\varepsilon_p\le\frac3{100}$ for $p\ge17$ ($\varepsilon_p=\frac1{2p}$ for $p\le47$, $\varepsilon_{53}=\frac3{106}$, and $\varepsilon^*+\frac1{118}<\frac3{100}$). So for $p\ge17$ Proposition~\ref{prop:wfamily}(1), (2) gives the value $\min(q_a(\varepsilon_p),q_b(\varepsilon_p))$ without the third digits. Both $q_a$ and $q_b$ are convex, so each lies above its tangent line at $\varepsilon^*$; there $q_a$ decreases and $q_b$ increases, and $q_b'(\varepsilon^*)<|q_a'(\varepsilon^*)|$; hence
\[
\delta N^{1/3}(\mathcal W^*_p)\ >\ \sqrt{\mu^{*2}-\frac Kp}-\frac{\sqrt3}p,\qquad K=\frac{|q_a'(\varepsilon^*)|}2=\frac{9\sqrt{31}-36}5=2.8219\ldots,
\]
for every prime $p\ge17$. With Proposition~\ref{prop:wfamily}(3), $\delta N^{1/3}(\mathcal W^*_p)\to\mu^*$: the optimum of the family is attained in the limit. The bound exceeds $1$ from $p=67$ on and $\sqrt{66}/8$ from $p=97$ on, as that of Proposition~\ref{prop:wnet}(2). It exceeds $\sqrt{281}/16$ from $p=9791$ on, and the upper bound $\sqrt{281}/16+5/(2p)$ of Proposition~\ref{prop:wnet}(3) from $p=17737$ on; only from there on is $\mathcal W^*_p$ provably ahead of $\mathcal W_p$.

The optimum is reached only asymptotically: at moderate $p$, $\mathcal W^*_p$ is not ahead of $\mathcal W_p$. $\mathcal W^*_p=\mathcal W_p$ for $p<193$ (computed), and $\mathcal W^*_p\ne\mathcal W_p$ for $p>8548$: as $\frac1{48}-\varepsilon^*=\frac{96\sqrt{31}-534}{4320}>\frac1{8549}$, for $p\ge8549$ we have $p/48-p\varepsilon^*>1$, so $\lfloor p/48\rfloor>\lfloor p\varepsilon^*\rfloor$, and the two nets have different $\rho<p$. At the primes $p\le1009$ where they differ, namely $193$, $241$, $337$, $433$, $577$, $673$, $769$ and $1009$, $\mathcal W_p$ has the larger $\delta N^{1/3}$ (computed; $1.0405\ldots$ against $1.0355\ldots$ at $p=193$, and $1.0463\ldots$ against $1.0458\ldots$ at $p=1009$). At these primes $\varepsilon_p<\varepsilon^*$ for $\mathcal W^*_p$, so without the third digits its minimum is the step (b) along $\e_z$, while $\mathcal W_p$ has $\varepsilon>\varepsilon^*$ and its minimum is the diagonal contact (a). A heuristic reading of the computed values is that the third rows of Proposition~\ref{prop:wnet} cost the step (b) more than the contact (a): at moderate $p$ the jitter, not the drift, decides, and the value that $\mathcal W_p$ gives up by its limit $\varepsilon\to\frac1{48}$ instead of $\varepsilon^*$ is small, $\mu^{*2}-\frac{281}{256}<7\cdot10^{-4}$.
\end{remark}

\begin{remark}[Digital shifts of $\mathcal T_p$ and $\mathcal W_p$]\label{rem:wshift}
\lean{QUNetsM3.W_shift, QUNetsM3.digitalPtS_eq_dshiftPt, QUNetsM3.TC_shift, QUNetsM3.W_wrap_one}\leanok
\uses{lem:shift, rem:shiftshared, prop:twocoset, prop:wnet, lem:wwrap}
As for the column nets (Remark~\ref{rem:shiftshared}), a digital shift translates the wrap windows of the coordinates independently, so bounds that use locked wraps need not survive it. The bound of Proposition~\ref{prop:twocoset}(2) does not extend to shifts whose second digits differ in $y$ and $z$, the two coordinates that share their geometric digit, for instance at $p=41$, $53$, $101$ (computed). If they agree, $\tau_{y,2}=\tau_{z,2}$ (all other shift digits arbitrary), the bound holds for the shifted net, for every prime $p\ge41$. Indeed, by Lemma~\ref{lem:shift} the shifted net, indexed by its own cells, is $\mathcal T_p$ with constants added to the offset digits. As $C_y$ and $C_z$ have the same second row, the geometric digits of $y$ and $z$ of each cell differ by the constant $\tau_{y,2}-\tau_{z,2}=0$, so $g_y=g_z$ at every cell; and the differences of the geometric digits of two cells satisfy the same congruences $\xi\equiv hS$, $\eta\equiv rE_x-T$ as for $\mathcal T_p$, since the constants cancel. The proof of Proposition~\ref{prop:twocoset}(2) uses the geometric digits only through these congruences and through $g_y=g_z$, and the third digits only through $|\Delta f_c|<p$; so it applies verbatim. The lock of $\mathcal W_p$ (Proposition~\ref{prop:wnet}) is more fragile, because $z$ reads half of the digit of $x$ and $y$. In the notation of Lemma~\ref{lem:wwrap} ($z$ reads $\zeta(n)=u\cdot n\bmod p$, and $x$ and $y$ read $2\zeta(n)\bmod p$), let the second shift digits be $\tau_{x,2}=\tau_{y,2}=t$ and $\tau_{z,2}=t'$. Re-indexing the cells by $n'=n+n_0$, for a cell $n_0$ with $u\cdot n_0\equiv t'$ (a translation of the torus by $-n_0$ cells, together with a shift of the third digits), turns the shifted net into the one in which $z$ reads $\zeta(n)$ and $x$ and $y$ read $(2\zeta(n)+t-2t')\bmod p$. For $t\equiv2t'$ this is $\mathcal W_p$ with shifted third digits. For $t\equiv2t'+1$ the digit $(2\zeta+1)\bmod p$ wraps, i.e.\ $2\zeta+1\ge p$, if and only if $\zeta\ge h$, one step earlier than $2\zeta$, and with $\omega=[2\zeta(n')+1\ge p]-[2\zeta(n)+1\ge p]$ the pairs $(\Delta\zeta,\omega)$ that occur are exactly those of Lemma~\ref{lem:wwrap}. Indeed, the proof of the lemma goes through with $h$ in place of $h+1$: the values of $\Delta\zeta$ and the reversal are unchanged, and for $\Delta\zeta=d>0$ and $\zeta=\zeta(n)\in[0,p-1-d]$ we have $\omega\in\{0,1\}$; $\omega=1$ if and only if $\zeta\le h-1<\zeta+d$, which holds for $\zeta=\max(0,h-d)$, so it occurs for every $d$; and $\omega=0$ if and only if $\zeta+d\le h-1$, possible if and only if $d\le h-1$ (both digits below the wrap), or $\zeta\ge h$, possible if and only if $h\le p-1-d$, i.e.\ $d\le h$ (both at or above it). So $\omega=0$ occurs if and only if $d\le h$, as in the lemma. The proof of Proposition~\ref{prop:wnet}(2) uses the second digits only through these pairs and allows any third digits, so its bound holds for the shifted net whenever $t\equiv2t'$ or $t\equiv2t'+1$; the exact value need not be kept, since the third digits move. Other second digits break the lock in every case computed: at $p=101$ the shifts $(\tau_{x,2},\tau_{y,2},\tau_{z,2})=(0,0,1)$ and $(7,7,7)$ give $\delta N^{1/3}=.1355\ldots$, against $1.0362\ldots$ without a shift, and every shift $(t,t,0)$ with $2\le t\le100$ gives $\delta N^{1/3}<0.1356$, while $(1,1,0)$, $(3,3,1)$ and $(5,5,2)$ keep $1.0362\ldots$ (computed). That only $t\equiv2t'$ and $t\equiv2t'+1$ keep the lock is computed, not proved. In the computed cases where the lock breaks, the step along $\e_z$ at a wrap of $z$ alone shrinks to $(4\varepsilon,4\varepsilon,2\varepsilon)$ up to the third digits. Proposition~\ref{prop:wnet} is a statement about unshifted nets.
\end{remark}


\section{No half-integer layer in the frame orbit}\label{sec:Pnohalf}

\begin{proposition}[No half-integer frame layer]\label{prop:P}
\lean{Z3NetF.P_half_mod, Z3NetF.P_half_mod_large, Z3NetF.P_half_rat}\leanok
Let $\tilde M$ be a layer of this Cauchy form over a field $\mathbb K$, and let $a,b,c,d$ be four distinct coordinates.
\begin{enumerate}
\item If $\mathbb K=\F_p$, $p$ prime, and $2\tilde M_{ij}\equiv h_{ij}\pmod p$ for all twelve ordered pairs $i\ne j$ in $\{a,b,c,d\}$, with odd integers $|h_{ij}|\le\Omega$, then $p$ divides a nonzero odd integer $Z$ with $|Z|\le13\Omega^{16}$. In particular this is impossible for $p>13\Omega^{16}$.
\item If $\mathbb K=\mathbb{Q}$, then the twelve entries $\tilde M_{ij}$, $i\ne j$ in $\{a,b,c,d\}$, do not all lie in $\frac12+\Z$.
\end{enumerate}
\end{proposition}

\begin{proof}
\leanok
\uses{lem:CH}
By Lemma~\ref{lem:CH}(P) over $\mathbb K$, with $\pi_{ij}=\tilde M_{ij}\tilde M_{ji}$, the numbers $\chi^2=\pi_{ab}\pi_{cd}/(\pi_{ac}\pi_{bd})$ and $\psi^2=\pi_{ab}\pi_{cd}/(\pi_{ad}\pi_{bc})$ satisfy $(\chi^2-\psi^2-1)^2=4\psi^2$.
(1) The entries of a layer of Cauchy form are units, so for odd $p$ the $h_{ij}\equiv2\tilde M_{ij}$ are units modulo $p$ (for $p=2$ the hypothesis fails, the $h_{ij}$ being odd), and $4\pi_{ij}\equiv h_{ij}h_{ji}$. Let $A_1=|h_{ab}h_{ba}h_{cd}h_{dc}|$, $A_2=|h_{ac}h_{ca}h_{bd}h_{db}|$, $A_3=|h_{ad}h_{da}h_{bc}h_{cb}|$, odd integers at most $\Omega^4$. Then $\chi^2\equiv\varepsilon_2A_1/A_2$ and $\psi^2\equiv\varepsilon_3A_1/A_3$ with signs $\varepsilon_2,\varepsilon_3\in\{\pm1\}$, and multiplying the last relation by $A_2^2A_3^2$ shows that $p$ divides
\[
Z=\bigl(\varepsilon_2A_1A_3-\varepsilon_3A_1A_2-A_2A_3\bigr)^2-4\varepsilon_3A_1A_2^2A_3 .
\]
The bracket is a sum of three odd numbers, hence odd, so $Z$ is odd and in particular nonzero, and $|Z|\le(3\Omega^8)^2+4\Omega^{16}=13\Omega^{16}$.
(2) Over $\mathbb{Q}$ the same computation with $h_{ij}=2\tilde M_{ij}$, odd integers by assumption, gives $Z=0$ exactly, contradicting that $Z$ is odd.
\end{proof}


\section{The split frame}\label{sec:split}

\begin{definition}\label{def:split}
\lean{Z3NetF.splitRead, Z3NetF.SplitOK, Z3NetF.splitGraph}\leanok
Let $v_0,\dots,v_{s-1}$ (\emph{cell points}) and $u_0,\dots,u_{s-1}$ (\emph{offset points}) be $2s$ pairwise distinct rational numbers, let $D$ be the product of their denominators and of the numerators of their pairwise differences, $R=\Z[1/D]$ and $m=sk$. For $V\in R[t]_{<m}$ let coordinate $c$ have the $m$ digits
\[
[(t-v_c)^j]\,V\quad(j<k),\qquad\text{followed by}\qquad[(t-u_c)^j]\,V\quad(j<m-k).
\]
Put $\vartheta_{cd}=u_c-v_d$, $\Pi(t)=\prod_d(t-v_d)^k$, $\Pi_d=\Pi/(t-v_d)^k$ and $\Theta_c=\prod_{d}\vartheta_{cd}$, so that $\Pi(u_c)=\Theta_c^k$.
\end{definition}

\begin{lemma}[Box minors of the split frame]\label{lem:Tsplit}
\lean{Z3NetF.Tsplit_minor, Z3NetF.Tsplit_net_iff, Z3NetF.Tsplit_prime_bound, Z3NetF.split_spec, Z3NetF.split_boxMinor, Z3NetHL.HL_frameMinor}\leanok
For depths $(d_c)$ with $\sum_cd_c=m$ the box matrix of the split frame consists of the $\min(d_c,k)$-jets of $V$ at $v_c$ and the $(d_c-k)_+$-jets at $u_c$, and its determinant is, up to sign, the confluent Vandermonde product $\prod_{\{P,Q\}}(x_Q-x_P)^{\alpha_P\alpha_Q}$ over the $2s$ points $x_P$ with these multiplicities $\alpha_P$. In particular, for $s\ge4$, every $k\ge1$ and $q$ coprime to the denominators of the points, the split frame is a $(0,sk,s)$-net over $\Z/q\Z$ if and only if the $2s$ points are pairwise distinct modulo every prime factor of $q$ (which needs every prime factor of $q$ to be at least $2s$).
\end{lemma}

\begin{proof}
\leanok
\uses{lem:nie, lem:hermite}
The box matrix is the Hermite matrix of Lemma~\ref{lem:hermite} for the $2s$ points with the multiplicities $\alpha_P$, $\sum_P\alpha_P=\sum_cd_c=m$ (no point at $\infty$); this gives the value, a unit modulo $p$ if the points are distinct modulo $p$. Conversely, a coincidence modulo $p$ puts two equal rows into some box: $V(v_c)$ and $V(v_d)$ for $v_c\equiv v_d$ (the box $(k,\dots,k)$); $V(u_c)$ and $V(v_d)$ for $u_c\equiv v_d$, $c=d$ allowed (depth $k+1$ on $c$, at least $1$ on $d$); $V(u_c)$ and $V(u_d)$ for $u_c\equiv u_d$, $c\ne d$ (depths $k+1$ on $c$ and $d$, $k-1$ on two other coordinates, $k$ elsewhere; here $s\ge4$ is used). Then Lemma~\ref{lem:nie}.
\end{proof}

\begin{lemma}[First offset rows of the split frame]\label{lem:Rsplit}
\lean{Z3NetF.Rsplit_row0}\leanok
For every $k\ge1$ and $c\ne d$, row $0$ of $\Phi^S_{cd}$ has, with $x^n$ marking carry depth $n$, the generating function
\[
\frac{\Theta_c^k}{\Pi_d(v_d+x)\,(\vartheta_{cd}-x)}.
\]
\end{lemma}

\begin{proof}
\leanok
As $\deg V<sk=\deg\Pi$, partial fractions give $V=\sum_d\Pi_dN_d$ with $\deg N_d<k$; as $V\equiv\Pi_dN_d\pmod{(t-v_d)^k}$ and $\Pi_d(v_d)$ is a unit, the cell of $d$ determines $N_d$ and conversely. The first offset digit of $c$ is $V(u_c)$. For the column (level $j$ of the cell of $d$, all other cells $0$) we have $V=\Pi_dN_d$ with $N_d\equiv(t-v_d)^j/\Pi_d\pmod{(t-v_d)^k}$, i.e. $N_d=\sum_{j\le i<k}\varphi_{i-j}(t-v_d)^i$ with $\varphi(x)=\sum_n\varphi_nx^n=1/\Pi_d(v_d+x)$. Hence, with $\vartheta=\vartheta_{cd}$, $V(u_c)=\Pi_d(u_c)\sum_{j\le i<k}\varphi_{i-j}\vartheta^i$. Weighting the level $j=k-1-n$ by $x^n$ gives $\Pi_d(u_c)\,\vartheta^{k-1}\varphi(x)\sum_{r\ge0}(x/\vartheta)^r=\Pi_d(u_c)\vartheta^k\varphi(x)/(\vartheta-x)$ modulo $x^k$, and $\Pi_d(u_c)\vartheta^k=\Pi(u_c)=\Theta_c^k$.
\end{proof}

\begin{proposition}[Normal form and reachable designs of the split frame]\label{prop:Gsplit}
\lean{Z3NetF.Gsplit_normal, Z3NetF.Gsplit_nodes, Z3NetF.Gsplit_beta, Z3NetF.Gsplit_gauge, Z3NetF.Gsplit_row0, Z3NetF.Gsplit_tsk_unit, Z3NetF.Gsplit_tsk_net, Z3NetF.Gsplit_reach_iff, Z3NetF.Gsplit_poly_iff, Z3NetHL.HL_propN, Z3NetHL.HL_row0, Z3NetHL.HL_net, Z3NetHL.HL_rowDesign, Z3NetHL.HL_row0_lev, Z3NetHL.HL_net_lev, Z3NetHL.HL_relax_lev, Z3NetHL.HL_dom_lev, Z3NetHL.HL_dom_lev_real}\leanok
\uses{lem:box, lem:Ts, prop:Gs, def:Ns, cor:tsk, prop:F, lem:Tsplit}
Apply the gauge move of Lemma~\ref{lem:box}(c) with $L_c=\operatorname{diag}(\Theta_c^{-k},1,\dots,1)$ and $U_d=\operatorname{Toep}(\Pi_d(v_d+x))$. Then row $0$ of every off-diagonal block becomes $w_{cd}/(1-y_{cd}x)$ with
\[
w_{cd}=y_{cd}=\frac1{\vartheta_{cd}}=\frac1{u_c-v_d}\qquad(c\ne d),
\]
independently of $k$, and within each column $d$ the $y_{cd}$ ($c\ne d$) are pairwise distinct with unit differences $y_{cd}-y_{c'd}=(u_{c'}-u_c)/(\vartheta_{cd}\vartheta_{c'd})$. Consequently Propositions~\ref{prop:Gs} and~\ref{prop:F}, Definition~\ref{def:Ns} and Corollary~\ref{cor:tsk} hold verbatim for the split frame, with $Q_c$ replaced by $\Theta_c$ and Lemma~\ref{lem:Ts} by Lemma~\ref{lem:Tsplit}. In particular a design is reachable with row scalars $\lambda'$ if and only if, for every column $d$ and every $j\le s-2$,
\[
M^{(j)}[c][d]=\lambda'_c\,\frac{g_{d,j}(y_{cd})}{u_c-v_d}\qquad(c\ne d)
\]
for polynomials $g_{d,j}$ of degree $j$ with unit leading coefficient; and then, for every $k\ge1$ and every $q$ coprime to $DD'$, where $D'$ collects the primes of the units $\lambda'_c$, of the leading coefficients of the $g_{d,j}$ and of the denominators, the graph form $\mathcal P^{u,v,F}_{k,q}$ of Definition~\ref{def:Ns} (on the split frame) is a digital $(0,sk,s)$-net over $\Z/q\Z$ with one point in every cell, whose row $0$ equals the design on the levels $k-1,\dots,k-\min(k,s-1)$ and vanishes below.
\end{proposition}

\begin{proof}
\leanok
\uses{prop:Ns, prop:Gs, cor:tsk, prop:F, lem:Tsplit, lem:Rsplit}
The normal form is Lemma~\ref{lem:Rsplit} times $\Theta_c^{-k}$ and $\Pi_d(v_d+x)$, as in Proposition~\ref{prop:Ns}, with unit diagonal entries of $L_c$ and $U_d$. The proofs of Propositions~\ref{prop:Gs} and~\ref{prop:F} use only that the $w_{cd}$ are units and that the $y_{cd}$ of a column are distinct with unit differences (the matrix $(w_{cd}y_{cd}^n)_{c\ne d,\,n\le s-2}$ is then a unit diagonal times an invertible Vandermonde matrix), and the proof of Corollary~\ref{cor:tsk} uses only the unit box minors of the frame, here Lemma~\ref{lem:Tsplit}.
\end{proof}

\begin{lemma}[Half-integer designs]\label{lem:Hs}
\lean{Z3NetHL.HL_halfDesign, Z3NetHL.HL_Hs, Z3NetHL.HL_main, Z3NetHL.HL_main_real}\leanok
Let the $u_c$ be even and the $v_d$ odd integers, so that every $\vartheta_{cd}=u_c-v_d$ is odd, put $\Lambda_c=\prod_{d\ne c}\vartheta_{cd}^{\,s-1}$, and take $\lambda'_c=\frac12\Lambda_c$, $g_{d,0}=1$ and $g_{d,j}(y)=1+2y^j$ for $1\le j\le s-2$. Then
\[
M^{(0)}[c][d]=\frac{\Lambda_c}{2\vartheta_{cd}},\qquad M^{(j)}[c][d]=\frac{\Lambda_c}{2\vartheta_{cd}}+\frac{\Lambda_c}{\vartheta_{cd}^{\,j+1}}\quad(1\le j\le s-2),
\]
all in $\frac12+\Z$ ($c\ne d$),
so the off-diagonal part of $F_0$ lies in $\frac12+\Z$ and $F_j=M^{(j)}-M^{(j-1)}$ is an integer matrix for $j\ge1$. Complete the design by $F_0[c][c]=-\alpha_c$, with $\alpha_c$ the least element of $\frac12+\Z$ that is at least $A_c=\sum_{d\ne c}\max_{j\le s-2}|M^{(j)}[c][d]|$, and $F_j[c][c]=0$ for $j\ge1$. The primes of $DD'$ are $2$ and the primes of the point differences and of the $\Lambda_c$. For $u_c=2c$, $v_d=2d+1$ they are all at most $2s-1$.
\end{lemma}

\begin{proof}
\leanok
\uses{prop:Gsplit}
By Proposition~\ref{prop:Gsplit}, $M^{(0)}[c][d]=\frac12\Lambda_c\vartheta_{cd}^{-1}$ and $M^{(j)}[c][d]=\frac12\Lambda_c\vartheta_{cd}^{-1}(1+2\vartheta_{cd}^{-j})$ for $1\le j\le s-2$. For $1\le i\le s-1$, $\Lambda_c/\vartheta_{cd}^{\,i}$ is a product of odd integers, hence an odd integer; so the first term lies in $\frac12+\Z$ and the second in $\Z$. The leading coefficients are $1$ and $2$, and $\mu_d=1$. For $H_s$ the differences $u_c-u_{c'}$, $v_d-v_{d'}$ and $\vartheta_{cd}=2(c-d)-1$ have absolute value at most $2s-1$.
\end{proof}


\section{Separation with carries}

\begin{proposition}[Half-integer layers with carries]\label{prop:Hk}
\lean{Z3NetHL.HL_relax, Z3NetHL.HL_dom, Z3NetHL.HL_dom_real, Z3NetHL.GF_relax, Z3NetHL.GF_dom, Z3NetHL.GF_dom_real, Z3NetHL.GF_upper, Z3NetHL.HL_domL_real, Z3NetHL.half_upper, Z3NetHL.half_limit, Z3NetHL.acct_LBHk}\leanok
Let $p$ be an odd prime, $s\ge3$, $k\ge1$, $L=\min(k,s-1)$, and let a graph form over $\F_p$ be given (a digital net or not; no part uses $t=0$) with depths $(k,\dots,k)$ whose row $0$ of every block, diagonal included, equals a design $(F_0,\dots,F_{L-1})$ on the levels $k-1,\dots,k-L$ and vanishes below (the rows $\ge1$ are arbitrary). Assume
\[
\text{(Half)}\qquad F_0\in(\tfrac12+\Z)^{s\times s},\qquad F_i\in\Z^{s\times s}\ (1\le i\le L-1),
\]
with $\frac12\equiv(p+1)/2$. For $0\ne E\in\{-2,\dots,2\}^s$ (steps of two cells are included, because an offset change of about $\pm\frac p2$ can bring the points of cells two apart to about $1.5$ cells apart, below the cap $2$ of $\mathrm{LB}$) and carry depths $j_d\in\{0,\dots,L-1\}$ on the moving axes put $\Sigma E=\sum_cE_c$,
\[
\sigma_c=\sum_{d:\,E_d\ne0}\Bigl(E_dF_0[c][d]+\sgn(E_d)\sum_{i=1}^{j_d}F_i[c][d]\Bigr)
\]
(so $\sigma_c=\sum_dE_dM^{(j_d)}[c][d]$ if $|E|_\infty\le1$), let the \emph{options} of coordinate $c$ be $O_c=\{\sigma_c+\frac p2,\sigma_c-\frac p2\}$ if $\Sigma E$ is odd, $O_c=\{\sigma_c,\sigma_c-p\sgn\sigma_c\}$ if $\Sigma E$ is even ($O_c=\{0\}$ if $\sigma_c=0$), $b_c=\min_{t\in O_c}\max(0,|pE_c+t|-1)$, and $\mathrm{LB}=\min\bigl(2,\min_{E,j}|b|/p\bigr)$.
\begin{enumerate}
\item If $p>2\max|\sigma_c|$ (over all $E$, $j$ and $c$), then $\delta N^{1/s}\ge\mathrm{LB}$ (toroidal distance), for every choice of the rows $\ge1$.
\item Suppose moreover that $F_i[c][c]=0$ for $i\ge1$ and $F_0[c][c]=-\alpha_c$ with
\[
\alpha_c\ \ge\ A_c:=\sum_{d\ne c}\max_{j<L}\bigl|M^{(j)}[c][d]\bigr|.
\]
Let $B=\max_c(2\alpha_c+2A_c)$ (the dominance constant) and $\epsilon_p=(B+1)/p$. If $p>2B$, then $\delta N^{1/s}\ge\min\bigl(\sqrt s(\frac12-\epsilon_p),\sqrt2(1-\epsilon_p)\bigr)$.
\item Let $C=\max_{c,d}|F_0[c][d]|$. If $p>4sC+2$, then $\delta N^{1/s}\le\min(\sqrt s/2,\sqrt2)+\sqrt s\,(2C+1)/p$.
\end{enumerate}
\end{proposition}

\begin{proof}
\leanok
(1) In cell units, coordinate $c$ of the point of cell $n$ is $n_c+(r_c+f_c)/p$ with $r_c$ the first offset digit and $f_c\in[0,1)$. Let $n\ne n'$ and let $E$ be a torus representative of $n'-n$ (modulo $p^k$ in each coordinate). If some $|E_c|\ge3$, then the two points differ by $|D_c|>|E_c|-1\ge2$ cells in coordinate $c$, with $D$ their difference in cell units, since $(r'_c+f'_c)-(r_c+f_c)\in(-p,p)$. Otherwise adding $E_d$ to the cell integer $n_d$ changes its level $k-1$ by $E_d$ and the levels $k-1-i$, $1\le i\le j_d$, by $\sgn(E_d)$ (carries or borrows; a carry is at most $1$ since $|E_d|\le2<p$, and the torus wrap is the full carry); since row $0$ vanishes below level $k-L$, capping $j_d$ at $L-1$ is exact. So $\Delta r_c=r'_c-r_c\equiv\sigma_c\pmod p$ with $\Delta r_c\in(-p,p)$. As $F_0\in(\frac12+\Z)^{s\times s}$ and $F_i$ is integral, $\sigma_c\in\Sigma E/2+\Z$. If $\Sigma E$ is odd, $\sigma_c\pm\frac p2$ are the integers congruent to $\sigma_c$ (with $\frac12\equiv(p+1)/2$), and as $|\sigma_c|<p/2$ they are the only ones in $(-p,p)$; if $\Sigma E$ is even, $\sigma_c$ is an integer with $|\sigma_c|<p$, and the integers in $(-p,p)$ congruent to it are $\sigma_c$ and $\sigma_c-p\sgn\sigma_c$ (only $0$ if $\sigma_c=0$). So $\Delta r_c\in O_c$, and $p|D_c|=|pE_c+\Delta r_c+\Delta f_c|>|pE_c+\Delta r_c|-1$, i.e. $p|D_c|\ge b_c$.

(2) First $|\sigma_c|\le2\alpha_c+2A_c\le B$, since the own term is $-E_c\alpha_c$ and the term of an axis $d\ne c$ is $E_dM^{(j_d)}[c][d]$ ($|E_d|=1$) or $\sgn(E_d)(M^{(0)}+M^{(j_d)})[c][d]$ ($|E_d|=2$); so (1) applies. Let $e=|E|_\infty\le2$ and $|E_c|=e$. Then $E_c\sigma_c\le-e^2\alpha_c+e\cdot eA_c\le0$: every coordinate of maximal size is backward. If $\Sigma E$ is odd, every coordinate has $|pE_c+t|\ge p|E_c\pm\frac12|-|\sigma_c|\ge p/2-|\sigma_c|$ for $t\in O_c$, so $|D_c|\ge\frac12-\epsilon_p$ for all $s$ coordinates. If $\Sigma E$ is even, a backward coordinate $c$ with $|E_c|=e$ has $\sgn\sigma_c\ne\sgn E_c$, so both options give $|pE_c+t|\ge pe-|\sigma_c|$ and $|D_c|\ge e-\epsilon_p$. For $e=2$ this is at least $2-\epsilon_p\ge\sqrt2(1-\epsilon_p)$; for $e=1$ an even $\Sigma E$ needs at least two nonzero coordinates, each of maximal size, so $|D|\ge\sqrt2(1-\epsilon_p)$.

(3) Since $|2F_0[c][d]|\le2C<p$ is odd, every entry of $F_0$ is a unit modulo $p$, and $C<p/2-1$. First take the step $E=\e_a$ from a cell in which the last digit of $n_a$ is not $p-1$ (no carry) and $r_a$ wraps, i.e. $\Delta r_a=F_0[a][a]-\frac p2$; such cells exist, because $r_a$ depends on the last digit of $n_d$, $d\ne a$, with the unit coefficient $F_0[a][d]$. Then $|D_a|\le\frac12+(C+1)/p$, and every other coordinate has $\Delta r_c=F_0[c][a]\pm\frac p2$ and $|D_c|\le\frac12+(C+1)/p$; so $\delta N^{1/s}\le\sqrt s(\frac12+(C+1)/p)$. Next take $E=\e_a-\e_b$, with the last digit of $n_a$ not $p-1$ and that of $n_b$ not $0$ (no carries); then $\sigma_c=\beta_c:=F_0[c][a]-F_0[c][b]$ is an integer with $|\beta_c|\le2C$. Fix these two digits. For every $c$, $r_c$ depends with a unit coefficient on the last digit of some $n_g$, $g\notin\{a,b\}$ (as $s\ge3$), so it is uniformly distributed on the remaining cells, and the fraction of them on which $r_c+\beta_c\notin[0,p)$ is $|\beta_c|/p$. As $\sum_c|\beta_c|\le2sC<p$, some cell has no wrap; there $|D_c-E_c|\le(|\beta_c|+1)/p\le(2C+1)/p$ for every $c$, so $|D|\le\sqrt2+\sqrt s(2C+1)/p$.
\end{proof}

\begin{theorem}[Beyond the cubic lattice]\label{thm:HL}
\lean{Z3NetHL.HL_main, Z3NetHL.HL_main_real, Z3NetHL.HL_upper, Z3NetHL.HL_limit, Z3NetHL.HL_Bbound}\leanok
\uses{lem:Hs}
Let $s\ge4$. For the explicit design $H_s$ on the split frame $u_c=2c$, $v_d=2d+1$ (Lemma~\ref{lem:Hs}) there are constants $B_s=e^{O(s^2\log s)}$ and $C_s$ such that for every prime $p\ge\max(2s,2B_s+2)$ and every $k\ge1$ the net $\mathcal H^s_{k,p}$ is a digital $(0,sk,s)$-net over $\F_p$ with one point in every cell, and
\[
\min\Bigl(\sqrt s\bigl(\tfrac12-\epsilon_p\bigr),\ \sqrt2\,(1-\epsilon_p)\Bigr)\ \le\ \delta N^{1/s}\qquad\Bigl(\epsilon_p=\frac{B_s+1}p\Bigr),
\]
toroidal, hence Euclidean; if also $p>4sC_s+2$, then $\delta N^{1/s}\le\min(\sqrt s/2,\sqrt2)+\sqrt s\,(2C_s+1)/p$. So $\delta N^{1/s}\to\min(\sqrt s/2,\sqrt2)$ as $p\to\infty$, uniformly in $k$.
\end{theorem}

\begin{proof}[Proof of Theorem~\ref{thm:HL}]
\leanok
\uses{prop:Gsplit, lem:Hs, prop:Hk}
Let $H_s$ be the design of Lemma~\ref{lem:Hs} on $u_c=2c$, $v_d=2d+1$, let $\mathcal H^s_{k,p}=\mathcal P^{u,v,H_s}_{k,p}$, $B_s=\max_c(2\alpha_c+2A_c)$ and $C_s=\max_c\alpha_c$, which is the largest entry of $F_0$ in absolute value, since $|F_0[c][d]|\le A_c\le\alpha_c$. For a prime $p\ge2s$ and every $k\ge1$, Lemma~\ref{lem:Hs} and Proposition~\ref{prop:Gsplit} make $\mathcal H^s_{k,p}$ a digital $(0,sk,s)$-net with one point in every cell, whose row $0$ is $H_s$ truncated to $L=\min(k,s-1)$ levels and $0$ below. It satisfies (Half) and the dominance hypothesis of Proposition~\ref{prop:Hk}(2) (with $A_c$ over all $j\le s-2\ge L-1$), so Proposition~\ref{prop:Hk}(2) gives the lower bound (as $p\ge2B_s+2>2B_s$; for $p\le2B_s+2$ it would not be positive) and Proposition~\ref{prop:Hk}(3) the upper bound, both independent of $k$. Finally $|\vartheta_{cd}|\le2s-1$, so $|\Lambda_c|\le(2s-1)^{(s-1)^2}$, $|M^{(j)}[c][d]|\le\frac32|\Lambda_c|$ (absolute values, as $\Lambda_c$ is negative for some $c$ when $s$ is even), $\alpha_c\le A_c+1$ and $B_s\le6(s-1)\max_c|\Lambda_c|+2=e^{O(s^2\log s)}$.
\end{proof}


\section{One resolution, and searched designs}\label{sec:halfms}

\begin{lemma}[Two-digit relaxation]\label{lem:R2}
\lean{Z3NetHL.R2_sep}\leanok
\uses{prop:Hk}
Let $p$ be an odd prime, $s\ge3$, and let a graph form over $\F_p$ (a digital net or not) with the depths $(1,\dots,1)$ have first offset rows $F_0\in(\frac12+\Z)^{s\times s}$ and second offset rows $D_1\in\Z^{s\times s}$; the deeper rows are arbitrary. For $0\ne E\in\{-2,\dots,2\}^s$ let $\beta=F_0E$, $v=D_1E$, let $O_c$ be the options of Proposition~\ref{prop:Hk} (with $\sigma_c=\beta_c$), $U_c=\{v_c,\,v_c-p\operatorname{sgn}v_c\}$ ($U_c=\{0\}$ if $v_c=0$), and
\[
b_c=\min_{t\in O_c,\ u\in U_c}\max\bigl(0,\ |p^2E_c+pt+u|-1\bigr),\qquad \mathrm{LB}_2(p)=\min\Bigl(2,\ \min_E|b|/p^2\Bigr).
\]
If $p>\max|2\beta_c|$ and $p>\max|v_c|$ over these $E$ and $c$, then $\delta N^{1/s}\ge\mathrm{LB}_2(p)$ for the toroidal distance.
\end{lemma}

\begin{proof}
\leanok
\uses{prop:Hk}
In cell units coordinate $c$ of the point of cell $n$ is $n_c+r_c/p+(d_c+g_c)/p^2$, with $r_c$ and $d_c$ the first and second offset digits and $g_c\in[0,1)$. For two distinct points and a torus representative $E$ of their cell difference with some $|E_c|\ge3$, coordinate $c$ differs by more than $2$, as in the proof of Proposition~\ref{prop:Hk}(1). Otherwise, as there (at $k=1$ there are no carries), $\Delta r_c\in O_c$. As $d=D_1n\bmod p$ is linear in the cell, $\Delta d_c\equiv v_c\pmod p$ with $\Delta d_c\in(-p,p)$ and $|v_c|<p$, so $\Delta d_c\in U_c$; and $|\Delta g_c|<1$. So $p^2|D_c|=|p^2E_c+p\Delta r_c+\Delta d_c+\Delta g_c|\ge b_c$. The comparison with the toroidal distance is that of the proof of Proposition~\ref{prop:Hk}(1).
\end{proof}

\begin{proposition}[Dimension five, one resolution; computer-assisted]\label{prop:half5}
\lean{Z3NetHL.halfPhi_row, Z3NetHL.half_sep, Z3NetHL.half_LB2_ge, Z3NetHL.half_lam_le, Z3NetHL.half_lam_dom, Z3NetHL.half_upper, Z3NetHL.half_limit, Z3NetHL.half_net_generic, Z3NetHL.half_M_exists, Z3NetHL.qstar_net, Z3NetHL.qstar_sep, Z3NetHL.qstar_bounds, Z3NetHL.qstar_limit, Z3NetHL.qstar_hk, Z3NetHL.ell2_eq_halfEllInf}\leanok
Let $F_0^*$, $D_1$, $D_2$, $D_3$ be the matrices
\[
2F_0^*=\left(\begin{smallmatrix}-3&-3&-3&-1&-1\\3&-3&-1&-3&1\\3&1&-3&1&-1\\1&3&-1&-3&-1\\1&-1&1&1&-1\end{smallmatrix}\right),\quad
D_1=\left(\begin{smallmatrix}-2&0&-1&-1&1\\0&-1&-1&-1&-1\\2&-1&-3&1&0\\-1&2&-1&-2&0\\-1&-1&3&2&-5\end{smallmatrix}\right),
\]
\[
D_2=\left(\begin{smallmatrix}0&-1&-1&0&0\\-1&0&0&0&-1\\-1&-1&0&0&0\\-1&-1&0&0&0\\-1&-1&0&0&0\end{smallmatrix}\right),\quad
D_3=\left(\begin{smallmatrix}0&-1&0&0&0\\-1&0&0&0&0\\-1&0&0&0&0\\-1&0&0&0&0\\-1&0&0&0&0\end{smallmatrix}\right),
\]
and, for an odd prime $p$, let $\mathcal Q^*_p$ be the digital net over $\F_p$ whose point with index $n\in\F_p^5$ has, in coordinate $c$, the digits $n_c$, $(F_0^*n)_c$, $(D_1n)_c$, $(D_2n)_c$, $(D_3n)_c$ (with $\frac12\equiv(p+1)/2$).
\begin{enumerate}
\item $\mathcal Q^*_p$ is a $(0,5,5)$-net for every prime $p\ge7$.
\item For every prime $p\ge37$, $\delta N^{1/5}>1$ (toroidal), and $\delta N^{1/5}\to\sqrt5/2=1.1180\ldots$ as $p\to\infty$.
\item For every prime $p\ge53$, every graph form over $\F_p$ with the depths $(1,\dots,1)$ and first offset rows $F_0^*$ (deeper rows arbitrary) has $\delta N^{1/5}>1$ (toroidal).
\end{enumerate}
\end{proposition}

\begin{proof}[Proof (computer-assisted)]
\leanok
\uses{lem:nie, lem:box, prop:Hk, lem:R2}
(1) The generator matrix of coordinate $c$ has the rows $\e_c$ and row $c$ of $F_0^*$, $D_1$, $D_2$, $D_3$. By Lemma~\ref{lem:box} its box minors are the minors of the stacked offset rows with disjoint row and column sets; each of the $126$ of them, computed in exact rational arithmetic, is a nonzero rational whose numerator and denominator have no prime factors other than $2$, $3$ and $5$. Lemma~\ref{lem:nie} concludes.

(2) Over $E\in\{-2,\dots,2\}^5$, $\max|2(F_0^*E)_c|=22$ and $\max|(D_1E)_c|=24$, so Lemma~\ref{lem:R2} applies for $p\ge29$. Evaluated exactly, $\mathrm{LB}_2(p)>1$ for every prime $37\le p<200$ (the least value is $1.0127\ldots$, at $p=37$; and $\mathrm{LB}_2(31)=0.9918\ldots$), and $\mathrm{LB}_2^*(200)^2=1.1877\ldots>1$, so $\mathrm{LB}_2(p)\ge\mathrm{LB}_2^*(p)\ge\mathrm{LB}_2^*(200)>1$ for every $p\ge200$. For the limit, with the row-sum norm $\|F_0^*\|_\infty=\max_c\sum_d|F_0^*[c][d]|=\frac{11}2$ every option of Proposition~\ref{prop:Hk}(1) is at least $p\,\ell_c(E)-(2\|F_0^*\|_\infty+1)$, where
\[
\ell_c(E)=\min\bigl(|E_c+\tfrac12|,|E_c-\tfrac12|\bigr)\ (\Sigma E\text{ odd}),\qquad\ell_c(E)=\min_{\omega}|E_c-\omega|\ (\Sigma E\text{ even}),
\]
$\omega$ over $\{0,\operatorname{sgn}\beta_c\}$ (over $\{0\}$ if $\beta_c=0$). So $\delta N^{1/5}\ge\lambda-\sqrt5(2\|F_0^*\|_\infty+1)/p$ with $\lambda=\min_E|\ell(E)|$. If $\Sigma E$ is odd, every coordinate gives at least $\frac12$, with equality exactly on the binding steps; if $\Sigma E$ is even, $|\ell(E)|^2$ is an integer, and a finite check over the $5^5-1$ vectors shows that it is at least $2$. So $\lambda=\sqrt5/2$, and Proposition~\ref{prop:Hk}(3) with $C=\frac32$ gives $\delta N^{1/5}\le\sqrt5/2+4\sqrt5/p$ for $p>32$.

(3) Proposition~\ref{prop:Hk}(1) with $k=1$ and $F_0=F_0^*$ applies for $p>22$. Evaluated exactly, $\mathrm{LB}>1$ for every prime $53\le p<200$ ($\mathrm{LB}(47)=0.9906\ldots$, $\mathrm{LB}(53)=1.0049\ldots$), and the nondecreasing bound $\mathrm{LB}^*$ after Proposition~\ref{prop:Hk} has $\mathrm{LB}^*(200)^2=1.1689\ldots>1$. (The monotone bounds already exceed $1$ from $p=50$ for $\mathrm{LB}_2^*$ and from $p=63$ for $\mathrm{LB}^*$, so the exact scans up to $200$ in (2) and (3) are more than needed.)
\end{proof}

\begin{remark}[First-order accounting]\label{rem:accounting}
\lean{Z3NetHL.sepN_spec, Z3NetHL.ell2_eq_halfEllInf, Z3NetHL.acct_upper, Z3NetHL.acct_lower, Z3NetHL.acct_main, Z3NetHL.acct_liminf, Z3NetHL.acct_exact, Z3NetHL.acct_qstar, Z3NetHL.acct_counterexample, Z3NetHL.acct_LB2, Z3NetHL.acct_LBHk}\leanok
\uses{prop:Hk, lem:R2, prop:half5, thm:D}
Let $F_0\in(\frac12+\Z)^{5\times5}$ with $\det F_0\ne0$ satisfy
\[
\text{(Gap)}\qquad |\ell(E)|>\sqrt5/2\ \text{for every } 0\ne E\in\{-2,\dots,2\}^5 \text{ that is not a binding step},
\]
with $\ell$ as in the proof of Proposition~\ref{prop:half5}, let $D_1,D_2,D_3$ be fixed integer matrices, and let $K_{\mathrm{eff}}(p)=\frac{5p}2\bigl(1-\delta N^{1/5}/\frac{\sqrt5}2\bigr)$ for the digital net with these rows over $\F_p$, normalized so that $K_{\mathrm{eff}}\to\mathrm{dr}$ when the second row is lean. Then
\[
\mathrm{dr}(F_0)\ \le\ \liminf_{p\to\infty}K_{\mathrm{eff}}(p)\ \le\ \limsup_{p\to\infty}K_{\mathrm{eff}}(p)\ \le\ \mathrm{dj}(F_0,D_1).
\]
In particular, if $\mathrm{dj}(F_0,D_1)=\mathrm{dr}(F_0)$, then $\delta N^{1/5}=\frac{\sqrt5}2\bigl(1-\frac{2\,\mathrm{dr}(F_0)}{5p}\bigr)+O(p^{-2})$. For $\mathcal Q^*_p$ the finite checks give $\mathrm{dr}=\mathrm{dj}=\frac{17}2$, $\det2F_0^*=-880$ and $|\ell(E)|^2\ge2$ on the non-binding steps, so $\delta N^{1/5}=\frac{\sqrt5}2\bigl(1-\frac{17}{5p}\bigr)+O(p^{-2})$.

\emph{Proof.} Upper bound for $K_{\mathrm{eff}}$: on a binding step each option of Lemma~\ref{lem:R2} is $p^2\bigl(\frac12-x/p\bigr)+O(1)$ in absolute value, and the smallest has $x=\mathrm{dr}_c(E)+[c\text{ harmful}]$: the shorter option of the geometric digit $r$ shortens by $\mathrm{dr}_c(E)$, a second-digit wrap of the sign $\varsigma_c(E)$ by $1$ more, a residue of the other sign by $O(1)$ only, and for $E_c=0$ the longer $r$-option combined with a harmful wrap is not shorter, as $|\beta_c|\ge\frac12$. So these steps give $|b|^2/p^4=\frac54-(\mathrm{dr}(E)+\#\text{harmful})/p+O(p^{-2})$. By (Gap) the other steps of $\{-2,\dots,2\}^5$ have $|b|/p^2\to|\ell(E)|>\sqrt5/2$. Hence $\mathrm{LB}_2(p)=\frac{\sqrt5}2\bigl(1-\frac{2\,\mathrm{dj}}{5p}\bigr)+O(p^{-2})$, and $\delta N^{1/5}\ge\mathrm{LB}_2(p)$ gives $\limsup K_{\mathrm{eff}}\le\mathrm{dj}$. Lower bound: fix a binding $E$ with $\mathrm{dr}(E)=\mathrm{dr}(F_0)$. For $p$ large ($p\nmid\det2F_0$, $p$ larger than the entries of $D_1$ and than $2|2\beta_c|+2$) $F_0$ is invertible modulo $p$, so $n\mapsto r=F_0n\bmod p$ is a bijection, and the cells whose five first offset digits all take the shorter option for $E$ form a box in the $r$-coordinates with at least $(\frac p2-O(1))^5$ cells. A second-digit wrap $\Delta d_c\ne v_c$ needs $d_c(n)$ in a set of $|v_c|$ residues, which happens on $|v_c|p^4$ cells if $v_c\ne0$ (as $d_c$ is a nonzero linear form) and never if $v_c=0$. So for $p$ large some cell has the shorter options and $\Delta d=v$, and its step $E$ gives $|D|^2=\frac54-\mathrm{dr}(F_0)/p+O(p^{-2})$, i.e. $\liminf K_{\mathrm{eff}}\ge\mathrm{dr}(F_0)$. In the same way Proposition~\ref{prop:Hk}(1) has $\mathrm{LB}=\frac{\sqrt5}2\bigl(1-\frac{2(\mathrm{dr}(F_0)+5)}{5p}\bigr)+O(p^{-2})$.

(Gap) is a genuine hypothesis. The layer
\[
2F_0=\left(\begin{smallmatrix}-3&1&-1&-3&-3\\-3&-3&3&1&-1\\-3&1&-3&3&1\\1&-1&3&-3&1\\1&1&1&1&-1\end{smallmatrix}\right)
\]
satisfies (B2) and (MIN), defined before Theorem~\ref{thm:D}, and has $\det2F_0=560$, but the even step $E=(1,-1,-2,-1,-1)$ has $\beta=(2,-3,-1,-1,-1)$ and $|\ell(E)|=1$: its coordinates with $E_c=\pm1$ move forward, so the step is short (length $1$). Each half of the remark uses only one hypothesis: (Gap) only the upper bound $\limsup K_{\mathrm{eff}}\le\mathrm{dj}$, and $\det F_0\ne0$ only the lower bound; the deeper rows are arbitrary in both. With (Gap) the accounting is an identity for the relaxations and a two-sided bound for the nets, exact when the second row is lean. Cauchy--Schwarz, $\sum_c(\frac12-x_c/p)^2\ge\frac15(\frac52-K/p)^2$, makes $\delta N^{1/5}>1$ plausible once $p\gtrsim\frac{2K}{5(1-2/\sqrt5)}\approx3.79\,K$; this is a first-order estimate, not a threshold. It predicts $51$ for $F_0^*$ with Proposition~\ref{prop:Hk}(1) ($K=\frac{17}2+5$) and $32$ for $\mathcal Q^*_p$ ($K=\frac{17}2$), against the certified $53$ and $37$.
\end{remark}

\begin{theorem}[The least drift]\label{thm:D}
\lean{Z3NetHL.thmD, Z3NetHL.thmD_star, Z3NetHL.thmD_skew, Z3NetHL.thmD_conj}\leanok
Every $F_0\in(\frac12+\Z)^{5\times5}$ with (B2) and (MIN), with entries of any size, has $\mathrm{dr}(F_0)\ge\frac{17}2$, and $\mathrm{dr}(F_0^*)=\frac{17}2$. The part of the proof for layers whose off-diagonal part is not skew-symmetric is computer-assisted.
\end{theorem}

\begin{proof}
\leanok
\emph{Diagonal.} $\mathrm{dr}(E)$ contains $\alpha_c$ exactly for the coordinates with $E_c\ne0$, since $-E_c\beta_c=\alpha_c-E_c\sum_{d\ne c}F_0[c][d]E_d$, and $\beta_c$ does not involve $\alpha_c$ when $E_c=0$. So $\mathrm{dr}(F_0)$ does not increase when every $\alpha_c$ is lowered to $m_c$, which keeps (B2) and (MIN); we assume $\alpha=m$.

\emph{Minor tax.} At most one row of $F_0$ has all its off-diagonal entries in $\{\pm\frac12\}$: if rows $a\ne b$ did, the three ratios $F_0[a][c]/F_0[b][c]$, $c\notin\{a,b\}$, would lie in $\{\pm1\}$, two would be equal, and the $2\times2$ minor on those two columns would vanish. So at most one $m_c$ is $\frac12$: at most one row is \emph{light}, with all off-diagonal entries $\pm\frac12$, and the others are \emph{heavy}, and $\sum_c\alpha_c\ge4\cdot\frac32+\frac12=\frac{13}2$.

\emph{Balance on the full steps.} A \emph{full} step has all $E_c=\pm1$. For $E\in\{\pm1\}^5$ (binding, as $5$ is odd), $\mathrm{dr}(E)=-E^{\!\top}F_0E=\sum_c\alpha_c+\phi(E)$ with $\phi(E)=\sum_{c<d}\gamma_{cd}E_cE_d$ and $\gamma_{cd}=-(F_0[c][d]+F_0[d][c])\in\Z$. The function $\phi$ has mean $0$ over $\{\pm1\}^5$, and it vanishes identically if and only if the off-diagonal part of $F_0$ is skew-symmetric; otherwise $\max\phi\ge1$, as $\phi$ is integer-valued. So $\mathrm{dr}(F_0)\ge\sum_c\alpha_c$, and $\mathrm{dr}(F_0)\ge\sum_c\alpha_c+1$ if the off-diagonal part is not skew.

\emph{Skew layers.} If some entry has $|F_0[c][d]|\ge\frac52$, it occurs in the rows $c$ and $d$, and with the minor tax $\sum\alpha\ge\frac52+\frac52+\frac32+\frac32+\frac12=\frac{17}2$. Otherwise all off-diagonal entries are $\pm\frac12$ or $\pm\frac32$; let $G$ be the graph on the five coordinates with the edges $\{c,d\}$ where $|F_0[c][d]|=\frac32$. For two coordinates $x,y$ and the set $S$ of the other three, the binding steps with $E_x=E_y=0$ and $E_S\in\{\pm1\}^3$ have, as the skew terms inside $S$ cancel, $\mathrm{dr}(E)=\sum_{c\in S}\alpha_c+|u\cdot E_S|+|w\cdot E_S|$ with $u=F_0[x,S]$, $w=F_0[y,S]$. The step with $E_S=\sgn u$ (coordinatewise) has $u\cdot E_S=\|u\|_1$ and $w\cdot E_S\in\frac12+\Z$, as $w$ has three entries in $\frac12+\Z$; so
\[
\max_{E_S}\bigl(|u\cdot E_S|+|w\cdot E_S|\bigr)\ \ge\ \|u\|_1+\tfrac12,\qquad \|u\|_1=\tfrac32+\deg_S(x).
\] (i) If there is a light row $\ell$, then $\sum\alpha=\frac{13}2$, the other four rows have degree at least $1$ in $G$, and $\ell$ is isolated. $G$ is not a perfect matching on the four heavy rows: for an edge $\{a,b\}$ of such a matching, the rows $a$ and $b$ would be $\pm\frac12$ on the three columns outside $\{a,b\}$, which (MIN) excludes as in the minor tax. So some heavy $h$ has degree at least $2$, and the pair $\{x,y\}=\{h,\ell\}$ gives $\mathrm{dr}(F_0)\ge\frac92+(\frac32+2+\frac12)=\frac{17}2$. (ii) If there is no light row, $\sum\alpha=\frac{15}2$ and every degree is at least $1$; as $5$ is odd, $G$ is not a perfect matching, so some $x$ has degree at least $2$, and with $y$ outside $x$ and its neighbours if $\deg x=2$ (any $y\ne x$ otherwise) we get $\deg_S(x)\ge2$ and again $\mathrm{dr}(F_0)\ge\frac92+4=\frac{17}2$.

\emph{Layers that are not skew.} If some entry has absolute value at least $\frac52$, then $\sum\alpha\ge\frac52+3\cdot\frac32+\frac12=\frac{15}2$, and balance gives $\mathrm{dr}\ge\frac{17}2$; the same holds if $\sum\alpha\ge\frac{15}2$. There remain the layers with entries in $\{\pm\frac12,\pm\frac32\}$, one light row and four heavy rows with $\alpha=\frac32$. Conjugation by permutation matrices and by diagonal sign matrices preserves the binding steps, $\mathrm{dr}$, (B2) and (MIN), so we may take the light row to be row $4$, equal to $(\frac12,\frac12,\frac12,\frac12,-\frac12)$. Each heavy row then has $240$ possibilities, and an exhaustive enumeration with the (MIN) conditions finds $1{,}004{,}544$ layers satisfying (MIN), none with $\mathrm{dr}\le\frac{15}2$. Finally $\mathrm{dr}(F_0^*)=\frac{17}2$ is a finite check over the $61$ binding steps up to sign.
\end{proof}

\begin{remark}[Accounting for a dominant layer]\label{rem:M5}
\uses{prop:Hk, thm:HL, thm:D}
For comparison, the layer
\[
2M_5=\left(\begin{smallmatrix}-5&3&1&-3&1\\-3&-3&3&-1&-1\\1&-1&-5&-3&3\\-1&-1&3&-7&1\\3&-1&1&-1&-3\end{smallmatrix}\right)
\]
satisfies (B2) with room to spare: $\alpha=\frac12(5,3,5,7,3)$, $\sum\alpha=\frac{23}2$, against $\sum m=\frac{15}2$; and its symmetric part has $\max\phi=6$. Its drift $\mathrm{dr}(M_5)=\frac{35}2$ is attained on full steps, and as an accounting identity
\[
\tfrac{35}2=\tfrac{13}2+1+4+6:
\]
$\frac{13}2$, the least possible $\sum_cm_c$ under (MIN) (the minor tax of the proof of Theorem~\ref{thm:D}); $1$, the excess of the actual $\sum_cm_c=\frac{15}2$ of $M_5$ over it; $4$, its diagonal above $m$; and $6$, the asymmetric part ($\max\phi$).
This is an accounting, not a causal decomposition: the drift of $M_5$ is attained on full steps, that of $F_0^*$ on steps of weight three. With the diagonal lowered to $m$ the drift of $M_5$ would be $\frac{27}2$. Proposition~\ref{prop:Hk}(1) applies to $M_5$ for $p>26$ and gives $\delta N^{1/5}>1$ only from $p=89$ on ($\mathrm{LB}(83)=0.9975\ldots$, $\mathrm{LB}(89)=1.0056\ldots$, computed), against $53$ for $F_0^*$ (Figure~\texttt{fig:half5}). Dominance, which Proposition~\ref{prop:Hk}(2) and Theorem~\ref{thm:HL} use to obtain a bound for all designs at once, is what the accounting says to avoid: by the diagonal step of the proof of Theorem~\ref{thm:D} it costs $\sum_c(\alpha_c-m_c)$ in drift.
\end{remark}

\begin{proposition}[Dimension five; computer-assisted]\label{prop:HL5}
\lean{Z3NetHL.B5_design, Z3NetHL.B5_net, Z3NetHL.B5_sep, Z3NetHL.B5_sep_real, Z3NetHL.B5_net_lev}\leanok
Let $\mathcal H^{5*}_{k,p}=\mathcal P^{u,v,F}_{k,p}$ for the split frame $u=(-2,0,2,4,8)$, $v=(5,7,-1,3,1)$ and the design $F=H_5^*$ of Appendix~\texttt{app:designs}. For every prime $p\notin\{2,3,5,7,97,101\}$ and every $k\ge1$ it is a digital $(0,5k,5)$-net over $\F_p$ with one point in every cell ($p=97,101$ enter only for $k\ge4$), and for every prime $p\ge1511$ and every $k\ge1$,
\[
\delta N^{1/5}>1.0265 .
\]
\end{proposition}

\begin{proof}[Proof (computer-assisted)]
\leanok
\uses{prop:Gsplit, prop:Hk}
The design $H_5^*$ satisfies (Half); its geometric layer $F_0$ has the row scalars $\lambda'=\frac12(-9,-3,-3,3,9)$ and the column scalars $\mu=(5,5,5,-5,1)$. The primes of the point differences, of the units $\lambda'_c\mu_d$ and of the leading coefficients of the column polynomials $g_{d,j}$ (exact interpolation) are $2,3,5,7$ and $97,101$, the last two only in leading coefficients at carry depth $j=3$; Proposition~\ref{prop:Gsplit} gives the net statement. For the separation, $\max|\sigma_c|=253$, so Proposition~\ref{prop:Hk}(1) applies for $p>506$; the monotone bound $\mathrm{LB}^*$, minimized over the number of design levels $L=\min(k,4)\in\{1,2,3,4\}$ (that is, over $k$), is $1.026529\ldots$ at $p=1511$, the same value as $\mathrm{LB}$ (attained at $L=4$, i.e. $k\ge4$). As $\mathrm{LB}^*$ is nondecreasing in $p$ for each $L$, $\delta N^{1/5}\ge\mathrm{LB}^*(p)\ge\mathrm{LB}^*(1511)=1.026529\ldots>1.0265$ for every prime $p\ge1511$ and every $k$.
\end{proof}


\chapter{Open problems}\label{sec:remarks}

\begin{remark}[Open problems]\label{rem:open}
\uses{thm:A, prop:F, lem:skellat, thm:S2, lem:frameid, cor:frameid, prop:korobov4, thm:col, lem:shared, thm:base23, prop:ceil, cor:ceil3, rem:sign, prop:4D2, prop:k1, prop:k1designs, thm:L, thm:O5, prop:E, thm:Lstar, thm:O, cor:linoff, prop:twocoset, prop:wnet, prop:wfamily, rem:wstar, prop:half5, prop:HL5, rem:S1Elim}
\emph{(I) Bounded layers.} (1) \emph{Growth of the deficit.} Let $K(s)$ be the best first-order deficit $K$ in $\delta N^{1/s}=1-K/q+O(q^{-2})$, $q\to\infty$, over families of digital $(0,sk,s)$-nets over $\Z/q\Z$ whose geometric layers have entries bounded independently of $q$ and no zero diagonal entry. Lower bounds: $K(s)\ge\sqrt{(s-3)/8}$ (Theorem~\ref{thm:L}). For designs on the places $\infty,0,\dots,s-2$: $K(F)\ge s^{1-o(1)}$ (Theorem~\ref{thm:Lstar}); if moreover the row scalars $\lambda'_c$ (Proposition~\ref{prop:F}) take one value on the finite places, $K(F)\ge\operatorname{lcm}(1,\dots,s-3)$ (Proposition~\ref{prop:E}); and the two-axis bound $K_{\mathrm{pair}}$ is at least $2$ for $s\ge5$ (Theorem~\ref{thm:O5}). Upper bounds: the constructions here achieve $1$ for $s=3$, $2$ for $s=4$ (Proposition~\ref{prop:4D2}), $5$ and $10$ for $s=5,6$ (Table~\texttt{tab:Kdesigns}), and a certified deficit at most $a^*_s\le e^{s\log s+O(s)}$ for every $s\ge4$ (Theorem~\ref{thm:A}). Does $K(s)$ grow polynomially or exponentially? On these places a certified deficit polynomial in $s$ would need row scalars that differ between the finite places (Proposition~\ref{prop:E}). Also open: $K=1$ for $s=4$; whether layers with zero diagonal do better (Remark~\ref{rem:diag0}); and the optimal second-order constant $C$ in $\delta N^{1/3}=1-1/p+C/p^2+o(p^{-2})$ in dimension three. For nets with the layer of $\mathcal P_{k,p}$ it is at most $4$ for $k\ge4$ (after Corollary~\ref{cor:ceil3}); for other bounded layers Proposition~\ref{prop:ceil}(i) bounds it by $\frac12\bigl((|\tilde M_{ba}|+1)^2+(|\tilde M_{b'a}|+1)^2\bigr)$ for a coordinate $a$ with $\tilde M_{aa}=-1$, and for zero diagonals (Remark~\ref{rem:diag0}) we have no bound. (2) \emph{Small bases.} Theorem~\ref{thm:A} needs $q\ge5a^*_s+4$, and the designs of Table~\texttt{tab:Kdesigns} certify smaller bases only by finite scans; for $s=4$ the bases $q\le15$ have only weaker proved constants, and for $s\ge4$ Theorem~\ref{thm:A} excludes even $q$ (it needs $\operatorname{lpf}(q)\ge s-1$; for $q=2$ no $(0,m,s)$-net with $s\ge4$ exists, as $s\le q+1$). In dimension three the bases $2$ and $3$ are covered for every $m$ (Theorem~\ref{thm:base23}, computer-assisted); at $m=3$ the column nets cover every base coprime to $6$ (Theorem~\ref{thm:col}), column nets do not exist in even bases (Theorem~\ref{thm:O}), and the other bases $q\ge4$ divisible by $2$ or $3$ are open at $m=3$.

\emph{(II) Lattice skeletons.} (3) For $s\ge5$ the residues $1\le r\le s-3$ ($\ell\ge3$ skeleton coordinates) are open for all $p$; per-prime nets exist for $s=5$, $r=2$ at seven primes (computed, after Corollary~\ref{cor:frameid}), and a family for all large $p$ is open. Lattices with the optimal minimum order $p^{\ell/s}$ exist by a counting argument (after Lemma~\ref{lem:skellat}); explicit ones are known here only for $\ell=1$, $\ell=2$ and $r=1$, and for $3\le\ell\le s-2$ an explicit family is open (a mixed radix in the ring of integers of a degree-$\ell$ field is a natural candidate). The bottleneck is to realize such a skeleton with bounded faults: inside the frame orbit Lemma~\ref{lem:frameid} excludes the Korobov skeleton digits of $\mathcal B^1$ with small faults (Corollary~\ref{cor:frameid}); other Korobov vectors, faults growing slowly with $p$, and completions outside the frame orbit are not excluded, and per prime frame-orbit skeletons with faults $\pm1$ work (Proposition~\ref{prop:korobov4}). For one skeleton coordinate it is open whether the certified Euclidean constants tend to $\sqrt{\gamma_s}$ (Remark~\ref{rem:S1Elim}) and which explicit $\kappa$ attain it; for two skeleton coordinates, Euclidean (rather than maximum-norm) bounds are open. (4) In dimension four: $m\equiv1\pmod4$ for primes $p>83$ (no closed form is known; per-prime nets are known at every prime $7\le p\le83$, after Corollary~\ref{cor:frameid} (computed), Proposition~\ref{prop:korobov4} works out $p=13$, and we know of no obstruction for larger $p$), $m=4$, and small primes ($m\equiv1$ at $p=5$, $m\equiv2$ at $p\le19$, $m\equiv3$ at $p\le13$). Also open: an explicit $p_0(s)$ in Theorem~\ref{thm:S2} for $s\ge9$, and a covering radius uniform in $p$ for $\ell\ge2$. (5) In dimension three: a mesh ratio uniform in $p$ for $m\equiv1$; $m\in\{4,7\}$ for $p\le31$ and $m=5$ for $p\le7$; and an explicit family for all large primes with limit $>1$ at the resolutions $m\ge4$ (at $m=3$ the nets $\mathcal T_p$ and $\mathcal W_p$ of Propositions~\ref{prop:twocoset} and~\ref{prop:wnet} are such families). At $m=3$, among constants below $1$: whether digital $(0,3,3)$-nets that are not column point sets have smaller deficits than $\frac13$ and $\frac23$ (Theorem~\ref{thm:O} settles this among column point sets; outside them the drift-and-jitter bookkeeping of Section~\ref{sec:drift} is only a heuristic).

\emph{(III) Large-entry layers.} (6) (a) Smaller bases for $s\ge5$: Proposition~\ref{prop:HL5} needs $p\ge1511$, and at one resolution it is open whether some $(0,5,5)$-net over $\F_p$ with $p\le31$ has $\delta N^{1/5}>1$ (Proposition~\ref{prop:half5}). (b) The drift-and-jitter bookkeeping for half-integer layers at $k\ge2$, where the split-frame layers are Cauchy matrices, and for $s\ge6$. (c) Split-frame designs with explicit constants of moderate size for $s=6,7,8$. (d) Half-integer layers at the resolutions $m\not\equiv0\pmod s$. (e) Other split-frame layers beyond $D_s^+$. (7) In dimension three: whether a geometric layer $\tilde M\approx qA$ with a fixed rational matrix $A$ can be realized for every $k$ with $\delta N^{1/3}$ bounded above $1$. At $k=1$, layers whose geometric digits wrap together in two or three coordinates (locked wraps, Section~\ref{sec:wnet}) give the limits $\sqrt{66}/8$ (Proposition~\ref{prop:twocoset}), $\sqrt{281}/16$ (Proposition~\ref{prop:wnet}) and $\mu^*=1.0480\ldots$ (Remark~\ref{rem:wstar}), the optimum of the layers of rank one with image $\F_p(2,2,1)$ (Proposition~\ref{prop:wfamily}); the optimum over the other images of rank one, over layers of rank two and over non-digital point sets is open. At $k\ge2$ the geometric digits of any two coordinates are jointly uniformly distributed (Lemma~\ref{lem:shared}, as $2(k+1)\le3k$), so no relation between the geometric digits of two coordinates survives, and neither the two-coordinate lock of $\mathcal T_p$ nor the three-coordinate lock of $\mathcal W_p$ extends to $k\ge2$. For digital nets the only linear relation among the three geometric digits of one index that this count leaves is one of full support, $h_xr_x+h_yr_y+h_zr_z\equiv0$ with all $h_c\ne0$; whether such a relation, or a correlation that involves digits of other indices, can raise the constant above $1$ is open.

\emph{The case $k=1$ and digital shifts.} (8) At $k=1$: whether $K_1=1$ is possible for $s=5$, or $K_1=2$ for $s=6$ (Proposition~\ref{prop:k1designs}); frameless layers ($F_0$ chosen directly, not from a frame, Proposition~\ref{prop:k1designs}(2)) for $s\ge7$; and the choice of the deeper rows, which are free at $k=1$ apart from $t=0$ and change the exact separation at small $q$ (Proposition~\ref{prop:k1} and the remark after it). (9) In the bases two and three the constants of Theorem~\ref{thm:base23}(3) are certified lower bounds, not exact values; for $b=3$ and $m\equiv2\pmod3$ the infimum over $m$ and $\tau$ for these nets lies between $0.0770$ and $0.0901$, the latter the exact value at $m=20$ (computed). Whether other seeds or lift parameters give larger every-shift constants in base three is open.

\emph{Sequences.} (10) All nets here depend on $m$; they are not prefixes of one sequence, and the question of \cite{DGS2} for sequences remains open in dimension $s\ge3$.

\emph{Beyond digital nets.} (11) Can nets that are not digital, for instance permutation nets, exceed the constants of the digital nets here, for instance by avoiding the slip planes that linear offsets force (Corollary~\ref{cor:linoff})?
\end{remark}

\appendix


\chapter{Deferred proofs and side results}\label{app:proofs}


\section{A conjectural limit for one skeleton coordinate (heuristic)}

\begin{remark}[The limit $\sqrt{\gamma_s}$]\label{rem:S1Elim}
\uses{prop:S1E, tab:S1E, thm:HL}
Table~\ref{tab:S1E} suggests that $\max_\kappa\lambda_{\mathrm{jit}}(\kappa)p^{-1/s}\to\sqrt{\gamma_s}$, the square root of the Hermite constant ($1.122$, $1.189$, $1.231$, $1.290$, $1.345$, $\sqrt2$ for $s=3,\dots,8$); here $\lambda_{\mathrm{jit}}(\kappa)$ is the Euclidean minimum of the skeleton lattice $\Lambda_\kappa$ after one fine cell of jitter is subtracted in each coordinate (Proposition~\ref{prop:S1E}). The upper bound holds by the definition of $\gamma_s$, since $\lambda_{\mathrm{jit}}\le\lambda_2(\Lambda_\kappa)\le\sqrt{\gamma_s}\,p^{1/s}$. For the lower bound it would suffice that the index-$p$ sublattices of $\Z^s$, scaled to covolume $1$, equidistribute in $\mathrm{SL}_s(\Z)\backslash\mathrm{SL}_s(\R)$ as $p\to\infty$. Indeed, the set of unimodular lattices with $\lambda_1>\sqrt{\gamma_s}-\varepsilon$ is open and nonempty, so for large $p$ it would contain a scaled $\Lambda_\kappa$ ($\lambda_1$ is continuous, and the sublattices containing a unit vector have normalized minimum at most $p^{-1/s}$), and then $\lambda_{\mathrm{jit}}\ge\lambda_2-\sqrt s$. Such equidistribution is a theorem of Clozel, Oh and Ullmo on Hecke points \cite{COU01}, but as usually stated it is for simply connected $G$ and Hecke operators $T_a$ with $a\in G(\mathbb Q)$, which does not literally contain the index-$p$ correspondence of $\mathrm{diag}(p,1,\dots,1)\notin\mathrm{SL}_s(\mathbb Q)$. The version of Eskin and Oh \cite[Thm.~1.2]{EO06} allows $a$ in the commensurator of the arithmetic group; we have not carried out the reduction in detail, so we do not state the limit as a theorem. The limit would hold at the residue $m\equiv-1\pmod s$, as $p\to\infty$ uniformly in $k$, and would exceed the limit $\min(\sqrt s/2,\sqrt2)$ of Theorem~\ref{thm:HL} (at the residue $m\equiv0$) for $4\le s\le7$ and equal it at $s=8$.
\end{remark}


\section{Size of the constants in Theorem~\ref{thm:S2} (heuristic)}

\begin{remark}[Constants]\label{rem:skelconst}
\uses{lem:G2N1}
$B(p)$ and $\rho(p)$ are nondecreasing, so the bounds are positive for every $p\ge p_0$. For $s\ge5$ a rough count gives $B\approx(p/\sqrt2)^{2/s}$, i.e. $\delta N^{1/s}\to2^{-1/s}$ (heuristic). For $s\ge9$ the route of Lemma~\ref{lem:G2N1}(c) needs $B$ of order $(2s/\pi)^2$, which suggests $p_0(s)=e^{O(s\log s)}$ (heuristic). Nothing is optimized: the faults cost $g^*$ cells at small $p$, and for even $r$ a \emph{block radix}, in which each of the two rows of $A$ is a radix chain on its own group of the full coordinates, has a better constant.
\end{remark}


\chapter{Per-prime designs in dimension three}\label{app:perprime}

\begin{definition}\label{def:ppdesign}
\lean{Z3NetKN.KN1_net, Z3NetKN.KN2_net}\leanok
A \emph{per-prime design} for $(p,r)$ assigns to each prescribed row $R$ and each coordinate $d$ integers $R|d^{(i)}$, $0\le i<L_d$. A graph form over $\F_p$ with depths $\mathbf k$ \emph{carries} it if (i) each prescribed row $R$ has the coefficient $R|d^{(i)}\bmod p$ on level $k_d-1-i$ of $d$ and $0$ on the higher levels of $d$, and (ii) each offset row of $c$ below its prescribed rows has the coefficient $0$ on the cell of $c$.
\end{definition}

\begin{lemma}[Invariants]\label{lem:ppinv}
\lean{Z3NetPI.ppinv_window, Z3NetPI.ppinv_a2, Z3NetPI.ppinv_a1, Z3NetPI.ppinv_b2, Z3NetPI.ppinv_b1}\leanok
\uses{lem:box}
(a) Where defined, $\chi_1,\chi_2$ and $J_1,\dots,J_5$ are unchanged by the moves (a), (c) of Lemma~\ref{lem:box} and the block scalings (b).
(b) For every prime $p$, $\Phi^F$ has $\chi_1=\chi_2=-1$ for every $k\ge1$ ($r=2$), and $(J_1,\dots,J_5)=(1,-1,2,2,1)$ for every $k\ge2$ ($r=1$).
\end{lemma}

\begin{proof}
\leanok
\uses{prop:faure}
(a) Diagonal blocks are not read. A move (c) scales the last level of a cell and adds to every other level multiples of deeper ones; so $\tilde x,\tilde y,\tilde\zeta,\tilde\eta,\tilde\xi$, each fixed by its leading level and a vanishing condition, are multiplied by units, and $d_x$ by the product of the factors of $\xi_0,\xi_1$. $L_c$ scales the first offset row of $c$ and replaces the second by a unit multiple of itself plus a multiple of the first; so $R_Y$, $R_Z$ (the combinations vanishing on $\xi_0$) and $Z1$ restricted to $\tilde x,\tilde y$ (where $Z0$ vanishes) are multiplied by units, and the bracket in $\chi_2$ is invariant under $Z1\mapsto Z1+\lambda Z0$ and scales like $Z1/Z0$. In each quantity every row, column and off-diagonal block occurs equally often in numerator and denominator, so all unit factors, including the block scalings, cancel.

(b) Let $\varphi_\infty(V)$, $\varphi_0(V)$, $\varphi_1(V)$ be the coefficient vectors of $V\in\F_p[t]_{<m}$ with respect to $t^{m-1},t^{m-2},\dots$, resp. $1,t,\dots$, resp. $1,(t-1),\dots$. The first $d$ digits of $[\Psi^2V]_m$ vanish iff $\deg V<m-d$ (as $\Psi^2V(t)=(-t)^{m-1}V((t-1)/t)$), i.e. iff the first $d$ entries of $\varphi_\infty(V)$ do; those of $[\Psi V]_m$ vanish iff $(t-1)^d\mid V$ (proof of Proposition~\ref{prop:faure}), i.e. iff those of $\varphi_1(V)$ do. Invertible matrices whose leading $d$ rows have equal kernels for every $d$ differ by a lower-triangular left factor; so the frame readings are $(T_x\varphi_\infty,\varphi_0,T_z\varphi_1)$, up to the unit $s_k$, with $T_x,T_z$ lower-triangular. Multiplying the reading of $c$ by such a $T_c$ changes $\Phi_{cd}$ into $T''_c\Phi_{cd}(T'_d)^{-1}$ ($d\ne c$; $T'_d$, $T''_c$ the leading $k_d\times k_d$ block of $T_d$ and the trailing diagonal block of $T_c$) and changes $\Phi_{cc}$, i.e. it is a combination of moves (a), (c). By (a) we may use the readings $(\varphi_\infty,\varphi_0,\varphi_1)$. Then the column of level $j$ of $d$ is the offset part of the unique $V$, $\deg V<m$, whose cell digits (the top $k_x$ coefficients, the coefficients of $t^0,\dots,t^{k_y-1}$, those of $(t-1)^0,\dots,(t-1)^{k_z-1}$) are $1$ at $(d,j)$ and $0$ elsewhere; $X0$ reads the coefficient of $t^{m-2-k}$, the first offset rows of $y$ those of $t^{k_y},t^{k_y+1}$, and $Z0,Z1$ those of $(t-1)^k,(t-1)^{k+1}$. Let $\bar t=t-1$, $\epsilon=(-1)^k$ and $\beta_i=\binom{k-1+i}{i}$, the coefficients of $(1-t)^{-k}$.

For $r=2$ the degree bound, the vanishing orders at $0$ and $1$ (for $\tilde x,\tilde y$ one more factor $\bar t$, as $Z0=0$) and one normalization give $x_k:\,t^{k+1}\bar t^k$, $\tilde x:\,t^{k+1}\bar t^{k+1}$, $y_k:\,\epsilon\,t^k\bar t^k$, $\tilde y:\,-\epsilon\,t^{k-1}\bar t^{k+1}$. So $(Y0,Z0,Z1)|x_k=(\epsilon,1,k+1)$, $(Y0,Z1)|\tilde x=(-\epsilon,1)$, $(X0,Z0,Z1)|y_k=(\epsilon,\epsilon,k\epsilon)$, $(X0,Z1)|\tilde y=(-\epsilon,-\epsilon)$, and
\[
\chi_1=\frac{(-\epsilon)(-\epsilon)\,\epsilon\cdot1}{1\cdot\epsilon\,(-\epsilon)\,\epsilon}=-1,\qquad \chi_2=\bigl((k+1)-k\bigr)\frac{1\cdot(-\epsilon)}{1\cdot\epsilon}=-1.
\]
For $r=1$, inverting $t^k$ resp. $\bar t^k$ as a power series at the other finite place, the columns are
\[
\xi_j:\ t^k\bar t^k\textstyle\sum_{i\le j}\beta_it^{j-i},\qquad \eta_j:\ \epsilon\,\bar t^kt^{k-1-j}\sum_{i\le j}\beta_it^i,\qquad \zeta_j:\ t^k\bar t^{k-1-j}\sum_{i\le j}(-1)^i\beta_i\bar t^i,
\]
with the entries
\begin{center}\small
\begin{tabular}{@{}lccccc@{}}
\toprule
 & $X0$ & $Y0$ & $Y1$ & $Z0$ & $Z1$\\
\midrule
$\xi_0$ & & $\epsilon$ & $-k\epsilon$ & $1$ & $k$\\
$\xi_1$ & & $k\epsilon$ & $-(k^2-1)\epsilon$ & $k+1$ & $k^2+k+1$\\
$\xi_2$ & & $\epsilon\binom{k+1}2$ & $-\epsilon k(k-1)(k+2)/2$ & $\binom{k+2}2$ & $(k+2)(k^2+k+2)/2$\\
$\eta_0$, $\eta_1$ & $\epsilon$, $k\epsilon$ & & & $\epsilon$, $(k+1)\epsilon$ & $(k-1)\epsilon$, $(k^2-2)\epsilon$\\
$\zeta_0$, $\zeta_1$ & $1$, $-k$ & $-\epsilon$, $(k+1)\epsilon$ & $(k-1)\epsilon$, $(2-k^2)\epsilon$ & & \\
\bottomrule
\end{tabular}
\end{center}
Hence $d_x=\epsilon$, $R_Y=Y1+kY0$, $R_Z=Z1-kZ0$, $\tilde\zeta=\zeta_1+k\zeta_0$, $\tilde\eta=\eta_1-k\eta_0$, $\tilde\xi=\xi_2-(k+1)\xi_1+\binom{k+1}2\xi_0$,
\[
R_Y|(\xi_1,\zeta_0,\tilde\zeta,\tilde\xi)=(\epsilon,-\epsilon,2\epsilon,-\epsilon),\qquad R_Z|(\xi_1,\eta_0,\tilde\eta,\tilde\xi)=(1,-\epsilon,-2\epsilon,1),\qquad Y0|\tilde\zeta=Z0|\tilde\eta=\epsilon,
\]
and substituting gives $(J_1,\dots,J_5)=(1,-1,2,2,1)$. All entries are integers and all divisors are $\pm1$, so this holds for every $k$ and every $p$.
\end{proof}

\begin{proposition}[Existence, $m\equiv2$]\label{prop:pp2}
\lean{Z3NetKN.KN1_net, Z3NetKN.KN2_net, Z3NetKN.pp2_net}\leanok
\uses{lem:box}
Let $D$ be a per-prime design for $(p,2)$ with
\begin{enumerate}

\item[(V2)] $Z0|x_{k-1}=Z0|x_{k-2}=Z0|y_{k-1}=Z0|y_{k-2}=0$ and $X0|z_{k-3}=Y0|z_{k-3}=0$;
\item[(U2)] the following are nonzero modulo $p$: the entries $Z0|x_k$, $Z0|y_k$, $Y0|x_k$, $Y0|x_{k-1}$, $Z1|x_{k-1}$, $X0|y_k$, $X0|y_{k-1}$, $X0|z_{k-1}$, $Y0|z_{k-1}$ and the brackets $[Y0,Z1\,|\,x_{k-2},x_{k-1}]$, $[X0,Z1\,|\,y_{k-2},y_{k-1}]$, $[X0,Y0\,|\,z_{k-2},z_{k-1}]$;
\item[(I2)] $\chi_1(D)=\chi_2(D)=-1$ in $\F_p$.
\end{enumerate}
Then for every $k\ge2$ the moves of Lemma~\ref{lem:box} turn $\Phi^F$ into a digital $(0,3k+2,3)$-net over $\F_p$ that carries $D$.
\end{proposition}

\begin{proof}
\leanok
\uses{lem:nie, def:ppdesign, lem:ppinv}
All moves are moves (a)--(c), so every intermediate graph form has unit box minors; each pivot below is nonzero because, given the entries already normalized, it is a box minor or a quotient of two (box in brackets). An entry is \emph{set} when it is made equal to its value in $D$.

\emph{Step 1 (windows).} $Y0,Z0,Z1$ on $x_{k-2},x_{k-1},x_k$ form an invertible matrix (box $(k-2,k+2,k+2)$); subtracting combinations of these levels from the levels $x_l$, $l<k-2$, clears the three rows there. Likewise $X0,Z0,Z1$ on $y_{k-2},y_{k-1},y_k$ (box $(k+2,k-2,k+2)$), and $X0,Y0$ on $z_l$, $l<k-2$, with $z_{k-2},z_{k-1}$ (box $(k+2,k+2,k-2)$). Then subtract multiples of $x_k$ from $x_{k-1},x_{k-2}$ to clear $Z0$ there (pivot $Z0|x_k$, box $(k,k+1,k+1)$), and likewise in $y$ (box $(k+1,k,k+1)$); now $x_{k-1}=\tilde x$, $y_{k-1}=\tilde y$.

\emph{Step 2.} Scale $\Phi_{yx}$ to set $Y0|x_k$ (box $(k,k+2,k)$) and the level $x_{k-1}$ to set $Y0|x_{k-1}$ (box $(k-1,k+2,k+1)$); apply $L_z$: $Z0\mapsto\ell Z0$, $Z1\mapsto\ell'Z1+\lambda Z0$, to set $Z0|x_k$, $Z1|x_{k-1}$ (box $(k-1,k+1,k+2)$) and $Z1|x_k$.

\emph{Step 3.} Scale $\Phi_{xy}$ to set $X0|y_k$ (box $(k+2,k,k)$), the level $y_{k-1}$ to set $X0|y_{k-1}$ (box $(k+2,k-1,k+1)$), and $\Phi_{zy}$ to set $Z0|y_k$ (box $(k+1,k,k+1)$).

\emph{Step 4 (forced entries).} This is where the invariants enter. All entries read by $\chi_1,\chi_2$ are now set except $Z1|y_{k-1}$, $Z1|y_k$. By Lemma~\ref{lem:ppinv} and (I2) the current $\chi_1,\chi_2$ equal those of $D$; as $\chi_1$ is a nonzero multiple of $Z1|y_{k-1}$ and $\chi_2$ is affine with nonzero slope in $Z1|y_k$, these entries already have their values in $D$, and no move is needed (or possible) to set them.

\emph{Step 5.} Replace $x_{k-2}$ by $a\,x_{k-2}+b\,x_{k-1}$ so as to set $Y0|x_{k-2}$, $Z1|x_{k-2}$; the matrix $[Y0,Z1\,|\,x_{k-2},x_{k-1}]$ is invertible (box $(k-2,k+2,k+2)$, divided by $Z0|x_k$), and $Z0$ stays $0$ on $x_{k-2}$. Afterwards the bracket is $a$ times its old value and equals that of $D$, nonzero by (U2); so $a\ne0$ and this is a move (c). The same for $y_{k-2}$ (box $(k+2,k-2,k+2)$).

\emph{Step 6.} Scale $\Phi_{xz}$, $\Phi_{yz}$ to set $X0|z_{k-1}$, $Y0|z_{k-1}$ (boxes $(k+2,k+1,k-1)$, $(k+1,k+2,k-1)$), and replace $z_{k-2}$ by $a\,z_{k-2}+b\,z_{k-1}$ to set $X0|z_{k-2}$, $Y0|z_{k-2}$ (box $(k+2,k+2,k-2)$; $a\ne0$ as in Step 5).

\emph{Step 7.} By move (a), give each coordinate's prescribed rows their values in $D$ on the window of its own cell ($0$ above it) and its other offset rows the value $0$ on its own cell.

Steps 2--6 only combine window levels, so the prescribed rows still vanish above the windows (as Definition~\ref{def:ppdesign}(i) requires), and the graph form carries $D$. Lemma~\ref{lem:nie} completes the proof.
\end{proof}

\begin{proposition}[Existence, $m\equiv1$]\label{prop:pp1}
\lean{Z3NetKN.KN1_net, Z3NetKN.KN2_net, Z3NetKN.pp1_net, Z3NetKN.pp1_net_all}\leanok
\uses{lem:box}
Let $D$ be a per-prime design for $(p,1)$ with
\begin{enumerate}

\item[(V1)] $Y0,Z0$ vanish on $x_{k-2},x_{k-3}$; $X0,Z0$ on $y_{k-3}$; $X0,Y0$ on $z_{k-3}$;
\item[(U1)] the following are nonzero modulo $p$: the entries $Y0|x_k$, $Z0|x_k$, $X0|y_{k-1}$, $Z0|y_{k-1}$, $X0|z_{k-1}$, $Y0|z_{k-1}$, $Y1|x_{k-2}$, $Z1|y_{k-3}$, $Y1|z_{k-3}$ and the brackets
\[[Y0,Z0\,|\,x_{k-1},x_k],\quad[X0,Z0\,|\,y_{k-2},y_{k-1}],\quad[X0,Y0\,|\,z_{k-2},z_{k-1}],\]
\[[Y1,Z1\,|\,x_{k-3},x_{k-2}],\quad[Y1,Y0\,|\,x_k,z_{k-1}],\quad[Z1,Z0\,|\,x_k,y_{k-1}];\]
\item[(I1)] $(J_1,\dots,J_5)(D)=(1,-1,2,2,1)$ in $\F_p$.
\end{enumerate}
Then for every $k\ge3$ the moves of Lemma~\ref{lem:box} turn $\Phi^F$ into a digital $(0,3k+1,3)$-net over $\F_p$ that carries $D$.
\end{proposition}

\begin{proof}
\leanok
\uses{lem:ppinv, prop:pp2}
As for Proposition~\ref{prop:pp2}; in each solve below, $a\ne0$ because the corresponding bracket of $D$ is nonzero (U1).

\emph{Step 1 (windows).} Clear $Y0,Y1,Z0,Z1$ on $x_l$, $l<k-3$ (box $(k-3,k+2,k+2)$), $X0,Z0,Z1$ on $y_l$, $l<k-3$ (box $(k+2,k-3,k+2)$), and $X0,Y0,Y1$ on $z_l$, $l<k-3$ (box $(k+2,k+2,k-3)$). Then clear $Y0,Z0$ on $x_{k-2},x_{k-3}$ with $x_{k-1},x_k$ (box $(k-1,k+1,k+1)$), $X0,Z0$ on $y_{k-3}$ with $y_{k-2},y_{k-1}$ (box $(k+2,k-2,k+1)$), and $X0,Y0$ on $z_{k-3}$ with $z_{k-2},z_{k-1}$ (box $(k+2,k+1,k-2)$).

\emph{Step 2 ($X0,Y0,Z0$).} Scale $\Phi_{yx}$, $\Phi_{zx}$ to set $Y0|x_k$, $Z0|x_k$ (boxes $(k,k+1,k)$, $(k,k,k+1)$) and replace $x_{k-1}$ by $a\,x_{k-1}+b\,x_k$ to set $Y0|x_{k-1}$, $Z0|x_{k-1}$ (box $(k-1,k+1,k+1)$). In the same way set $X0|y_{k-1}$, $Z0|y_{k-1}$ by $\Phi_{xy}$, $\Phi_{zy}$ and $X0|y_{k-2}$, $Z0|y_{k-2}$ by $y_{k-2}\mapsto a\,y_{k-2}+b\,y_{k-1}$ (boxes $(k+2,k-1,k)$, $(k+1,k-1,k+1)$, $(k+2,k-2,k+1)$), and $X0|z_{k-1}$, $Y0|z_{k-1}$, $X0|z_{k-2}$, $Y0|z_{k-2}$ by $\Phi_{xz}$, $\Phi_{yz}$ and $z_{k-2}$ (boxes $(k+2,k,k-1)$, $(k+1,k+1,k-1)$, $(k+2,k+1,k-2)$).

\emph{Step 3 ($Y1,Z1$).} Apply $L_y$: $Y1\mapsto aY1+bY0$ to set $Y1|x_k$, $Y1|z_{k-1}$ (matrix $[Y1,Y0\,|\,x_k,z_{k-1}]$, box $(k,k+2,k-1)$), and $L_z$: $Z1\mapsto cZ1+dZ0$ to set $Z1|x_k$, $Z1|y_{k-1}$ (box $(k,k-1,k+2)$).

\emph{Step 4 (forced entries).} All entries read by $J_1,\dots,J_4$ are now set except $Y1|x_{k-1}$, $Z1|x_{k-1}$, $Y1|z_{k-2}$, $Z1|y_{k-2}$, and $J_1,\dots,J_4$ are affine with nonzero slope in one of them each (through $R_Y|\xi_1$, $R_Z|\xi_1$, $R_Y|\tilde\zeta$, $R_Z|\tilde\eta$). By Lemma~\ref{lem:ppinv} and (I1) they already have their values in $D$.

\emph{Step 5.} Scale the levels $z_{k-3}$, $y_{k-3}$, $x_{k-2}$ to set $Y1|z_{k-3}$, $Z1|y_{k-3}$, $Y1|x_{k-2}$ (boxes $(k+2,k+2,k-3)$, $(k+2,k-3,k+2)$, $(k-2,k+2,k+1)$). Now $\tilde\xi=\xi_2$, as $Y0,Z0$ vanish on $x_{k-2}$, and $J_5$ is affine with nonzero slope in $Z1|x_{k-2}$; by Lemma~\ref{lem:ppinv} and (I1) this entry has its value in $D$. Finally replace $x_{k-3}$ by $a\,x_{k-3}+b\,x_{k-2}$ to set $Y1|x_{k-3}$, $Z1|x_{k-3}$ (box $(k-3,k+2,k+2)$).

\emph{Step 6} is Step 7 of Proposition~\ref{prop:pp2}.
\end{proof}

\begin{definition}[Per-prime relaxation]\label{def:pprelax}
\lean{Z3NetKN.ppsep_general, Z3NetKN.rlxLin, Z3NetKN.rlxRes, Z3NetKN.rlxT, Z3NetKN.rlxTerm, Z3NetKN.IsPPCert}\leanok
Let $D$ be a per-prime design for $(p,r)$ and $B\in\Z_{>0}^3$ a search radius. A \emph{configuration} is a pair $(E,j)$ of a cell difference $E\in\Z^3\setminus\{0\}$, $|E_c|\le B_c$, and carry depths $j_c\in\{0,\dots,L_c-1\}$, with $j_c=0$ if $E_c=0$. Put $\Delta_{d,0}=E_d$, $\Delta_{d,i}=\sgn(E_d)\,[i\le j_d]$ for $1\le i<L_d$, and let $\sigma_c$ resp. $\iota_c\in\{0,\dots,p-1\}$ be the residue of $\sum_{d,i}R|d^{(i)}\Delta_{d,i}$, where $R$ is the skeleton row resp. the row of the inner digit of $c$. With $\mathcal A(v)=\{v,v-p\}$ for $v\ne0$, $\mathcal A(0)=\{0\}$, and $\varepsilon_c=0$ if $E_d=0$ for all $d\ne c$, $\varepsilon_c=1$ otherwise, let
\[
t_c=\min_{\beta\in\mathcal A(\iota_c)}\bigl(|pE_c+\beta|-\varepsilon_c\bigr)^+\ \ (c\text{ full}),\qquad
t_c=\min_{\alpha\in\mathcal A(\sigma_c),\,\beta\in\mathcal A(\iota_c)}\bigl(|p(pE_c+\alpha)+\beta|-\varepsilon_c\bigr)^+\ \ (c\text{ skeleton}).
\]
A \emph{certificate} is a pair $(T,B)$ with $T\le t_x^2+t_y^2+t_z^2$ for every configuration, $B_c<p$ for all $c$, $T\le p^2B_c^2$ for full $c$, and $B_c=1$, $T\le p^4$ for skeleton $c$.
\end{definition}

\begin{proposition}[Separation]\label{prop:ppsep}
\lean{Z3NetKN.KN1_sep, Z3NetKN.KN2_sep, Z3NetKN.ppsep_general}\leanok
Let $p$ be odd, $D$ a per-prime design for $(p,r)$ with a certificate $(T,B)$, and $k\ge3$ ($r=1$) resp. $k\ge2$ ($r=2$). Every graph form over $\F_p$ with depths $\mathbf k$ carrying $D$, net or not, has $\|X(w)-X(v)\|^2_{\mathrm{torus}}\ge Tg_0^2$ for all indices $v\ne w$ (integer coordinates on the torus of side $N$). Hence $\delta N^{1/3}\ge\sqrt T\,p^{-5/3}$ ($r=1$) resp. $\sqrt T\,p^{-4/3}$ ($r=2$).
\end{proposition}

\begin{proof}
\leanok
\uses{lem:skelred, def:ppdesign}
Let $E_c$ be the centred residue of $n_c(w)-n_c(v)$ modulo $p^{k_c}$; as the cells are the index, $E\ne0$. The proof of Lemma~\ref{lem:skelred}(i) only uses $X_c=g\,n_c+\theta$, $0\le\theta<g$, $N=g\,p^{k_c}$ and the centring; with $g=pg_0$ (full) resp. $p^2g_0$ (skeleton) the toroidal distance in coordinate $c$ is at least $(|E_c|-1)g$. So if $|E_c|>B_c$ for a full $c$, or $|E_c|\ge2$ for a skeleton $c$, the claim follows from $T\le p^2B_c^2$ resp. $T\le p^4$.

Otherwise $|E_c|\le B_c<p$ for all $c$. For $E_c\ne0$ the digits of $n_c(w)$ and $n_c(v)$ differ, modulo $p$, by $E_c$ on the last level, by $\sgn E_c$ on the next $j'_c\ge0$ levels (a carry chain, cf. Lemma~\ref{lem:skelred}(ii)), and not above; with $j_c=\min(j'_c,L_c-1)$, $(E,j)$ is a configuration. As the prescribed rows have the coefficients of $D$ on the windows and $0$ above, the skeleton and inner digits of $c$ change by $\sigma_c$ resp. $\iota_c$ modulo $p$, and since they lie in $[0,p)$, $S_c(w)-S_c(v)\in\mathcal A(\sigma_c)$ and $\rho_c(w)-\rho_c(v)\in\mathcal A(\iota_c)$. Further $|\theta_c(w)-\theta_c(v)|<g_0$, with equality $\theta_c(w)=\theta_c(v)$ if $\varepsilon_c=0$: the deeper digits of $c$ have coefficient $0$ on the cell of $c$ (Definition~\ref{def:ppdesign}(ii)), and no other cell changes. Hence $X_c(w)-X_c(v)\equiv g_0y_c+\theta_c(w)-\theta_c(v)\pmod N$ with $y_c=pE_c+\beta$ resp. $p(pE_c+\alpha)+\beta$ for some $\alpha\in\mathcal A(\sigma_c)$, $\beta\in\mathcal A(\iota_c)$; as $|y_c|<3p^2$ and $N/g_0=p^{k+2}\ge p^4$, the toroidal distance is at least $g_0(|y_c|-\varepsilon_c)^+\ge g_0t_c$. Summing squares gives $\|X(w)-X(v)\|^2_{\mathrm{torus}}\ge g_0^2\sum_ct_c^2\ge Tg_0^2$. Finally $\delta N^{1/3}\ge\sqrt T\,g_0N^{-2/3}=\sqrt T\,p^{(m-3k-6)/3}$.
\end{proof}

\begin{remark}[Digital shifts]\label{rem:ppshift}
\lean{Z3NetKN.ppsep_shift, Z3NetKN.KN1_sep_shift, Z3NetKN.KN2_sep_shift}\leanok
\uses{def:dshift, prop:ppsep}
Proposition~\ref{prop:ppsep} holds verbatim, with the same certificate $(T,B)$, for every digital shift (Definition~\ref{def:dshift}) of every graph form carrying $D$. So the constants of Table~\texttt{tab:perprime} hold for every digitally shifted per-prime net.
\end{remark}

\begin{proof}
\leanok
\uses{lem:shift, def:ppdesign, prop:ppsep}
By Lemma~\ref{lem:shift} the shifted point set is indexed by its own cells $n$, one point per cell, and its offsets are $\Phi n+\gamma$, with the blocks $\Phi$ of the unshifted graph form and a constant $\gamma$. Each coordinate still has the form $X_c=g_0(pn_c+\rho_c)+\theta_c$ resp.\ $g_0(p(pn_c+S_c)+\rho_c)+\theta_c$ with $S_c,\rho_c\in[0,p)$ and $0\le\theta_c<g_0$. So the first paragraph of the proof of Proposition~\ref{prop:ppsep} applies, and $E$ and the carry chains are those of the shifted cells, so $(E,j)$ is a configuration. The prescribed rows applied to $n(w)-n(v)$ give $S_c(w)-S_c(v)\equiv\sigma_c$ and $\rho_c(w)-\rho_c(v)\equiv\iota_c\pmod p$, since $\gamma$ cancels. As both digits lie in $[0,p)$, the differences lie in $\mathcal A(\sigma_c)$ resp.\ $\mathcal A(\iota_c)$. A shift may change which of the two representatives occurs, but $t_c$ is the minimum over both. If $\varepsilon_c=0$, the deeper digits of $c$ are given by rows with coefficient $0$ on the cell of $c$ (Definition~\ref{def:ppdesign}(ii)) plus the same constants, so $\theta_c(w)=\theta_c(v)$; otherwise $|\theta_c(w)-\theta_c(v)|<g_0$ as before. The rest of the proof is unchanged.
\end{proof}

\begin{remark}[Formalized statements of this section without a result label]\label{bp:app:perprime}
\lean{Z3NetKN.pp2_table, Z3NetKN.pp1_table}\leanok
The 17 resp.\ 20 per-prime designs of Appendix~\ref{app:perprime} satisfy (S2), (N2), (K) resp.\ (S1), (N1), (J).
\end{remark}


\begin{thebibliography}{99}
\bibitem{BHMN20} J.~Brakensiek, M.~Heule, J.~Mackey, D.~Narv\'aez, \emph{The resolution of Keller's conjecture}, J. Automat. Reason. 66 (2022) 277--300, doi:10.1007/s10817-022-09623-5; conference version in: IJCAR 2020, Lecture Notes in Comput. Sci. 12166, Springer, 2020, 48--65, doi:10.1007/978-3-030-51074-9\_4; arXiv:1910.03740.
\bibitem{BS83} E.~S.~Barnes, N.~J.~A.~Sloane, \emph{New lattice packings of spheres}, Canad. J. Math. 35 (1983) 117--130, doi:10.4153/cjm-1983-008-1.
\bibitem{CE03} H.~Cohn, N.~Elkies, \emph{New upper bounds on sphere packings I}, Ann. of Math. 157 (2003) 689--714, doi:10.4007/annals.2003.157.689; arXiv:math/0110009.
\bibitem{CKMRV17} H.~Cohn, A.~Kumar, S.~D.~Miller, D.~Radchenko, M.~Viazovska, \emph{The sphere packing problem in dimension $24$}, Ann. of Math. 185 (2017) 1017--1033, doi:10.4007/annals.2017.185.3.8; arXiv:1603.06518.
\bibitem{CM67} R.~R.~Coveyou, R.~D.~MacPherson, \emph{Fourier analysis of uniform random number generators}, J. ACM 14 (1967) 100--119, doi:10.1145/321371.321379.
\bibitem{Coo86} R.~L.~Cook, \emph{Stochastic sampling in computer graphics}, ACM Trans. Graph. 5 (1986) 51--72, doi:10.1145/7529.8927.
\bibitem{COU01} L.~Clozel, H.~Oh, E.~Ullmo, \emph{Hecke operators and equidistribution of Hecke points}, Invent. Math. 144 (2001) 327--351, doi:10.1007/s002220100126.
\bibitem{CP76} R.~Cranley, T.~N.~L.~Patterson, \emph{Randomization of number theoretic methods for multiple integration}, SIAM J. Numer. Anal. 13 (1976) 904--914, doi:10.1137/0713071.
\bibitem{CS99} J.~H.~Conway, N.~J.~A.~Sloane, \emph{Sphere Packings, Lattices and Groups}, 3rd ed., Springer, 1999.
\bibitem{DGLPS1} J.~Dick, T.~Goda, G.~Larcher, F.~Pillichshammer, K.~Suzuki, \emph{On the quasi-uniformity properties of quasi-Monte Carlo point sets and sequences, Part I: lattices and Kronecker sequences}, Math. Comp. (2026), doi:10.1090/mcom/4264; arXiv:2502.06202 (theorem numbers cited from arXiv v2).
\bibitem{DGS2} J.~Dick, T.~Goda, K.~Suzuki, \emph{On the quasi-uniformity properties of quasi-Monte Carlo point sets and sequences, Part II: digital nets and sequences}, Math. Comp. (2026), doi:10.1090/mcom/4263; arXiv:2501.18226 (theorem numbers cited from arXiv v2).
\bibitem{dlVP96} C.-J.~de la Vall\'ee Poussin, \emph{Recherches analytiques sur la th\'eorie des nombres premiers}, Ann. Soc. Sci. Bruxelles 20 B (1896) 183--256.
\bibitem{dMU21} L.~de Moura, S.~Ullrich, \emph{The Lean~4 theorem prover and programming language}, in: Automated Deduction -- CADE 28, Lecture Notes in Comput. Sci. 12699, Springer, 2021, 625--635, doi:10.1007/978-3-030-79876-5\_37.
\bibitem{DP10} J.~Dick, F.~Pillichshammer, \emph{Digital Nets and Sequences: Discrepancy Theory and Quasi-Monte Carlo Integration}, Cambridge Univ. Press, 2010.
\bibitem{EO06} A.~Eskin, H.~Oh, \emph{Ergodic theoretic proof of equidistribution of Hecke points}, Ergodic Theory Dynam. Systems 26 (2006) 163--167, doi:10.1017/s0143385705000428; arXiv:math/0409507.
\bibitem{Fau82} H.~Faure, \emph{Discr\'epance de suites associ\'ees \`a un syst\`eme de num\'eration (en dimension $s$)}, Acta Arith. 41 (1982) 337--351, doi:10.4064/aa-41-4-337-351.
\bibitem{FL09} H.~Faure, C.~Lemieux, \emph{Generalized Halton sequences in 2008: a comparative study}, ACM Trans. Model. Comput. Simul. 19 (2009), Art.~15, 1--31, doi:10.1145/1596519.1596520.
\bibitem{GHS26} T.~Goda, R.~Hofer, K.~Suzuki, \emph{Disproving the quasi-uniformity of the Halton sequences and of some Halton-type sequences}, J. Complexity 95 (2026) 102047, doi:10.1016/j.jco.2026.102047; arXiv:2511.21153 (theorem numbers cited from arXiv v2).
\bibitem{GHSK08} L.~Gr\"unschlo{\ss}, J.~Hanika, R.~Schwede, A.~Keller, \emph{$(t,m,s)$-nets and maximized minimum distance}, in: A.~Keller, S.~Heinrich, H.~Niederreiter (eds.), Monte Carlo and Quasi-Monte Carlo Methods 2006, Springer, 2008, 397--412, doi:10.1007/978-3-540-74496-2\_23.
\bibitem{GK09} L.~Gr\"unschlo{\ss}, A.~Keller, \emph{$(t,m,s)$-nets and maximized minimum distance, Part II}, in: P.~L'Ecuyer, A.~B.~Owen (eds.), Monte Carlo and Quasi-Monte Carlo Methods 2008, Springer, 2009, 395--409, doi:10.1007/978-3-642-04107-5\_25.
\bibitem{Had96} J.~Hadamard, \emph{Sur la distribution des z\'eros de la fonction $\zeta(s)$ et ses cons\'equences arithm\'etiques}, Bull. Soc. Math. France 24 (1896) 199--220, doi:10.24033/bsmf.545.
\bibitem{Haj42} G.~Haj\'os, \emph{\"Uber einfache und mehrfache Bedeckung des $n$-dimensionalen Raumes mit einem W\"urfelgitter}, Math. Z. 47 (1942) 427--467, doi:10.1007/bf01180974.
\bibitem{Hal05} T.~C.~Hales, \emph{A proof of the Kepler conjecture}, Ann. of Math. 162 (2005) 1065--1185, doi:10.4007/annals.2005.162.1065.
\bibitem{Ham60} J.~M.~Hammersley, \emph{Monte Carlo methods for solving multivariable problems}, Ann. New York Acad. Sci. 86 (1960) 844--874, doi:10.1111/j.1749-6632.1960.tb42846.x.
\bibitem{HW08} G.~H.~Hardy, E.~M.~Wright, \emph{An Introduction to the Theory of Numbers}, 6th ed., revised by D.~R.~Heath-Brown and J.~H.~Silverman, Oxford Univ. Press, 2008.
\bibitem{He17} X.~He, \emph{Rotated sphere packing designs}, J. Amer. Statist. Assoc. 112 (2017) 1612--1622, doi:10.1080/01621459.2016.1222289.
\bibitem{He19} X.~He, \emph{Interleaved lattice-based maximin distance designs}, Biometrika 106 (2019) 453--464, doi:10.1093/biomet/asy069; arXiv:1807.02289.
\bibitem{He21} X.~He, \emph{Lattice-based designs with quasi-optimal separation distance on all projections}, Biometrika 108 (2021) 443--454, doi:10.1093/biomet/asaa057; arXiv:1709.02062.
\bibitem{HP55} M.~Hall, L.~J.~Paige, \emph{Complete mappings of finite groups}, Pacific J. Math. 5 (1955) 541--549, doi:10.2140/pjm.1955.5.541.
\bibitem{HS19} R.~Hofer, K.~Suzuki, \emph{A complete classification of digital $(0,m,3)$-nets and digital $(0,2)$-sequences in base $2$}, Unif. Distrib. Theory 14 (2019) 43--52, doi:10.2478/udt-2019-0004; arXiv:1807.04383.
\bibitem{KB79} R.~Kannan, A.~Bachem, \emph{Polynomial algorithms for computing the Smith and Hermite normal forms of an integer matrix}, SIAM J. Comput. 8 (1979) 499--507, doi:10.1137/0208040.
\bibitem{Knu97} D.~E.~Knuth, \emph{The Art of Computer Programming, Vol.~2: Seminumerical Algorithms}, 3rd ed., Addison-Wesley, 1997.
\bibitem{KO14} W.~Kositwattanarerk, F.~Oggier, \emph{Connections between Construction D and related constructions of lattices}, Des. Codes Cryptogr. 73 (2014) 441--455, doi:10.1007/s10623-014-9939-3; arXiv:1308.6175.
\bibitem{Kor59} N.~M.~Korobov, \emph{The approximate computation of multiple integrals}, Dokl. Akad. Nauk SSSR 124 (1959) 1207--1210.
\bibitem{Kri21} R.~Kritzinger, \emph{Dispersion of digital $(0,m,2)$-nets}, Monatsh. Math. 195 (2021) 155--171, doi:10.1007/s00605-021-01525-9; arXiv:2004.14760.
\bibitem{Len77} H.~W.~Lenstra, Jr., \emph{Euclidean number fields of large degree}, Invent. Math. 38 (1976/77) 237--254, doi:10.1007/bf01403131.
\bibitem{L15} C.~Lemieux, \emph{Tractability using periodized generalized Faure sequences}, J. Complexity 31 (2015) 42--56, doi:10.1016/j.jco.2014.08.001.
\bibitem{LF09} C.~Lemieux, H.~Faure, \emph{New perspectives on $(0,s)$-sequences}, in: P.~L'Ecuyer, A.~B.~Owen (eds.), Monte Carlo and Quasi-Monte Carlo Methods 2008, Springer, 2009, 113--130, doi:10.1007/978-3-642-04107-5\_7.
\bibitem{LNS96} G.~Larcher, H.~Niederreiter, W.~Ch.~Schmid, \emph{Digital nets and sequences constructed over finite rings and their application to quasi-Monte Carlo integration}, Monatsh. Math. 121 (1996) 231--253, doi:10.1007/bf01298952.
\bibitem{LP01} G.~Larcher, F.~Pillichshammer, \emph{Walsh series analysis of the $L_2$-discrepancy of symmetrisized point sets}, Monatsh. Math. 132 (2001) 1--18, doi:10.1007/s006050170054.
\bibitem{LS92} J.~C.~Lagarias, P.~W.~Shor, \emph{Keller's cube-tiling conjecture is false in high dimensions}, Bull. Amer. Math. Soc. (N.S.) 27 (1992) 279--283, doi:10.1090/s0273-0979-1992-00318-x; arXiv:math/9210222.
\bibitem{Mac02} J.~Mackey, \emph{A cube tiling of dimension eight with no facesharing}, Discrete Comput. Geom. 28 (2002) 275--279, doi:10.1007/s00454-002-2801-9.
\bibitem{Mat98} J.~Matou\v sek, \emph{On the $L_2$-discrepancy for anchored boxes}, J. Complexity 14 (1998) 527--556, doi:10.1006/jcom.1998.0489.
\bibitem{Mathlib20} The mathlib Community, \emph{The Lean mathematical library}, in: Proc. 9th ACM SIGPLAN Int. Conf. on Certified Programs and Proofs (CPP 2020), ACM, 2020, 367--381, doi:10.1145/3372885.3373824.
\bibitem{MS77} F.~J.~MacWilliams, N.~J.~A.~Sloane, \emph{The Theory of Error-Correcting Codes}, North-Holland, 1977.
\bibitem{Nie87} H.~Niederreiter, \emph{Point sets and sequences with small discrepancy}, Monatsh. Math. 104 (1987) 273--337, doi:10.1007/bf01294651.
\bibitem{Owe95} A.~B.~Owen, \emph{Randomly permuted $(t,m,s)$-nets and $(t,s)$-sequences}, in: H.~Niederreiter, P.~J.-S.~Shiue (eds.), Monte Carlo and Quasi-Monte Carlo Methods in Scientific Computing, Lecture Notes in Statist. 106, Springer, 1995, 299--317, doi:10.1007/978-1-4612-2552-2\_19.
\bibitem{Pai51} L.~J.~Paige, \emph{Complete mappings of finite groups}, Pacific J. Math. 1 (1951) 111--116, doi:10.2140/pjm.1951.1.111.
\bibitem{Per40} O.~Perron, \emph{\"Uber l\"uckenlose Ausf\"ullung des $n$-dimensionalen Raumes durch kongruente W\"urfel}, Math. Z. 46 (1940) 1--26, doi:10.1007/bf01181421; \emph{II}, Math. Z. 46 (1940) 161--180, doi:10.1007/bf01181436.
\bibitem{PFA09} L.~Panek, M.~Firer, M.~M.~S.~Alves, \emph{Symmetry groups of Rosenbloom--Tsfasman spaces}, Discrete Math. 309 (2009) 763--771, doi:10.1016/j.disc.2008.01.013.
\bibitem{RS85} R.~M.~Roth, G.~Seroussi, \emph{On generator matrices of MDS codes (Corresp.)}, IEEE Trans. Inform. Theory 31 (1985) 826--830, doi:10.1109/tit.1985.1057113.
\bibitem{RT97} M.~Yu.~Rosenbloom, M.~A.~Tsfasman, \emph{Codes for the $m$-metric}, Probl. Inf. Transm. 33 (1997) 45--52.
\bibitem{Rud26} L.~Rudolph, \emph{The even-order Gr\"unschlo{\ss}--Keller permutation nets are $(0,m,2)$-nets}, preprint, Zenodo, version 1.0.0, 2026, doi:10.5281/zenodo.22264675.
\bibitem{SG26} N.~Sakai, T.~Goda, \emph{Space-filling lattice designs for computer experiments}, arXiv:2602.15390 (2026).
\bibitem{SJ94} I.~H.~Sloan, S.~Joe, \emph{Lattice Methods for Multiple Integration}, Oxford Univ. Press, 1994.
\bibitem{Skr01} M.~M.~Skriganov, \emph{Coding theory and uniform distributions}, Algebra i Analiz 13 (2001) 191--239; English transl. St. Petersburg Math. J. 13 (2002) 301--337; arXiv:math/9909163 (theorem numbers cited from the arXiv version).
\bibitem{Sob67} I.~M.~Sobol', \emph{On the distribution of points in a cube and the approximate evaluation of integrals}, Zh. Vychisl. Mat. Mat. Fiz. 7 (1967) 784--802; English transl. USSR Comput. Math. Math. Phys. 7 (1967) 86--112, doi:10.1016/0041-5553(67)90144-9.
\bibitem{SS07} I.~M.~Sobol', B.~V.~Shukhman, \emph{Quasi-random points keep their distance}, Math. Comput. Simulation 75 (2007) 80--86, doi:10.1016/j.matcom.2006.09.004.
\bibitem{Suz25} K.~Suzuki, \emph{Exact $\ell^\infty$-separation radius of Sobol' sequences in dimension~$2$}, arXiv:2508.14803 (2025).
\bibitem{Suz26} K.~Suzuki, \emph{Separation properties of scrambled digital nets and related random point sets}, arXiv:2607.29063 (2026).
\bibitem{Tez93} S.~Tezuka, \emph{Polynomial arithmetic analogue of Halton sequences}, ACM Trans. Model. Comput. Simul. 3 (1993) 99--107, doi:10.1145/169702.169694.
\bibitem{Tez95} S.~Tezuka, \emph{Uniform Random Numbers: Theory and Practice}, Kluwer, 1995.
\bibitem{TV06} R.~Tavildar, P.~Viswanath, \emph{Approximately universal codes over slow-fading channels}, IEEE Trans. Inform. Theory 52 (2006) 3233--3258, doi:10.1109/tit.2006.876226; arXiv:cs/0512017.
\bibitem{vDHHM07} E.~R.~van Dam, B.~Husslage, D.~den Hertog, H.~Melissen, \emph{Maximin Latin hypercube designs in two dimensions}, Oper. Res. 55 (2007) 158--169, doi:10.1287/opre.1060.0317.
\bibitem{vDRH09} E.~R.~van Dam, G.~Rennen, B.~Husslage, \emph{Bounds for maximin Latin hypercube designs}, Oper. Res. 57 (2009) 595--608, doi:10.1287/opre.1080.0604.
\bibitem{VG06} P.~O.~Vontobel, A.~Ganesan, \emph{On universally decodable matrices for space-time coding}, Des. Codes Cryptogr. 41 (2006) 325--342, doi:10.1007/s10623-006-9019-4; arXiv:cs/0601075.
\bibitem{Via17} M.~S.~Viazovska, \emph{The sphere packing problem in dimension $8$}, Ann. of Math. 185 (2017) 991--1015, doi:10.4007/annals.2017.185.3.7.
\end{thebibliography}
